% Generated mechanically from frozen input inputs/p05/ja/perfect.tex.
% Do not edit this generated file; edit the bound locale source instead.

% OK, start here.
%

\setcounter{chapter}{35}
\chapter{スキームの導来圏}



\phantomsection
\label{perfect-section-phantom}


\section{序論}
\label{perfect-section-introduction}

\noindent
本章では、スキーム上の加群の導来圏について論じる。
ここで扱う内容の大部分は
\cite{TT}, \cite{Bokstedt-Neeman}, \cite{BvdB}, および \cite{LN}
に見いだすことができる。もちろん、ほかにも多くの参考文献がある。


\section{規約}
\label{perfect-section-conventions}

\noindent
$\mathcal{A}$ をアーベル圏とし、$M$ を $\mathcal{A}$ の対象とする。このとき、
次数 $0$ に $M$ をもち、それ以外のすべての次数で零である複体に対応する
$K(\mathcal{A})$ および／または $D(\mathcal{A})$ の対象も $M$ と記す。

\medskip\noindent
環 $A$ が与えられたとき、$K(A)$ は $A$-加群の複体のホモトピー圏を表し、
$D(A)$ はそれに付随する導来圏を表す。同様に、環付き空間
$(X, \mathcal{O}_X)$ が与えられたとき、記号 $K(\mathcal{O}_X)$ は
$\mathcal{O}_X$-加群の複体のホモトピー圏を表し、$D(\mathcal{O}_X)$ は
それに付随する導来圏を表す。










\section{準連接加群の導来圏}
\label{perfect-section-derived-quasi-coherent}

\noindent
本節では、スキーム $X$ 上の準連接加群とすべての加群との関係を論じる。
参考文献として \cite[Appendix B]{TT} がある。
Schemes, Section \ref{schemes-section-quasi-coherent} の議論により、埋め込み
$\QCoh(\mathcal{O}_X) \subset \textit{Mod}(\mathcal{O}_X)$
は $\QCoh(\mathcal{O}_X)$ を $\mathcal{O}_X$-加群の圏の弱 Serre 部分圏として
表示する。コホモロジー層が準連接である複体からなる部分圏を
$$
D_\QCoh(\mathcal{O}_X) \subset D(\mathcal{O}_X)
$$
と記す。Derived Categories, Section \ref{derived-section-triangulated-sub}
を参照せよ。したがって、標準的な関手
\begin{equation}
\label{perfect-equation-compare}
D(\QCoh(\mathcal{O}_X))
\longrightarrow
D_\QCoh(\mathcal{O}_X)
\end{equation}
を得る。Derived Categories, Equation (\ref{derived-equation-compare})
を参照せよ。

\begin{lemma}
\label{perfect-lemma-quasi-coherence-direct-sums}
スキーム $X$ をとる。このとき $D_\QCoh(\mathcal{O}_X)$ は直和をもつ。
\end{lemma}

\begin{proof}
Injectives, Lemma \ref{injectives-lemma-derived-products} により、導来圏
$D(\mathcal{O}_X)$ は直和をもち、それは任意の代表複体の項ごとの直和を
とることで計算される。任意の Grothendieck アーベル圏において直和をとる
関手は完全であるから、直和のコホモロジー層がコホモロジー層の直和である
ことは明らかである。準連接層の直和は準連接であるため、補題が従う。
Schemes, Section \ref{schemes-section-quasi-coherent} を参照せよ。
\end{proof}

\noindent
導来極限についていくらかの情報が必要になる。以下の補題に現れる導来極限は、
一般には $D_\QCoh$ の対象ではないことに注意されたい。

\begin{lemma}
\label{perfect-lemma-Rlim-quasi-coherent}
スキーム $X$ をとる。$(K_n)$ を $D_\QCoh(\mathcal{O}_X)$ の逆系とし、
$D(\mathcal{O}_X)$ におけるその導来極限を $K = R\lim K_n$ とする。
すべての $q \in \mathbf{Z}$ および $n \geq 1$ に対して
$H^q(K_{n + 1}) \to H^q(K_n)$ が全射であると仮定する。このとき
\begin{enumerate}
\item $H^q(K) = \lim H^q(K_n)$,
\item $R\lim H^q(K_n) = \lim H^q(K_n)$ であり、
\item 任意のアフィン開集合 $U \subset X$ に対して
$H^p(U, \lim H^q(K_n)) = 0$ が $p > 0$ のとき成り立つ。
\end{enumerate}
\end{lemma}

\begin{proof}
$\mathcal{B}$ を $X$ のアフィン開集合全体とする。$H^q(K_n)$ は準連接なので、
Cohomology of Schemes, Lemma
\ref{coherent-lemma-quasi-coherent-affine-cohomology-zero} により、
$U \in \mathcal{B}$ に対して $H^p(U, H^q(K_n)) = 0$ である。
さらに Schemes, Lemma \ref{schemes-lemma-equivalence-quasi-coherent}
により、$U \in \mathcal{B}$ に対して写像
$H^0(U, H^q(K_{n + 1})) \to H^0(U, H^q(K_n))$ は全射である。
いま確認した条件を満たすので、(1) は Cohomology, Lemma
\ref{cohomology-lemma-derived-limit-suitable-system} から従う。
(2) と (3) は Cohomology, Lemma
\ref{cohomology-lemma-inverse-limit-is-derived-limit} から従う。
\end{proof}

\noindent
次の補題は、$D_\QCoh(\mathcal{O}_X)$ の対象上で右導来関手を「計算」する際に
役立つ。

\begin{lemma}
\label{perfect-lemma-nice-K-injective}
スキーム $X$ をとり、$E$ を $D_\QCoh(\mathcal{O}_X)$ の対象とする。
このとき Derived Categories, Remark
\ref{derived-remark-map-into-derived-limit-truncations} の写像
$E \to R\lim \tau_{\geq -n}E$ は同型である\footnote{特に、
Derived Categories, Lemma \ref{derived-lemma-difficulty-K-injectives}
により、$E$ は K-入射的な代表をもつ。}。
\end{lemma}

\begin{proof}
$E$ の第 $i$ コホモロジー層を $\mathcal{H}^i = H^i(E)$ と記す。
$\mathcal{B}$ を $X$ のアフィン開部分集合全体とする。このとき、すべての
$p > 0$、$i \in \mathbf{Z}$、および $U \in \mathcal{B}$ に対して
$H^p(U, \mathcal{H}^i) = 0$ である。Cohomology of Schemes, Lemma
\ref{coherent-lemma-quasi-coherent-affine-cohomology-zero} を参照せよ。
したがって補題は Cohomology, Lemma
\ref{cohomology-lemma-is-limit-dimension} から従う。
\end{proof}

\begin{lemma}
\label{perfect-lemma-application-nice-K-injective}
スキーム $X$ をとる。$F : \textit{Mod}(\mathcal{O}_X) \to \textit{Ab}$
を加法的関手とし、$N \geq 0$ を整数とする。次を仮定する。
\begin{enumerate}
\item $F$ は可算直積と可換である。
\item すべての $p \geq N$ および準連接な $\mathcal{F}$ に対して
$R^pF(\mathcal{F}) = 0$ である。
\end{enumerate}
このとき $E \in D_\QCoh(\mathcal{O}_X)$ に対して
\begin{enumerate}
\item $i \leq a$ ならば $H^i(RF(\tau_{\leq a}E)) \to H^i(RF(E))$
は同型である。
\item $i \geq b$ ならば $H^i(RF(E)) \to H^i(RF(\tau_{\geq b - N + 1}E))$
は同型である。
\item ある $-\infty \leq a \leq b \leq \infty$ に対し、
$i \not \in [a, b]$ ならば $H^i(E) = 0$ であるとする。このとき
$i \not \in [a, b + N - 1]$ ならば $H^i(RF(E)) = 0$ である。
\end{enumerate}
\end{lemma}

\begin{proof}
(1) は Derived Categories, Lemma
\ref{derived-lemma-negative-vanishing} である。

\medskip\noindent
(2) を証明する。$E_n = \tau_{\geq -n}E$ と書く。Lemma
\ref{perfect-lemma-nice-K-injective} により $E = R\lim E_n$ である。
したがって Injectives, Lemma
\ref{injectives-lemma-RF-commutes-with-Rlim} により、$D(\textit{Ab})$ において
$RF(E) = R\lim RF(E_n)$ である。よって、すべての $i \in \mathbf{Z}$ に対して
短完全列
$$
0 \to R^1\lim H^{i - 1}(RF(E_n)) \to H^i(RF(E)) \to \lim H^i(RF(E_n)) \to 0
$$
を得る。More on Algebra, Remark
\ref{more-algebra-remark-compare-derived-limit} を参照せよ。
(2) を示すため、左辺の項が零であり、$i \geq b$ を満たす任意の $b$ に対して
右辺の項が $H^i(RF(E_{-b + N - 1}))$ に等しいことを示す。

\medskip\noindent
すべての $n$ に対して完全三角形
$$
H^{-n}(E)[n] \to E_n \to E_{n - 1} \to H^{-n}(E)[n + 1]
$$
が $D(\mathcal{O}_X)$ に存在する（Derived Categories, Remark
\ref{derived-remark-truncation-distinguished-triangle}）。
$H^{-n}(E)$ は準連接なので
$$
H^i(RF(H^{-n}(E)[n])) = R^{i + n}F(H^{-n}(E)) = 0
$$
$i + n \geq N$ に対して成り立ち、また
$$
H^i(RF(H^{-n}(E)[n + 1])) = R^{i + n + 1}F(H^{-n}(E)) = 0
$$
$i + n + 1 \geq N$ に対して成り立つ。したがって
$$
H^i(RF(E_n)) \to H^i(RF(E_{n - 1}))
$$
は $n \geq N - i$ に対して同型である。ゆえに系 $H^i(RF(E_n))$ は
いずれも ML 条件を満たし、上の短完全列の $R^1\lim$ の項は零である
（More on Algebra, Section \ref{more-algebra-section-Rlim} の議論を参照せよ）。
さらに系 $H^i(RF(E_n))$ は、望みどおり $n = N - i - 1$ から定値となる。

\medskip\noindent
(3) を証明する。$E$ に関する仮定のもとで $\tau_{\leq a - 1}E = 0$
であり、(1) から $i \leq a - 1$ に対する $H^i(RF(E))$ の消滅を得る。
同様に $\tau_{\geq b + 1}E = 0$ であるから、(2) により
$i \geq b + N$ に対する $H^i(RF(E))$ の消滅を得る。
\end{proof}

\noindent
次の補題は、本章の多くの結果にとって鍵となる材料である。

\begin{lemma}
\label{perfect-lemma-affine-compare-bounded}
アフィン・スキーム $X = \Spec(A)$ をとる。図式
$$
\xymatrix{
D(\QCoh(\mathcal{O}_X)) \ar[rr]_{(\ref{perfect-equation-compare})}
& &
D_\QCoh(\mathcal{O}_X) \ar[ld]^{R\Gamma(X, -)} \\
& D(A) \ar[lu]^{\widetilde{\ \ }}
}
$$
に現れるすべての関手は三角圏の同値である。さらに、
$D_\QCoh(\mathcal{O}_X)$ の $E$ に対して
$H^0(X, E) = H^0(X, H^0(E))$ が成り立つ。
\end{lemma}

\begin{proof}
関手 $R\Gamma(X, -)$ は関手 $D(\mathcal{O}_X) \to D(A)$ を与えるので、
制限することにより関手
\begin{equation}
\label{perfect-equation-back}
R\Gamma(X, -) : D_\QCoh(\mathcal{O}_X) \longrightarrow D(A).
\end{equation}
を得る。$X$ 上の準連接加群の圏と $A$-加群の圏との同値を介して、この関手が
(\ref{perfect-equation-compare}) の擬逆であることを示す。

\medskip\noindent
補足説明。$A$ によって与えられる環の層を備えた一点空間を
$(Y, \mathcal{O}_Y)$ と記す。また、環付き空間の自然な射を
$\pi : (X, \mathcal{O}_X) \to (Y, \mathcal{O}_Y)$
と記す。このとき $R\Gamma(X, -)$ は $R\pi_*$ と同一視できる。また、同値
$\textit{Mod}(\mathcal{O}_Y) = \text{Mod}_A = \QCoh(\mathcal{O}_X)$
を介して、関手 (\ref{perfect-equation-compare}) は
$L\pi^* = \pi^* = \widetilde{\ }$ と同一視できる（Modules, Lemma
\ref{modules-lemma-construct-quasi-coherent-sheaves} および Schemes,
Lemmas \ref{schemes-lemma-compare-constructions} および
\ref{schemes-lemma-equivalence-quasi-coherent} を参照せよ）。したがって関手
$$
\xymatrix{
D(A) \ar@<1ex>[r] & D(\mathcal{O}_X) \ar@<1ex>[l]
}
$$
は随伴する（Cohomology, Lemma \ref{cohomology-lemma-adjoint} による）。
特に、標準的な随伴射
$$
a : \widetilde{R\Gamma(X, E)} \longrightarrow E
$$
$D(\mathcal{O}_X)$ の $E$ に対するもの、および
$$
b : M^\bullet \longrightarrow R\Gamma(X, \widetilde{M^\bullet})
$$
$A$-加群の複体 $M^\bullet$ に対するものを得る。

\medskip\noindent
$E$ を $D_\QCoh(\mathcal{O}_X)$ の対象とする。Cohomology of Schemes,
Lemma \ref{coherent-lemma-quasi-coherent-affine-cohomology-zero} により、
$N = 1$ として関手 $F(-) = \Gamma(X, -)$ に Lemma
\ref{perfect-lemma-application-nice-K-injective} を適用できる。したがって
$$
H^0(R\Gamma(X, E)) = H^0(R\Gamma(X, \tau_{\geq 0}E)) = \Gamma(X, H^0(E))
$$
である（最後の等式は標準切断の定義による）。これを用いて、随伴射 $a$ と
$b$ が同型 $H^0(a)$ および $H^0(b)$ を誘導することを示す。
この主張はシフトに関して不変なので、$a$ と $b$ は擬同型となり、補題が示される。

\medskip\noindent
いずれの場合にも、$\widetilde{\ }$ が完全関手であることを用いる
（Schemes, Lemma \ref{schemes-lemma-spec-sheaves}）。すなわち、これは
$$
H^0\left(\widetilde{R\Gamma(X, E)}\right) =
\widetilde{H^0(R\Gamma(X, E))} =
\widetilde{\Gamma(X, H^0(E))}
$$
を意味する。$H^0(E)$ は準連接なので、これは $H^0(E)$ に等しい。
したがって $H^0(a)$ は同型である。もう一方の向きについては
$$
H^0(R\Gamma(X, \widetilde{M^\bullet})) =
\Gamma(X, H^0(\widetilde{M^\bullet})) =
\Gamma(X, \widetilde{H^0(M^\bullet)}) =
H^0(M^\bullet)
$$
となり、$H^0(b)$ が同型であることが示される。
\end{proof}

\begin{lemma}
\label{perfect-lemma-affine-K-flat}
アフィン・スキーム $X = \Spec(A)$ をとる。$K^\bullet$ が $A$-加群の
K-平坦複体ならば、$\widetilde{K^\bullet}$ は $\mathcal{O}_X$-加群の
K-平坦複体である。
\end{lemma}

\begin{proof}
More on Algebra, Lemma \ref{more-algebra-lemma-base-change-K-flat}
により、すべての $\mathfrak p \in \Spec(A)$ に対して
$K^\bullet \otimes_A A_\mathfrak p$ は $A_\mathfrak p$-加群の K-平坦複体である。
したがって Cohomology, Lemma \ref{cohomology-lemma-check-K-flat-stalks}
（および Schemes, Lemma \ref{schemes-lemma-spec-sheaves}）から、
$\widetilde{K^\bullet}$ は K-平坦であると結論できる。
\end{proof}

\begin{lemma}
\label{perfect-lemma-quasi-coherence-pushforward}
$f : X \to Y$ を環準同型 $A \to B$ によって与えられるアフィン・スキームの
射とする。このとき図式
$$
\xymatrix{
D(B) \ar[d] \ar[r] & D_\QCoh(\mathcal{O}_X) \ar[d]^{Rf_*} \\
D(A) \ar[r] & D_\QCoh(\mathcal{O}_Y)
}
$$
は可換である。
\end{lemma}

\begin{proof}
Cohomology, Lemma \ref{cohomology-lemma-Leray-unbounded} による等式
$R\Gamma(Y, Rf_*K) = R\Gamma(X, K)$ を用いれば、Lemma
\ref{perfect-lemma-affine-compare-bounded} から従う。
\end{proof}

\begin{lemma}
\label{perfect-lemma-quasi-coherence-pullback}
スキームの射 $f : Y \to X$ をとる。
\begin{enumerate}
\item 関手 $Lf^*$ は $D_\QCoh(\mathcal{O}_X)$ を
$D_\QCoh(\mathcal{O}_Y)$ の中へ送る。
\item $X$ と $Y$ がアフィンであり、$f$ が環準同型 $A \to B$ によって
与えられるならば、図式
$$
\xymatrix{
D(B) \ar[r] & D_\QCoh(\mathcal{O}_Y) \\
D(A) \ar[r] \ar[u]^{- \otimes_A^\mathbf{L} B} &
D_\QCoh(\mathcal{O}_X) \ar[u]_{Lf^*}
}
$$
は可換である。
\end{enumerate}
\end{lemma}

\begin{proof}
まず図式
$$
\xymatrix{
D(B) \ar[r] & D(\mathcal{O}_Y) \\
D(A) \ar[r] \ar[u]^{- \otimes_A^\mathbf{L} B} &
D(\mathcal{O}_X) \ar[u]_{Lf^*}
}
$$
が可換であることを示す。これは Lemma \ref{perfect-lemma-affine-K-flat} と、問題の
関手の構成から明らかである。(1) を示すため、$E$ を
$D_\QCoh(\mathcal{O}_X)$ の対象とする。$Lf^*E$ が準連接なコホモロジー層を
もつことは $X$ 上局所的に確かめられる。$Lf^*$ は開部分スキームへの制限と
両立することに注意せよ。したがって、$f$ は (2) のようなアフィン・スキームの
射であると仮定してよい。Lemma \ref{perfect-lemma-affine-compare-bounded} により、
$E$ は $A$-加群の複体から来る。証明の最初の図式の可換性により $Lf^*E$
についても同じことが成り立つので、(1) が従う。
\end{proof}

\begin{lemma}
\label{perfect-lemma-quasi-coherence-tensor-product}
スキーム $X$ をとる。
\begin{enumerate}
\item $D_\QCoh(\mathcal{O}_X)$ の対象 $K, L$ に対して、導来テンソル積
$K \otimes^\mathbf{L}_{\mathcal{O}_X} L$ は $D_\QCoh(\mathcal{O}_X)$
に属する。
\item $X = \Spec(A)$ がアフィンならば
$$
\widetilde{M^\bullet} \otimes_{\mathcal{O}_X}^\mathbf{L} \widetilde{K^\bullet}
=
\widetilde{M^\bullet \otimes_A^\mathbf{L} K^\bullet}
$$
が任意の $A$-加群の複体の対 $K^\bullet$, $M^\bullet$ に対して成り立つ。
\end{enumerate}
\end{lemma}

\begin{proof}
(2) の等式は Lemma \ref{perfect-lemma-affine-K-flat} と導来テンソル積の構成から
直ちに従う。(1) を示すため、$K, L$ を $D_\QCoh(\mathcal{O}_X)$ の対象とする。
$K \otimes^\mathbf{L} L$ が $D_\QCoh(\mathcal{O}_X)$ に属することは $X$ 上
局所的に確かめられるので、$X = \Spec(A)$ はアフィンであると仮定してよい。
Lemma \ref{perfect-lemma-affine-compare-bounded} により $K$ と $L$ を $A$-加群の複体で
表すことができ、(2) から結果が従う。
\end{proof}





\section{全直像}
\label{perfect-section-total-direct-image}

\noindent
次の補題は Cohomology of Schemes, Lemma
\ref{coherent-lemma-quasi-coherence-higher-direct-images} の類似である。

\begin{lemma}
\label{perfect-lemma-quasi-coherence-direct-image}
スキームの射 $f : X \to S$ をとる。$f$ は準分離かつ準コンパクトであると
仮定する。
\begin{enumerate}
\item 関手 $Rf_*$ は $D_\QCoh(\mathcal{O}_X)$ を
$D_\QCoh(\mathcal{O}_S)$ の中へ送る。
\item $S$ が準コンパクトならば、次を満たす整数 $N = N(X, S, f)$ が存在する。
$D_\QCoh(\mathcal{O}_X)$ の対象 $E$ が $m > 0$ に対して
$H^m(E) = 0$ を満たすならば、$m \geq N$ に対して
$H^m(Rf_*E) = 0$ である。
\item 実際、$S$ が準コンパクトならば、スキームの任意の射 $S' \to S$ に対して
同じ結論が関手 $R(f')_*$ について成り立つような $N = N(X, S, f)$ をとれる。
ここで $f' : X' \to S'$ は $f$ の基底変換である。
\end{enumerate}
\end{lemma}

\begin{proof}
$E$ を $D_\QCoh(\mathcal{O}_X)$ の対象とする。(1) を示すには、$Rf_*E$ が
準連接なコホモロジー層をもつことを示さなければならない。この問題は $S$ 上
局所的なので、$S$ は準コンパクトであると仮定してよい。Cohomology of Schemes,
Lemma \ref{coherent-lemma-quasi-coherence-higher-direct-images} のように
$N = N(X, S, f)$ をとる。すると、すべての準連接 $\mathcal{O}_X$-加群
$\mathcal{F}$ およびすべての $p \geq N$ に対して
$R^pf_*\mathcal{F} = 0$ であり、このことは基底変換後にも成り立つ。

\medskip\noindent
まず $E$ が下に有界であると仮定する。このような $E$ に対し、上で選んだ
$N$ について (1), (2), (3) が成り立つことを示す。この場合、例えばスペクトル系列
$$
R^pf_*H^q(E) \Rightarrow R^{p + q}f_*E
$$
と、$R^pf_*H^q(E)$ の準連接性、および $p \geq N$ に対する
$R^pf_*H^q(E)$ の消滅を用いればよい（Derived Categories, Lemma
\ref{derived-lemma-two-ss-complex-functor}）。これにより、この場合の
(1), (2), (3) が従う。

\medskip\noindent
次に (2) と (3) を示す。$m > 0$ に対して $H^m(E) = 0$ であるとする。
$U \subset S$ をアフィン開集合とする。Cohomology of Schemes, Lemma
\ref{coherent-lemma-quasi-coherence-higher-direct-images-application}
と $N$ の選び方により、任意の準連接 $\mathcal{O}_X$-加群 $\mathcal{F}$ と
$p \geq N$ に対して $H^p(f^{-1}(U), \mathcal{F}) = 0$ である。
したがって関手 $\Gamma(f^{-1}(U), -)$ に Lemma
\ref{perfect-lemma-application-nice-K-injective} を適用すると、
$$
R\Gamma(U, Rf_*E) = R\Gamma(f^{-1}(U), E)
$$
の次数 $\geq N$ におけるコホモロジーは消滅する。これはすべてのアフィン開集合
$U \subset S$ について成り立つので、$m \geq N$ に対して
$H^m(Rf_*E) = 0$ と結論できる。

\medskip\noindent
次に一般の場合の (1) を示す。$D(\mathcal{O}_X)$ に完全三角形
$$
\tau_{\leq -n - 1}E \to E \to \tau_{\geq -n}E \to
(\tau_{\leq -n - 1}E)[1]
$$
が存在することを思い出そう。Derived Categories, Remark
\ref{derived-remark-truncation-distinguished-triangle} を参照せよ。
(2) により、$Rf_*\tau_{\leq -n - 1}E$ の次数 $\geq -n + N$ における
コホモロジー層は消滅する。したがって、整数 $q$ が与えられたとき、ある $n$
について $R^qf_*E$ は $R^qf_*\tau_{\geq -n}E$ に等しく、上の結果を適用できる。
\end{proof}

\begin{lemma}
\label{perfect-lemma-acyclicity-lemma}
$f : X \to S$ を準分離かつ準コンパクトなスキームの射とする。
$\mathcal{F}^\bullet$ を準連接 $\mathcal{O}_X$-加群の複体とし、その各項は
$f_*$ に関して右非輪状であるとする。このとき $D(\mathcal{O}_S)$ において
$f_*\mathcal{F}^\bullet$ は $Rf_*\mathcal{F}^\bullet$ を表す。
\end{lemma}

\begin{proof}
標準的な写像 $f_*\mathcal{F}^\bullet \to Rf_*\mathcal{F}^\bullet$ は常に存在する。
これがコホモロジー層上で同型であることを示せばよい。この主張はシフトに関して
不変なので、$H^0(f_*(\mathcal{F}^\bullet)) \to R^0f_*\mathcal{F}^\bullet$
が同型であることを示せば十分である。主張は $S$ 上局所的なので、$S$ はアフィン
であると仮定してよい。Lemma \ref{perfect-lemma-quasi-coherence-direct-image} により、
十分大きいすべての $n$ に対して
$R^0f_*\mathcal{F}^\bullet = R^0f_*\tau_{\geq -n}\mathcal{F}^\bullet$
である。したがって $\mathcal{F}^\bullet$ は下に有界であると仮定してよい。
仮定により各 $\mathcal{F}^n$ は右 $f_*$-非輪状なので、Leray の非輪状性補題
（Derived Categories, Lemma \ref{derived-lemma-leray-acyclicity}）から
$f_*\mathcal{F}^\bullet \to Rf_*\mathcal{F}^\bullet$ は擬同型である。
\end{proof}

\begin{lemma}
\label{perfect-lemma-acyclicity-lemma-global}
$X$ を準分離かつ準コンパクトなスキームとする。$\mathcal{F}^\bullet$ を
準連接 $\mathcal{O}_X$-加群の複体とし、その各項は $\Gamma(X, -)$ に関して
右非輪状であるとする。このとき $D(\Gamma(X, \mathcal{O}_X)$ において
$\Gamma(X, \mathcal{F}^\bullet)$ は $R\Gamma(X, \mathcal{F}^\bullet)$ を表す。
\end{lemma}

\begin{proof}
標準的な射 $X \to \Spec(\Gamma(X, \mathcal{O}_X))$ に Lemma
\ref{perfect-lemma-acyclicity-lemma} を適用する。細部は省略する。
\end{proof}

\begin{lemma}
\label{perfect-lemma-spectral-sequence}
$X$ を準分離かつ準コンパクトなスキームとする。
$D_\QCoh(\mathcal{O}_X)$ の任意の対象 $K$ に対して、Cohomology, Example
\ref{cohomology-example-spectral-sequence} のスペクトル系列
$$
E_2^{i, j} = H^i(X, H^j(K)) \Rightarrow H^{i + j}(X, K)
$$
は有界で収束する。
\end{lemma}

\begin{proof}
Cohomology, Lemma \ref{cohomology-lemma-spectral-sequence-filtered-object}
を介し、$\tau_{\leq -p}K$ が与えるフィルトレーションを用いるスペクトル系列の
構成によれば、与えられた $n \in \mathbf{Z}$ に対して
$$
H^n(X, \tau_{\leq -p}K) = 0 \text{（ただし）} p \gg 0
$$
および
$$
H^n(X, K) = H^n(X, \tau_{\leq -p}K) \text{（ただし）} p \ll 0
$$
を示せば十分である。前者は、$F = \Gamma(X, -)$ として Lemma
\ref{perfect-lemma-application-nice-K-injective} を適用し、Cohomology of Schemes,
Lemma \ref{coherent-lemma-quasi-coherence-higher-direct-images} の上界を
用いれば従う。後者は、任意の環付き空間 $(X, \mathcal{O}_X)$ と任意の
$K \in D(\mathcal{O}_X)$ に対し、$-p \leq n$ ならば成り立つ。
\end{proof}

\begin{lemma}
\label{perfect-lemma-quasi-coherence-pushforward-direct-sums}
$f : X \to S$ を準分離かつ準コンパクトなスキームの射とする。このとき
$Rf_* : D_\QCoh(\mathcal{O}_X) \to D_\QCoh(\mathcal{O}_S)$
は直和と可換である。
\end{lemma}

\begin{proof}
$E_i$ を $D_\QCoh(\mathcal{O}_X)$ の対象の族とし、$E = \bigoplus E_i$
とおく。写像
$$
\bigoplus Rf_*E_i \longrightarrow Rf_*E
$$
が同型であることを示したい。これが次数 $0$ のコホモロジー層上で同型を誘導する
ことを示せば、補題が従う。これは $S$ 上局所的に証明できるので、$S$ は
準コンパクトであると仮定する。Lemma
\ref{perfect-lemma-quasi-coherence-direct-image} のように整数 $N$ を選ぶ。
引用した補題により $R^0f_*E = R^0f_*\tau_{\geq -N}E$ および
$R^0f_*E_i = R^0f_*\tau_{\geq -N}E_i$ である。また、
$\tau_{\geq -N}E = \bigoplus \tau_{\geq -N}E_i$ であることに注意せよ。
したがって、すべての $E_i$ の次数 $< -N$ におけるコホモロジー層は消滅すると
仮定してよい。次にスペクトル系列
$$
R^pf_*H^q(E) \Rightarrow R^{p + q}f_*E
\quad\text{および}\quad
R^pf_*H^q(E_i) \Rightarrow R^{p + q}f_*E_i
$$
を用い、準連接層の直和の場合に帰着する（Derived Categories, Lemma
\ref{derived-lemma-two-ss-complex-functor}）。この場合は Cohomology of Schemes,
Lemma \ref{coherent-lemma-colimit-cohomology} で扱われている。
\end{proof}









\section{アフィン射}
\label{perfect-section-affine-morphisms}

\noindent
本節では、スキームのアフィン射に沿う直像についての情報をまとめる。

\begin{lemma}
\label{perfect-lemma-pushforward-affine-morphism}
$f : X \to S$ をスキームのアフィン射とする。$\mathcal{F}^\bullet$ を
準連接 $\mathcal{O}_X$-加群の複体とする。このとき
$f_*\mathcal{F}^\bullet = Rf_*\mathcal{F}^\bullet$ である。
\end{lemma}

\begin{proof}
Lemma \ref{perfect-lemma-acyclicity-lemma} と Cohomology of Schemes, Lemma
\ref{coherent-lemma-relative-affine-vanishing} を組み合わせる。
別の証明として、$S$ 上アフィン局所的に議論し、Lemma
\ref{perfect-lemma-quasi-coherence-pushforward} を用いてもよい。
\end{proof}

\begin{lemma}
\label{perfect-lemma-affine-morphism}
$f : X \to S$ をスキームのアフィン射とする。このとき
$Rf_* : D_\QCoh(\mathcal{O}_X) \to D_\QCoh(\mathcal{O}_S)$
は同型を反映する。
\end{lemma}

\begin{proof}
この主張は、$D_\QCoh(\mathcal{O}_X)$ の射 $\alpha : E \to F$ について、
$Rf_*\alpha$ が同型ならば元の射も同型であることを意味する。
これはコホモロジー層上で確かめられる。特に、この問題は $S$ 上局所的である。
したがって $S$ はアフィンであり、それゆえ $X$ もアフィンであると仮定してよい。
この場合、Lemma \ref{perfect-lemma-affine-compare-bounded} による導来圏
$D_\QCoh(\mathcal{O}_X)$ と $D_\QCoh(\mathcal{O}_S)$ の記述から主張は明らかである。
細部は省略する。
\end{proof}

\begin{lemma}
\label{perfect-lemma-affine-morphism-pull-push}
$f : X \to S$ をスキームのアフィン射とする。
$D_\QCoh(\mathcal{O}_S)$ の $E$ に対して
$Rf_* Lf^* E = E \otimes^\mathbf{L}_{\mathcal{O}_S} f_*\mathcal{O}_X$.
\end{lemma}

\begin{proof}
$f$ はアフィンなので、写像 $f_*\mathcal{O}_X \to Rf_*\mathcal{O}_X$ は同型である
（Cohomology of Schemes, Lemma
\ref{coherent-lemma-relative-affine-vanishing}）。写像
$E \otimes^\mathbf{L} f_*\mathcal{O}_X =
E \otimes^\mathbf{L} Rf_*\mathcal{O}_X \to Rf_* Lf^* E$ が標準的に存在し、
これは写像
$$
Lf^*(E \otimes^\mathbf{L} Rf_*\mathcal{O}_X) =
Lf^*E \otimes^\mathbf{L} Lf^*Rf_*\mathcal{O}_X \longrightarrow
Lf^* E \otimes^\mathbf{L} \mathcal{O}_X = Lf^* E
$$
に随伴する。この写像は $1 : Lf^*E \to Lf^*E$ と標準写像
$Lf^*Rf_*\mathcal{O}_X \to \mathcal{O}_X$ から生じる。こうして構成した写像が
同型であることは $S$ 上局所的に確かめられる。したがって $S$ はアフィンであり、
それゆえ $X$ もアフィンであると仮定してよい。この場合、Lemmas
\ref{perfect-lemma-affine-compare-bounded} および
\ref{perfect-lemma-quasi-coherence-pullback} による導来圏
$D_\QCoh(\mathcal{O}_X)$ と $D_\QCoh(\mathcal{O}_S)$、および関手 $Lf^*$ の
記述から主張は明らかである。細部は省略する。
\end{proof}

\noindent
$Y$ をスキームとし、$\mathcal{A}$ を $\mathcal{O}_Y$-代数の層とする。
制限関手 $D(\mathcal{A}) \to D(\mathcal{O}_X)$ による
$D_\QCoh(\mathcal{O}_X)$ の逆像を $D_\QCoh(\mathcal{A})$ と記す。
言い換えれば、$K \in D(\mathcal{A})$ が $D_\QCoh(\mathcal{A})$ に属することと、
そのコホモロジー層が $\mathcal{O}_X$-加群として準連接であることとは同値である。
$\mathcal{A}$ 自身が準連接ならば、これはコホモロジー層が $\mathcal{A}$-加群として
準連接であることと同じである。Morphisms, Lemma
\ref{morphisms-lemma-affine-equivalence-modules} を参照せよ。

\begin{lemma}
\label{perfect-lemma-affine-morphism-equivalence}
$f : X \to Y$ をスキームのアフィン射とする。このとき $f_*$ は同値
$$
\Phi : D_\QCoh(\mathcal{O}_X) \longrightarrow D_\QCoh(f_*\mathcal{O}_X)
$$
を誘導する。これと $D_\QCoh(f_*\mathcal{O}_X) \to D_\QCoh(\mathcal{O}_Y)$
との合成は $Rf_* : D_\QCoh(\mathcal{O}_X) \to D_\QCoh(\mathcal{O}_Y)$
である。
\end{lemma}

\begin{proof}
$Rf_*$ の対象 $K \in D_\QCoh(\mathcal{O}_X)$ における値は、$K$ を表す
$\mathcal{O}_X$-加群の K-入射的複体 $\mathcal{I}^\bullet$ を選び、
$f_*\mathcal{I}^\bullet$ をとることで計算されることを思い出そう。
そこで $\Phi(K)$ を、$f_*\mathcal{O}_X$-加群の複体とみなした複体
$f_*\mathcal{I}^\bullet$ とする。環付き空間の自然な射を
$g : (X, \mathcal{O}_X) \to (Y, f_*\mathcal{O}_X)$ と記す。
すると $g$ は環付き空間の平坦射であり（茎の記述は以下を参照）、$\Phi$ は
$Rg_*$ の $D_\QCoh(\mathcal{O}_X)$ への制限である。$Lg^*$ が擬逆であると主張する。
まず、$g^*$ は準連接加群を準連接加群へ移すので（Modules, Lemma
\ref{modules-lemma-pullback-quasi-coherent}）、$Lg^*$ は
$D_\QCoh(f_*\mathcal{O}_X)$ を $D_\QCoh(\mathcal{O}_X)$ の中へ送る。
証明を完了するには、随伴射
$$
Lg^*\Phi(K) = Lg^*Rg_*K \to K
\quad\text{および}\quad
M \to Rg_*Lg^*M = \Phi(Lg^*M)
$$
が $K \in D_\QCoh(\mathcal{O}_X)$ および
$M \in D_\QCoh(f_*\mathcal{O}_X)$ に対して同型であることを示せば十分である。
これは局所的な問題なので、$Y$ はアフィンであり、それゆえ $X$ もアフィンであると
仮定してよい。

\medskip\noindent
$Y = \Spec(B)$ および $X = \Spec(A)$ と仮定する。
$\mathfrak q = y \in \Spec(B) = Y$ に写る点を
$\mathfrak p = x \in \Spec(A) = X$ とする。このとき
$(f_*\mathcal{O}_X)_y = A_\mathfrak q$ かつ $\mathcal{O}_{X, x} = A_\mathfrak p$
なので、$g$ は平坦である。したがって $g^*$ は完全であり、
$D(f_*\mathcal{O}_X)$ の任意の $M$ に対して
$H^i(Lg^*M) = g^*H^i(M)$ である。
$K \in D_\QCoh(\mathcal{O}_X)$ に対しては
$$
H^i(\Phi(K)) = H^i(Rf_*K) = f_*H^i(K)
$$
である。これは高次直像の消滅（Cohomology of Schemes, Lemma
\ref{coherent-lemma-relative-affine-vanishing}）と Lemma
\ref{perfect-lemma-application-nice-K-injective} による（細部を一つ省略した）。
したがって、
$$
g^*g_*\mathcal{F} \to \mathcal{F}
\quad\text{および}\quad
\mathcal{G} \to g_*g^*\mathcal{G}
$$
が同型であることを示せば十分である。ここで $\mathcal{F}$ は準連接
$\mathcal{O}_X$-加群、$\mathcal{G}$ は準連接 $f_*\mathcal{O}_X$-加群である。
これは Morphisms, Lemma \ref{morphisms-lemma-affine-equivalence-modules}
から従う。
\end{proof}





\section{閉部分集合に台をもつコホモロジー}
\label{perfect-section-cohomology-support}

\noindent
スキームと準連接加群について、Cohomology, Sections
\ref{cohomology-section-cohomology-support} および
\ref{cohomology-section-cohomology-support-bis} の内容を詳しく述べる。

\begin{definition}
\label{perfect-definition-supported-on}
スキーム $X$、$D(\mathcal{O}_X)$ の対象 $E$、および閉部分集合
$T \subset X$ をとる。コホモロジー層 $H^i(E)$ が $T$ に台をもつとき、
$E$ は {\it $T$ に台をもつ}という。
\end{definition}

\noindent
この定義の状況で、Cohomology, Section
\ref{cohomology-section-cohomology-support-bis} の議論の一部を繰り返す。
$X$ をスキーム、$T \subset X$ を閉部分集合とする。台が $T$ に含まれる
$\mathcal{O}_X$-加群の圏は、すべての $\mathcal{O}_X$-加群の圏の Serre 部分圏で
ある。Homology, Definition \ref{homology-definition-serre-subcategory}
および Modules, Lemma \ref{modules-lemma-support-section-closed} を参照せよ。
したがって Derived Categories, Section
\ref{derived-section-triangulated-sub} の議論を適用でき、$T$ に台をもつ対象からなる
$D(\mathcal{O}_X)$ の厳密充満・飽和三角部分圏を得る。以下、この圏を
$D_T(\mathcal{O}_X)$ と記す。

\medskip\noindent
Definition \ref{perfect-definition-supported-on} の状況で、包含写像を
$i : T \to X$ と記す。Cohomology, Section
\ref{cohomology-section-cohomology-support-bis} により、この状況では関手
$R\mathcal{H}_T : D(\mathcal{O}_X) \to D(i^{-1}\mathcal{O}_X)$
があり、これは $i_* : D(i^{-1}\mathcal{O}_X) \to D(\mathcal{O}_X)$ の
右随伴である。

\begin{lemma}
\label{perfect-lemma-support-quasi-coherent}
スキーム $X$ と閉部分集合 $T \subset X$ をとり、$X \setminus T$ は $X$ の
逆コンパクトな開集合であるとする。包含を $i : T \to X$ と記す。
\begin{enumerate}
\item $D_\QCoh(\mathcal{O}_X)$ の $E$ に対して、
$i_*R\mathcal{H}_T(E)$ は $D_{\QCoh, T}(\mathcal{O}_X)$ に属する。
\item 関手 $i_* \circ R\mathcal{H}_T : D_\QCoh(\mathcal{O}_X) \to
D_{\QCoh, T}(\mathcal{O}_X)$ は、包含関手
$D_{\QCoh, T}(\mathcal{O}_X) \to D_\QCoh(\mathcal{O}_X)$ の右随伴である。
\end{enumerate}
\end{lemma}

\begin{proof}
$U = X \setminus T$ とおき、包含を $j : U \to X$ と記す。Cohomology,
Lemma \ref{cohomology-lemma-triangle-sections-with-support-sheaves} により、
$D(\mathcal{O}_X)$ に完全三角形
$$
i_*R\mathcal{H}_T(E) \to E \to Rj_*(E|_U) \to i_*R\mathcal{H}_Z(E)[1]
$$
が存在する。Lemma \ref{perfect-lemma-quasi-coherence-direct-image} により、複体
$Rj_*(E|_U)$ は準連接なコホモロジー層をもつ（ここで $U$ が $X$ 内で
逆コンパクトであることを用いる）。したがって (1) が成り立つ。
(2) はこれと Cohomology, Lemma
\ref{cohomology-lemma-complexes-with-support-on-closed} における関手の随伴性から
従う。
\end{proof}

\begin{lemma}
\label{perfect-lemma-support-direct-sums}
スキーム $X$ と閉部分集合 $T \subset X$ をとり、$X \setminus T$ は $X$ の
逆コンパクトな開集合であるとする。このとき、$D_\QCoh(\mathcal{O}_X)$ の
対象の族 $E_i$, $i \in I$ に対して
$R\mathcal{H}_T(\bigoplus E_i) = \bigoplus R\mathcal{H}_T(E_i)$.
\end{lemma}

\begin{proof}
$U = X \setminus T$ とおき、包含を $j : U \to X$ と記す。Cohomology,
Lemma \ref{cohomology-lemma-triangle-sections-with-support-sheaves} により、
任意の $D(\mathcal{O}_X)$ の $E$ に対して完全三角形
$$
i_*R\mathcal{H}_T(E) \to E \to Rj_*(E|_U) \to i_*R\mathcal{H}_Z(E)[1]
$$
が $D(\mathcal{O}_X)$ に存在する。Lemma
\ref{perfect-lemma-quasi-coherence-pushforward-direct-sums} により、
関手 $E \mapsto Rj_*(E|_U)$ は $D_\QCoh(\mathcal{O}_X)$ 上で直和と可換である。
したがって関手 $i_* \circ R\mathcal{H}_T$ についても同じことが成り立つ
（細部は省略する）。$i_* : D(i^{-1}\mathcal{O}_X) \to D_T(\mathcal{O}_X)$ は
同値なので（Cohomology, Lemma
\ref{cohomology-lemma-complexes-with-support-on-closed}）、結論を得る。
\end{proof}

\begin{remark}
\label{perfect-remark-support-c-equations}
スキーム $X$ と $f_1, \ldots, f_c \in \Gamma(X, \mathcal{O}_X)$ をとる。
$f_1, \ldots, f_c$ によって切り出される閉部分スキームを $Z \subset X$ と記す。
$0 \leq p < c$ および $1 \leq i_0 < \ldots < i_p \leq c$ に対して、
$f_{i_0} \ldots f_{i_p}$ が可逆となる開部分スキームを
$U_{i_0 \ldots i_p} \subset X$ と記す。任意の $\mathcal{O}_X$-加群
$\mathcal{F}$ に対して
$$
\mathcal{F}_{i_0 \ldots i_p} =
(U_{i_0 \ldots i_p} \to X)_*(\mathcal{F}|_{U_{i_0 \ldots i_p}})
$$
とおく。この状況における {\it 拡張交代 {\v C}ech 複体}とは、
$\mathcal{O}_X$-加群の複体
\begin{equation}
\label{perfect-equation-extended-alternating}
0 \to \mathcal{F} \to
\bigoplus\nolimits_{i_0} \mathcal{F}_{i_0} \to
\ldots \to
\bigoplus\nolimits_{i_0 < \ldots < i_p} \mathcal{F}_{i_0 \ldots i_p} \to
\ldots \to \mathcal{F}_{1 \ldots c} \to 0
\end{equation}
のことであり、$\mathcal{F}$ は次数 $0$ に置かれる。写像は次のように構成する。
$1 \leq i_0 < \ldots < i_{p + 1} \leq c$ および $0 \leq j \leq p + 1$
が与えられたとき、標準写像
$$
\mathcal{F}_{i_0 \ldots \hat i_j \ldots i_{p + 1}} \to
\mathcal{F}_{i_0 \ldots i_p}
$$
があり、これは包含
$U_{i_0 \ldots i_p} \subset U_{i_0 \ldots \hat i_j \ldots i_{p + 1}}$
から生じる。
拡張交代複体の微分は、これらの標準写像に符号 $(-1)^j$ を付けて用いる。
\end{remark}

\begin{lemma}
\label{perfect-lemma-extended-alternating-zero}
$X$, $f_1, \ldots, f_c \in \Gamma(X, \mathcal{O}_X)$, および
$\mathcal{F}$ は Remark \ref{perfect-remark-support-c-equations} のとおりとする。
このとき複体 (\ref{perfect-equation-extended-alternating}) を $X \setminus Z$ 上に
制限すると非輪状複体となる。
\end{lemma}

\noindent
より一般に、有限開被覆 $X \setminus Z = U_1 \cup \ldots \cup U_c$ から始めて
Remark \ref{perfect-remark-support-c-equations} のように定義される任意の拡張交代
{\v C}ech 複体についても、この補題が成り立つことに注意せよ。

\begin{proof}
$W \subset X \setminus Z$ を開部分集合とする。層の複体
(\ref{perfect-equation-extended-alternating}) を $W$ 上で評価すると、複体
$$
\mathcal{F}(W) \to \bigoplus\nolimits_{i_0} \mathcal{F}(U_{i_0} \cap W) \to
\bigoplus\nolimits_{i_0 < i_1} \mathcal{F}(U_{i_0i_1} \cap W) \to \ldots
$$
を得る。言い換えれば、被覆 $W = \bigcup U_i \cap W$ と
$\{1, \ldots, c\}$ 上の標準順序に関する拡張順序付き {\v C}ech 複体を得る。
Cohomology, Section \ref{cohomology-section-alternating-cech} を参照せよ。
Cohomology, Lemma \ref{cohomology-lemma-alternating-cech-trivial} により、
ある $1 \leq i \leq c$ について $W$ が $V(f_i)$ に含まれれば、この複体は
零とホモトピー同値である。これで証明が完了する。
\end{proof}

\begin{remark}
\label{perfect-remark-extended-alternating-map-to-support}
$X$, $f_1, \ldots, f_c \in \Gamma(X, \mathcal{O}_X)$, および
$\mathcal{F}$ は Remark \ref{perfect-remark-support-c-equations} のとおりとする。
複体 (\ref{perfect-equation-extended-alternating}) を $\mathcal{F}^\bullet$ と記す。
Lemma \ref{perfect-lemma-extended-alternating-zero} により、
$\mathcal{F}^\bullet$ のコホモロジー層は $Z$ に台をもつので、
$\mathcal{F}^\bullet$ は $D_Z(\mathcal{O}_X)$ の対象である。一方、等式
$\mathcal{F}^0 = \mathcal{F}$ は $D(\mathcal{O}_X)$ における標準写像
$\mathcal{F}^\bullet \to \mathcal{F}$ を定める。
$i_* \circ R\mathcal{H}_Z$ は包含関手
$D_Z(\mathcal{O}_X) \to D(\mathcal{O}_X)$ の右随伴なので（Cohomology,
Lemma \ref{cohomology-lemma-complexes-with-support-on-closed}）、標準可換図式
$$
\xymatrix{
\mathcal{F}^\bullet \ar[rd] \ar[rr] & & \mathcal{F} \\
& i_*R\mathcal{H}_Z(\mathcal{F}) \ar[ru]
}
$$
を $D(\mathcal{O}_X)$ に得る。これは $\mathcal{O}_X$-加群 $\mathcal{F}$
について関手的である。
\end{remark}

\begin{lemma}
\label{perfect-lemma-extended-alternating-represented}
$X$, $f_1, \ldots, f_c \in \Gamma(X, \mathcal{O}_X)$, および
$\mathcal{F}$ は Remark \ref{perfect-remark-support-c-equations} のとおりとする。
$\mathcal{F}$ が準連接ならば、複体 (\ref{perfect-equation-extended-alternating}) は
$D_Z(\mathcal{O}_X)$ において $i_* R\mathcal{H}_Z(\mathcal{F})$ を表す。
\end{lemma}

\begin{proof}
複体 (\ref{perfect-equation-extended-alternating}) を $\mathcal{F}^\bullet$ と記す。
補題の主張は、Remark \ref{perfect-remark-extended-alternating-map-to-support} の写像
$\mathcal{F}^\bullet \to i_*R\mathcal{H}_Z(\mathcal{F})$ が同型であることを
意味する。$\mathcal{F}^\bullet$ は $D_Z(\mathcal{O}_X)$ に属するので
（引用した注意を参照）、Cohomology, Lemma
\ref{cohomology-lemma-complexes-with-support-on-closed} により
$i_*R\mathcal{H}_Z(\mathcal{F}^\bullet) = \mathcal{F}^\bullet$ である。
射 $U_{i_0 \ldots i_p} \to X$ はアフィンである。実際、$X$ のアフィン開集合上では
関数 $f_{i_0} \ldots f_{i_p}$ を可逆化することで与えられる。したがって
$$
\mathcal{F}_{i_0 \ldots i_p} =
(U_{i_0 \ldots i_p} \to X)_*\mathcal{F}|_{U_{i_0 \ldots i_p}} =
R(U_{i_0 \ldots i_p} \to X)_*\mathcal{F}|_{U_{i_0 \ldots i_p}}
$$
である。これは Cohomology of Schemes, Lemma
\ref{coherent-lemma-relative-affine-vanishing} と $\mathcal{F}$ の準連接性による。
Cohomology, Lemma
\ref{cohomology-lemma-sections-support-in-closed-disjoint-open} から
$R\mathcal{H}_Z(\mathcal{F}_{i_0 \ldots i_p}) = 0$ と結論できる。
したがって $p > 0$ に対して $i_*R\mathcal{H}_Z(\mathcal{F}^p) = 0$ である。
すべてを合わせると
$$
\mathcal{F}^\bullet = i_*R\mathcal{H}_Z(\mathcal{F}^\bullet) =
i_*R\mathcal{H}_Z(\mathcal{F})
$$
を得る。これは望んだ等式である。
\end{proof}

\begin{lemma}
\label{perfect-lemma-supported-trivial-vanishing}
スキーム $X$ をとる。$T \subset X$ を、構造層の高々 $c$ 個の元で局所的に
切り出せる閉部分集合とする。このとき、任意の準連接 $\mathcal{O}_X$-加群
$\mathcal{F}$ と $i > c$ に対して $\mathcal{H}^i_Z(\mathcal{F}) = 0$ である。
\end{lemma}

\begin{proof}
これは Lemma \ref{perfect-lemma-extended-alternating-represented} による
$R\mathcal{H}_T(\mathcal{F})$ の局所的記述から直ちに従う。
\end{proof}

\begin{lemma}
\label{perfect-lemma-supported-vanishing}
スキーム $X$ をとる。$Z \subset X$ を、$c$ 個の元からなる Koszul 正則列で
局所的に切り出せる閉部分集合とする。このとき、任意の平坦な準連接
$\mathcal{O}_X$-加群 $\mathcal{F}$ と $i \not = c$ に対して
$\mathcal{H}^i_Z(\mathcal{F}) = 0$ である。
\end{lemma}

\begin{proof}
Lemma \ref{perfect-lemma-extended-alternating-represented} による
$R\mathcal{H}_Z(\mathcal{F})$ の記述から、これは次の代数的主張に帰着する。
環 $R$、Koszul 正則列 $f_1, \ldots, f_c \in R$、および平坦 $R$-加群 $M$
が与えられたとき、More on Algebra, Section
\ref{more-algebra-section-alternating-cech} の拡張交代 {\v C}ech 複体
$M \to \bigoplus\nolimits_{i_0} M_{f_{i_0}} \to
\bigoplus\nolimits_{i_0 < i_1} M_{f_{i_0}f_{i_1}} \to
\ldots \to M_{f_1 \ldots f_c}$
は次数 $c$ にしかコホモロジーをもたない。More on Algebra, Lemma
\ref{more-algebra-lemma-vanishing-extended-alternating-koszul} により、
$R$ の拡張交代 {\v C}ech 複体について所望の消滅を得る。$M$ に対する複体は、
これと平坦 $R$-加群 $M$ とのテンソル積をとることで得られるので
（More on Algebra, Lemma
\ref{more-algebra-lemma-extended-alternating-form-module}）、結論が従う。
\end{proof}

\begin{remark}
\label{perfect-remark-supported-map-c-equations}
$X$, $f_1, \ldots, f_c \in \Gamma(X, \mathcal{O}_X)$, および
$\mathcal{F}$ は Remark \ref{perfect-remark-support-c-equations} のとおりとする。
このとき標準的な $\mathcal{O}_X|_Z$-線形写像
$$
c_{f_1, \ldots, f_c} :
i^*\mathcal{F}
\longrightarrow
\mathcal{H}^c_Z(\mathcal{F})
$$
があり、$\mathcal{F}$ について関手的である。実際、拡張交代 {\v C}ech 複体
(\ref{perfect-equation-extended-alternating}) を $\mathcal{F}^\bullet$ と記すと、
Remark \ref{perfect-remark-extended-alternating-map-to-support} の標準写像
$\mathcal{F}^\bullet \to i_*R\mathcal{H}_Z(\mathcal{F})$ がある。
これは次数 $c$ のコホモロジー層上で標準写像
$$
\Coker\left(\bigoplus \mathcal{F}_{1 \ldots \hat i \ldots c} \to
\mathcal{F}_{1 \ldots c}\right)
\longrightarrow
i_*\mathcal{H}^c_Z(\mathcal{F})
$$
を定める。$\mathcal{F}$ の局所切断 $s$ が与えられたとき、
$\mathcal{F}_{1 \ldots c}$ の局所切断
$$
\frac{s}{f_1 \ldots f_c}
$$
を考えられる。上に表示した余核におけるこの切断の類は、写像
$(f_1, \ldots, f_c) : \mathcal{F}^{\oplus c} \to \mathcal{F}$ の像を法とした
$s$ のみに依存する。$i_*i^*\mathcal{F}$ は
$(f_1, \ldots, f_c) : \mathcal{F}^{\oplus c} \to \mathcal{F}$ の余核に等しいので、
$\mathcal{O}_X$-加群の写像
$i_*i^*\mathcal{F} \to i_*\mathcal{H}_Z^c(\mathcal{F})$ を得る。
$i_*$ は充満忠実なので、写像 $c_{f_1, \ldots, f_c}$ を得る。
\end{remark}

\begin{example}
\label{perfect-example-affine-supported-map-c-equations}
$X = \Spec(A)$ をアフィンとし、$f_1, \ldots, f_c \in A$ とする。また、ある
$A$-加群 $M$ に対して $\mathcal{F} = \widetilde{M}$ とする。Remark
\ref{perfect-remark-supported-map-c-equations} の写像 $c_{f_1, \ldots, f_c}$ は、写像
$$
M/(f_1, \ldots, f_c)M
\longrightarrow
\Coker\left(
\bigoplus M_{f_1 \ldots \hat f_i \ldots f_c} \to M_{f_1 \ldots f_c}
\right)
$$
として記述できる。これは $s \in M$ の類を、余核における
$s/f_1 \ldots f_c$ の類へ送る。
\end{example}

\begin{lemma}
\label{perfect-lemma-supported-map-determinant}
$X$, $f_1, \ldots, f_c \in \Gamma(X, \mathcal{O}_X)$, および
$\mathcal{F}$ は Remark \ref{perfect-remark-support-c-equations} のとおりとする。
$1 \leq i, j \leq c$ に対して $a_{ji} \in \Gamma(X, \mathcal{O}_X)$ をとり、
$g_j = \sum_{i = 1, \ldots, c} a_{ji}f_i$ とおく。$g_1, \ldots, g_c$ が
スキーム論的に $Z$ を切り出すと仮定する。$\mathcal{F}$ が準連接ならば
$$
c_{f_1, \ldots, f_c} = \det(a_{ji}) c_{g_1, \ldots, g_c}
$$
である。ここで $c_{f_1, \ldots, f_c}$ と $c_{g_1, \ldots, g_c}$ は Remark
\ref{perfect-remark-supported-map-c-equations} のとおりである。
\end{lemma}

\begin{proof}
$s \in \mathcal{F}(X)$ を任意にとるとき、$\mathcal{H}_Z(\mathcal{F})$ の
大域切断として $c_{f_1, \ldots, f_c}(s) =
\det(a_{ij}) c_{g_1, \ldots, g_c}(s)$ であることを示す。
そうすれば $X$ の任意の開部分スキーム上の切断についても同じ結果が得られるので、
これで十分である。そのため、$1 \leq i_0 < \ldots < i_p \leq c$ および
$1 \leq j_0 < \ldots < j_q \leq c$ に対して、
$U_{i_0 \ldots i_p} \subset X$,
$V_{j_0 \ldots j_q} \subset X$, および
$W_{i_0 \ldots i_p, j_0 \ldots j_q} \subset X$
を、それぞれ $f_{i_0} \ldots f_{i_p}$、$g_{j_0} \ldots g_{j_q}$、および
$f_{i_0} \ldots f_{i_p}g_{j_0} \ldots g_{j_q}$ が可逆となる開部分スキームとする。
$\mathcal{F}$ の $U_{i_0 \ldots i_p}$、resp.\ $V_{j_0 \ldots j_q}$、
resp.\ $W_{i_0 \ldots i_p, j_0 \ldots j_q}$ への制限を $X$ へ直像したものを、
$\mathcal{F}_{i_0 \ldots i_p}$、resp.\ $\mathcal{F}'_{j_0 \ldots j_q}$、
resp.\ $\mathcal{F}''_{i_0 \ldots i_p, j_0 \ldots j_q}$ と記す。
すると三つの拡張交代 {\v C}ech 複体
$$
\mathcal{F}^\bullet :
\mathcal{F} \to \bigoplus\nolimits_{i_0} \mathcal{F}_{i_0} \to
\bigoplus\nolimits_{i_0 < i_1} \mathcal{F}_{i_0i_1} \to \ldots
$$
および
$$
(\mathcal{F}')^\bullet :
\mathcal{F} \to \bigoplus\nolimits_{j_0} \mathcal{F}'_{j_0} \to
\bigoplus\nolimits_{j_0 < j_1} \mathcal{F}'_{j_0j_1} \to \ldots
$$
および
$$
(\mathcal{F}'')^\bullet :
\mathcal{F} \to
\bigoplus\nolimits_{i_0} \mathcal{F}_{i_0} \oplus
\bigoplus\nolimits_{j_0} \mathcal{F}'_{j_0} \to
\bigoplus\nolimits_{i_0 < i_1} \mathcal{F}_{i_0i_1} \oplus
\bigoplus\nolimits_{i_0, j_0} \mathcal{F}''_{i_0, j_0} \oplus
\bigoplus\nolimits_{j_0 < j_1} \mathcal{F}'_{j_0j_1} \to
\ldots
$$
を得る。その微分は (\ref{perfect-equation-extended-alternating}) の定義で用いたものである。
複体の写像
$$
(\mathcal{F}'')^\bullet \to \mathcal{F}^\bullet
\quad\text{および}\quad
(\mathcal{F}'')^\bullet \to (\mathcal{F}')^\bullet
$$
があり、各項上の射影写像によって与えられる（したがって次数 $0$ では恒等写像を
誘導する）。Lemma \ref{perfect-lemma-extended-alternating-represented} により、これらの
複体はいずれも $i_*R\mathcal{H}_Z(\mathcal{F})$ を表し、これらの写像はこの対象上の
恒等写像を表すことに注意せよ。したがって、二つの写像により
$c_{f_1, \ldots, f_c}(s)$ および $\det(a_{ji})c_{g_1, \ldots, g_c}(s)$
へ写る元
$$
\sigma \in H^c((\mathcal{F}'')^\bullet(X))
$$
を見つければ十分である。$\sigma$ を表すコサイクルは明示的に与えられる。
すなわち
$$
\sigma_{1 \ldots c} = \frac{s}{f_1 \ldots f_c} \in \mathcal{F}_{1 \ldots c}(X)
\quad\text{および}\quad
\sigma'_{1 \ldots c} = \frac{\det(a_{ji})s}{g_1 \ldots g_c} \in
\mathcal{F}'_{1 \ldots c}(X)
$$
とおき、さらに
$$
\sigma_{i_0 \ldots i_p, j_0 \ldots j_{c - p - 2}} =
\frac{\lambda(i_0 \ldots i_p, j_0 \ldots j_{c - p - 2})s}%
{f_{i_0} \ldots f_{i_p}g_{j_0} \ldots g_{j_{c - p - 2}}}
\in \mathcal{F}''_{i_0 \ldots i_p, j_0 \ldots j_{c - p - 2}}(X)
$$
とおく。ここで $\lambda(i_0 \ldots i_p, j_0 \ldots j_{c - p - 2})$ は、形式的な式
$$
e_{i_0} \wedge \ldots \wedge e_{i_p} \wedge
(a_{j_01} e_1 + \ldots + a_{j_0c}e_c) \wedge \ldots \wedge
(a_{j_{c - p - 2}1} e_1 + \ldots + a_{j_{c - p - 2}c}e_c)
$$
における $e_1 \wedge \ldots \wedge e_c$ の係数である。$\sigma$ がコサイクルで
あることを確かめるには、$1 \leq i_0 < \ldots < i_p \leq c$ および
$1 \leq j_0 < \ldots < j_{c - p - 1} \leq c$ に対して
\begin{align*}
0 & =
\sum\nolimits_{a = 0, \ldots, p} (-1)^a
f_{i_a} \lambda(i_0 \ldots \hat i_a \ldots i_p, j_0 \ldots j_{c - p - 1}) \\
& +
\sum\nolimits_{b = 0, \ldots, c - p - 1} (-1)^{p + b + 1}g_{j_b}
\lambda(i_0 \ldots i_p, j_0 \ldots \hat j_b \ldots j_{c - p - 1})
\end{align*}
が成り立つことを示さなければならない。これを見る最も簡単な方法は、おそらく
形式的な式
$$
\xi = e_{i_0} \wedge \ldots \wedge e_{i_p} \wedge
(a_{j_01} e_1 + \ldots + a_{j_0c}e_c) \wedge \ldots \wedge
(a_{j_{c - p - 1}1} e_1 + \ldots + a_{j_{c - p - 1}c}e_c)
$$
が $0$ であると論じることである。実際、これは $e_1, \ldots, e_c$ 上の自由加群の
第 $(c + 1)$ 外冪の元であり、上の式は $e_i \to f_i$ と送る Koszul 微分による
$\xi$ の像である。細部は省略する。
\end{proof}

\begin{lemma}
\label{perfect-lemma-supported-map-global}
スキーム $X$ をとる。$Z \to X$ を有限表示の閉埋め込みとし、その余法層
$\mathcal{C}_{Z/X}$ は階数 $c$ の局所自由層であるとする。このとき標準写像
$$
c :
\wedge^c(\mathcal{C}_{Z/X})^\vee \otimes_{\mathcal{O}_Z} i^*\mathcal{F}
\longrightarrow
\mathcal{H}_Z^c(\mathcal{F})
$$
があり、準連接加群 $\mathcal{F}$ について関手的である。
\end{lemma}

\begin{proof}
Remark \ref{perfect-remark-supported-map-c-equations} の構成と、Lemma
\ref{perfect-lemma-supported-map-determinant} で示されたイデアル層の生成元の選び方への
非依存性から従う。細部は省略する。
\end{proof}

\begin{remark}
\label{perfect-remark-supported-functorial}
$g : X' \to X$ をスキームの射とし、
$f_1, \ldots, f_c \in \Gamma(X, \mathcal{O}_X)$ とする。
$f'_i = g^\sharp(f_i) \in \Gamma(X', \mathcal{O}_{X'})$ とおく。
$f_1, \ldots, f_c$、resp.\ $f'_1, \ldots, f'_c$ によって切り出される閉部分
スキームを $Z \subset X$、resp.\ $Z' \subset X'$ と記す。このとき
$Z' = Z \times_X X'$ である。誘導されるスキームの射を $h : Z' \to Z$ と記す。
$\mathcal{F}$ を $\mathcal{O}_X$-加群とし、$\mathcal{F}' = g^*\mathcal{F}$
とおく。この状況で $\mathcal{F}$ が準連接ならば、図式
$$
\xymatrix{
(i')^{-1}\mathcal{O}_{X'} \otimes_{h^{-1}i^{-1}\mathcal{O}_X}
h^{-1}\mathcal{H}^c_Z(\mathcal{F}) \ar[r] &
\mathcal{H}_{Z'}^c(\mathcal{F}') \\
h^*i^*\mathcal{F} \ar[r] \ar[u]_-{c_{f_1, \ldots, f_c}} &
(i')^*\mathcal{F}' \ar[u]^-{c_{f'_1, \ldots, f'_c}}
}
$$
は可換である。ここで上の水平射は Cohomology, Remark
\ref{cohomology-remark-support-functorial} の写像の次数 $c$ のコホモロジー層上の
写像である。具体的には、$\mathcal{F}, f_1, \ldots, f_c$、resp.\
$\mathcal{F}', f'_1, \ldots, f'_c$ を用いて Remark
\ref{perfect-remark-support-c-equations} で構成した拡張交代 {\v C}ech 複体を
$\mathcal{F}^\bullet$、resp.\ $(\mathcal{F}')^\bullet$ と記す。
$(\mathcal{F}')^\bullet = g^*\mathcal{F}^\bullet$ であることに注意せよ。
このとき $\mathcal{F}$ の準連接性を仮定しなくても、図式
$$
\xymatrix{
i'_* L(g|_{Z'})^* R\mathcal{H}_Z(\mathcal{F}) \ar[r] \ar@{=}[d] &
i'_*R\mathcal{H}_{Z'}(Lg^*\mathcal{F}) \ar[d] \\
Lg^*i_*R\mathcal{H}_Z(\mathcal{F}) &
i'_*R\mathcal{H}_{Z'}(\mathcal{F}') \\
Lg^*(\mathcal{F}^\bullet) \ar[u] \ar[r] &
(\mathcal{F}')^\bullet \ar[u]
}
$$
は可換である。ここで $g|_{Z'} : (Z', (i')^{-1}\mathcal{O}_{X'}) \to
(Z, i^{-1}\mathcal{O}_X)$ は誘導される環付き空間の射である。
上の水平射は Cohomology, Remark
\ref{cohomology-remark-support-functorial} で与えられ、等号の説明もそこにある。
上向きの射は Remark \ref{perfect-remark-extended-alternating-map-to-support} から来る。
下の水平射は写像 $Lg^*\mathcal{F}^\bullet
\to g^*\mathcal{F}^\bullet = (\mathcal{F}')^\bullet$ であり、下向きの射は
$Lg^*\mathcal{F} \to g^*\mathcal{F} = \mathcal{F}'$ によって誘導される。
この図式は可換である。実際、図式をどちら向きに回っても二つの射
$Lg^*\mathcal{F}^\bullet \to i'_*R\mathcal{H}_{Z'}(\mathcal{F}')$ を得るが、
それぞれを $i'_*R\mathcal{H}_{Z'}(\mathcal{F}') \to \mathcal{F}'$ と合成すると、
標準写像 $Lg^*\mathcal{F}^\bullet \to \mathcal{F}'$ となる。細部は省略する。
ここで、この図式の次数 $c$ のコホモロジー層を調べ、Remark
\ref{perfect-remark-supported-map-c-equations} で用いた写像
$i^*\mathcal{F} \to
\Coker(\bigoplus \mathcal{F}_{1 \ldots \hat i \ldots c} \to
\mathcal{F}_{1 \ldots c})$
の構成が引き戻しと両立することを用いれば、最初の図式の可換性が従う。
\end{remark}








\section{コヒーレータ}
\label{perfect-section-coherator}

\noindent
スキーム $X$ をとる。{\it コヒーレータ}とは、関手
$$
Q_X :
\textit{Mod}(\mathcal{O}_X)
\longrightarrow
\QCoh(\mathcal{O}_X)
$$
であって、包含関手
$\QCoh(\mathcal{O}_X) \to \textit{Mod}(\mathcal{O}_X)$
の右随伴となるものをいう。これは任意のスキーム $X$ に対して存在し、
さらに随伴射 $Q_X(\mathcal{F}) \to \mathcal{F}$ は任意の
準連接加群 $\mathcal{F}$ に対して同型である。Properties,
Proposition \ref{properties-proposition-coherator} を参照せよ。
$Q_X$ は右随伴であるから左完全であり、その右導来拡張
$$
RQ_X :
D(\mathcal{O}_X)
\longrightarrow
D(\QCoh(\mathcal{O}_X)).
$$
を考えることができる。$Q_X$ は包含関手
$\QCoh(\mathcal{O}_X) \to \textit{Mod}(\mathcal{O}_X)$
の右随伴なので、Derived Categories, Lemma
\ref{derived-lemma-derived-adjoint-functors} により、
$RQ_X$ は標準関手
$D(\QCoh(\mathcal{O}_X)) \to D(\mathcal{O}_X)$ の右随伴である。

\medskip\noindent
本節では関手 $RQ_X$ を調べる。Section \ref{perfect-section-better-coherator}
では、包含関手
$D_\QCoh(\mathcal{O}_X) \to D(\mathcal{O}_X)$
の右随伴（存在する場合）を調べる。これは前述の関手と密接に関係する。

\begin{lemma}
\label{perfect-lemma-affine-pushforward}
スキームのアフィン射 $f : X \to Y$ をとる。このとき $f_*$ は導来関手
$f_* : D(\QCoh(\mathcal{O}_X)) \to D(\QCoh(\mathcal{O}_Y))$
を定め、この関手について図式
$$
\xymatrix{
D(\QCoh(\mathcal{O}_X)) \ar[d]_{f_*} \ar[r] &
D_\QCoh(\mathcal{O}_X) \ar[d]^{Rf_*} \\
D(\QCoh(\mathcal{O}_Y)) \ar[r] &
D_\QCoh(\mathcal{O}_Y)
}
$$
は可換である。
\end{lemma}

\begin{proof}
関手
$f_* : \QCoh(\mathcal{O}_X) \to \QCoh(\mathcal{O}_Y)$
は完全である。Cohomology of Schemes, Lemma
\ref{coherent-lemma-relative-affine-vanishing} を参照せよ。
したがって $f_*$ は、Derived Categories, Lemma
\ref{derived-lemma-right-derived-exact-functor} により、
任意の代表複体に単に $f_*$ を適用することで導来関手
$f_* : D(\QCoh(\mathcal{O}_X)) \to D(\QCoh(\mathcal{O}_Y))$
が定まる。図式の可換性は Lemma \ref{perfect-lemma-pushforward-affine-morphism}
による。
\end{proof}

\begin{lemma}
\label{perfect-lemma-flat-pushforward-coherator}
スキームの射 $f : X \to Y$ をとる。$f$ は準コンパクト、
準分離かつ平坦であるとする。
左完全関手
$f_* : \QCoh(\mathcal{O}_X) \to \QCoh(\mathcal{O}_Y)$
の右導来関手を
$$
\Phi : D(\QCoh(\mathcal{O}_X)) \to D(\QCoh(\mathcal{O}_Y))
$$
と記す。このとき $RQ_Y \circ Rf_* = \Phi \circ RQ_X$ である。
\end{lemma}

\begin{proof}
$RQ_Y \circ Rf_*$ と $\Phi \circ RQ_X$ が同じ関手
$D(\QCoh(\mathcal{O}_Y)) \to D(\mathcal{O}_X)$
の右随伴であることを示す。

\medskip\noindent
$f$ は準コンパクトかつ準分離なので、$f_*$ は準連接性を保つ。
Schemes, Lemma \ref{schemes-lemma-push-forward-quasi-coherent}
を参照せよ。$\QCoh(\mathcal{O}_X)$ はグロタンディーク・アーベル圏である
ことを思い出そう（Properties, Proposition
\ref{properties-proposition-coherator}）。したがって
$D(\QCoh(\mathcal{O}_X))$ の任意の $K$ は、
$\QCoh(\mathcal{O}_X)$ の K-入射的複体 $\mathcal{I}^\bullet$
によって表される。Injectives, Theorem
\ref{injectives-theorem-K-injective-embedding-grothendieck}
を参照せよ。このとき $\Phi(K) = f_*\mathcal{I}^\bullet$
と定めることができる。

\medskip\noindent
$f$ は平坦なので、関手 $f^*$ は完全である。したがって $f^*$ は
$f^* : D(\mathcal{O}_Y) \to D(\mathcal{O}_X)$ および
$f^* : D(\QCoh(\mathcal{O}_Y)) \to D(\QCoh(\mathcal{O}_X))$
を定める。関手
$f^* = Lf^* : D(\mathcal{O}_Y) \to D(\mathcal{O}_X)$ は
$Rf_* : D(\mathcal{O}_X) \to D(\mathcal{O}_Y)$,
の左随伴である。Cohomology, Lemma \ref{cohomology-lemma-adjoint}
を参照せよ。同様に、Derived Categories, Lemma
\ref{derived-lemma-derived-adjoint-functors} により、関手
$f^* : D(\QCoh(\mathcal{O}_Y)) \to D(\QCoh(\mathcal{O}_X))$ は
$\Phi : D(\QCoh(\mathcal{O}_X)) \to D(\QCoh(\mathcal{O}_Y))$
の左随伴である。

\medskip\noindent
$A$ を $D(\QCoh(\mathcal{O}_Y))$ の対象、
$E$ を $D(\mathcal{O}_X)$ の対象とする。このとき
\begin{align*}
\Hom_{D(\QCoh(\mathcal{O}_Y))}(A, RQ_Y(Rf_*E))
& =
\Hom_{D(\mathcal{O}_Y)}(A, Rf_*E) \\
& =
\Hom_{D(\mathcal{O}_X)}(f^*A, E) \\
& =
\Hom_{D(\QCoh(\mathcal{O}_X))}(f^*A, RQ_X(E)) \\
& =
\Hom_{D(\QCoh(\mathcal{O}_Y))}(A, \Phi(RQ_X(E)))
\end{align*}
となり、主張が従う。
\end{proof}

\begin{lemma}
\label{perfect-lemma-affine-coherator}
アフィン・スキーム $X = \Spec(A)$ に対して、次が成り立つ。
\begin{enumerate}
\item $Q_X : \textit{Mod}(\mathcal{O}_X) \to \QCoh(\mathcal{O}_X)$
は、$\mathcal{F}$ を $A$-加群 $\Gamma(X, \mathcal{F})$
に付随する準連接 $\mathcal{O}_X$-加群へ送る関手である。
\item $RQ_X : D(\mathcal{O}_X) \to D(\QCoh(\mathcal{O}_X))$
は、$E$ を $D(A)$ の対象 $R\Gamma(X, E)$ に付随する
準連接 $\mathcal{O}_X$-加群の複体へ送る関手である。
\item $D_\QCoh(\mathcal{O}_X)$ に制限した関手 $RQ_X$ は、
(\ref{perfect-equation-compare}) の擬逆関手を定める。
\end{enumerate}
\end{lemma}

\begin{proof}
Schemes, Lemma \ref{schemes-lemma-compare-constructions} により、
関手 $Q_X$ は
$$
\mathcal{F} \mapsto \widetilde{\Gamma(X, \mathcal{F})}
$$
である。これから直ちに (1) と (2) が従う。第三の主張は
Lemma \ref{perfect-lemma-affine-compare-bounded}（の証明）から従う。
\end{proof}

\noindent
ここで、関手
$D(\QCoh(\mathcal{O}_X)) \to D_\QCoh(\mathcal{O}_X)$
が圏同値となるための判定条件を証明する準備が整った。

\begin{lemma}
\label{perfect-lemma-argument-proves}
準コンパクトかつ準分離なスキーム $X$ をとる。任意のアフィン開集合
$U \subset X$ に対し、左完全関手
$j_* : \QCoh(\mathcal{O}_U) \to \QCoh(\mathcal{O}_X)$
の右導来関手
$$
\Phi : D(\QCoh(\mathcal{O}_U)) \to D(\QCoh(\mathcal{O}_X))
$$
が可換図式
$$
\xymatrix{
D(\QCoh(\mathcal{O}_U)) \ar[d]_\Phi \ar[r]_{i_U} &
D_\QCoh(\mathcal{O}_U) \ar[d]^{Rj_*} \\
D(\QCoh(\mathcal{O}_X)) \ar[r]^{i_X} &
D_\QCoh(\mathcal{O}_X)
}
$$
に組み込まれると仮定する。このとき関手 (\ref{perfect-equation-compare})
$$
D(\QCoh(\mathcal{O}_X))
\longrightarrow
D_\QCoh(\mathcal{O}_X)
$$
は圏同値であり、その擬逆関手は $RQ_X$ で与えられる。
\end{lemma}

\begin{proof}
$E$ を $D_\QCoh(\mathcal{O}_X)$ の対象、
$A$ を $D(\QCoh(\mathcal{O}_X))$ の対象とする。随伴射
$$
A \to RQ_X(i_X(A))
\quad\text{および}\quad
i_X(RQ_X(E)) \to E
$$
が同型であることを示せばよい。次の仮説 $H_n$ を考える：
$E$ と $i_X(A)$ が、$X$ の $n$ 個のアフィン開集合の合併に含まれる
$X$ の閉部分集合に支持される
（Definition \ref{perfect-definition-supported-on}）ならば、
上の随伴射は同型である。$n$ に関する帰納法で $H_n$ を証明する。

\medskip\noindent
基底の場合：$n = 0$ とする。このとき $E = 0$ なので、写像
$i_X(RQ_X(E)) \to E$ は同型である。同様に $i_X(A) = 0$ である。
したがって $i_X(A)$ のコホモロジー層は零である。包含関手
$\QCoh(\mathcal{O}_X) \to \textit{Mod}(\mathcal{O}_X)$
は充満忠実かつ完全なので、$A$ のコホモロジー対象は零である。
すなわち $A = 0$ であり、$A \to RQ_X(i_X(A))$ も同型である。

\medskip\noindent
帰納段階。$U_i \subset X$ をアフィン開集合とし、$E$ と $i_X(A)$ が
$U_1 \cup \ldots \cup U_n$ に含まれる $X$ の閉部分集合 $T$
に支持されると仮定する。$U = U_n$ とおく。完全三角形
$$
A \to \Phi(A|_U) \to A' \to A[1]
\quad\text{および}\quad
E \to Rj_*(E|_U) \to E' \to E[1]
$$
を考える。ここで $\Phi$ は補題の主張に現れる関手である。
$E \to Rj_*(E|_U)$ は $U = U_n$ 上で擬同型であることに注意する。
仮定より $i_X \circ \Phi = Rj_* \circ i_U$ であり、
$i_X(A)|_U = i_U(A|_U)$ なので、
$i_X(A) \to i_X(\Phi(A|_U))$ は $U$ 上の擬同型である。
したがって $i_X(A')$ と $E'$ は、$U_1 \cup \ldots \cup U_{n - 1}$
に含まれる $X$ の閉部分集合 $T \setminus U$ に支持される。
帰納法の仮定により、主張は $A'$ と $E'$ に対して成り立つ。
Derived Categories, Lemma
\ref{derived-lemma-third-isomorphism-triangle} により、写像
$$
\Phi(A|_U) \to RQ_X(i_X(\Phi(A|_U)))
\quad\text{および}\quad
i_X(RQ_X(Rj_*E|_U)) \to Rj_*(E|_U)
$$
が同型であることを示せば十分である。仮定と Lemma
\ref{perfect-lemma-flat-pushforward-coherator} により（包含射
$j : U \to X$ は平坦、準コンパクトかつ準分離である）、
$$
RQ_X(i_X(\Phi(A|_U))) = RQ_X(Rj_*(i_U(A|_U))) = \Phi(RQ_U(i_U(A|_U)))
$$
および
$$
i_X(RQ_X(Rj_*(E|_U))) = i_X(\Phi(RQ_U(E|_U))) = Rj_*(i_U(RQ_U(E|_U)))
$$
を得る。最後に、写像
$$
A|_U \to RQ_U(i_U(A|_U))
\quad\text{および}\quad
i_U(RQ_U(E|_U)) \to E|_U
$$
は Lemma \ref{perfect-lemma-affine-coherator} により同型である。以上で結論が従う。
\end{proof}

\begin{proposition}
\label{perfect-proposition-quasi-compact-affine-diagonal}
対角射がアフィンである準コンパクト・スキーム $X$ をとる。
このとき関手 (\ref{perfect-equation-compare})
$$
D(\QCoh(\mathcal{O}_X))
\longrightarrow
D_\QCoh(\mathcal{O}_X)
$$
は圏同値であり、その擬逆関手は $RQ_X$ で与えられる。
\end{proposition}

\begin{proof}
$U \subset X$ をアフィン開集合とする。Morphisms, Lemma
\ref{morphisms-lemma-affine-permanence} により、射 $U \to X$
はアフィンである。したがって Lemma
\ref{perfect-lemma-affine-pushforward} により Lemma
\ref{perfect-lemma-argument-proves} の仮定が成り立ち、結論を得る。
\end{proof}

\begin{lemma}
\label{perfect-lemma-direct-image-coherator}
スキームの射 $f : X \to Y$ をとる。$X$ と $Y$ は準コンパクトであり、
その対角射はアフィンであると仮定する。
左完全関手
$f_* : \QCoh(\mathcal{O}_X) \to \QCoh(\mathcal{O}_Y)$
の右導来関手を
$$
\Phi : D(\QCoh(\mathcal{O}_X)) \to D(\QCoh(\mathcal{O}_Y))
$$
と記す。このとき図式
$$
\xymatrix{
D(\QCoh(\mathcal{O}_X)) \ar[d]_\Phi \ar[r] &
D_\QCoh(\mathcal{O}_X) \ar[d]^{Rf_*} \\
D(\QCoh(\mathcal{O}_Y)) \ar[r] &
D_\QCoh(\mathcal{O}_Y)
}
$$
は可換である。
\end{lemma}

\begin{proof}
Proposition \ref{perfect-proposition-quasi-compact-affine-diagonal} により、
図式の横矢印は圏同値であることに注意する。したがってこれらの圏を
同一視できる（対角射がアフィンである他の準コンパクト・スキームに
ついても同様である）。補題の主張は、
$D(\QCoh(\mathcal{O}_X))$ のすべての $K$ に対して標準写像
$\Phi(K) \to Rf_*(K)$ が同型になるということである。
$K_1 \to K_2 \to K_3 \to K_1[1]$ が
$D(\QCoh(\mathcal{O}_X))$ の完全三角形であり、
三項のうち二項について主張が成り立つなら、残りの一項についても
成り立つことに注意する。

\medskip\noindent
$U \subset X$ をアフィン開集合とする。$X$ の対角射はアフィンなので、
包含射 $j : U \to X$ はアフィンである（Morphisms, Lemma
\ref{morphisms-lemma-affine-permanence}）。同様に、合成
$g = f \circ j : U \to Y$ もアフィンである。
$\mathcal{I}^\bullet$ を $\QCoh(\mathcal{O}_U)$ の
K-入射的複体とする。関手
$j_* : \QCoh(\mathcal{O}_U) \to \QCoh(\mathcal{O}_X)$
は完全な左随伴
$j^* : \QCoh(\mathcal{O}_X) \to \QCoh(\mathcal{O}_U)$
をもつので、$j_*\mathcal{I}^\bullet$ は
$\QCoh(\mathcal{O}_X)$ の K-入射的複体である。Derived Categories,
Lemma \ref{derived-lemma-adjoint-preserve-K-injectives} を参照せよ。
したがって
$$
\Phi(j_*\mathcal{I}^\bullet) =
f_*j_*\mathcal{I}^\bullet =
g_*\mathcal{I}^\bullet
$$
を得る。Lemma \ref{perfect-lemma-affine-pushforward} により、
$j_*\mathcal{I}^\bullet$ は $Rj_*\mathcal{I}^\bullet$ を表し、
$g_*\mathcal{I}^\bullet$ は $Rg_*\mathcal{I}^\bullet$ を表す。
一方、$Rf_* \circ Rj_* = Rg_*$ である。ゆえに
$f_*j_*\mathcal{I}^\bullet$ は
$Rf_*(j_*\mathcal{I}^\bullet)$ を表す。
したがって、$\mathcal{G}^\bullet$ が $U$ 上の準連接加群の複体ならば、
$j_*\mathcal{G}^\bullet$ の形の任意の複体について補題が成り立つ。
（実際、$\mathcal{G}^\bullet \to \mathcal{I}^\bullet$ が擬同型なら、
$j_*$ は準連接加群上で完全な関手なので、
$j_*\mathcal{G}^\bullet \to j_*\mathcal{I}^\bullet$
も擬同型である。）

\medskip\noindent
$\mathcal{F}^\bullet$ を準連接 $\mathcal{O}_X$-加群の複体とする。
$T \subset X$ を、すべての $p$ について $\mathcal{F}^p$
の台が $T$ に含まれるような閉部分集合とする。
$T \subset U_1 \cup \ldots \cup U_n$ となるアフィン開集合
$U_1, \ldots, U_n$ の最小個数 $n$ に関する帰納法を用いる。
基底の場合 $n = 0$ は自明である。$n \geq 1$ なら $U = U_1$
とおき、上と同じ開埋め込みを $j : U \to X$ と記す。
複体の写像
$c : \mathcal{F}^\bullet \to j_*j^*\mathcal{F}^\bullet$
を考える。複体の二つの短完全列
$$
0 \to \Ker(c) \to \mathcal{F}^\bullet \to \Im(c) \to 0
$$
および
$$
0 \to \Im(c) \to j_*j^*\mathcal{F}^\bullet \to \Coker(c) \to 0
$$
を得る。複体 $\Ker(c)$ と $\Coker(c)$ は
$T \setminus U \subset U_2 \cup \ldots \cup U_n$ に支持されるので、
帰納法によりこれらについて結論が成り立つ。
前段落の議論により $j_*j^*\mathcal{F}^\bullet$ についても成り立つ。
短完全列に付随する完全三角形を調べ、証明冒頭の注意を用いれば結論を得る。
\end{proof}

\begin{remark}[Warning]
\label{perfect-remark-warning-coherator}
対角射がアフィンである準コンパクト・スキーム $X$ をとる。
Proposition \ref{perfect-proposition-quasi-compact-affine-diagonal} により
$D(\QCoh(\mathcal{O}_X)) = D_\QCoh(\mathcal{O}_X)$
であることが分かっていても、奇妙な現象が起こりうるため、
この題材では誤りを犯しやすい。注意すべき点の一つは、
この等式が導来関手まで同じであることを意味すると軽率に仮定することである。
例えば、準コンパクトな開集合 $U \subset X$ があるとする。このとき、
左完全関手 $\Gamma(U, -)$ の高次右導来関手
$$
R^i(\QCoh)\Gamma(U, -) : \QCoh(\mathcal{O}_X) \to \textit{Ab}
$$
を考えることができる。これは普遍 $\delta$-関手であり、また
$X$ 上のすべてのアーベル層に対して定義される関手 $H^i(U, -)$
を $\QCoh(\mathcal{O}_X)$ に制限したものは $\delta$-関手をなすので、
標準変換
$$
t^i : R^i(\QCoh)\Gamma(U, -) \to H^i(U, -).
$$
を得る。$X = \Spec(A)$ がアフィンであっても、これらの変換は一般には
同型ではない。実際、$I$ が入射的 $A$-加群なら、右導来関手の構成と
$\QCoh(\mathcal{O}_X)$ と $\text{Mod}_A$ の圏同値により、
$R^1(\QCoh)\Gamma(U, \widetilde{I}) = 0$ である。
しかし Examples, Lemma \ref{examples-lemma-nonvanishing} は、
$H^1(U, \widetilde{I}) \not = 0$ となる $A$, $I$, $U$ が存在することを示す。
\end{remark}




\section{ネーター・スキームに対するコヒーレータ}
\label{perfect-section-coherator-Noetherian}

\noindent
ネーター・スキームの場合には、次の補題を用いることができる。

\begin{lemma}
\label{perfect-lemma-injective-quasi-coherent-sheaf-Noetherian}
$X$ をネーター・スキームとし、$\mathcal{J}$ を
$\QCoh(\mathcal{O}_X)$ の入射的対象とする。このとき $\mathcal{J}$
は $\mathcal{O}_X$-加群の脆弱層である。
\end{lemma}

\begin{proof}
$U \subset X$ を開部分集合、$s \in \mathcal{J}(U)$ を切断とする。
$X \setminus U$ 上の誘導された被約スキーム構造を定める
準連接イデアル層を $\mathcal{I} \subset \mathcal{O}_X$ とする
（Schemes, Definition \ref{schemes-definition-reduced-induced-scheme}
を参照せよ）。Cohomology of Schemes, Lemma
\ref{coherent-lemma-homs-over-open} により、切断 $s$ はある $n$
に対する写像 $\sigma : \mathcal{I}^n \to \mathcal{J}$ に対応する。
$\mathcal{J}$ は $\QCoh(\mathcal{O}_X)$ の入射的対象なので、
$\sigma$ を写像 $\tilde s : \mathcal{O}_X \to \mathcal{J}$
へ延長できる。このとき $\tilde s$ は、$s$ に制限される
$\mathcal{J}$ の大域切断に対応する。
\end{proof}

\begin{lemma}
\label{perfect-lemma-Noetherian-pushforward}
ネーター・スキームの射 $f : X \to Y$ をとる。
準連接層上の $f_*$ は右導来拡張
$\Phi : D(\QCoh(\mathcal{O}_X)) \to D(\QCoh(\mathcal{O}_Y))$
をもち、図式
$$
\xymatrix{
D(\QCoh(\mathcal{O}_X)) \ar[d]_{\Phi} \ar[r] &
D_\QCoh(\mathcal{O}_X) \ar[d]^{Rf_*} \\
D(\QCoh(\mathcal{O}_Y)) \ar[r] &
D_\QCoh(\mathcal{O}_Y)
}
$$
は可換である。
\end{lemma}

\begin{proof}
$X$ と $Y$ はネーター・スキームなので、この射は準コンパクトかつ
準分離である（Properties, Lemma
\ref{properties-lemma-locally-Noetherian-quasi-separated} および
Schemes, Remark \ref{schemes-remark-quasi-compact-and-quasi-separated}
を参照せよ）。したがって $f_*$ は準連接性を保つ。Schemes, Lemma
\ref{schemes-lemma-push-forward-quasi-coherent} を参照せよ。
次に、$K$ を $D(\QCoh(\mathcal{O}_X))$ の対象とする。
$\QCoh(\mathcal{O}_X)$ はグロタンディーク・アーベル圏なので
（Properties, Proposition \ref{properties-proposition-coherator}）、
各 $\mathcal{I}^n$ が $\QCoh(\mathcal{O}_X)$ の入射的対象となる
K-入射的複体 $\mathcal{I}^\bullet$ によって $K$ を表すことができる。
Injectives, Theorem
\ref{injectives-theorem-K-injective-embedding-grothendieck}
を参照せよ。したがって関手 $\Phi$ は
$$
\Phi(K) = f_*\mathcal{I}^\bullet
$$
とおくことで定義される。ここで右辺は
$D(\QCoh(\mathcal{O}_Y))$ の対象とみなす。
補題の証明を終えるには、標準写像
$$
f_*\mathcal{I}^\bullet \longrightarrow Rf_*\mathcal{I}^\bullet
$$
が $D(\mathcal{O}_Y)$ における同型であることを示せば十分である。
Lemma \ref{perfect-lemma-acyclicity-lemma} により、これを示すには
すべての $n \in \mathbf{Z}$ に対して $\mathcal{I}^n$ が右
$f_*$-非輪状であることを示せばよい。Lemma
\ref{perfect-lemma-injective-quasi-coherent-sheaf-Noetherian} により
$\mathcal{I}^n$ は脆弱であり、Cohomology, Lemma
\ref{cohomology-lemma-flasque-acyclic-pushforward} により
脆弱加群は右 $f_*$-非輪状なので、これは成り立つ。
\end{proof}

\begin{proposition}
\label{perfect-proposition-Noetherian}
$X$ をネーター・スキームとする。このとき関手 (\ref{perfect-equation-compare})
$$
D(\QCoh(\mathcal{O}_X))
\longrightarrow
D_\QCoh(\mathcal{O}_X)
$$
は圏同値であり、その擬逆関手は $RQ_X$ で与えられる。
\end{proposition}

\begin{proof}
Lemma \ref{perfect-lemma-argument-proves} と
Lemma \ref{perfect-lemma-Noetherian-pushforward} から従う。
\end{proof}






\section{コズル複体}
\label{perfect-section-koszul}

\noindent
$A$ を環、$f_1, \ldots, f_r$ を $A$ の元の列とする。
コズル複体 $K_\bullet(f_1, \ldots, f_r)$ は More on Algebra,
Definition \ref{more-algebra-definition-koszul-complex} で定義した。
これは次数 $r, \ldots, 0$ に位置する鎖複体である。
$K^{-n}(f_1, \ldots, f_r) = K_n(f_1, \ldots, f_r)$ とおき、
同じ微分を用いることで、これを上鎖複体
$K^\bullet(f_1, \ldots, f_r)$ にする。本節の残りでは、
すべての複体を上鎖複体とする。

\medskip\noindent
完全三角形
$$
I^\bullet(f_1, \ldots, f_r) \to
A \to
K^\bullet(f_1, \ldots, f_r) \to
I^\bullet(f_1, \ldots, f_r)[1]
$$
が $K(A)$ において得られるような複体
$I^\bullet(f_1, \ldots, f_r)$ を定義する。すなわち
$$
I^i(f_1, \ldots, f_r) =
\left\{
\begin{matrix}
K^{i - 1}(f_1, \ldots, f_r) & \text{ただし } i \leq 0 \\
0 & \text{それ以外}
\end{matrix}
\right.
$$
とおき、$K^\bullet(f_1, \ldots, f_r)$ 上の微分の負を用いる。
完全三角形の写像は明らかなものである。
$I^0(f_1, \ldots, f_r) = A^{\oplus r} \to A$ は、
$i$ 番目の因子上では $f_i$ による乗法で与えられることに注意する。
したがって $I^\bullet(f_1, \ldots, f_r) \to A$ は
$$
I^\bullet(f_1, \ldots, f_r) \to I \to A
$$
と分解する。ここで $I = (f_1, \ldots, f_r)$ である。実際、短完全列
$$
0 \to H^{-1}(K^\bullet(f_1, \ldots, f_s)) \to
H^0(I^\bullet(f_1, \ldots, f_s)) \to I \to 0
$$
があり、すべての $i < 0$ に対して
$H^i(I^\bullet(f_1, \ldots, f_r)) = H^{i - 1}(K^\bullet(f_1, \ldots, f_r))$
である。$A$ の元の第二の列 $g_1, \ldots, g_r$ が与えられると、
標準写像
$$
I^\bullet(f_1g_1, \ldots, f_rg_r) \to I^\bullet(f_1, \ldots, f_r)
\quad\text{および}\quad
K^\bullet(f_1g_1, \ldots, f_rg_r) \to K^\bullet(f_1, \ldots, f_r)
$$
があり、これらは上記の写像と両立する。第一の写像は
$I^0(f_1g_1, \ldots, f_rg_r) = A^{\oplus r}$ の $i$ 番目の直和因子上で
$g_i$ による乗法として与えられる。特に、$f_1, \ldots, f_r$
が与えられると、複体の逆系
\begin{equation}
\label{perfect-equation-system}
I^\bullet(f_1, \ldots, f_r) \leftarrow
I^\bullet(f_1^2, \ldots, f_r^2) \leftarrow
I^\bullet(f_1^3, \ldots, f_r^3) \leftarrow \ldots
\end{equation}
を得る。これは以下で重要な役割を果たす。次の補題を簡潔に述べるため、
いくつかの記法を固定する。

\begin{situation}
\label{perfect-situation-complex}
ここで $A$ は環、$f_1, \ldots, f_r$ は $A$ の元の列である。
$X = \Spec(A)$ および
$U = D(f_1) \cup \ldots \cup D(f_r) \subset X$ とおく。
$U$ のこの開被覆を
$\mathcal{U} : U = \bigcup_{i = 1, \ldots, r} D(f_i)$ と記す。
\end{situation}

\noindent
最初の補題は、上の複体を用いて $U$ 上の準連接層の
コホモロジーを計算できるというものである。

\begin{lemma}
\label{perfect-lemma-alternating-cech-complex}
Situation \ref{perfect-situation-complex} において、$M$ を $A$-加群とし、
付随する $\mathcal{O}_X$-加群を $\mathcal{F}$ と記す。このとき、
複体の標準同型
$$
\Psi : \colim_e \Hom_A(I^\bullet(f_1^e, \ldots, f_r^e), M)
\longrightarrow
\check{\mathcal{C}}_{alt}^\bullet(\mathcal{U}, \mathcal{F})
$$
があり、これは $M$ に関して関手的である。ここで $\Hom$-複体の微分は
$I^\bullet(f_1^e, \ldots, f_r^e)$ 上の微分の反変写像である。
\end{lemma}

\begin{proof}
交代 {\v C}ech 複体とは、通常の {\v C}ech 複体のうち交代上鎖で
与えられる部分複体であることを思い出そう。Cohomology, Section
\ref{cohomology-section-alternating-cech} を参照せよ。通常どおり、
$\check{\mathcal{C}}_{alt}^\bullet(\mathcal{U}, \mathcal{F})$
の $p$-上鎖を、$\{1, \ldots, r\}^{p + 1}$ 上の交代関数 $s$
とみなす。その $(i_0, \ldots, i_p)$ における値
$s_{i_0\ldots i_p}$ は
$M_{f_{i_0}\ldots f_{i_p}} = \mathcal{F}(U_{i_0\ldots i_p})$
に属する。一方、
$\Hom^\bullet(I^\bullet(f_1^e, \ldots, f_r^e), M)$
の $p$-上鎖 $t$ は写像
$t : \wedge^{p + 1}(A^{\oplus r}) \to M$ である。
$i$ 番目の基底元を $[i] \in A^{\oplus r}$ と書き、
$$
[i_0, \ldots, i_p] = [i_0] \wedge \ldots \wedge [i_p]
\in \wedge^{p + 1}(A^{\oplus r})
$$
と書く。上の $t$ に対して
$$
\Psi(t)_{i_0 \ldots i_p} =
(-1)^p \frac{t([i_0, \ldots, i_p])}{f_{i_0}^e\ldots f_{i_p}^e}
$$
$\Psi(t)$ は明らかに交代上鎖である。系 (\ref{perfect-equation-system})
の遷移写像
$$
I^\bullet(f_1^e, \ldots, f_r^e) \leftarrow
I^\bullet(f_1^{e + 1}, \ldots, f_r^{e + 1}),
$$
は $[i_0, \ldots, i_p]$ を
$f_{i_0}\ldots f_{i_p}[i_0, \ldots, i_p]$ へ送るので、
上の規則は系の遷移写像と両立する。局所化
$M_{f_{i_0} \ldots f_{i_p}}$ の Algebra, Lemma
\ref{algebra-lemma-localization-colimit} における記述から、
規則 $\Psi$ が余極限において次数 $p$ の上鎖加群の同型を定めることは
明らかである。証明を終えるには、この写像が微分と両立することを
示せばよい。そのため、上の $t$ に対して次を計算する。
\begin{align*}
d(\Psi(t))_{i_0 \ldots i_{p + 1}}
& =
\sum\nolimits_{j = 0}^{p + 1} (-1)^j
\Psi(t)_{i_0\ldots \hat i_j \ldots i_{p + 1}} \\
& =
(-1)^p
\sum\nolimits_{j = 0}^{p + 1} (-1)^j t([i_0 \ldots \hat i_j \ldots i_{p + 1}])
(f_{i_0} \ldots \hat f_{i_j} \ldots f_{i_p})^{-e}
\end{align*}
$I^\bullet(f_1^e, \ldots, f_r^e)$ 上の微分は
$K^\bullet(f_1, \ldots, f_r)$ 上の微分の負であることを思い出そう。
したがって
\begin{align*}
\Psi(d(t))_{i_0 \ldots i_{p + 1}} & =
(-1)^{p + 1}
d(t)([i_0, \ldots, i_{p + 1}]) (f_{i_0} \ldots f_{i_{p + 1}})^{-e} \\
& =
(-1)^{p + 1}
t(d([i_0, \ldots, i_{p + 1}])) (f_{i_0} \ldots f_{i_{p + 1}})^{-e} \\
& =
(-1)^{p + 1}
t(-\sum\nolimits_{j = 0}^{p + 1}
(-1)^j f_{i_j}^e [i_0, \ldots, \hat i_j, \ldots i_{p + 1}])
(f_{i_0} \ldots f_{i_{p + 1}})^{-e} \\
& =
-(-1)^{p + 1}
\sum\nolimits_{j = 0}^{p + 1}
(-1)^j t([i_0, \ldots, \hat i_j, \ldots i_{p + 1}])
(f_{i_0} \ldots \hat f_{i_j} \ldots f_{i_p})^{-e}
\end{align*}
二つの式は一致し、証明が完了する。
\end{proof}

\noindent
$A$-加群の有限複体 $I^\bullet$ と $A$-加群の複体 $M^\bullet$
が与えられたとする。このとき、対応する Hom 複体
$\Hom^\bullet(I^\bullet, M^\bullet)$ がある。これは次数 $n$ の項が
$$
\bigoplus\nolimits_{p + q = n} \Hom_A(I^{-q}, M^p)
$$
で与えられ、微分は More on Algebra, Section
\ref{more-algebra-section-hom-complexes} に記述されたものである。
複体 $I^\bullet$ の非零項は有限個しかないので、上に表示した直和は有限である。
複体の {\v C}ech 複体に付随する全複体をとる際の規約は、
Cohomology, Section
\ref{cohomology-section-cech-cohomology-of-complexes} のとおりである。

\begin{lemma}
\label{perfect-lemma-alternating-cech-complex-complex}
Situation \ref{perfect-situation-complex} において、$M^\bullet$ を
$A$-加群の複体とし、付随する $\mathcal{O}_X$-加群の複体を
$\mathcal{F}^\bullet$ と記す。このとき、複体の標準同型
$$
\colim_e \Hom^\bullet(I^\bullet(f_1^e, \ldots, f_r^e), M^\bullet)
\longrightarrow
\text{Tot}(\check{\mathcal{C}}_{alt}^\bullet(\mathcal{U}, \mathcal{F}^\bullet))
$$
があり、これは $M^\bullet$ に関して関手的である。
\end{lemma}

\begin{proof}
項
$F^{p, q} = \mathcal{C}_{alt}^p(\mathcal{U}, \mathcal{F}^q)$
をもつ二重複体 $F^{\bullet, \bullet}$ を考える。これは Cohomology,
Section \ref{cohomology-section-cech-cohomology-of-complexes}
で論じたものである。また、項
$G^{p, q} = \colim_e \Hom_A(I^{-p}(f_1^e, \ldots, f_r^e), M^q)$
をもち、関手性によって（符号を介さずに）与えられる微分を備えた
二重複体 $G^{\bullet, \bullet}$ を考える。Lemma
\ref{perfect-lemma-alternating-cech-complex} の証明で構成した写像
$\psi^{p, q} : G^{p, q} \to F^{p, q}$ は同型であり、
微分 $d_1$（補題による）および $d_2$（これは明らか）と両立する。
しかし、補題の矢印の左辺と右辺にある複体の微分 $d$ は符号が異なる。
すなわち、$g \in G^{p, q}$ に対して
$$
d(g) =  d_2(g) - (-1)^{p + q} d_1(g) 
$$
であり（More on Algebra, Section
\ref{more-algebra-section-hom-complexes} を参照せよ）、
$f \in F^{p, q}$ に対する微分は
$$
d(f) = d_1(f) + (-1)^p d_2(f)
$$
である。したがって $\psi^{p, q}$ に
$(-1)^{pq + p(p - 1)/2}$ を掛けることで符号を補正できる。
\end{proof}

\begin{lemma}
\label{perfect-lemma-alternating-cech-complex-complex-computes-cohomology}
Situation \ref{perfect-situation-complex} において、$\mathcal{F}^\bullet$
を準連接 $\mathcal{O}_X$-加群の複体とする。このとき、標準同型
$$
\text{Tot}(\check{\mathcal{C}}_{alt}^\bullet(\mathcal{U}, \mathcal{F}^\bullet))
\longrightarrow
R\Gamma(U, \mathcal{F}^\bullet)
$$
が $D(A)$ において存在し、$\mathcal{F}^\bullet$ に関して関手的である。
\end{lemma}

\begin{proof}
$\mathcal{B}$ を $U$ のアフィン開集合全体の集合とする。
アフィン・スキーム上の準連接加群の高次コホモロジー群は零なので
（Cohomology of Schemes, Lemma
\ref{coherent-lemma-quasi-coherent-affine-cohomology-zero}）、
これは Cohomology, Lemma
\ref{cohomology-lemma-alternating-cech-complex-complex-ss} の特別な場合である。
\end{proof}

\noindent
Situation \ref{perfect-situation-complex} において、Lemma
\ref{perfect-lemma-affine-compare-bounded} の圏同値を通じて
$A$-加群の複体 $I^\bullet(f_1^e, \ldots, f_r^e)$ に対応する
$D(\mathcal{O}_X)$ の対象を $I_e$ と記す。写像
(\ref{perfect-equation-system}) は系
$$
I_1 \leftarrow
I_2 \leftarrow
I_3 \leftarrow \ldots
$$
を与える。さらに、$U$ に制限すると同型になる写像
$I_e \to \mathcal{O}_X$ の両立する系がある。したがって、
$D(\mathcal{O}_X)$ の任意の対象 $E$ に対して標準写像
\begin{equation}
\label{perfect-equation-comparison}
\colim_e \Hom_{D(\mathcal{O}_X)}(I_e, E) \longrightarrow H^0(U, E)
\end{equation}
がある。これは写像 $I_e \to E$ をその $U$ への制限へ送り、
$\Hom_{D(\mathcal{O}_U)}(\mathcal{O}_U, E|_U) = H^0(U, E)$.
を用いて構成される。

\begin{proposition}
\label{perfect-proposition-represent-cohomology-class-on-open}
Situation \ref{perfect-situation-complex} において、
$D_\QCoh(\mathcal{O}_X)$ の任意の対象 $E$ に対し、
写像 (\ref{perfect-equation-comparison}) は同型である。
\end{proposition}

\begin{proof}
Lemma \ref{perfect-lemma-affine-compare-bounded} により、$E$ は
準連接層の複体 $\mathcal{F}^\bullet$ で与えられると仮定してよい。
対応する $A$-加群の複体を
$M^\bullet = \Gamma(X, \mathcal{F}^\bullet)$ とおく。Lemmas
\ref{perfect-lemma-alternating-cech-complex-complex} および
\ref{perfect-lemma-alternating-cech-complex-complex-computes-cohomology}
により、擬同型
$$
\colim_e \Hom^\bullet(I^\bullet(f_1^e, \ldots, f_r^e), M^\bullet)
\longrightarrow
\text{Tot}(\check{\mathcal{C}}_{alt}^\bullet(\mathcal{U}, \mathcal{F}^\bullet))
\longrightarrow
R\Gamma(U, \mathcal{F}^\bullet)
$$
を得る。More on Algebra, Lemma
\ref{more-algebra-lemma-RHom-out-of-projective} および
Equation (\ref{more-algebra-equation-h0-RHom}) により、複体
$\Hom^\bullet(I^\bullet(f_1^e, \ldots, f_r^e), M^\bullet)$
の $H^0$ をとることで $D(A)$ における $\Hom$ が計算される。
したがって両辺の $H^0$ をとると
$$
\colim_e \Hom_{D(A)}(I^\bullet(f_1^e, \ldots, f_r^e), M^\bullet)
=
H^0(U, E)
$$
Lemma \ref{perfect-lemma-affine-compare-bounded} により
$\Hom_{D(A)}(I^\bullet(f_1^e, \ldots, f_r^e), M^\bullet) =
\Hom_{D(\mathcal{O}_X)}(I_e, E)$ なので、補題が従う。
\end{proof}

\noindent
Situation \ref{perfect-situation-complex} において、Lemma
\ref{perfect-lemma-affine-compare-bounded} の圏同値を通じて
$A$-加群の複体 $K^\bullet(f_1^e, \ldots, f_r^e)$ に対応する
$D(\mathcal{O}_X)$ の対象を $K_e$ と記す。このとき完全三角形
$$
I_e \to \mathcal{O}_X \to K_e \to I_e[1]
$$
および系
$$
K_1 \leftarrow
K_2 \leftarrow
K_3 \leftarrow \ldots
$$
があり、後者は系 $(I_e)$ と両立する。さらに、写像の両立する系
$$
K_e \to H^0(K_e) = \mathcal{O}_X/(f_1^e, \ldots, f_r^e)
$$

\begin{lemma}
\label{perfect-lemma-represent-cohomology-class-on-closed}
Situation \ref{perfect-situation-complex} において、$E$ を
$D_\QCoh(\mathcal{O}_X)$ の対象とする。
$i = - r + 1, \ldots, 0$ に対して $H^i(E)|_U = 0$ と仮定する。
このとき、$s \in H^0(X, E)$ が与えられると、
$H^0(X, K_e) \to H^0(X, E)$ の像に $s$ が含まれるような
$e \geq 0$ と射 $K_e \to E$ が存在する。
\end{lemma}

\begin{proof}
$U$ は $r$ 個のアフィン開集合で被覆されるので、任意の準連接加群に対し
$j \geq r$ なら $H^j(U, \mathcal{F}) = 0$ である
（Cohomology of Schemes, Lemma
\ref{coherent-lemma-vanishing-nr-affines}）。Lemma
\ref{perfect-lemma-application-nice-K-injective} により、$H^0(U, E)$ は
$H^0(U, \tau_{\geq -r + 1}E)$ に等しい。スペクトル系列
$$
H^j(U, H^i(\tau_{\geq -r + 1}E)) \Rightarrow H^{i + j}(U, \tau_{\geq -N}E)
$$
がある。Derived Categories, Lemma
\ref{derived-lemma-two-ss-complex-functor} を参照せよ。
$E$ のコホモロジー層に仮定した消滅性により、$H^0(U, E) = 0$ である。
したがって $s|_U = 0$ である。$s$ を $D(\mathcal{O}_X)$ における射
$\mathcal{O}_X \to E$ とみなす。Proposition
\ref{perfect-proposition-represent-cohomology-class-on-open} により、
ある $e$ に対して合成 $I_e \to \mathcal{O}_X \to E$ は零である。
完全三角形 $I_e \to \mathcal{O}_X \to K_e \to I_e[1]$ により、
$s$ が合成 $\mathcal{O}_X \to K_e \to E$ となるような射
$K_e \to E$ を得る。
\end{proof}


\section{擬連接複体と完全複体}
\label{perfect-section-spell-out}

\noindent
本節では、Cohomology, Sections
\ref{cohomology-section-strictly-perfect},
\ref{cohomology-section-pseudo-coherent},
\ref{cohomology-section-tor}, および
\ref{cohomology-section-perfect}
で定義した一般概念と、More on Algebra, Sections
\ref{more-algebra-section-pseudo-coherent},
\ref{more-algebra-section-tor}, および
\ref{more-algebra-section-perfect}
における加群の複体についての対応する概念を結び付ける。

\begin{lemma}
\label{perfect-lemma-pseudo-coherent}
$X$ をスキームとする。$E$ が $D(\mathcal{O}_X)$ の
$m$-擬連接対象ならば、$i > m$ に対して $H^i(E)$ は
準連接 $\mathcal{O}_X$-加群であり、$H^m(E)$ は準連接
$\mathcal{O}_X$-加群の商である。$E$ が擬連接ならば、
$E$ は $D_\QCoh(\mathcal{O}_X)$ の対象である。
\end{lemma}

\begin{proof}
$X$ 上局所的に、$i > m$ に対して $H^i(E)$ が
$H^i(\mathcal{E}^\bullet)$ と同型であり、$H^m(E)$ が
$H^m(\mathcal{E}^\bullet)$ の商となるような厳密完全複体
$\mathcal{E}^\bullet$ が存在する。層 $\mathcal{E}^i$ は有限自由加群の
直和因子なので準連接である。補題が従う。
\end{proof}

\begin{lemma}
\label{perfect-lemma-pseudo-coherent-affine}
$X = \Spec(A)$ をアフィン・スキームとする。$M^\bullet$ を
$A$-加群の複体、$E$ を対応する $D(\mathcal{O}_X)$ の対象とする。
$E$ が $D(\mathcal{O}_X)$ の対象として $m$-擬連接
（resp.\ 擬連接）であるための必要十分条件は、
$M^\bullet$ が $A$-加群の複体として $m$-擬連接
（resp.\ 擬連接）であることである。
\end{lemma}

\begin{proof}
定義から、$M^\bullet$ が $m$-擬連接なら $E$ もそうであることは
直ちに分かる。逆を証明するため、$E$ は $m$-擬連接であると仮定する。
$X = \Spec(A)$ は準コンパクトであり、その位相は標準開集合からなる
基底をもつので、標準開被覆
$X = D(f_1) \cup \ldots \cup D(f_n)$、各 $D(f_i)$ 上の厳密完全複体
$\mathcal{E}_i^\bullet$、および $j > m$ では $H^j$ 上の同型を、
$H^m$ 上では全射を誘導する写像
$\alpha_i : \mathcal{E}_i^\bullet \to E|_{U_i}$ をとることができる。
Cohomology, Lemma \ref{cohomology-lemma-local-actual} により、
開被覆を細分した後、各 $i$ について $\alpha_i$ は複体の写像
$\mathcal{E}_i^\bullet \to \widetilde{M^\bullet}|_{U_i}$
で与えられると仮定してよい。Modules, Lemma
\ref{modules-lemma-direct-summand-of-locally-free-is-locally-free}
により、項 $\mathcal{E}_i^n$ は有限局所自由加群である。
したがって開被覆をさらに細分し、各 $\mathcal{E}_i^n$ は
有限自由 $\mathcal{O}_{U_i}$-加群であると仮定してよい。
定義から、$M^\bullet_{f_i}$ は $A_{f_i}$-加群の
$m$-擬連接複体である。More on Algebra, Lemma
\ref{more-algebra-lemma-glue-pseudo-coherent} を適用して結論を得る。

\medskip\noindent
「擬連接」の場合は次の事実から従う：定義により、$E$ が擬連接である
ための必要十分条件は、すべての $m$ に対して $E$ が $m$-擬連接で
あることであり、More on Algebra, Lemma
\ref{more-algebra-lemma-pseudo-coherent} により $M^\bullet$
についても同じことが成り立つ。
\end{proof}

\begin{lemma}
\label{perfect-lemma-identify-pseudo-coherent-noetherian}
$X$ をネーター・スキーム、$E$ を
$D_\QCoh(\mathcal{O}_X)$ の対象とする。$m \in \mathbf{Z}$
に対して次は同値である。
\begin{enumerate}
\item $i \geq m$ に対して $H^i(E)$ は連接であり、
$i \gg 0$ に対して零である。
\item $E$ は $m$-擬連接である。
\end{enumerate}
特に、$E$ が擬連接であるための必要十分条件は、$E$ が
$D^-_{\textit{Coh}}(\mathcal{O}_X)$ の対象であることである。
\end{lemma}

\begin{proof}
$X$ は準コンパクトなので、(1) と (2) のいずれにおいても対象 $E$
は上に有界である。したがって問題は $X$ 上局所的であり、$X$
はアフィンであると仮定してよい。あるネーター環 $A$ に対して
$X = \Spec(A)$ と書く。この場合、Lemma
\ref{perfect-lemma-affine-compare-bounded} により、$E$ は $A$-加群の複体
$M^\bullet$ に対応する。Lemma \ref{perfect-lemma-pseudo-coherent-affine}
により、$E$ が $m$-擬連接であるための必要十分条件は、
$M^\bullet$ が $m$-擬連接であることである。一方、$H^i(E)$ が
連接であるための必要十分条件は、$H^i(M^\bullet)$ が有限
$A$-加群であることである（Properties, Lemma
\ref{properties-lemma-finite-type-module}）。ゆえに結果は
More on Algebra, Lemma
\ref{more-algebra-lemma-Noetherian-pseudo-coherent} から従う。
\end{proof}

\begin{lemma}
\label{perfect-lemma-tor-dimension-affine}
$X = \Spec(A)$ をアフィン・スキームとする。$M^\bullet$ を
$A$-加群の複体、$E$ を対応する $D(\mathcal{O}_X)$ の対象とする。
このとき
\begin{enumerate}
\item $E$ が $[a, b]$ に Tor 振幅をもつための必要十分条件は、
$M^\bullet$ が $[a, b]$ に Tor 振幅をもつことである。
\item $E$ が有限 Tor 次元をもつための必要十分条件は、
$M^\bullet$ が有限 Tor 次元をもつことである。
\end{enumerate}
\end{lemma}

\begin{proof}
(2) は (1) から自明に従う。(1) の証明では、Lemma
\ref{perfect-lemma-affine-compare-bounded} の圏同値
$D(A) = D_\QCoh(X)$ を断りなく用いる。
$M^\bullet$ が $[a, b]$ に Tor 振幅をもつと仮定する。このとき
$K^\bullet$ は $D(A)$ において、$i \not \in [a, b]$ では
$K^i = 0$ となる平坦 $A$-加群の複体 $K^\bullet$ と同型である。
More on Algebra, Lemma \ref{more-algebra-lemma-tor-amplitude}
を参照せよ。このとき $E$ は $\widetilde{K^\bullet}$ と同型である。
各 $\widetilde{K^i}$ は平坦 $\mathcal{O}_X$-加群なので、
Cohomology, Lemma \ref{cohomology-lemma-tor-amplitude} により
$E$ は $[a, b]$ に Tor 振幅をもつ。

\medskip\noindent
$E$ が $[a, b]$ に Tor 振幅をもつと仮定する。このとき $E$ は
有界であり、したがって $M^\bullet$ は $K^-(A)$ に属する。
ゆえに $M^\bullet$ を上に有界な $A$-加群の複体で置き換えてよい。
さらに射影分解を選び、$M^\bullet$ は自由 $A$-加群からなる
上に有界な複体であると仮定してよい。このとき任意の $A$-加群 $N$ に対し
$$
E \otimes_{\mathcal{O}_X}^\mathbf{L} \widetilde{N}
\cong
\widetilde{M^\bullet} \otimes_{\mathcal{O}_X}^\mathbf{L} \widetilde{N}
\cong
\widetilde{M^\bullet \otimes_A N}
$$
が $D(\mathcal{O}_X)$ において成り立つ。したがって左辺の
コホモロジー層の消滅性から、$M^\bullet$ は $[a, b]$
に Tor 振幅をもつことが従う。
\end{proof}

\begin{lemma}
\label{perfect-lemma-tor-dimension-rel-affine}
環準同型 $R \to A$ に対応するアフィン・スキームの射
$f : X \to S$ をとる。$M^\bullet$ を $A$-加群の複体、
$E$ を対応する $D(\mathcal{O}_X)$ の対象とする。このとき
\begin{enumerate}
\item $E$ が $D(f^{-1}\mathcal{O}_S)$ の対象として
$[a, b]$ に Tor 振幅をもつための必要十分条件は、
$M^\bullet$ が $D(R)$ の対象として $[a, b]$ に Tor 振幅を
もつことである。
\item $E$ が $D(f^{-1}\mathcal{O}_S)$ の対象として局所的に
有限 Tor 次元をもつための必要十分条件は、$M^\bullet$ が
$D(R)$ の対象として有限 Tor 次元をもつことである。
\end{enumerate}
\end{lemma}

\begin{proof}
$\mathfrak p \subset R$ の上にある素イデアル
$\mathfrak q \subset A$ を考える。対応する点を
$x \in X$ および $s = f(x) \in S$ とする。このとき
$(f^{-1}\mathcal{O}_S)_x = \mathcal{O}_{S, s} = R_\mathfrak p$
かつ $E_x = M^\bullet_\mathfrak q$ である。これを念頭に置けば、
同値は次のように分かる。

\medskip\noindent
$M^\bullet$ が $R$-加群の複体として $[a, b]$ に Tor 振幅を
もつなら、$A$ の任意の素イデアルにおける $M^\bullet$ の局所化に
ついても同じことが成り立つ。したがって Cohomology, Lemma
\ref{cohomology-lemma-tor-amplitude-stalk} により、$E$ は
$f^{-1}\mathcal{O}_S$-加群の層の複体として $[a, b]$
に Tor 振幅をもつ。逆に、$E$ が $D(f^{-1}\mathcal{O}_S)$
の対象として $[a, b]$ に Tor 振幅をもつと仮定する。
先ほど引用した補題により、$\mathfrak p \subset R$ の上にある
すべての素イデアル $\mathfrak q \subset A$ に対して、
$M^\bullet_\mathfrak q$ は $R_\mathfrak p$-加群の複体として
$[a, b]$ に Tor 振幅をもつ。More on Algebra, Lemma
\ref{more-algebra-lemma-tor-amplitude-localization} により、
$M^\bullet$ は $R$-加群の複体として $[a, b]$ に Tor 振幅をもつ。
これで (1) の証明が完了する。

\medskip\noindent
$X$ は準コンパクトなので、$E$ が $f^{-1}\mathcal{O}_S$-加群の
複体として局所的に有限 Tor 次元をもつなら、実際にはある $a, b$
に対して $E$ は $f^{-1}\mathcal{O}_S$-加群の複体として
$[a, b]$ に Tor 振幅をもつ。したがって (2) は (1) から従う。
\end{proof}

\begin{lemma}
\label{perfect-lemma-tor-qc-qs}
$X$ を準分離スキーム、$E$ を $D_\QCoh(\mathcal{O}_X)$ の対象とし、
$a \leq b$ とする。次は同値である。
\begin{enumerate}
\item $E$ は $[a, b]$ に Tor 振幅をもつ。
\item $\QCoh(\mathcal{O}_X)$ のすべての $\mathcal{F}$ に対して、
$i \not \in [a, b]$ なら
$H^i(E \otimes_{\mathcal{O}_X}^\mathbf{L} \mathcal{F}) = 0$ である。
\end{enumerate}
\end{lemma}

\begin{proof}
(1) が (2) を含意することは明らかである。(2) を仮定する。
$U \subset X$ をアフィン開集合とする。$X$ は準分離なので、射
$j : U \to X$ は準コンパクトかつ分離である。したがって
$j_*$ は準連接加群を準連接加群へ移す
（Schemes, Lemma \ref{schemes-lemma-push-forward-quasi-coherent}）。
ゆえに関手
$\QCoh(\mathcal{O}_X) \to \QCoh(\mathcal{O}_U)$
は本質的全射である。したがって条件 (2) から、すべての準連接
$\mathcal{O}_U$-加群 $\mathcal{G}$ に対して、
$i \not \in [a, b]$ なら
$H^i(E|_U \otimes_{\mathcal{O}_U}^\mathbf{L} \mathcal{G})$
が消滅することが従う。$U = \Spec(A)$ と書き、Lemma
\ref{perfect-lemma-affine-compare-bounded} により $E|_U$ に対応する
$A$-加群の複体を $M^\bullet$ とする。いま示したことから、
$M^\bullet \otimes_A^\mathbf{L} N$ のコホモロジー群は
$[a, b]$ の外で消滅する。言い換えると、$M^\bullet$ は
$[a, b]$ に Tor 振幅をもつ。Lemma
\ref{perfect-lemma-tor-dimension-affine} により、$E|_U$ は
$[a, b]$ に Tor 振幅をもつ。これで補題が証明された。
\end{proof}

\begin{lemma}
\label{perfect-lemma-perfect-affine}
$X = \Spec(A)$ をアフィン・スキームとする。$M^\bullet$ を
$A$-加群の複体、$E$ を対応する $D(\mathcal{O}_X)$ の対象とする。
$E$ が $D(\mathcal{O}_X)$ の完全対象であるための必要十分条件は、
$M^\bullet$ が $D(A)$ の対象として完全であることである。
\end{lemma}

\begin{proof}
これは Lemmas \ref{perfect-lemma-pseudo-coherent-affine} および
\ref{perfect-lemma-tor-dimension-affine},
Cohomology, Lemma \ref{cohomology-lemma-perfect}、ならびに
More on Algebra, Lemma \ref{more-algebra-lemma-perfect}
の論理的帰結である。
\end{proof}

\noindent
スキーム上の擬連接複体についての記述から、ある種の内部 Hom
が準連接であることを証明できる。

\begin{lemma}
\label{perfect-lemma-quasi-coherence-internal-hom}
$X$ をスキームとする。
\begin{enumerate}
\item $L$ が $D^+_\QCoh(\mathcal{O}_X)$ に属し、
$D(\mathcal{O}_X)$ の $K$ が擬連接ならば、
$R\SheafHom(K, L)$ は $D_\QCoh(\mathcal{O}_X)$ に属し、
局所的に下に有界である。
\item $L$ が $D_\QCoh(\mathcal{O}_X)$ に属し、
$D(\mathcal{O}_X)$ の $K$ が完全ならば、
$R\SheafHom(K, L)$ は $D_\QCoh(\mathcal{O}_X)$ に属する。
\item $X = \Spec(A)$ がアフィンで $K, L \in D(A)$ ならば、
$$
R\SheafHom(\widetilde{K}, \widetilde{L}) = \widetilde{R\Hom_A(K, L)}
$$
次の二つの場合に
\begin{enumerate}
\item $K$ が擬連接で $L$ が下に有界である。
\item $K$ が完全で $L$ は任意である。
\end{enumerate}
成り立つ。
\item $X = \Spec(A)$ かつ $K, L$ が $D(A)$ に属するなら、
$R\SheafHom(\widetilde{K}, \widetilde{L})$ の第 $n$ コホモロジー層は、
前層
$$
X \supset D(f) \longmapsto \Ext^n_{A_f}(K \otimes_A A_f, L \otimes_A A_f)
$$
に付随する層である。ここで $f \in A$ である。
\end{enumerate}
\end{lemma}

\begin{proof}
導来圏における $\mathcal{O}_X$ の内部 Hom の構成は局所化と可換である
（Cohomology, Section \ref{cohomology-section-internal-hom}
を参照せよ）。したがって (1) と (2) を証明する際、$X$
をアフィン開集合で置き換えてよい。Lemmas
\ref{perfect-lemma-affine-compare-bounded},
\ref{perfect-lemma-pseudo-coherent-affine}, および
\ref{perfect-lemma-perfect-affine}
により、(1) と (2) を証明するには (3) を証明すれば十分である。

\medskip\noindent
Cohomology, Lemma
\ref{cohomology-lemma-Rhom-complex-of-direct-summands-finite-free}
における内部 Hom の計算で、$K$ を有限射影 $A$-加群からなる
上に有界な（resp.\ 有限）複体で表し、$L$ を $A$-加群からなる
下に有界な（resp.\ 任意の）複体で表せば、(3) が従う。

\medskip\noindent
(4) を証明するため、任意の環付き空間上で
$R\SheafHom(A, B)$ の第 $n$ コホモロジー層は、前層
$$
U \mapsto \Hom_{D(U)}(A|_U, B|_U[n]) =
\Ext^n_{D(\mathcal{O}_U)}(A|_U, B|_U)
$$
に付随する層であることを思い出そう。Cohomology, Section
\ref{cohomology-section-internal-hom} を参照せよ。一方、
$\widetilde{K}$ の主開集合 $D(f)$ への制限は
$K \otimes_A A_f$ の像であり、$L$ についても同様である。
したがって (4) は Lemma \ref{perfect-lemma-affine-compare-bounded}
の圏同値から従う。
\end{proof}

\begin{lemma}
\label{perfect-lemma-internal-hom-evaluate-tensor-isomorphism}
$X$ をスキームとし、$K, L, M$ を
$D_\QCoh(\mathcal{O}_X)$ の対象とする。Cohomology, Lemma
\ref{cohomology-lemma-internal-hom-diagonal-better} の写像
$$
K \otimes_{\mathcal{O}_X}^\mathbf{L} R\SheafHom(M, L)
\longrightarrow
R\SheafHom(M, K \otimes_{\mathcal{O}_X}^\mathbf{L} L)
$$
は次の場合に同型である。
\begin{enumerate}
\item $M$ は完全である。
\item $K$ は完全である。
\item $M$ は擬連接であり、$L \in D^+(\mathcal{O}_X)$ であり、
$K$ は有限 Tor 次元をもつ。
\end{enumerate}
\end{lemma}

\begin{proof}
Lemma \ref{perfect-lemma-quasi-coherence-internal-hom}
により、場合 (1) と (3) はアフィンの場合に帰着され、これは
More on Algebra, Lemma
\ref{more-algebra-lemma-internal-hom-evaluate-tensor-isomorphism}
で扱われている。（代数へ翻訳するには、Lemmas
\ref{perfect-lemma-pseudo-coherent-affine},
\ref{perfect-lemma-perfect-affine}, および \ref{perfect-lemma-tor-dimension-affine}
も用いる必要がある。）
$K$ が完全である以外に仮定がない場合、矢印のいずれの辺も
$D_\QCoh(\mathcal{O}_X)$ に属するとは限らない。しかし局所的に作業し、
$K$ が有限自由加群からなる有限複体である場合に帰着でき、
その場合には明らかなので、結果は依然として成り立つ。
\end{proof}





\section{連接加群の導来圏}
\label{perfect-section-derived-coherent}

\noindent
$X$ を局所ネーター・スキームとする。この場合、連接
$\mathcal{O}_X$-加群の圏
$\textit{Coh}(\mathcal{O}_X) \subset \textit{Mod}(\mathcal{O}_X)$
は弱 Serre 部分圏である。Homology, Section
\ref{homology-section-serre-subcategories} および
Cohomology of Schemes, Lemma
\ref{coherent-lemma-coherent-abelian-Noetherian} を参照せよ。
コホモロジー層が連接である複体の部分圏を
$$
D_{\textit{Coh}}(\mathcal{O}_X) \subset D(\mathcal{O}_X)
$$
と記す。Derived Categories, Section
\ref{derived-section-triangulated-sub} を参照せよ。
こうして標準関手
\begin{equation}
\label{perfect-equation-compare-coherent}
D(\textit{Coh}(\mathcal{O}_X))
\longrightarrow
D_{\textit{Coh}}(\mathcal{O}_X)
\end{equation}
を得る。Derived Categories, Equation
(\ref{derived-equation-compare}) を参照せよ。

\begin{lemma}
\label{perfect-lemma-coh-to-qcoh}
$X$ をネーター・スキームとする。このとき関手
$$
D^-(\textit{Coh}(\mathcal{O}_X))
\longrightarrow
D^-_{\textit{Coh}(\mathcal{O}_X)}(\QCoh(\mathcal{O}_X))
$$
は圏同値である。
\end{lemma}

\begin{proof}
$\textit{Coh}(\mathcal{O}_X) \subset \QCoh(\mathcal{O}_X)$
は Serre 部分圏であることに注意する。Homology, Definition
\ref{homology-definition-serre-subcategory}、Lemma
\ref{homology-lemma-characterize-serre-subcategory}、および
Cohomology of Schemes, Lemmas
\ref{coherent-lemma-coherent-abelian-Noetherian}、
\ref{coherent-lemma-coherent-Noetherian-quasi-coherent-sub-quotient}
を参照せよ。一方、$\mathcal{G} \to \mathcal{F}$ が準連接
$\mathcal{O}_X$-加群から連接 $\mathcal{O}_X$-加群への全射なら、
$\mathcal{F}$ へ全射する連接部分加群
$\mathcal{G}' \subset \mathcal{G}$ が存在する。実際、Properties,
Lemma \ref{properties-lemma-quasi-coherent-colimit-finite-type}
により、$\mathcal{G}$ をその連接部分加群のフィルター付き合併として
書けるので、それらの一つで十分である。したがって補題は
Derived Categories, Lemma
\ref{derived-lemma-fully-faithful-embedding} から従う。
\end{proof}

\begin{proposition}
\label{perfect-proposition-DCoh}
$X$ をネーター・スキームとする。このとき関手
$$
D^-(\textit{Coh}(\mathcal{O}_X))
\longrightarrow
D^-_{\textit{Coh}}(\mathcal{O}_X)
\quad\text{および}\quad
D^b(\textit{Coh}(\mathcal{O}_X))
\longrightarrow
D^b_{\textit{Coh}}(\mathcal{O}_X)
$$
は圏同値である。
\end{proposition}

\begin{proof}
可換図式
$$
\xymatrix{
D^-(\textit{Coh}(\mathcal{O}_X)) \ar[r] \ar[d] &
D^-_{\textit{Coh}}(\mathcal{O}_X) \ar[d] \\
D^-(\QCoh(\mathcal{O}_X)) \ar[r] &
D^-_\QCoh(\mathcal{O}_X)
}
$$
を考える。Lemma \ref{perfect-lemma-coh-to-qcoh} により左の縦矢印は充満忠実である。
Proposition \ref{perfect-proposition-Noetherian} により下の矢印は圏同値である。
構成により右の縦矢印は充満忠実である。したがって上の横矢印は
充満忠実である。$K$ が $D^-_{\textit{Coh}}(\mathcal{O}_X)$
の対象なら、Proposition \ref{perfect-proposition-Noetherian} によって
これに対応する $D^-(\QCoh(\mathcal{O}_X))$ の対象 $K'$ は
連接なコホモロジー層をもつ。したがって Lemma
\ref{perfect-lemma-coh-to-qcoh} により、$K'$ は左の縦矢印の本質像に属する。
ゆえに上の横矢印は本質的全射である。
これで上に有界な場合の証明が完了する。有界な場合は
上に有界な場合から直ちに従う。
\end{proof}

\begin{lemma}
\label{perfect-lemma-direct-image-coherent}
$S$ をネーター・スキームとし、$f : X \to S$ を局所有限型な
スキームの射とする。$E$ を $D^b_{\textit{Coh}}(\mathcal{O}_X)$
の対象とし、すべての $i$ について $H^i(E)$ の台は $S$
上固有であると仮定する。このとき $Rf_*E$ は
$D^b_{\textit{Coh}}(\mathcal{O}_S)$ の対象である。
\end{lemma}

\begin{proof}
スペクトル系列
$$
R^pf_*H^q(E) \Rightarrow R^{p + q}f_*E
$$
を考える。Derived Categories, Lemma
\ref{derived-lemma-two-ss-complex-functor} を参照せよ。
仮定と Cohomology of Schemes, Lemma
\ref{coherent-lemma-support-proper-over-base-pushforward} により、
層 $R^pf_*H^q(E)$ は連接である。したがって
$R^{p + q}f_*E$ は連接、すなわち
$Rf_*E \in D_{\textit{Coh}}(\mathcal{O}_S)$ である。
下からの有界性は自明である。上からの有界性は
Cohomology of Schemes, Lemma
\ref{coherent-lemma-quasi-coherence-higher-direct-images}
または Lemma \ref{perfect-lemma-quasi-coherence-direct-image} から従う。
\end{proof}

\begin{lemma}
\label{perfect-lemma-direct-image-coherent-bdd-below}
$S$ をネーター・スキームとし、$f : X \to S$ を局所有限型な
スキームの射とする。$E$ を $D^+_{\textit{Coh}}(\mathcal{O}_X)$
の対象とし、すべての $i$ について $H^i(E)$ の台は $S$
上固有であると仮定する。このとき $Rf_*E$ は
$D^+_{\textit{Coh}}(\mathcal{O}_S)$ の対象である。
\end{lemma}

\begin{proof}
証明は Lemma \ref{perfect-lemma-direct-image-coherent} の証明と同じである。
また、完全関手 $Rf_*$ が完全三角形
$\tau_{\leq a}E \to E \to \tau_{\geq a + 1}E \to \tau_{\leq a}E[1]$.
にどう作用するかを考えれば、Lemma
\ref{perfect-lemma-direct-image-coherent} から導くこともできる。
\end{proof}

\begin{lemma}
\label{perfect-lemma-coherent-internal-hom}
$X$ を局所ネーター・スキームとする。$L$ が
$D^+_{\textit{Coh}}(\mathcal{O}_X)$ に属し、$K$ が
$D^-_{\textit{Coh}}(\mathcal{O}_X)$ に属するなら、
$R\SheafHom(K, L)$ は $D^+_{\textit{Coh}}(\mathcal{O}_X)$ に属する。
\end{lemma}

\begin{proof}
$X$ がネーター環 $A$ のスペクトルである場合を証明すれば十分である。
Lemma \ref{perfect-lemma-identify-pseudo-coherent-noetherian} により、
$K$ は擬連接である。すると Lemma
\ref{perfect-lemma-quasi-coherence-internal-hom} を用いて、問題を次の代数的問題へ
翻訳できる：$L \in D^+_{\textit{Coh}}(A)$ かつ
$K$ が $D^-_{\textit{Coh}}(A)$ に属するなら、
$R\Hom_A(K, L)$ は $D^+_{\textit{Coh}}(A)$ に属する。
$L$ は下に有界で $K$ は上に有界なので、収束するスペクトル系列
$$
\Ext^p_A(K, H^q(L)) \Rightarrow \text{Ext}^{p + q}_A(K, L)
$$
があり、さらに収束するスペクトル系列
$$
\Ext^i_A(H^{-j}(K), H^q(L)) \Rightarrow \text{Ext}^{i + j}_A(K, H^q(L))
$$
がある。Injectives, Remarks
\ref{injectives-remark-spectral-sequences-ext} および
\ref{injectives-remark-spectral-sequences-ext-variant} を参照せよ。
Algebra, Lemma \ref{algebra-lemma-ext-noetherian} により、
有限 $A$-加群 $M$, $N$ に対して加群 $\Ext^p_A(M, N)$ は有限なので、
これで証明が完了する。
\end{proof}

\begin{lemma}
\label{perfect-lemma-perfect-on-noetherian}
$X$ をネーター・スキームとし、$D(\mathcal{O}_X)$ の $E$
は完全であるとする。このとき
\begin{enumerate}
\item $E$ は $D^b_{\textit{Coh}}(\mathcal{O}_X)$ に属する。
\item $L$ が $D_{\textit{Coh}}(\mathcal{O}_X)$ に属するなら、
$E \otimes_{\mathcal{O}_X}^\mathbf{L} L$ と
$R\SheafHom_{\mathcal{O}_X}(E, L)$ は
$D_{\textit{Coh}}(\mathcal{O}_X)$ に属する。
\item $L$ が $D^b_{\textit{Coh}}(\mathcal{O}_X)$ に属するなら、
$E \otimes_{\mathcal{O}_X}^\mathbf{L} L$ と
$R\SheafHom_{\mathcal{O}_X}(E, L)$ は
$D^b_{\textit{Coh}}(\mathcal{O}_X)$ に属する。
\item $L$ が $D^+_{\textit{Coh}}(\mathcal{O}_X)$ に属するなら、
$E \otimes_{\mathcal{O}_X}^\mathbf{L} L$ と
$R\SheafHom_{\mathcal{O}_X}(E, L)$ は
$D^+_{\textit{Coh}}(\mathcal{O}_X)$ に属する。
\item $L$ が $D^-_{\textit{Coh}}(\mathcal{O}_X)$ に属するなら、
$E \otimes_{\mathcal{O}_X}^\mathbf{L} L$ と
$R\SheafHom_{\mathcal{O}_X}(E, L)$ は
$D^-_{\textit{Coh}}(\mathcal{O}_X)$ に属する。
\end{enumerate}
\end{lemma}

\begin{proof}
$X$ は準コンパクトなので、各主張は $X$ の任意の開被覆の
各要素上で確認できる。したがって $E$ は有限自由加群からなる
有界複体 $\mathcal{E}^\bullet$ によって表されると仮定してよい。
Cohomology, Lemma
\ref{cohomology-lemma-perfect-on-locally-ringed} を参照せよ。
この場合、
$R\SheafHom_{\mathcal{O}_X}(E, L)$
と $E \otimes_{\mathcal{O}_X}^\mathbf{L} L$ はともに
$\mathcal{E}^\bullet$ を用いて複体のレベルで計算できるので、
各主張は明らかである。Cohomology, Lemmas
\ref{cohomology-lemma-Rhom-strictly-perfect} および
\ref{cohomology-lemma-bounded-flat-K-flat} を参照せよ。細部は省略する。
\end{proof}

\begin{lemma}
\label{perfect-lemma-ext-finite}
$A$ をネーター環、$X$ を $A$ 上の固有スキームとする。
$D^+_{\textit{Coh}}(\mathcal{O}_X)$ の $L$ と
$D^-_{\textit{Coh}}(\mathcal{O}_X)$ の $K$ に対して、
$A$-加群 $\Ext_{\mathcal{O}_X}^n(K, L)$ は有限である。
\end{lemma}

\begin{proof}
等式
$$
\Ext_{\mathcal{O}_X}^n(K, L) =
H^n(X, R\SheafHom_{\mathcal{O}_X}(K, L)) =
H^n(\Spec(A), Rf_*R\SheafHom_{\mathcal{O}_X}(K, L))
$$
を思い出そう。Cohomology, Lemma
\ref{cohomology-lemma-section-RHom-over-U} および Cohomology,
Section \ref{cohomology-section-Leray} を参照せよ。
したがって結果は Lemmas \ref{perfect-lemma-coherent-internal-hom}
および \ref{perfect-lemma-direct-image-coherent-bdd-below} から従う。
\end{proof}

\begin{lemma}
\label{perfect-lemma-perfect-on-regular}
$X$ を正則スキームとする。このとき
$D^b_{\textit{Coh}}(\mathcal{O}_X)$ のすべての対象は完全である。
$X$ が準コンパクト、すなわちネーターかつ正則なら、逆に
$D(\mathcal{O}_X)$ のすべての完全対象は
$D^b_{\textit{Coh}}(\mathcal{O}_X)$ に属する。
\end{lemma}

\begin{proof}
$K$ を $D^b_{\textit{Coh}}(\mathcal{O}_X)$ の対象とする。
$K$ が完全であることを確認するには $X$ 上アフィン局所的に
作業してよい（Cohomology, Section
\ref{cohomology-section-perfect} を参照せよ）。
このとき Lemma \ref{perfect-lemma-perfect-affine} と More on Algebra,
Lemma \ref{more-algebra-lemma-regular-perfect} により $K$ は完全である。
逆は Lemma \ref{perfect-lemma-perfect-on-noetherian} である。
\end{proof}




\section{複体に関する有限性条件の降下}
\label{perfect-section-descent-finiteness}

\noindent
本節は、スキームの導来圏の対象についての
Descent, Section \ref{descent-section-descent-finiteness}
の類似である。この種の結果のうち最も容易なものは、おそらく次である。

\begin{lemma}
\label{perfect-lemma-tor-amplitude-descends}
$f : X \to Y$ をスキーム（より一般には局所環付き空間）の
全射平坦射とする。
$E \in D(\mathcal{O}_Y)$ とし、$a, b \in \mathbf{Z}$ とする。
このとき、$E$ の Tor 振幅が $[a, b]$ に含まれることと、
$Lf^*E$ の Tor 振幅が $[a, b]$ に含まれることは同値である。
\end{lemma}

\begin{proof}
引き戻しは常に Tor 振幅を保つ。Cohomology, Lemma
\ref{cohomology-lemma-tor-amplitude-pullback} を参照せよ。
Tor 振幅が $[a, b]$ に含まれることは茎で確認できる。Cohomology,
Lemma \ref{cohomology-lemma-tor-amplitude-stalk} を参照せよ。
平坦な局所環準同型は忠実平坦である（Algebra, Lemma
\ref{algebra-lemma-local-flat-ff}）。したがって結果は More on Algebra,
Lemma \ref{more-algebra-lemma-flat-descent-tor-amplitude} から従う。
\end{proof}

\begin{lemma}
\label{perfect-lemma-pseudo-coherent-descends-fpqc}
$\{f_i : X_i \to X\}$ をスキームの fpqc 被覆とする。
$E \in D_\QCoh(\mathcal{O}_X)$ とし、$m \in \mathbf{Z}$ とする。
このとき、$E$ が $m$-擬連接であることと、各
$Lf_i^*E$ が $m$-擬連接であることは同値である。
\end{lemma}

\begin{proof}
引き戻しは常に $m$-擬連接性を保つ。Cohomology, Lemma
\ref{cohomology-lemma-pseudo-coherent-pullback} を参照せよ。
逆に、すべての $i$ について $Lf_i^*E$ が $m$-擬連接であると仮定する。
$U \subset X$ をアフィン開集合とする。$E|_U$ が $m$-擬連接であることを
示せば十分である。$\{f_i : X_i \to X\}$ は fpqc 被覆なので、
有限個のアフィン開集合 $V_j \subset X_{a(j)}$ で、
$f_{a(j)}(V_j) \subset U$ かつ $U = \bigcup f_{a(j)}(V_j)$
となるものをとれる。$V = \coprod V_i$ とおく。
したがって $X$ を $U$ で、$\{f_i : X_i \to X\}$ を
$\{V \to U\}$ で置き換え、$X$ はアフィンであり、被覆は
アフィン・スキームの単一の全射平坦射 $\{f : Y \to X\}$
によって与えられると仮定してよい。この場合、Lemmas
\ref{perfect-lemma-affine-compare-bounded} および
\ref{perfect-lemma-pseudo-coherent-affine} を介して、結果は More on Algebra,
Lemma \ref{more-algebra-lemma-flat-descent-pseudo-coherent} から従う。
\end{proof}

\begin{lemma}
\label{perfect-lemma-pseudo-coherent-descends-fppf}
$\{f_i : X_i \to X\}$ をスキームの fppf 被覆とする。
$E \in D(\mathcal{O}_X)$ とし、$m \in \mathbf{Z}$ とする。
このとき、$E$ が $m$-擬連接であることと、各
$Lf_i^*E$ が $m$-擬連接であることは同値である。
\end{lemma}

\begin{proof}
引き戻しは常に $m$-擬連接性を保つ。Cohomology, Lemma
\ref{cohomology-lemma-pseudo-coherent-pullback} を参照せよ。
逆に、すべての $i$ について $Lf_i^*E$ が $m$-擬連接であると仮定する。
$U \subset X$ をアフィン開集合とする。$E|_U$ が $m$-擬連接であることを
示せば十分である。$\{f_i : X_i \to X\}$ は fppf 被覆なので、
有限個のアフィン開集合 $V_j \subset X_{a(j)}$ で、
$f_{a(j)}(V_j) \subset U$ かつ $U = \bigcup f_{a(j)}(V_j)$
となるものをとれる。$V = \coprod V_i$ とおく。
したがって $X$ を $U$ で、$\{f_i : X_i \to X\}$ を
$\{V \to U\}$ で置き換え、$X$ はアフィンであり、被覆は
有限表示の単一の全射平坦射 $\{f : Y \to X\}$
によって与えられると仮定してよい。

\medskip\noindent
$f$ は平坦なので、導来関手 $Lf^*$ は単に $f^*$ で与えられ、
$f^*$ は完全である。したがって $H^i(Lf^*E) = f^*H^i(E)$ である。
$Lf^*E$ は $m$-擬連接なので、
$Lf^*E \in D^-(\mathcal{O}_Y)$ である。$f$ は全射かつ平坦なので、
$E \in D^-(\mathcal{O}_X)$ でもある。$H^i(E)$ が非零となる最大の
整数を $i \in \mathbf{Z}$ とする。$i < m$ なら証明は終わる。
そうでなければ、Cohomology, Lemma
\ref{cohomology-lemma-finite-cohomology} により、
$f^*H^i(E)$ は有限型 $\mathcal{O}_Y$-加群である。
すると Descent, Lemma \ref{descent-lemma-finite-type-descends-fppf}
により、$\mathcal{O}_X$-加群 $H^i(E)$ は有限型である。
したがって、$X$ を有限アフィン開被覆の各要素で置き換えることにより、
写像
$$
\alpha : \mathcal{O}_X^{\oplus n}[-i] \longrightarrow E
$$
で $H^i(\alpha)$ が全射となるものが存在すると仮定してよい。
$C$ を $D(\mathcal{O}_X)$ における $\alpha$ の錐とする。
$Y$ へ引き戻し、Cohomology, Lemma
\ref{cohomology-lemma-cone-pseudo-coherent} を用いると、
$Lf^*C$ は $m$-擬連接である。さらに、$j \geq i$ に対して
$H^j(C) = 0$ である。したがって $i$ に関する帰納法により、
$C$ は $m$-擬連接である。もう一度 Cohomology, Lemma
\ref{cohomology-lemma-cone-pseudo-coherent} を用いれば結論を得る。
\end{proof}

\begin{lemma}
\label{perfect-lemma-perfect-descends-fpqc}
$\{f_i : X_i \to X\}$ をスキームの fpqc 被覆とする。
$E \in D(\mathcal{O}_X)$ とする。このとき、$E$ が完全であることと、
各 $Lf_i^*E$ が完全であることは同値である。
\end{lemma}

\begin{proof}
引き戻しは常に完全複体を保つ。Cohomology, Lemma
\ref{cohomology-lemma-perfect-pullback} を参照せよ。
逆に、すべての $i$ について $Lf_i^*E$ が完全であると仮定する。
すると、各 $Lf_i^*E$ のコホモロジー層は準連接である。Lemma
\ref{perfect-lemma-pseudo-coherent} および Cohomology, Lemma
\ref{cohomology-lemma-perfect} を参照せよ。
射 $f_i$ は平坦なので、$H^p(Lf_i^*E) = f_i^*H^p(E)$ である。
したがって Descent, Proposition
\ref{descent-proposition-fpqc-descent-quasi-coherent} により、
$E$ のコホモロジー層は準連接である。以上から、この補題は
Cohomology, Lemma \ref{cohomology-lemma-perfect} と Lemmas
\ref{perfect-lemma-tor-amplitude-descends} および
\ref{perfect-lemma-pseudo-coherent-descends-fpqc} から形式的に従う。
\end{proof}

\begin{lemma}
\label{perfect-lemma-closed-push-pseudo-coherent}
$i : Z \to X$ を環付き空間の射で、$i$ が基礎位相空間の閉埋め込みであり、
$i_*\mathcal{O}_Z$ が $\mathcal{O}_X$-加群として擬連接であるものとする。
$E \in D(\mathcal{O}_Z)$ とする。このとき、$E$ が $m$-擬連接で
あることと、$Ri_*E$ が $m$-擬連接であることは同値である。
\end{lemma}

\begin{proof}
この証明では一貫して、$i_*$ が完全関手であり、したがって
$Ri_* = i_*$ であることを用いる。Modules, Lemma
\ref{modules-lemma-i-star-exact} を参照せよ。

\medskip\noindent
$E$ が $m$-擬連接であると仮定する。$x \in X$ とする。
その上で $i_*E$ が $m$-擬連接となる $x$ の近傍を見つける。
$x \not \in Z$ ならこれは明らかである。したがって $x \in Z$ と
仮定してよい。$x \in U \subset X$ を満たす開集合に対する
$U \cap Z$ が、$Z$ における $x$ の近傍の基本系をなすことを用いる。
$X$ を縮小することにより、$E$ は上に有界であると仮定してよい。
$H^p(E)$ が非零となる最大の整数 $p$ に関する帰納法を用いる。
$p < m$ なら示すことはない。$p \geq m$ なら、
Cohomology, Lemma \ref{cohomology-lemma-finite-cohomology} により、
$H^p(E)$ は有限型 $\mathcal{O}_Z$-加群である。
したがって $X$ を縮小した後、全射
$\mathcal{O}_Z^{\oplus n} \to H^p(E)$ を誘導する写像
$\mathcal{O}_Z^{\oplus n}[-p] \to E$ を選べる。完全三角形
$$
\mathcal{O}_Z^{\oplus n}[-p] \to E \to C \to \mathcal{O}_Z^{\oplus n}[-p + 1]
$$
を選ぶ。$j \geq p$ に対して $H^j(C) = 0$ であり、Cohomology,
Lemma \ref{cohomology-lemma-cone-pseudo-coherent} により、
$C$ は $m$-擬連接である。帰納法により、$i_*C$ は $X$ 上
$m$-擬連接である。また $i_*\mathcal{O}_Z$ も $X$ 上
$m$-擬連接なので、完全三角形
$$
i_*\mathcal{O}_Z^{\oplus n}[-p] \to i_*E \to i_*C \to
i_*\mathcal{O}_Z^{\oplus n}[-p + 1]
$$
と Cohomology, Lemma \ref{cohomology-lemma-cone-pseudo-coherent}
から、$i_*E$ は $m$-擬連接であると結論する。

\medskip\noindent
$i_*E$ が $m$-擬連接であると仮定する。$z \in Z$ とする。
その上で $E$ が $m$-擬連接となる $z$ の近傍を見つける。
$z \in U \subset X$ を満たす開集合に対する $U \cap Z$ が、
$Z$ における $z$ の近傍の基本系をなすことを用いる。
$X$ を縮小することにより、$i_*E$、したがって $E$ は上に有界で
あると仮定してよい。$H^p(E)$ が非零となる最大の整数 $p$ に
関する帰納法を用いる。$p < m$ なら示すことはない。$p \geq m$ なら、
$H^p(i_*E) = i_*H^p(E)$ は有限型 $\mathcal{O}_X$-加群である。
Cohomology, Lemma \ref{cohomology-lemma-finite-cohomology} を参照せよ。
$E$ を表す $\mathcal{O}_Z$-加群の複体
$\mathcal{E}^\bullet$ を選ぶ。$X$ を縮小した後、全射
$\mathcal{O}_X^{\oplus n} \to i_*H^p(\mathcal{E}^\bullet)$
を誘導する写像
$\alpha : \mathcal{O}_X^{\oplus n}[-p] \to i_*\mathcal{E}^\bullet$
を選べる。随伴により、全射
$\mathcal{O}_Z^{\oplus n} \to H^p(\mathcal{E}^\bullet)$
を誘導する写像
$\alpha : \mathcal{O}_Z^{\oplus n}[-p] \to \mathcal{E}^\bullet$
を得る。完全三角形
$$
\mathcal{O}_Z^{\oplus n}[-p] \to E \to C \to \mathcal{O}_Z^{\oplus n}[-p + 1]
$$
を選ぶ。$j \geq p$ に対して $H^j(C) = 0$ である。完全三角形
$$
i_*\mathcal{O}_Z^{\oplus n}[-p] \to i_*E \to i_*C \to
i_*\mathcal{O}_Z^{\oplus n}[-p + 1]
$$
と、$i_*\mathcal{O}_Z$ が擬連接であること、および
Cohomology, Lemma \ref{cohomology-lemma-cone-pseudo-coherent}
から、$i_*C$ は $m$-擬連接であると結論する。
帰納法により $C$ は $m$-擬連接である。もう一度 Cohomology,
Lemma \ref{cohomology-lemma-cone-pseudo-coherent} を用いると、
$E$ は $m$-擬連接であると結論できる。
\end{proof}

\begin{lemma}
\label{perfect-lemma-finite-push-pseudo-coherent}
$f : X \to Y$ をスキームの有限射で、
$f_*\mathcal{O}_X$ が $\mathcal{O}_Y$-加群として擬連接で
あるものとする\footnote{これは $f$ が擬連接であることを意味する。
More on Morphisms, Lemma
\ref{more-morphisms-lemma-finite-pseudo-coherent} を参照せよ。}。
$E \in D_\QCoh(\mathcal{O}_X)$ とする。このとき、$E$ が
$m$-擬連接であることと、$Rf_*E$ が $m$-擬連接であることは同値である。
\end{lemma}

\begin{proof}
これは More on Algebra, Lemma
\ref{more-algebra-lemma-finite-push-pseudo-coherent}
をスキームの言葉に翻訳したものである。この翻訳には Lemmas
\ref{perfect-lemma-affine-compare-bounded} および
\ref{perfect-lemma-pseudo-coherent-affine} を用いる。
\end{proof}


\section{複体の持ち上げ}
\label{perfect-section-lift}

\noindent
$U \subset X$ を環付き空間の開部分空間とし、包含射を
$j : U \to X$ と記す。$Rj_*$ は制限の右逆なので、関手
$D(\mathcal{O}_X) \to D(\mathcal{O}_U)$ は本質的に全射である。
本節では、このことを準連接コホモロジー層をもつ複体などへ拡張する。

\begin{lemma}
\label{perfect-lemma-lift-quasi-coherent}
$X$ をスキームとし、$j : U \to X$ を準コンパクト開埋め込みとする。
関手
$$
D_\QCoh(\mathcal{O}_X) \to D_\QCoh(\mathcal{O}_U)
\quad\text{かつ}\quad
D^+_\QCoh(\mathcal{O}_X) \to D^+_\QCoh(\mathcal{O}_U)
$$
は本質的に全射である。$X$ が準コンパクトなら、関手
$$
D^-_\QCoh(\mathcal{O}_X) \to D^-_\QCoh(\mathcal{O}_U)
\quad\text{かつ}\quad
D^b_\QCoh(\mathcal{O}_X) \to D^b_\QCoh(\mathcal{O}_U)
$$
も本質的に全射である。
\end{lemma}

\begin{proof}
Lemma \ref{perfect-lemma-quasi-coherence-direct-image} により、$Rj_*$ は
$D_\QCoh(\mathcal{O}_U)$ を $D_\QCoh(\mathcal{O}_X)$ に写すので、
補題の直前の議論を最初の場合に適用できる。
$Rj_*$ が $D^+_\QCoh(\mathcal{O}_U)$ を
$D^+_\QCoh(\mathcal{O}_X)$ に写すことは明らかであり、
これから下に有界な複体についての主張が従う。
最後に、$X$ が準コンパクトなら、Lemma
\ref{perfect-lemma-quasi-coherence-direct-image} により $Rj_*$ は
$D^-_\QCoh(\mathcal{O}_U)$ を $D^-_\QCoh(\mathcal{O}_X)$ に写す。
これら二つを合わせると最後の主張を得る。
\end{proof}

\begin{lemma}
\label{perfect-lemma-lift-coherent}
$X$ をネーター・スキームとし、$j : U \to X$ を開埋め込みとする。
関手
$D^b_{\textit{Coh}}(\mathcal{O}_X) \to D^b_{\textit{Coh}}(\mathcal{O}_U)$
は本質的に全射である。
\end{lemma}

\begin{proof}
$K$ を $D^b_{\textit{Coh}}(\mathcal{O}_U)$ の対象とする。
Proposition \ref{perfect-proposition-DCoh} により、$K$ は連接
$\mathcal{O}_U$-加群からなる有界複体 $\mathcal{F}^\bullet$
で表せる。ある $a \leq b$ に対して、$i \not \in [a, b]$ なら
$\mathcal{F}^i = 0$ であるとする。$j$ は準コンパクトかつ分離的なので、
有界複体 $j_*\mathcal{F}^\bullet$ の各項は $X$ 上の準連接加群である。
Schemes, Lemma \ref{schemes-lemma-push-forward-quasi-coherent} を参照せよ。
連接部分加群 $\mathcal{G}^i \subset j_*\mathcal{F}^i$ を
次のように帰納的に選ぶ。$i = a$ のとき、$U$ への制限が
$\mathcal{F}^a$ となる任意の連接部分加群
$\mathcal{G}^a \subset j_*\mathcal{F}^a$ を選ぶ。これは
Properties, Lemma \ref{properties-lemma-extend} により可能である。
$i > a$ のとき、まず $U$ への制限が $\mathcal{F}^i$ となる
任意の連接部分加群 $\mathcal{H}^i \subset j_*\mathcal{F}^i$ を選び、
次に
$\mathcal{G}^i = \Im(\mathcal{H}^i \oplus \mathcal{G}^{i - 1}
\to j_*\mathcal{F}^i)$ とおく。すると
$\mathcal{G}^\bullet \subset j_*\mathcal{F}^\bullet$ は連接
$\mathcal{O}_X$-加群からなる有界複体で、その $U$ への制限は
所望どおり $\mathcal{F}^\bullet$ である。
\end{proof}

\begin{lemma}
\label{perfect-lemma-lift-pseudo-coherent}
$X$ をアフィン・スキームとし、$U \subset X$ を準コンパクト開部分
スキームとする。$D(\mathcal{O}_U)$ の任意の擬連接対象 $E$ に対して、
有限自由 $\mathcal{O}_X$-加群からなる上に有界な複体で、その $U$ への
制限が $E$ と同型になるものが存在する。
\end{lemma}

\begin{proof}
Lemma \ref{perfect-lemma-pseudo-coherent} により、$E$ は
$D_\QCoh(\mathcal{O}_U)$ の対象である。Lemma
\ref{perfect-lemma-lift-quasi-coherent} により、
$D_\QCoh(\mathcal{O}_X)$ のある対象 $E'$ に対して
$E = E'|U$ であると仮定してよい。
$X = \Spec(A)$ と書く。Lemma \ref{perfect-lemma-affine-compare-bounded}
により、付随する $\mathcal{O}_X$-加群の複体が $E'$ を表すような
$A$-加群の複体 $M^\bullet$ をとれる。

\medskip\noindent
$U = D(f_1) \cup \ldots \cup D(f_r)$ となる
$f_1, \ldots, f_r \in A$ を選ぶ。Lemma
\ref{perfect-lemma-pseudo-coherent-affine} により、複体
$M^\bullet_{f_j}$ は $A_{f_j}$-加群の擬連接複体である。
$n$ を整数とする。複体の写像
$\alpha : F^\bullet \to M^\bullet$ で、$F^\bullet$ は上に有界、
$i < n$ に対して $F^i = 0$、各 $F^i$ は有限自由
$R$-加群であり、かつ
$$
H^i(\alpha_{f_j}) : H^i(F^\bullet_{f_j}) \to H^i(M^\bullet_{f_j})
$$
が $i > n$ では同型、$i = n$ では全射となるものが与えられたと
仮定する。図式で表すと
$$
\xymatrix{
& F^n \ar[r] \ar[d]^\alpha & F^{n + 1} \ar[d]^\alpha \ar[r] & \ldots \\
M^{n-1} \ar[r] & M^n \ar[r] & M^{n + 1} \ar[r] & \ldots
}
$$
である。各 $M^\bullet_{f_j}$ の十分大きな次数のコホモロジーは
消えるので、$n \gg 0$ に対してこのような写像をとれる。
$n$ に関する帰納法により、これを複体の写像
$F^\bullet \to M^\bullet$ で、すべての $i$ について
$H^i(\alpha_{f_j})$ が同型となるものへ拡張する。
$\widetilde{F^\bullet}$ をとれば補題が従う。

\medskip\noindent
帰納段階では、上の図式に $F^{n - 1}$ を加えて拡張する。
$C^\bullet$ を $\alpha$ の錐とする（Derived Categories, Definition
\ref{derived-definition-cone}）。コホモロジーの長完全列から、
$i \geq n$ に対して $H^i(C^\bullet_{f_j}) = 0$ である。
More on Algebra, Lemma
\ref{more-algebra-lemma-cone-pseudo-coherent} により、
$C^\bullet_{f_j}$ は $(n - 1)$-擬連接である。More on Algebra,
Lemma \ref{more-algebra-lemma-finite-cohomology} により、
$H^{n - 1}(C^\bullet_{f_j})$ は有限 $A_{f_j}$-加群である。
有限自由 $A$-加群 $F^{n - 1}$ と $A$-加群準同型
$\beta : F^{n - 1} \to C^{n - 1}$ を、合成
$F^{n - 1} \to C^{n - 1} \to C^n$ が零であり、かつ
$F^{n - 1}_{f_j}$ が $H^{n - 1}(C^\bullet_{f_j})$ へ全射となるように
選ぶ。（細部は省略する。ヒント：分母を払え。）
$C^{n - 1} = M^{n - 1} \oplus F^n$ なので、
$\beta = (\alpha^{n - 1}, -d^{n - 1})$ と書ける。
合成 $F^{n - 1} \to C^{n - 1} \to C^n$ が零であることから、
これらの写像は複体の射
$$
\xymatrix{
& F^{n - 1} \ar[d]^{\alpha^{n - 1}} \ar[r]_{d^{n - 1}} &
F^n \ar[r] \ar[d]^\alpha &
F^{n + 1} \ar[d]^\alpha \ar[r] & \ldots \\
\ldots \ar[r] &
M^{n - 1} \ar[r] & M^n \ar[r] & M^{n + 1} \ar[r] & \ldots
}
$$
をなす。さらに、これらの写像は完全三角形の射
$$
\xymatrix{
(F^n \to \ldots) \ar[r] \ar[d] &
(F^{n-1} \to \ldots) \ar[r] \ar[d] &
F^{n-1} \ar[r] \ar[d]_\beta &
(F^n \to \ldots)[1] \ar[d] \\
(F^n \to \ldots) \ar[r] &
M^\bullet \ar[r] &
C^\bullet \ar[r] &
(F^n \to \ldots)[1]
}
$$
を定める。したがって $\beta$ の選び方により、複体の写像
$(F^{-1} \to \ldots) \to M^\bullet$ は、$f_j$ で局所化した
コホモロジー上、次数 $\geq n$ では同型を、次数 $n - 1$ では
全射を誘導する。これで補題の証明が完了する。
\end{proof}

\noindent
次の二つの補題は、おそらく別の場所に置くべきである。

\begin{lemma}
\label{perfect-lemma-vanishing-ext}
$X$ を準コンパクトかつ準分離的なスキームとする。
$E \in D^b_\QCoh(\mathcal{O}_X)$ とする。
ある整数 $n_0 > 0$ が存在して、
$\Ext^n_{D(\mathcal{O}_X)}(\mathcal{E}, E) = 0$
が、任意の有限局所自由 $\mathcal{O}_X$-加群 $\mathcal{E}$ と
任意の $n \geq n_0$ に対して成り立つ。
\end{lemma}

\begin{proof}
$\Ext^n_{D(\mathcal{O}_X)}(\mathcal{E}, E) =
\Hom_{D(\mathcal{O}_X)}(\mathcal{E}, E[n])$ であることを思い出そう。
導来圏の射に対する Mayer--Vietoris がある。Cohomology, Lemma
\ref{cohomology-lemma-mayer-vietoris-hom} を参照せよ。
したがって $X = U \cup V$ で、ある上界 $n_0$ に対して補題の結果が
$E|_U$, $E|_V$, および $E|_{U \cap V}$ に成り立つなら、
$E$ に対して上界 $n_0 + 1$ で成り立つ。
ゆえに $X$ がアフィンの場合を証明すれば十分である。Cohomology of
Schemes, Lemma \ref{coherent-lemma-induction-principle} を参照せよ。

\medskip\noindent
$X = \Spec(A)$ がアフィンであると仮定する。付随する準連接加群の
複体が $E$ を表すような $A$-加群の複体 $M^\bullet$ を選ぶ。
Lemma \ref{perfect-lemma-affine-compare-bounded} を参照せよ。
ある $A$-加群 $P$ に対して $\mathcal{E} = \widetilde{P}$ と書く。
$\mathcal{E}$ は有限局所自由なので、$P$ は有限射影
$A$-加群である。ここで
\begin{align*}
\Hom_{D(\mathcal{O}_X)}(\mathcal{E}, E[n])
& = 
\Hom_{D(A)}(P, M^\bullet[n]) \\
& =
\Hom_{K(A)}(P, M^\bullet[n]) \\
& =
\Hom_A(P, H^n(M^\bullet))
\end{align*}
最初の等号は Lemma \ref{perfect-lemma-affine-compare-bounded} により、
二番目の等号は Derived Categories, Lemma
\ref{derived-lemma-morphisms-from-projective-complex} により、
最後の等号は $\Hom_A(P, -)$ が完全関手であることによる。
$E$、したがって $M^\bullet$ は有界なので、十分大きなすべての
$n$ に対して零を得る。
\end{proof}

\noindent
次の補題は強めることができる（非零コホモロジー層が固定された
範囲内にしかないすべての $L$ にわたり、消滅は一様である）。

\begin{lemma}
\label{perfect-lemma-ext-from-perfect-into-bounded-QCoh}
$X$ を準コンパクトかつ準分離的なスキームとする。
$K$ を $D(\mathcal{O}_X)$ の完全対象とする。このとき
\begin{enumerate}
\item ある整数 $a \leq b$ が存在し、$i \in [a, b]$ に対して
$H^i(L) = 0$ となる任意の $L \in D_\QCoh(\mathcal{O}_X)$ に対して
$\Hom_{D(\mathcal{O}_X)}(K, L) = 0$ である。
\item $L$ が有界なら、有限個を除くすべての $n$ に対して
$\Ext^n_{D(\mathcal{O}_X)}(K, L)$ は零である。
\end{enumerate}
\end{lemma}

\begin{proof}
(2) は (1) から従う。実際、
$\Ext^n_{D(\mathcal{O}_X)}(K, L) =
\Hom_{D(\mathcal{O}_X)}(K, L[n])$ である。(1) を証明する。
$K$ は完全なので
$$
\Hom_{D(\mathcal{O}_X)}(K, L) =
H^0(X, K^\vee \otimes_{\mathcal{O}_X}^\mathbf{L} L)
$$
である。ここで $K^\vee$ は $K$ の「双対」完全複体である。
Cohomology, Lemma \ref{cohomology-lemma-dual-perfect-complex} を参照せよ。
Lemmas \ref{perfect-lemma-quasi-coherence-tensor-product} および
\ref{perfect-lemma-pseudo-coherent}（完全複体が準連接コホモロジー層を
もつことを見るため）により、
$K^\vee \otimes_{\mathcal{O}_X}^\mathbf{L} L$ は $D_\QCoh(X)$ に属する。
$K^\vee$ の Tor 振幅が $[a, b]$ に含まれるとする。このとき
スペクトル系列
$$
E_1^{p, q} = H^p(K^\vee \otimes_{\mathcal{O}_X}^\mathbf{L} H^q(L))
\Rightarrow
H^{p + q}(K^\vee \otimes_{\mathcal{O}_X}^\mathbf{L} L)
$$
から、$q \in [j - b, j - a]$ に対して $H^q(L) = 0$ なら
$H^j(K^\vee \otimes_{\mathcal{O}_X}^\mathbf{L} L)$ は零である。
$N$ を Cohomology of Schemes, Lemma
\ref{coherent-lemma-vanishing-nr-affines-quasi-separated} の整数 $d$ とする。
このとき、コホモロジー層
$$
H^{-N}(K^\vee \otimes_{\mathcal{O}_X}^\mathbf{L} L),
\ H^{-N + 1}(K^\vee \otimes_{\mathcal{O}_X}^\mathbf{L} L),
\ \ldots,
\ H^0(K^\vee \otimes_{\mathcal{O}_X}^\mathbf{L} L)
$$
が零なら $H^0(X, K^\vee \otimes_{\mathcal{O}_X}^\mathbf{L} L)$ は
消える。実際、引用した補題と Lemma
\ref{perfect-lemma-application-nice-K-injective} により
$$
H^0(X, K^\vee \otimes_{\mathcal{O}_X}^\mathbf{L} L) =
H^0(X, \tau_{\geq -N}(K^\vee \otimes_{\mathcal{O}_X}^\mathbf{L} L))
$$
であり、コホモロジー層の消滅により、これは
$H^0(X, \tau_{\geq 1}(K^\vee \otimes_{\mathcal{O}_X}^\mathbf{L} L))$
に等しい。後者は Derived Categories, Lemma
\ref{derived-lemma-negative-vanishing} により零である。
したがって $i \in [-b - N, -a]$ に対して $H^i(L) = 0$ なら、
$\Hom_{D(\mathcal{O}_X)}(K, L)$ は零である。
\end{proof}

\begin{lemma}
\label{perfect-lemma-lift-perfect-complex-plus-locally-free}
$X$ をアフィン・スキームとし、$U \subset X$ を準コンパクト開集合とする。
$D(\mathcal{O}_U)$ の任意の完全対象 $E$ に対して、整数 $r$ と
$U$ 上の有限局所自由層 $\mathcal{F}$ で、
$\mathcal{F}[-r] \oplus E$ が $D(\mathcal{O}_X)$ の完全対象の
制限となるものが存在する。
\end{lemma}

\begin{proof}
$X = \Spec(A)$ とする。完全複体は擬連接であることを思い出そう。
Cohomology, Lemma \ref{cohomology-lemma-perfect} を参照せよ。
Lemma \ref{perfect-lemma-lift-pseudo-coherent} により、有限自由
$A$-加群からなる上に有界な複体 $\mathcal{F}^\bullet$ で、
$E$ が $D(\mathcal{O}_U)$ において $\mathcal{F}^\bullet|_U$
と同型になるものをとれる。Cohomology, Lemma
\ref{cohomology-lemma-perfect} と $U$ の準コンパクト性により、
$E$ は有限 Tor 次元をもつ。$E$ の Tor 振幅が $[a, b]$ に
含まれるとする。$r < a$ を選び、
$$
\mathcal{K} = \Ker(\mathcal{F}^{r} \to \mathcal{F}^{r + 1})
= \Im(\mathcal{F}^{r - 1} \to \mathcal{F}^r).
$$
とおく。$E$ の Tor 振幅は $[a, b]$ に含まれるので、
$\mathcal{F} = \mathcal{K}|_U$ は平坦である（Cohomology, Lemma
\ref{cohomology-lemma-last-one-flat}）。したがって $\mathcal{F}$ は
平坦かつ有限表示であり、ゆえに有限局所自由である（Properties,
Lemma \ref{properties-lemma-finite-locally-free}）。
このことから
$$
\mathcal{F} \to \mathcal{F}^r|_U \to \mathcal{F}^{r + 1}|_U \to \ldots
$$
は $E$ を表す $U$ 上の厳密完全複体である。一方、複体
$P = (\mathcal{F}^r \to \mathcal{F}^{r + 1} \to \ldots )$
は $X$ 上の完全複体である。愚直切断を用いると完全三角形
$$
P|_U \to E \to \mathcal{F}[-r - 1] \to (P|_U)[1]
$$
を得る。写像 $E \to \mathcal{F}[-r - 1]$ が $D(\mathcal{O}_U)$
において零なら、$P|_U = \mathcal{F}[-r - 2] \oplus E$ である。
Derived Categories, Lemma \ref{derived-lemma-split} を参照せよ。
たとえば Lemma \ref{perfect-lemma-ext-from-perfect-into-bounded-QCoh} により、
$r \ll 0$ ならこれは成り立つ。
\end{proof}

\begin{lemma}
\label{perfect-lemma-lift-map}
$X$ をアフィン・スキームとし、$U \subset X$ を準コンパクト開集合とする。
$E, E'$ を $D_\QCoh(\mathcal{O}_X)$ の対象とし、$E$ は完全であるとする。
任意の写像 $\alpha : E|_U \to E'|_U$ に対して、$X$ 上の複体の写像
$$
E \xleftarrow{\beta} E_1 \xrightarrow{\gamma} E'
$$
で、$E_1$ は完全、$\beta : E_1 \to E$ の $U$ への制限は同型であり、
$\alpha = \gamma|_U \circ \beta|_U^{-1}$ となるものが存在する。
さらに、$X$ 上のある完全複体 $I$ に対して
$E_1 = E \otimes_{\mathcal{O}_X}^\mathbf{L} I$ と仮定できる。
\end{lemma}

\begin{proof}
$X = \Spec(A)$ と書き、$U = D(f_1) \cup \ldots \cup D(f_r)$ と書く。
$E$ を表す有限射影 $A$-加群の有限複体 $M^\bullet$ を選ぶ
（Lemma \ref{perfect-lemma-perfect-affine}）。$E'$ を表す $A$-加群の複体
$(M')^\bullet$ を選ぶ（Lemma \ref{perfect-lemma-affine-compare-bounded}）。
この場合、複体 $H^\bullet = \Hom_A(M^\bullet, (M')^\bullet)$ は
$A$-加群の複体であり、付随する準連接 $\mathcal{O}_X$-加群の複体は
$R\SheafHom(E, E')$ を表す。Cohomology, Lemma
\ref{cohomology-lemma-Rhom-strictly-perfect} を参照せよ。
すると $\alpha$ は $H^0(U, R\SheafHom(E, E'))$ の元 $s$ を定める。
Cohomology, Lemma \ref{cohomology-lemma-section-RHom-over-U} を参照せよ。
Proposition \ref{perfect-proposition-represent-cohomology-class-on-open} により、
ある $e$ と写像
$$
\xi : I^\bullet(f_1^e, \ldots, f_r^e) \to \Hom_A(M^\bullet, (M')^\bullet)
$$
で $s$ に対応するものが存在する。次に、複体
$$
\text{Tot}(I^\bullet(f_1^e, \ldots, f_r^e) \otimes_A M^\bullet)
$$
に付随する準連接 $\mathcal{O}_X$-加群の複体に対応する対象を
$E_1$ とする。標準写像
$I^\bullet(f_1^e, \ldots, f_r^e) \to A$ を用いて $E_1 \to E$ を得、
$\xi$ と Cohomology, Lemma
\ref{cohomology-lemma-section-RHom-over-U} を用いて $E_1 \to E'$ を得る。
\end{proof}

\begin{lemma}
\label{perfect-lemma-lift-perfect-complex-plus-shift}
$X$ をアフィン・スキームとし、$U \subset X$ を準コンパクト開集合とする。
$D(\mathcal{O}_U)$ の任意の完全対象 $F$ に対して、
対象 $F \oplus F[1]$ は $D(\mathcal{O}_X)$ の完全対象の制限である。
\end{lemma}

\begin{proof}
Lemma \ref{perfect-lemma-lift-perfect-complex-plus-locally-free} により、
$D(\mathcal{O}_X)$ の完全対象 $E$ で、ある有限局所自由
$\mathcal{O}_U$-加群 $\mathcal{F}$ に対して
$E|_U = \mathcal{F}[r] \oplus F$ となるものをとれる。
Lemma \ref{perfect-lemma-lift-map} により、完全複体の射
$\alpha : E_1 \to E$ で、$(E_1)|_U \cong E|_U$ かつ
$\alpha|_U$ が写像
$$
\left(
\begin{matrix}
\text{id}_{\mathcal{F}[r]} & 0 \\
0 & 0
\end{matrix}
\right)
:
\mathcal{F}[r] \oplus F \to \mathcal{F}[r] \oplus F
$$
となるものをとれる。このとき $\alpha$ の錐が解である。
\end{proof}

\begin{lemma}
\label{perfect-lemma-perfect-into-support-on-T}
$X$ を準コンパクトかつ準分離的なスキームとする。
$f \in \Gamma(X, \mathcal{O}_X)$ とする。
$D_\QCoh(\mathcal{O}_X)$ の任意の射 $\alpha : E \to E'$ で、
\begin{enumerate}
\item $E$ は完全であり、
\item $E'$ は $T = V(f)$ に支持される
\end{enumerate}
ものに対して、$f^n \alpha  = 0$ となる $n \geq 0$ が存在する。
\end{lemma}

\begin{proof}
導来圏の射に対する Mayer--Vietoris がある。Cohomology, Lemma
\ref{cohomology-lemma-mayer-vietoris-hom} を参照せよ。
したがって $X = U \cup V$ で、補題の結果が
$f|_U$, $f|_V$, および $f|_{U \cap V}$ に対して成り立つなら、
$f$ に対しても成り立つ。ゆえに $X$ がアフィンの場合を証明すれば
十分である。Cohomology of Schemes, Lemma
\ref{coherent-lemma-induction-principle} を参照せよ。

\medskip\noindent
$X = \Spec(A)$ とする。このとき $f \in A$ である。以下では
Lemma \ref{perfect-lemma-affine-compare-bounded} の同値
$D(A) = D_\QCoh(X)$ を断りなく用いる。
$E$ を有限射影 $A$-加群の有限複体 $P^\bullet$ で表す。
これは Lemma \ref{perfect-lemma-perfect-affine} により可能である。
$P^t$ が非零となる最大の整数を $t$ とする。完全三角形
$$
P^t[-t] \to P^\bullet \to \sigma_{\leq t - 1}P^\bullet \to P^t[-t + 1]
$$
により、複体 $P^\bullet$ の長さに関する帰納法を用いて、
$P^\bullet$ が非零な項を一つしかもたない場合に帰着できる。
これとシフト関手により、$P^\bullet$ が次数 $0$ に置かれた
単一の有限射影 $A$-加群 $P$ からなる場合に帰着する。
$E'$ を $A$-加群の複体 $M^\bullet$ で表す。
すると $\alpha$ は写像 $P \to H^0(M^\bullet)$ に対応する。
仮定 (2) により加群 $H^0(M^\bullet)$ は $V(f)$ に支持されるので、
$H^0(M^\bullet)$ の各元は $f$ のある冪で零化される。
$P$ は有限 $A$-加群なので、所望どおり、ある $n$ に対して写像
$f^n\alpha : P \to H^0(M^\bullet)$ は零である。
\end{proof}

\begin{lemma}
\label{perfect-lemma-lift-perfect-complex-plus-shift-support}
$X$ をアフィン・スキームとし、$T \subset X$ を
$X \setminus T$ が準コンパクトとなる閉部分集合とする。
$U \subset X$ を準コンパクト開集合とする。$T \cap U$ に支持される
$D(\mathcal{O}_U)$ の任意の完全対象 $F$ に対して、
対象 $F \oplus F[1]$ は、$T$ に支持される
$D(\mathcal{O}_X)$ の完全対象 $E$ の制限である。
\end{lemma}

\begin{proof}
$T = V(g_1, \ldots, g_s)$ とする。$g_j$ をある冪で置き換えることにより、
$F$ 上の $g_j$ 倍が零であると仮定してよい。Lemma
\ref{perfect-lemma-perfect-into-support-on-T} を参照せよ。Lemma
\ref{perfect-lemma-lift-perfect-complex-plus-shift} のように $E$ を選ぶ。
$g_j : E \to E$ の $U$ への制限は零であることに注意する。
完全三角形
$$
E \xrightarrow{g_1} E \to C_1 \to E[1]
$$
を選ぶ。Derived Categories, Lemma \ref{derived-lemma-split} により、
対象 $C_1$ の $U$ への制限は
$F \oplus F[1] \oplus F[1] \oplus F[2]$ である。
さらに Derived Categories, Lemma
\ref{derived-lemma-third-map-square-zero} により、
$g_1 : C_1 \to C_1$ は二乗零である。実際、図式
$$
\xymatrix{
E \ar[r] \ar[d]_0 & C_1 \ar[d]_{g_1} \ar[r] & E[1] \ar[d]_0 \\
E \ar[r] & C_1 \ar[r] & E[1]
}
$$
は可換である。なぜなら合成 $E \xrightarrow{g_1} E \to C_1$ と
$C_1 \to E[1] \xrightarrow{g_1} E[1]$ は零だからである。
この操作を続け、$C_{i + 1}$ を写像 $g_i : C_i \to C_i$ の錐とおくと、
$T$ に支持される $X$ 上の完全複体 $C_s$ で、その $U$ への制限が
$$
F \oplus F[1]^{\oplus s} \oplus F[2]^{\oplus {s \choose 2}}
\oplus \ldots \oplus F[s]
$$
となるものを得る。Lemma \ref{perfect-lemma-lift-map} のように、
完全複体の射 $\beta : C' \to C_s$ と $\gamma : C' \to C_s$ を、
$\beta|_U$ が同型であり、$\gamma|_U \circ \beta|_U^{-1}$ が写像
$$
F \oplus F[1]^{\oplus s} \oplus F[2]^{\oplus {s \choose 2}}
\oplus \ldots \oplus F[s]
\to
F \oplus F[1]^{\oplus s} \oplus F[2]^{\oplus {s \choose 2}}
\oplus \ldots \oplus F[s]
$$
で、$F$ 上では零、それ以外のすべての直和因子上では恒等写像となる
ものとして選ぶ。Lemma \ref{perfect-lemma-lift-map} により、さらに
$X$ 上のある完全複体 $I$ に対して
$C' = C_s \otimes^\mathbf{L} I$ である。したがって
$g_j^2\text{id}_{C_s}$ が零であることから、$C'$ についても同じことが
従う。ゆえに $C'$ も $T$ に支持される。このとき
$\text{Cone}(\gamma)$ が解である。
\end{proof}

\noindent
次の補題の特別な場合は \cite{Neeman-Grothendieck} に見られる。

\begin{lemma}
\label{perfect-lemma-lift-map-from-perfect-complex-with-support}
$X$ を準コンパクトかつ準分離的なスキームとする。
$U \subset X$ を準コンパクト開集合とする。$T \subset X$ を、
$X \setminus T$ が $X$ においてレトロコンパクトとなる閉部分集合とする。
$E$ を $D_\QCoh(\mathcal{O}_X)$ の対象とする。
$\alpha : P \to E|_U$ を写像とし、$P$ は $T \cap U$ に支持される
$D(\mathcal{O}_U)$ の完全対象であるとする。このとき写像
$\beta : R \to E$ で、$R$ は $T$ に支持される
$D(\mathcal{O}_X)$ の完全対象であり、$P$ が $D(\mathcal{O}_U)$
において $R|_U$ の直和因子となり、この分解が $\alpha$ と
$\beta|_U$ に適合するものが存在する。
\end{lemma}

\begin{proof}
$X$ は準コンパクトなので、$X$ のあるアフィン開集合 $V_j$ と
ある整数 $m$ に対して
$X = U \cup V_1 \cup \ldots \cup V_m$ である。
$m$ に関する帰納法により、$m = 1$ と仮定してよい。
すなわち、$V$ がアフィンで $X = U \cup V$ と仮定してよい。
Lemma \ref{perfect-lemma-lift-perfect-complex-plus-shift-support} により、
$T \cap V$ に支持される $D(\mathcal{O}_V)$ の完全対象 $Q$ と同型
$Q|_{U \cap V} \to (P \oplus P[1])|_{U \cap V}$ を選べる。
Lemma \ref{perfect-lemma-lift-map} により、$Q$ を
$Q \otimes^\mathbf{L} I$ で置き換え（これは依然として
$T \cap V$ に支持される）、写像
$$
Q|_{U \cap V} \to (P \oplus P[1])|_{U \cap V}
\longrightarrow P|_{U \cap V}
\longrightarrow
E|_{U \cap V}
$$
が $Q \to E|_V$ に持ち上がると仮定できる。
Cohomology, Lemma \ref{cohomology-lemma-glue} により、
$D(\mathcal{O}_X)$ の射 $a : R \to E$ で、
$a|_U$ は $P \oplus P[1] \to E|_U$ と同型であり、
$a|_V$ は $Q \to E|_V$ と同型であるものを得る。
したがって所望どおり、$R$ は完全で $T$ に支持される。
\end{proof}

\begin{remark}
\label{perfect-remark-addendum}
Lemma \ref{perfect-lemma-lift-map-from-perfect-complex-with-support} の証明から、
$$
R|_U = P \oplus P^{\oplus n_1}[1] \oplus \ldots \oplus P^{\oplus n_m}[m]
$$
がある $m \geq 0$ と $n_j \geq 0$ に対して成り立つことが分かる。
したがって $R|_U$ の最高次数のコホモロジー層は $P$ のそれに等しい。
写像
$P^{\oplus n_1}[1] \oplus \ldots \oplus P^{\oplus n_m}[m] \to R|_U$
に対してこの構成を繰り返し、錐をとり、帰納法を用いると、
任意に与えた次数より上で $R|_U$ と $P$ のコホモロジー層を
等しくできる。
\end{remark}



\section{完全複体による近似}
\label{perfect-section-approximation}

\noindent
本節では、擬連接複体を完全複体によって「近似」できるという、
Neeman と Lipman による観察を論じる。

\begin{definition}
\label{perfect-definition-approximation-holds}
$X$ をスキームとする。次を満たす三つ組 $(T, E, m)$ を考える。
\begin{enumerate}
\item $T \subset X$ は閉部分集合である。
\item $E$ は $D_\QCoh(\mathcal{O}_X)$ の対象である。
\item $m \in \mathbf{Z}$ である。
\end{enumerate}
$T$ に支持される $D(\mathcal{O}_X)$ の完全対象 $P$ と写像
$\alpha : P \to E$ で、$i > m$ に対して同型
$H^i(P) \to H^i(E)$ を、また全射 $H^m(P) \to H^m(E)$ を
誘導するものが存在するとき、{\it 三つ組 $(T, E, m)$ に対して
近似が成立する}という。
\end{definition}

\noindent
近似がすべての三つ組に対して成立することはない。実際、
三つ組 $(T, E, m)$ に対して近似が成立するなら、明らかに
\begin{enumerate}
\item $E$ は $m$-擬連接である。Cohomology, Definition
\ref{cohomology-definition-pseudo-coherent} を参照せよ。
\item $i \geq m$ に対してコホモロジー層 $H^i(E)$ は $T$ に支持される。
\end{enumerate}
さらに、完全複体の「支持」は、その補集合が $X$ において
レトロコンパクトとなる閉部分スキームである（細部は省略する）。
したがって $T$ に関するこの仮定なしに近似が成立するとは期待できない。
これは次の定義の条件を部分的に説明する。

\begin{definition}
\label{perfect-definition-approximation}
$X$ をスキームとする。$X \setminus T$ が $X$ において
レトロコンパクトとなる任意の閉部分集合 $T \subset X$ に対して
ある整数 $r$ が存在し、Definition
\ref{perfect-definition-approximation-holds} の任意の三つ組 $(T, E, m)$ で
\begin{enumerate}
\item $E$ は $(m - r)$-擬連接であり、
\item $i \geq m - r$ に対して $H^i(E)$ は $T$ に支持される
\end{enumerate}
ものについて近似が成立するとき、$X$ 上で{\it 完全複体による近似が
成立する}という。
\end{definition}

\noindent
準コンパクトかつ準分離的なスキームでは完全複体による近似が
成立することを証明する。第二の条件は我々の証明法に必要と思われる。
以下の議論を注意深く解析すれば、第一の条件を「$E$ は
$m$-擬連接である」へ弱められる可能性がある。

\begin{lemma}
\label{perfect-lemma-open}
$X$ をスキームとし、$U \subset X$ を開部分スキームとする。
$(T, E, m)$ を Definition \ref{perfect-definition-approximation-holds}
の三つ組とする。次を仮定する。
\begin{enumerate}
\item $T \subset U$ である。
\item $(T, E|_U, m)$ に対して近似が成立する。
\item $i \geq m$ に対して層 $H^i(E)$ は $T$ に支持される。
\end{enumerate}
このとき $(T, E, m)$ に対して近似が成立する。
\end{lemma}

\begin{proof}
$j : U \to X$ を包含射とする。$P \to E|_U$ が $U$ 上で三つ組
$(T, E|_U, m)$ の近似なら、
$j_!P = Rj_*P \to j_!(E|_U) \to E$ は $X$ 上で
$(T, E, m)$ の近似である。Cohomology, Lemmas
\ref{cohomology-lemma-pushforward-restriction} および
\ref{cohomology-lemma-pushforward-perfect} を参照せよ。
\end{proof}

\begin{lemma}
\label{perfect-lemma-approximation-affine}
$X$ をアフィン・スキームとする。このとき、ある整数 $r \geq 0$
が存在して次を満たす、Definition \ref{perfect-definition-approximation-holds}
の任意の三つ組 $(T, E, m)$ に対して近似が成立する。
\begin{enumerate}
\item $E$ は $m$-擬連接である。
\item $i \geq m - r + 1$ に対して $H^i(E)$ は $T$ に支持される。
\item $X \setminus T$ は $r$ 個のアフィン開集合の和である。
\end{enumerate}
特に、アフィン・スキームでは完全複体による近似が成立する。
\end{lemma}

\begin{proof}
$X = \Spec(A)$ とし、$T = V(f_1, \ldots, f_r)$ と書く。
（$r = 0$、すなわち $T = X$ の場合は、Lemma
\ref{perfect-lemma-pseudo-coherent-affine} と定義から直ちに従う。）
$(T, E, m)$ を補題の三つ組とする。$H^t(E)$ が非零となる
最大の整数を $t$ とする。$t$ に関する帰納法を用いる。
基底の場合は $t < m$ であり、この場合、結果は自明である。
次に $t \geq m$ と仮定する。Cohomology, Lemma
\ref{cohomology-lemma-finite-cohomology} により、層 $H^t(E)$ は
有限型である。これは準連接なので、有限個の切断で生成される
（Properties, Lemma \ref{properties-lemma-finite-type-module}）。
任意の $s \in \Gamma(X, H^t(E)) = H^t(X, E)$ に対して
（Lemma \ref{perfect-lemma-affine-compare-bounded} の証明を参照せよ）、
ある $e > 0$ と射 $K_e[-t] \to E$ で、$s$ が
$H^0(K_e) = H^t(K_e[-t]) \to H^t(E)$ の像に含まれるものをとれる。
Lemma \ref{perfect-lemma-represent-cohomology-class-on-closed} を参照せよ。
これらの写像の有限直和をとると写像 $P \to E$ を得る。ここで
$P$ は $T$ に支持される完全複体であり、$i > t$ に対して
$H^i(P) = 0$ であり、$H^t(P) \to E$ は全射である。完全三角形
$$
P \to E \to E' \to P[1]
$$
を選ぶ。このとき $E'$ は $m$-擬連接であり
（Cohomology, Lemma \ref{cohomology-lemma-cone-pseudo-coherent}）、
$i \geq t$ に対して $H^i(E') = 0$ であり、
$i \geq m - r + 1$ に対して $H^i(E')$ は $T$ に支持される。
帰納法により、$(T, E', m)$ の近似 $P' \to E'$ を得る。
合成 $P' \to E' \to P[1]$ を完全三角形
$P \to P'' \to P' \to P[1]$ に組み込み、射 $P' \to E'$ と
$P[1] \to P[1]$ を完全三角形の射
$$
\xymatrix{
P \ar[r] \ar[d] & P'' \ar[d] \ar[r] & P' \ar[d] \ar[r] & P[1] \ar[d] \\
P \ar[r] &  E \ar[r] & E' \ar[r] & P[1]
}
$$
へ TR3 を用いて拡張する。このとき $P''$ は $T$ に支持される
完全複体である（Cohomology, Lemma
\ref{cohomology-lemma-two-out-of-three-perfect}）。
容易な図式追跡により、$P'' \to E$ は所望の近似である。
\end{proof}

\begin{lemma}
\label{perfect-lemma-induction-step}
$X$ をスキームとする。$X = U \cup V$ を開被覆とし、
$U$ は準コンパクト、$V$ はアフィン、$U \cap V$ は
準コンパクトであるとする。$U$ 上で完全複体による近似が成立するなら、
$X$ 上で近似が成立する。
\end{lemma}

\begin{proof}
$T \subset X$ を、$X \setminus T$ が $X$ においてレトロコンパクト
となる閉部分集合とする。$r_U$ を、組 $(U, T \cap U)$ に適合する
Definition \ref{perfect-definition-approximation} の整数とする。
$T' = T \setminus U$ とおく。$T' \subset V$ であり、また
$V \setminus T' = (X \setminus T) \cap U \cap V$ は $T$ に関する
仮定により準コンパクトであることに注意する。
$V \setminus T'$ を被覆するのに必要なアフィンの個数を $r'$ とする。
$r = \max(r_U, r')$ が組 $(X, T)$ に対して使えると主張する。

\medskip\noindent
これを見るため、$E$ は $(m - r)$-擬連接であり、
$i \geq m - r$ に対して $H^i(E)$ は $T$ に支持されるような
三つ組 $(T, E, m)$ を選ぶ。$H^t(E)|_U$ が非零となる最大の整数を
$t$ とする。（$U$ は準コンパクトで $E|_U$ は
$(m - r)$-擬連接なので、このような整数は存在する。）
$t$ に関する帰納法により $E$ を近似できることを証明する。

\medskip\noindent
基底の場合：$t \leq m - r'$ とする。これは $i \geq m - r'$
に対して $H^i(E)$ が $T'$ に支持されることを意味する。
したがって Lemma \ref{perfect-lemma-approximation-affine} により、
$V$ 上で $(T', E|_V, m)$ の近似 $P \to E|_V$ が存在する。
Lemma \ref{perfect-lemma-open} を適用すると、$(T', E, m)$ を近似できる。
この近似は $(T, E, m)$ の近似でもある。

\medskip\noindent
帰納段階。$(T \cap U, E|_U, m)$ の近似 $P \to E|_U$ を選ぶ。
これは特に全射 $H^t(P) \to H^t(E|_U)$ を与える。Lemma
\ref{perfect-lemma-lift-perfect-complex-plus-shift-support} により、
$T \cap V$ に支持される $D(\mathcal{O}_V)$ の完全対象 $Q$ と同型
$Q|_{U \cap V} \to (P \oplus P[1])|_{U \cap V}$ を選べる。
Lemma \ref{perfect-lemma-lift-map} により、$Q$ を
$Q \otimes^\mathbf{L} I$ で置き換え、写像
$$
Q|_{U \cap V} \to (P \oplus P[1])|_{U \cap V}
\longrightarrow P|_{U \cap V}
\longrightarrow
E|_{U \cap V}
$$
が $Q \to E|_V$ に持ち上がると仮定できる。Cohomology, Lemma
\ref{cohomology-lemma-glue} により、$D(\mathcal{O}_X)$ の射
$a : R \to E$ で、$a|_U$ は $P \oplus P[1] \to E|_U$ と同型であり、
$a|_V$ は $Q \to E|_V$ と同型であるものを得る。
したがって $R$ は完全で $T$ に支持され、写像
$H^t(R) \to H^t(E)$ は $U$ への制限上で全射である。完全三角形
$$
R \to E \to E' \to R[1]
$$
を選ぶ。このとき $E'$ は $(m - r)$-擬連接であり
（Cohomology, Lemma \ref{cohomology-lemma-cone-pseudo-coherent}）、
$i \geq t$ に対して $H^i(E')|_U = 0$ であり、
$i \geq m - r$ に対して $H^i(E')$ は $T$ に支持される。
帰納法により、$(T, E', m)$ の近似 $R' \to E'$ を得る。
合成 $R' \to E' \to R[1]$ を完全三角形
$R \to R'' \to R' \to R[1]$ に組み込み、射 $R' \to E'$ と
$R[1] \to R[1]$ を完全三角形の射
$$
\xymatrix{
R \ar[r] \ar[d] & R'' \ar[d] \ar[r] & R' \ar[d] \ar[r] & R[1] \ar[d] \\
R \ar[r] &  E \ar[r] & E' \ar[r] & R[1]
}
$$
へ TR3 を用いて拡張する。このとき $R''$ は $T$ に支持される
完全複体である（Cohomology, Lemma
\ref{cohomology-lemma-two-out-of-three-perfect}）。
容易な図式追跡により、$R'' \to E$ は所望の近似である。
\end{proof}

\begin{theorem}
\label{perfect-theorem-approximation}
$X$ を準コンパクトかつ準分離的なスキームとする。
このとき $X$ 上で完全複体による近似が成立する。
\end{theorem}

\begin{proof}
これは Cohomology of Schemes, Lemma
\ref{coherent-lemma-induction-principle} の帰納原理と Lemmas
\ref{perfect-lemma-induction-step} および \ref{perfect-lemma-approximation-affine} から従う。
\end{proof}






\section{導来圏の生成}
\label{perfect-section-generating}

\noindent
本節では、準コンパクトかつ準分離的なスキームの導来圏
$D_\QCoh(\mathcal{O}_X)$ が単一の完全対象で生成できることを証明する。
この結果の証明については、Bondal と van den Bergh の優れた論文
\cite{BvdB} をぜひ読んでほしい。

\begin{lemma}
\label{perfect-lemma-direct-summand-of-a-restriction}
$X$ を準コンパクトかつ準分離的なスキームとする。
$U$ を準コンパクト開部分スキームとする。
$P$ を $D(\mathcal{O}_U)$ の完全対象とする。
このとき $P$ は、$D(\mathcal{O}_X)$ の完全対象の制限の直和因子である。
\end{lemma}

\begin{proof}
Lemma \ref{perfect-lemma-lift-map-from-perfect-complex-with-support} の特別な場合である。
\end{proof}

\begin{lemma}
\label{perfect-lemma-orthogonal-koszul-complex}
\begin{reference}
\cite[Proposition 6.1]{Bokstedt-Neeman}
\end{reference}
Situation \ref{perfect-situation-complex} において、開埋め込みを
$j : U \to X$ と記し、$K$ を $A$ 上の $f_1, \ldots, f_r$ に関する
Koszul 複体に対応する $D(\mathcal{O}_X)$ の完全対象とする。
$E \in D_\QCoh(\mathcal{O}_X)$ に対して、次は同値である。
\begin{enumerate}
\item $E = Rj_*(E|_U)$ である。
\item すべての $n \in \mathbf{Z}$ に対して
$\Hom_{D(\mathcal{O}_X)}(K[n], E) = 0$ である。
\end{enumerate}
\end{lemma}

\begin{proof}
完全三角形 $E \to Rj_*(E|_U) \to N \to E[1]$ を選ぶ。
$$
\Hom_{D(\mathcal{O}_X)}(K[n], Rj_*(E|_U)) =
\Hom_{D(\mathcal{O}_U)}(K|_U[n], E) = 0
$$
であることに注意する。実際、$K|_U = 0$ なので、これはすべての
$n$ に対して成り立つ。したがって $N$ に対して結果を証明すれば
十分である。言い換えると、$E$ の $U$ への制限は零であると
仮定してよい。また、Koszul 複体の完全三角形
$$
K^\bullet(f_1^{e_1}, \ldots, f_i^{e'_i}, \ldots, f_r^{e_r}) \to
K^\bullet(f_1^{e_1}, \ldots, f_i^{e'_i + e''_i}, \ldots, f_r^{e_r}) \to
K^\bullet(f_1^{e_1}, \ldots, f_i^{e''_i}, \ldots, f_r^{e_r}) \to \ldots
$$
がある。More on Algebra, Lemma \ref{more-algebra-lemma-koszul-mult}
を参照せよ。したがって、すべての $n \in \mathbf{Z}$ に対して
$\Hom_{D(\mathcal{O}_X)}(K[n], E) = 0$ なら、Lemma
\ref{perfect-lemma-represent-cohomology-class-on-closed} のように $K$ を
$K_e$ で置き換えても同じことが成り立つ。ゆえに、本補題はその補題と、
$E$ が $A$-加群の複体 $R\Gamma(X, E)$ によって決定されることから
直ちに従う。Lemma \ref{perfect-lemma-affine-compare-bounded} を参照せよ。
\end{proof}

\begin{theorem}
\label{perfect-theorem-bondal-van-den-Bergh}
$X$ を準コンパクトかつ準分離的なスキームとする。圏
$D_\QCoh(\mathcal{O}_X)$ は単一の完全対象で生成できる。
より正確には、$D(\mathcal{O}_X)$ の完全対象 $P$ で、
$E \in D_\QCoh(\mathcal{O}_X)$ に対して次が同値となるものが存在する。
\begin{enumerate}
\item $E = 0$ である。
\item すべての $n \in \mathbf{Z}$ に対して
$\Hom_{D(\mathcal{O}_X)}(P[n], E) = 0$ である。
\end{enumerate}
\end{theorem}

\begin{proof}
Cohomology of Schemes, Lemma
\ref{coherent-lemma-induction-principle} の帰納原理を用いて証明する。

\medskip\noindent
$X$ がアフィンなら、$\mathcal{O}_X$ は完全生成対象である。
これは Lemma \ref{perfect-lemma-affine-compare-bounded} から従う。

\medskip\noindent
$X = U \cup V$ を開被覆とし、$U$ は準コンパクトで $U$ 上では
定理が成り立ち、$V$ はアフィン開集合であると仮定する。
$P$ を $D(\mathcal{O}_U)$ の完全対象で
$D_\QCoh(\mathcal{O}_U)$ の生成対象となるものとする。Lemma
\ref{perfect-lemma-direct-summand-of-a-restriction} により、
$D(\mathcal{O}_X)$ の完全対象 $Q$ で、その $U$ への制限が
$P$ を直和因子の一つとして含む直和となるものを選べる。
$V = \Spec(A)$ とし、$Z = X \setminus U$ とする。
これは $V$ の閉部分集合で、$V \setminus Z$ は準コンパクトである。
$Z = V(f_1, \ldots, f_r)$ となる
$f_1, \ldots, f_r \in A$ を選ぶ。$K \in D(\mathcal{O}_V)$ を、
$A$ 上の $f_1, \ldots, f_r$ に関する Koszul 複体に対応する
完全対象とする。$K$ は閉集合 $Z \subset V$ に支持されるので、
直像 $K' = R(V \to X)_*K$ は $D(\mathcal{O}_X)$ の完全対象であり、
その $V$ への制限は $K$ であることに注意する（Cohomology,
Lemma \ref{cohomology-lemma-pushforward-perfect}）。
$Q \oplus K'$ は $D_\QCoh(\mathcal{O}_X)$ の生成対象であると主張する。

\medskip\noindent
$E$ を $D_\QCoh(\mathcal{O}_X)$ の対象とし、
$Q \oplus K'$ のどのシフトから $E$ へも非自明な写像がないとする。
Cohomology, Lemma \ref{cohomology-lemma-pushforward-restriction} により
$K' =  R(V \to X)_! K$ であり、したがって
$$
\Hom_{D(\mathcal{O}_X)}(K'[n], E) = \Hom_{D(\mathcal{O}_V)}(K[n], E|_V)
$$
である。Lemma \ref{perfect-lemma-orthogonal-koszul-complex} により、
これらの群が消えることから $E|_V$ は
$R(U \cap V \to V)_*E|_{U \cap V}$ と同型である。
これは $E = R(U \to X)_*E|_U$ を含意する（細部を一つ省略した）。
この場合
$$
\Hom_{D(\mathcal{O}_X)}(Q[n], E) = \Hom_{D(\mathcal{O}_U)}(Q|_U[n], E|_U)
$$
であり、これは $\Hom_{D(\mathcal{O}_U)}(P[n], E|_U)$ を
直和因子として含む。したがって $P$ の選び方により、
これらの群の消滅は $E|_U$ が零であることを含意する。
ゆえに $E$ は零である。
\end{proof}

\noindent
次の結果は Theorem \ref{perfect-theorem-bondal-van-den-Bergh} の強化であり、
まったく同じ方法で証明される。
スキーム $X$ の閉部分集合 $T$ に対して、
$D_T(\mathcal{O}_X)$ は $T$ に支持される対象からなる
$D(\mathcal{O}_X)$ の厳密充満・飽和・三角部分圏を表すことを思い出そう
（Definition \ref{perfect-definition-supported-on}）。
同様に、$D_{\QCoh, T}(\mathcal{O}_X)$ は、コホモロジー層が準連接で
$T$ に支持される複体からなる $D(\mathcal{O}_X)$ の
厳密充満・飽和・三角部分圏を表す。

\begin{lemma}
\label{perfect-lemma-generator-with-support}
\begin{reference}
\cite[Theorem 6.8]{Rouquier-dimensions}
\end{reference}
$X$ を準コンパクトかつ準分離的なスキームとする。
$T \subset X$ を、$X \setminus T$ が準コンパクトとなる閉部分集合とする。
上の記法のもとで、圏 $D_{\QCoh, T}(\mathcal{O}_X)$ は
単一の完全対象で生成される。
\end{lemma}

\begin{proof}
Cohomology of Schemes, Lemma
\ref{coherent-lemma-induction-principle} の帰納原理を用いて証明する。

\medskip\noindent
$X = \Spec(A)$ がアフィンであると仮定する。この場合、
$T = V(f_1, \ldots, f_r)$ となる $f_1, \ldots, f_r \in A$ が存在する。
$K$ を Lemma \ref{perfect-lemma-orthogonal-koszul-complex} のような
$f_1, \ldots, f_r$ に関する Koszul 複体とする。
すると $K$ はコホモロジーが $T$ に支持される完全対象であり、
したがって $D_{\QCoh, T}(\mathcal{O}_X)$ の完全対象である。
一方、$E \in D_{\QCoh, T}(\mathcal{O}_X)$ かつすべての $n$ に対して
$\Hom(K, E[n]) = 0$ なら、Lemma
\ref{perfect-lemma-orthogonal-koszul-complex} により
$E = Rj_*(E|_{X \setminus T}) = 0$ である。
ゆえに $K$ は $D_{\QCoh, T}(\mathcal{O}_X)$ を生成する
（Derived Categories, Definition \ref{derived-definition-generators}
における三角圏の生成対象の定義による）。

\medskip\noindent
$X = U \cup V$ を開被覆とし、$V$ はアフィン、$U$ は準コンパクトで、
$U$ 上では補題が成り立つと仮定する。
$P$ を $T \cap U$ に支持される $D(\mathcal{O}_U)$ の完全対象で、
$D_{\QCoh, T \cap U}(\mathcal{O}_U)$ の生成対象となるものとする。
Lemma \ref{perfect-lemma-lift-map-from-perfect-complex-with-support} により、
$T$ に支持される $D(\mathcal{O}_X)$ の完全対象 $Q$ で、
その $U$ への制限が $P$ を直和因子の一つとして含む直和となるものを
選べる。$V = \Spec(B)$ と書き、$Z = X \setminus U$ とする。
すると $Z$ は $V$ の閉部分集合で、$V \setminus Z$ は準コンパクトである。
$X$ は準分離的なので、$Z \cap T$ は $V$ の閉部分集合で、
$W = V \setminus (Z \cap T)$ は準コンパクトである。したがって
$Z \cap T = V(g_1, \ldots, g_r)$ となる
$g_1, \ldots, g_s \in B$ を選べる。
$K \in D(\mathcal{O}_V)$ を、$B$ 上の $g_1, \ldots, g_s$ に関する
Koszul 複体に対応する完全対象とする。$K$ は閉集合
$(Z \cap T) \subset V$ に支持されるので、直像
$K' = R(V \to X)_*K$ は $D(\mathcal{O}_X)$ の完全対象であり、
その $V$ への制限は $K$ であることに注意する（Cohomology,
Lemma \ref{cohomology-lemma-pushforward-perfect}）。
$Q \oplus K'$ は $D_{\QCoh, T}(\mathcal{O}_X)$ の生成対象であると主張する。

\medskip\noindent
$E$ を $D_{\QCoh, T}(\mathcal{O}_X)$ の対象とし、
$Q \oplus K'$ のどのシフトから $E$ へも非自明な写像がないとする。
Cohomology, Lemma \ref{cohomology-lemma-pushforward-restriction} により
$K' =  R(V \to X)_! K$ であり、したがって
$$
\Hom_{D(\mathcal{O}_X)}(K'[n], E) = \Hom_{D(\mathcal{O}_V)}(K[n], E|_V)
$$
である。したがって Lemma \ref{perfect-lemma-orthogonal-koszul-complex} により、
$E|_V = Rj_*E|_W$ である。ここで $j : W \to V$ は包含である。図式は
$$
\xymatrix{
W \ar[r]_j & V & Z \cap T \ar[l] \ar[d] \\
U \cap V \ar[u]^{j'} \ar[ru]_{j''} & & Z \ar[lu]
}
$$
である。$E$ は $T$ に支持されるので、$E|_W$ は
$T \cap W = T \cap U \cap V$ に支持され、これは $W$ で閉である。
したがって
$$
E|_V = Rj_*(E|_W) = Rj_*(Rj'_*(E|_{U \cap V})) = Rj''_*(E|_{U \cap V})
$$
である。ここで二番目の等号は Cohomology, Lemma
\ref{cohomology-lemma-pushforward-restriction} の (1) である。
これは $E = R(U \to X)_*E|_U$ を含意する（細部を一つ省略した）。
この場合
$$
\Hom_{D(\mathcal{O}_X)}(Q[n], E) = \Hom_{D(\mathcal{O}_U)}(Q|_U[n], E|_U)
$$
であり、これは $\Hom_{D(\mathcal{O}_U)}(P[n], E|_U)$ を
直和因子として含む。したがって $P$ の選び方により、
これらの群の消滅は $E|_U$ が零であることを含意する。
ゆえに $E$ は零である。
\end{proof}



\section{生成対象の例}
\label{perfect-section-example-generator}

\noindent
本節では、環上の射影空間の導来圏がベクトル束、実際には
構造層のシフトの直和によって生成されることを証明する。

\medskip\noindent
次の補題は、$\mathcal{L}$ が豊富なら
$\bigoplus_{n \geq 0} \mathcal{L}^{\otimes -n}$ が生成対象であることをいう。

\begin{lemma}
\label{perfect-lemma-nonzero-some-cohomology}
$X$ をスキームとし、$\mathcal{L}$ を豊富な可逆
$\mathcal{O}_X$-加群とする。$K$ が $D_\QCoh(\mathcal{O}_X)$ の
非零対象なら、ある $n \geq 0$ と $p \in \mathbf{Z}$ に対して
コホモロジー群
$H^p(X, K \otimes_{\mathcal{O}_X}^\mathbf{L} \mathcal{L}^{\otimes n})$
は非零である。
\end{lemma}

\begin{proof}
豊富な可逆層をもつ $X$ は準コンパクトかつ分離的であることを
思い出そう（Properties, Definition \ref{properties-definition-ample}
および Lemma \ref{properties-lemma-affine-s-opens-cover-quasi-separated}）。
したがって Proposition \ref{perfect-proposition-quasi-compact-affine-diagonal}
を適用し、$K$ を準連接加群の複体 $\mathcal{F}^\bullet$ で表せる。
次が非零となる任意の $p$ を選ぶ。
$\mathcal{H}^p = \Ker(\mathcal{F}^p \to \mathcal{F}^{p + 1})/
\Im(\mathcal{F}^{p - 1} \to \mathcal{F}^p)$ は非零である。
茎 $\mathcal{H}^p_x$ が非零となる点 $x \in X$ を選ぶ。
$X_s$ が $x$ のアフィン開近傍となるような $n \geq 0$ と
$s \in \Gamma(X, \mathcal{L}^{\otimes n})$ を選ぶ。
茎 $\mathcal{H}^p_x$ の非零元へ写る
$\tau \in \mathcal{H}^p(X_s)$ を選ぶ。$\mathcal{H}^p$ は準連接で
$X_s$ はアフィンなので、これは可能である。$X_s$ 上の切断をとることは
準連接加群上の完全関手なので、$\mathcal{F}^{p + 1}(X_s)$ では零へ、
$\mathcal{H}^p(X_s)$ では $\tau$ へ写る切断
$\tau' \in \mathcal{F}^p(X_s)$ をとれる。Properties, Lemma
\ref{properties-lemma-invert-s-sections} により、
$\tau' \otimes s^{\otimes m}$ が切断
$\tau'' \in \Gamma(X, \mathcal{F}^p \otimes \mathcal{L}^{\otimes mn})$
の像となるような $m$ が存在する。同じ補題をもう一度適用すると、
$\tau'' \otimes s^{\otimes l}$ が
$\mathcal{F}^{p + 1} \otimes \mathcal{L}^{\otimes (m + l)n}$ で
零へ写るような $l \geq 0$ を得る。このとき $\tau''$ は
$H^p(X, K \otimes^\mathbf{L}_{\mathcal{O}_X} \mathcal{L}^{(m + l)n})$
の非零類を与え、所望の結論を得る。
\end{proof}

\begin{lemma}
\label{perfect-lemma-construct-the-next-one}
$A$ を環とし、$X = \mathbf{P}^n_A$ とする。任意の $a \in \mathbf{Z}$
に対して、$X$ 上のベクトル束の完全複体
$$
0 \to \mathcal{O}_X(a) \to \ldots
\to \mathcal{O}_X(a + i)^{\oplus {n + 1 \choose i}} \to
\ldots \to \mathcal{O}_X(a + n + 1) \to 0
$$
が存在する。
\end{lemma}

\begin{proof}
$\mathbf{P}^n_A$ は $\text{Proj}(A[X_0, \ldots, X_n])$ であることを
思い出そう。Constructions, Definition
\ref{constructions-definition-projective-space} を参照せよ。
Koszul 複体
$$
K_\bullet = K_\bullet(A[X_0, \ldots, X_n], X_0, \ldots, X_n)
$$
を、$S = A[X_0, \ldots, X_n]$ 上の $X_0, \ldots, X_n$ に関して考える。
$X_0, \ldots, X_n$ は多項式環 $S$ の正則列であることが明らかなので
（More on Algebra, Lemma
\ref{more-algebra-lemma-regular-koszul-regular}）、Koszul 複体
$K_\bullet$ は次数 $0$ を除いて完全であり、その次数の
コホモロジーは $S/(X_0, \ldots, X_n)$ である。
$K_i$ の生成元を次数 $+i$ に置けば、$K_\bullet$ は次数付き加群の
複体となることに注意する。言い換えると、完全複体
$$
0 \to S(-n - 1) \to \ldots \to S(-n - 1 + i)^{\oplus {n \choose i}} \to \ldots
\to S \to S/(X_0, \ldots, X_n) \to 0
$$
を得る。Constructions, Lemma \ref{constructions-lemma-proj-sheaves}
の完全関手 $\tilde{\ }$ を適用し、最後の項がこの関手の核に
属することを用いると、完全複体
$$
0 \to \mathcal{O}_X(-n - 1) \to \ldots
\to \mathcal{O}_X(-n - 1 + i)^{\oplus {n + 1 \choose i}} \to
\ldots \to \mathcal{O}_X \to 0
$$
を得る。可逆層 $\mathcal{O}_X(n + a + 1)$ で捻ると、
補題の完全複体を得る。
\end{proof}

\begin{lemma}
\label{perfect-lemma-generator-P1}
$A$ を環とし、$X = \mathbf{P}^n_A$ とする。このとき
$$
E =
\mathcal{O}_X \oplus \mathcal{O}_X(-1) \oplus \ldots \oplus \mathcal{O}_X(-n)
$$
は $D_\QCoh(X)$ の生成対象である
（Derived Categories, Definition \ref{derived-definition-generators}）。
\end{lemma}

\begin{proof}
$K \in D_\QCoh(\mathcal{O}_X)$ とし、すべての
$p \in \mathbf{Z}$ に対して $\Hom(E, K[p]) = 0$ と仮定する。
$K = 0$ を示さなければならない。Derived Categories, Lemma
\ref{derived-lemma-right-orthogonal} により、すべての
$E' \in \langle E \rangle$ と $p \in \mathbf{Z}$ に対して
$\Hom(E', K[p])$ は零である。Lemma
\ref{perfect-lemma-construct-the-next-one} を $a = -n - 1$ に適用すると、
$\mathcal{O}_X(-n - 1) \in \langle E \rangle$ である。
実際、これは各項が $E$ の直和因子の有限直和である有限複体と
擬同型である。$a = -n - 2$ として議論を繰り返すと、
$\mathcal{O}_X(-n - 2) \in \langle E \rangle$ である。
帰納法により、すべての $m \geq 0$ に対して
$\mathcal{O}_X(-m) \in \langle E \rangle$ を得る。また
$$
\Hom(\mathcal{O}_X(-m), K[p]) =
H^p(X, K \otimes_{\mathcal{O}_X}^\mathbf{L} \mathcal{O}_X(m)) =
H^p(X, K \otimes_{\mathcal{O}_X}^\mathbf{L} \mathcal{O}_X(1)^{\otimes m})
$$
なので、Lemma \ref{perfect-lemma-nonzero-some-cohomology} により $K = 0$
と結論する。（ここでは $\mathcal{O}_X(1)$ が $X$ 上の豊富な可逆層で
あることも用いた。これは Properties, Lemma
\ref{properties-lemma-open-in-proj-ample} から従う。）
\end{proof}

\begin{remark}
\label{perfect-remark-pullback-generator}
$f : X \to Y$ を準コンパクトかつ準分離的なスキームの射とする。
$E \in D_\QCoh(\mathcal{O}_Y)$ を生成対象とする
（Theorem \ref{perfect-theorem-bondal-van-den-Bergh} を参照せよ）。
このとき次は同値である。
\begin{enumerate}
\item $K \in D_\QCoh(\mathcal{O}_X)$ に対して
$Rf_*K = 0$ であることと $K = 0$ であることは同値である。
\item $Rf_* : D_\QCoh(\mathcal{O}_X) \to D_\QCoh(\mathcal{O}_Y)$
は同型を反映する。
\item $Lf^*E$ は $D_\QCoh(\mathcal{O}_X)$ の生成対象である。
\end{enumerate}
(1) と (2) の同値性は、
$Rf_* : D_\QCoh(\mathcal{O}_X) \to D_\QCoh(\mathcal{O}_Y)$ が
三角圏の完全関手であることから形式的に従う。同様に、(1) と (3) の
同値性は $Lf^*$ が $Rf_*$ の左随伴であることから形式的に従う。
$f$ がアフィンである場合（Lemma \ref{perfect-lemma-affine-morphism}）、
$f$ が開埋め込みである場合、または $f$ がそのような射の合成である
場合に、これらの条件は成り立つ。したがって
\begin{enumerate}
\item $X$ が準アフィン・スキームなら、$\mathcal{O}_X$ は
$D_\QCoh(\mathcal{O}_X)$ の生成対象である。
\item $X \subset \mathbf{P}^n_A$ が準コンパクトな局所閉部分スキームなら、
$\mathcal{O}_X \oplus \mathcal{O}_X(-1) \oplus \ldots \oplus \mathcal{O}_X(-n)$
は Lemma \ref{perfect-lemma-generator-P1} により
$D_\QCoh(\mathcal{O}_X)$ の生成対象である。
\end{enumerate}
\end{remark}




\section{コンパクト対象と完全対象}
\label{perfect-section-compact}

\noindent
$X$ を有限次元ネーター・スキームとする。Cohomology, Proposition
\ref{cohomology-proposition-vanishing-Noetherian} と Cohomology on Sites,
Lemma \ref{sites-cohomology-lemma-when-jshriek-compact} により、
任意の開集合 $U \subset X$ に対して加群の層
$j_!\mathcal{O}_U$ は $D(\mathcal{O}_X)$ のコンパクト対象である。
これらの層は一般には準連接でないため、導来圏
$D(\mathcal{O}_X)$ の完全対象を与えない。しかし、準連接
コホモロジー層をもつ複体に制限すれば、この現象は起こらない。
正確な主張は次である。

\begin{proposition}
\label{perfect-proposition-compact-is-perfect}
$X$ を準コンパクトかつ準分離的なスキームとする。
$D_\QCoh(\mathcal{O}_X)$ の対象がコンパクトであることと
完全であることは同値である。
\end{proposition}

\begin{proof}
$K$ を $D(\mathcal{O}_X)$ の完全対象とし、その双対を
$K^\vee$ とする（Cohomology, Lemma
\ref{cohomology-lemma-dual-perfect-complex}）。このとき
$$
\Hom_{D(\mathcal{O}_X)}(K, M) =
H^0(X, K^\vee \otimes_{\mathcal{O}_X}^\mathbf{L} M)
$$
が $M$ に関して関手的に成り立つ。
$K^\vee \otimes_{\mathcal{O}_X}^\mathbf{L} -$ は直和と可換であり、
Lemma \ref{perfect-lemma-quasi-coherence-pushforward-direct-sums} により
$H^0(X, -)$ は $D_\QCoh(\mathcal{O}_X)$ 上で直和と可換なので、
$K$ は $D_\QCoh(\mathcal{O}_X)$ においてコンパクトである。

\medskip\noindent
逆に、$K$ を $D_\QCoh(\mathcal{O}_X)$ のコンパクト対象とする。
$K$ が完全であることを示すには、任意のアフィン開集合
$U \subset X$ に対して $K|_U$ が完全であることを示せば十分である。
Cohomology, Lemma
\ref{cohomology-lemma-perfect-independent-representative} を参照せよ。
$j : U \to X$ は準コンパクトかつ分離的な射であることに注意する。
したがって
$Rj_* : D_\QCoh(\mathcal{O}_U) \to D_\QCoh(\mathcal{O}_X)$
は直和と可換する。Lemma
\ref{perfect-lemma-quasi-coherence-pushforward-direct-sums} を参照せよ。
$U$ への制限と $Rj_*$ の随伴性により、$K|_U$ は
$D_\QCoh(\mathcal{O}_U)$ のコンパクト対象である。
ゆえに $X$ がアフィンの場合に帰着する。

\medskip\noindent
$X = \Spec(A)$ がアフィンであると仮定する。Lemma
\ref{perfect-lemma-affine-compare-bounded} により、問題は $D(A)$ に対する
同じ問題へ翻訳される。$D(A)$ に対しては、結果は More on Algebra,
Proposition \ref{more-algebra-proposition-perfect-is-compact} である。
\end{proof}

\begin{remark}
\label{perfect-remark-classical-generator}
$X$ を準コンパクトかつ準分離的なスキームとする。
$G$ を $D(\mathcal{O}_X)$ の完全対象で
$D_\QCoh(\mathcal{O}_X)$ の生成対象となるものとする。
Theorem \ref{perfect-theorem-bondal-van-den-Bergh} により、このような対象は
少なくとも一つ存在する。Lemma \ref{perfect-lemma-quasi-coherence-direct-sums},
Proposition \ref{perfect-proposition-compact-is-perfect}, および
Derived Categories, Proposition
\ref{derived-proposition-generator-versus-classical-generator} を合わせると、
$G$ は $D_{perf}(\mathcal{O}_X)$ の古典的生成対象である。
\end{remark}

\noindent
次の結果は Proposition \ref{perfect-proposition-compact-is-perfect} の強化である。
$T \subset X$ をスキーム $X$ の閉部分集合とする。前と同様、
$D_T(\mathcal{O}_X)$ は $T$ に支持される対象からなる
$D(\mathcal{O}_X)$ の厳密充満・飽和・三角部分圏を表す
（Definition \ref{perfect-definition-supported-on}）。
直和をとることはコホモロジー層をとることと可換するので、
$D_T(\mathcal{O}_X)$ は直和をもち、それらは
$D(\mathcal{O}_X)$ における直和と等しい。

\begin{lemma}
\label{perfect-lemma-compact-is-perfect-with-support}
$X$ を準コンパクトかつ準分離的なスキームとする。
$T \subset X$ を、$X \setminus T$ が準コンパクトとなる閉部分集合とする。
$D_{\QCoh, T}(\mathcal{O}_X)$ の対象がコンパクトであることと、
$D(\mathcal{O}_X)$ の対象として完全であることは同値である。
\end{lemma}

\begin{proof}
補題の直前の注意により、$D_{\QCoh, T}(\mathcal{O}_X)$ は直和をもつ
三角圏である。Proposition \ref{perfect-proposition-compact-is-perfect} により、
完全対象は $D(\mathcal{O}_X)$ のコンパクト対象を定め、したがって
まして、直和をとることで保たれる任意の部分圏のコンパクト対象を定める。
逆については、$\mathcal{O}_X$-加群の完全複体である生成対象
$E \in D_{\QCoh, T}(\mathcal{O}_X)$ が存在することを用いる。
Lemma \ref{perfect-lemma-generator-with-support} を参照せよ。
したがって上の議論により $E$ はコンパクトである。すると
Derived Categories, Proposition
\ref{derived-proposition-generator-versus-classical-generator} により、
$E$ は $D_{\QCoh, T}(\mathcal{O}_X)$ のコンパクト対象からなる
充満部分圏の古典的生成対象である。ゆえに任意のコンパクト対象は、
$E$ から (a) シフト、(b) 有限直和、(c) 錐、(d) 直和因子をとる、
という有限列の操作によって構成できる。これらの各操作は完全性を保つので、
結果が従う。
\end{proof}

\begin{remark}
\label{perfect-remark-classical-generator-with-support}
$X$ を準コンパクトかつ準分離的なスキームとする。
$T \subset X$ を、$X \setminus T$ が準コンパクトとなる閉部分集合とする。
$G$ を $D_{\QCoh, T}(\mathcal{O}_X)$ の完全対象で、
$D_{\QCoh, T}(\mathcal{O}_X)$ の
生成対象となるものとする。Lemma \ref{perfect-lemma-generator-with-support}
により、このような対象は少なくとも一つ存在する。
$D_{\QCoh, T}(\mathcal{O}_X)$ が直和をもつことと、Lemma
\ref{perfect-lemma-compact-is-perfect-with-support} および Derived Categories,
Proposition \ref{derived-proposition-generator-versus-classical-generator}
を合わせると、$G$ は $D_{perf, T}(\mathcal{O}_X)$ の
古典的生成対象である。
\end{remark}

\noindent
次の補題は、直前の補題の証明に用いた考えの応用である。

\begin{lemma}
\label{perfect-lemma-map-from-pseudo-coherent-to-complex-with-support}
$X$ を準コンパクトかつ準分離的なスキームとする。$T \subset X$ を、
$U = X \setminus T$ が準コンパクトとなる閉部分集合とする。
$\alpha : P \to E$ を $D_\QCoh(\mathcal{O}_X)$ の射で、
次のいずれかを満たすものとする。
\begin{enumerate}
\item $P$ は完全で、$E$ は $T$ に支持される。
\item $P$ は擬連接、$E$ は $T$ に支持され、$E$ は下に有界である。
\end{enumerate}
このとき、$\mathcal{O}_X$-加群の完全複体 $I$ と写像
$I \to \mathcal{O}_X[0]$ で、$I \otimes^\mathbf{L} P \to E$ は零、
$I|_U \to \mathcal{O}_U[0]$ は同型となるものが存在する。
\end{lemma}

\begin{proof}
$\mathcal{D} = D_{\QCoh, T}(\mathcal{O}_X)$ とおく。どちらの場合も、
複体 $K = R\SheafHom(P, E)$ は $\mathcal{D}$ の対象である。
準連接性については Lemma \ref{perfect-lemma-quasi-coherence-internal-hom}
を参照せよ。$R\SheafHom$ の形成は開集合への制限と可換するので、
$K$ が $T$ に支持されることは明らかである。写像 $\alpha$ は
$H^0(K) = \Hom_{D(\mathcal{O}_X)}(\mathcal{O}_X[0], K)$ の元を定める。
したがって写像 $\alpha : \mathcal{O}_X[0] \to K$ に対して
結果を証明すれば十分である。

\medskip\noindent
$E \in \mathcal{D}$ を完全生成対象とする。Lemma
\ref{perfect-lemma-generator-with-support} を参照せよ。Derived Categories,
Lemma \ref{derived-lemma-write-as-colimit} において生成対象 $E$ を用いた
ときのように
$$
K = \text{hocolim} K_n
$$
と書く。関手 $\mathcal{D} \to D(\mathcal{O}_X)$ は直和と可換するので、
$K = \text{hocolim} K_n$ は $D(\mathcal{O}_X)$ においても成り立つ。
$\mathcal{O}_X$ は $D(\mathcal{O}_X)$ のコンパクト対象なので、
$\alpha$ を誘導する整数 $n$ と射
$\alpha_n : \mathcal{O}_X \to K_n$ を得る。Derived Categories, Lemma
\ref{derived-lemma-commutes-with-countable-sums} を参照せよ。
周囲の圏 $D(\mathcal{O}_X)$ における射
$\mathcal{O}_X[0] \to K_n$ に Derived Categories, Lemma
\ref{derived-lemma-factor-through} を適用すると、$\alpha_n$ は
$\mathcal{O}_X[0] \to Q \to K_n$ と分解する。ここで $Q$ は
$\langle E \rangle$ の対象である。したがって $Q$ は $T$ に支持される
完全複体である。

\medskip\noindent
完全三角形
$$
I \to \mathcal{O}_X[0] \to Q \to I[1]
$$
を選ぶ。構成により $I$ は完全で、写像 $I \to \mathcal{O}_X[0]$
は $U$ 上で同型に制限され、$\alpha$ は $Q$ を経由して分解するので
合成 $I \to K$ は零である。これで補題が証明された。
\end{proof}






\section{加群圏としての導来圏}
\label{perfect-section-derived-is-dga}

\noindent
本節では、これまでの議論からいくつかの帰結を導く。
その前に、さらに二つほど補題が必要である。

\begin{lemma}
\label{perfect-lemma-tensor-with-QCoh-complex}
$X$ をスキームとする。$K^\bullet$ を $\mathcal{O}_X$-加群の複体で、
そのコホモロジー層が擬連接であるものとする。
$(E, d) = \Hom_{\text{Comp}^{dg}(\mathcal{O}_X)}(K^\bullet, K^\bullet)$
を自己準同型微分次数付き代数とする。このとき
$$
- \otimes_E^\mathbf{L} K^\bullet :
D(E, \text{d}) \longrightarrow D(\mathcal{O}_X)
$$
という Differential Graded Algebra, Lemma
\ref{dga-lemma-tensor-with-complex-derived}
の関手の像は $D_\QCoh(\mathcal{O}_X)$ に含まれる。
\end{lemma}

\begin{proof}
$P$ を性質 (P) をもつ微分次数付き $E$-加群とし、
$F_\bullet$ を
Differential Graded Algebra, Section \ref{dga-section-P-resolutions}
における $P$ のフィルトレーションとする。このとき
$$
P \otimes_E K^\bullet = \text{hocolim}\ F_iP \otimes_E K^\bullet
$$
である。各 $F_iP$ は、その次数付き商が $E[k]$ の直和であるような
有限フィルトレーションをもつ。したがって結果は直ちに従う。
\end{proof}

\noindent
次の結果は \cite{BvdB} による。

\begin{theorem}
\label{perfect-theorem-DQCoh-is-Ddga}
$X$ を準コンパクトかつ準分離なスキームとする。
このとき、非零なコホモロジー群 $H^i(E)$ が有限個しかなく、
$D_\QCoh(\mathcal{O}_X)$ が $D(E, \text{d})$ と同値になるような
微分次数付き代数 $(E, \text{d})$ が存在する。
\end{theorem}

\begin{proof}
$K^\bullet$ を、完全で $D_\QCoh(\mathcal{O}_X)$ を生成する
$\mathcal{O}$-加群の K-入射的複体とする。このような複体は
定理 \ref{perfect-theorem-bondal-van-den-Bergh} と K-入射分解の存在により
存在する。以下では
$$
(E, \text{d}) = \Hom_{\text{Comp}^{dg}(\mathcal{O}_X)}(K^\bullet, K^\bullet)
$$
とすれば定理が成り立つことを示す。ここで
$\text{Comp}^{dg}(\mathcal{O}_X)$ は $\mathcal{O}$-加群の複体からなる
微分次数付き圏である。
Differential Graded Algebra, Section \ref{dga-section-variant-base-change}
を参照されたい。$K^\bullet$ は K-入射的であるから、すべての
$n \in \mathbf{Z}$ に対して
\begin{equation}
\label{perfect-equation-E-is-OK}
H^n(E) = \Ext^n_{D(\mathcal{O}_X)}(K^\bullet, K^\bullet)
\end{equation}
である。補題 \ref{perfect-lemma-ext-from-perfect-into-bounded-QCoh} により、
これらの Ext のうち非零なものは有限個しかない。関手
$$
- \otimes_E^\mathbf{L} K^\bullet :
D(E, \text{d}) \longrightarrow D(\mathcal{O}_X)
$$
を考える。これは Differential Graded Algebra, Lemma
\ref{dga-lemma-tensor-with-complex-derived}
の関手である。$K^\bullet$ は完全であるから
$D(\mathcal{O}_X)$ のコンパクト対象を定める。命題
\ref{perfect-proposition-compact-is-perfect} を参照されたい。
これと (\ref{perfect-equation-E-is-OK}) を合わせると、
Differential Graded Algebra, Lemmas
\ref{dga-lemma-fully-faithful-in-compact-case} により、上の関手は
充満忠実である。この関手は
$$
R\Hom(K^\bullet, - ) : D(\mathcal{O}_X) \longrightarrow D(E, \text{d})
$$
という右随伴をもつ。これは Differential Graded Algebra, Lemmas
\ref{dga-lemma-tensor-with-complex-hom-adjoint} による。
随伴関手に関する一般論により、この右随伴は左擬逆関手である。
一方、補題 \ref{perfect-lemma-tensor-with-QCoh-complex} により
$$
- \otimes_E^\mathbf{L} K^\bullet :
D(E, \text{d}) \longrightarrow D_\QCoh(\mathcal{O}_X)
$$
を得る。また $K^\bullet$ を $D_\QCoh(\mathcal{O}_X)$ の生成対象として
選んだので、この随伴を $D_\QCoh(\mathcal{O}_X)$ に制限したときの
核は零である。したがって形式的な議論により所望の同値を得る。
Derived Categories, Lemma
\ref{derived-lemma-fully-faithful-adjoint-kernel-zero}
を参照されたい。
\end{proof}

\begin{remark}[台を指定した変種]
\label{perfect-remark-DQCoh-is-Ddga-with-support}
$X$ を準コンパクトかつ準分離なスキームとし、$T \subset X$ を
$X \setminus T$ が準コンパクトとなる閉部分集合とする。
定理 \ref{perfect-theorem-DQCoh-is-Ddga} の類似は
$D_{\QCoh, T}(\mathcal{O}_X)$ に対しても成り立つ。
実際、補題 \ref{perfect-lemma-generator-with-support}、
\ref{perfect-lemma-compact-is-perfect-with-support}、および台を指定した
補題 \ref{perfect-lemma-tensor-with-QCoh-complex} の変種を用いれば、
定理の証明とまったく同じ議論から従う。
これが必要になれば、ここに結果を精密に述べて詳しい証明を与える。
\end{remark}

\begin{remark}[微分次数付き代数の一意性]
\label{perfect-remark-independence-choice}
$X$ を環 $R$ 上の準コンパクトかつ準分離なスキームとする。
定理 \ref{perfect-theorem-DQCoh-is-Ddga} の証明の構成により、
$D_\QCoh(X)$ が三角圏として $D(A, \text{d})$ と $R$-線型同値になる
ような $R$ 上の微分次数付き代数 $(A, \text{d})$ が存在する。
そこで、$(A, \text{d})$ はどの程度一意であるかを問うことができる。
答えは、$(A, \text{d})$ が導来同値を除いて一意に定まるというだけの
主張より（わずかに）強い。実際、$(B, \text{d})$ をそのような
第二の組としよう。このとき
$$
(A, \text{d}) = \Hom_{\text{Comp}^{dg}(\mathcal{O}_X)}(K^\bullet, K^\bullet)
$$
および
$$
(B, \text{d}) = \Hom_{\text{Comp}^{dg}(\mathcal{O}_X)}(L^\bullet, L^\bullet)
$$
である。ただし $K^\bullet$ と $L^\bullet$ は
$D_\QCoh(\mathcal{O}_X)$ の完全生成対象に対応する
$\mathcal{O}_X$-加群の K-入射的複体である。ここで
$$
\Omega = \Hom_{\text{Comp}^{dg}(\mathcal{O}_X)}(K^\bullet, L^\bullet)
\quad
\Omega' = \Hom_{\text{Comp}^{dg}(\mathcal{O}_X)}(L^\bullet, K^\bullet)
$$
とおく。このとき $\Omega$ は微分次数付き
$B^{opp} \otimes_R A$-加群であり、$\Omega'$ は微分次数付き
$A^{opp} \otimes_R B$-加群である。さらに、同値
$$
D(A, \text{d}) \to D_\QCoh(\mathcal{O}_X) \to
D(B, \text{d})
$$
は関手 $- \otimes_A^\mathbf{L} \Omega'$ によって与えられ、
擬逆についても同様である。したがってこれは
Differential Graded Algebra, Remark \ref{dga-remark-hochschild-cohomology}
の状況である。この注意が必要になれば、ここに精密な主張と
詳しい証明を与える。
\end{remark}






\section{擬連接複体の特徴付け I}
\label{perfect-section-pseudo-coherent-hocolim}

\noindent
上の方法を用いると、擬連接対象を完全対象による近似の
導来ホモトピー極限として特徴付けることができる。

\begin{lemma}
\label{perfect-lemma-pseudo-coherent-hocolim}
$X$ を準コンパクトかつ準分離なスキームとし、
$K \in D(\mathcal{O}_X)$ とする。次は同値である。
\begin{enumerate}
\item $K$ は擬連接である。
\item $K = \text{hocolim} K_n$ であり、$K_n$ は完全で、
すべての $n$ に対して
$\tau_{\geq -n}K_n \to \tau_{\geq -n}K$ は同型である。
\end{enumerate}
\end{lemma}

\begin{proof}
含意 (2) $\Rightarrow$ (1) は任意の環付き空間上で成り立つ。
実際、(2) を仮定する。導来圏の完全対象は擬連接であることを
思い出そう。Cohomology, Lemma \ref{cohomology-lemma-perfect}
を参照されたい。すると定義から、
$\tau_{\geq -n}K_n$ は $(-n + 1)$-擬連接
であるから $\tau_{\geq -n}K$ も $(-n + 1)$-擬連接であり、
したがって $K$ は $(-n + 1)$-擬連接である。これはすべての $n$
について成り立つので、$K$ は擬連接である。
Cohomology, Definition \ref{cohomology-definition-pseudo-coherent}
を参照されたい。

\medskip\noindent
(1) を仮定する。まず、完全複体 $K_1$ による
$(X, K, -2)$ の近似 $K_1 \to K$ を選ぶ。定義
\ref{perfect-definition-approximation-holds}、\ref{perfect-definition-approximation}
および定理 \ref{perfect-theorem-approximation} を参照されたい。
帰納的に
$$
K_1 \to K_2 \to \ldots \to K_n \to K
$$
を得ており、各 $K_i$ は完全で、すべての $1 \leq i \leq n$ に対して
$\tau_{\geq -i}K_i \to \tau_{\geq -i}K$ は同型であると仮定する。
完全対象 $K_n$ に対して、補題
\ref{perfect-lemma-ext-from-perfect-into-bounded-QCoh} のように
$a \leq b$ を選ぶ。$(X, K, \min(a - 1, -n - 1))$ の近似
$K_{n + 1} \to K$ を選び、完全三角形
$$
K_{n + 1} \to K \to C \to K_{n + 1}[1]
$$
を選ぶ。このとき $C \in D_\QCoh(\mathcal{O}_X)$ であり、
$i \geq a$ に対して $H^i(C) = 0$ である。したがって $a, b$ の
選び方により $\Hom_{D(\mathcal{O}_X)}(K_n, C) = 0$ である。
ゆえに合成 $K_n \to K \to C$ は零である。そこで
Derived Categories, Lemma \ref{derived-lemma-representable-homological}
により $K_n \to K$ は $K_{n + 1}$ を経由して分解し、
帰納法の一段階が示される。

\medskip\noindent
残るのは $K = \text{hocolim} K_n$ を示すことである。これは関手
$H^i( - ) : D(\mathcal{O}_X) \to \textit{Mod}(\mathcal{O}_X)$
と $K_n$ の選び方に
Derived Categories, Lemma \ref{derived-lemma-cohomology-of-hocolim}
を適用すれば従う。
\end{proof}

\begin{lemma}
\label{perfect-lemma-pseudo-coherent-hocolim-with-support}
$X$ を準コンパクトかつ準分離なスキームとする。
$T \subset X$ を $X \setminus T$ が準コンパクトとなる閉部分集合とし、
$K \in D(\mathcal{O}_X)$ を $T$ 上に台をもつ対象とする。
次は同値である。
\begin{enumerate}
\item $K$ は擬連接である。
\item $K = \text{hocolim} K_n$ であり、$K_n$ は完全で $T$ 上に台をもち、
すべての $n$ に対して
$\tau_{\geq -n}K_n \to \tau_{\geq -n}K$ は同型である。
\end{enumerate}
\end{lemma}

\begin{proof}
近似を選ぶ際に三つ組 $(T, K, m)$ を用いる点を除けば、
この補題の証明は補題 \ref{perfect-lemma-pseudo-coherent-hocolim}
の証明とまったく同じである。
\end{proof}









\section{同値の一例}
\label{perfect-section-example-equivalence}

\noindent
節 \ref{perfect-section-example-generator} では、環 $A$ 上の射影空間
$\mathbf{P}^n_A$ の導来圏がベクトル束、実際には構造層の
シフトの直和によって生成されることを証明した。本節では、
これが $D_\QCoh(\mathcal{O}_{\mathbf{P}^n_A})$ と
ある $A$-代数の導来圏との同値を定めることを証明する。

\medskip\noindent
結果を述べる前に記法を準備する。$A$ を環とし、
$X = \mathbf{P}^n_A = \text{Proj}(S)$ とおく。ただし
$S = A[X_0, \ldots, X_n]$ である。補題 \ref{perfect-lemma-generator-P1}
により
\begin{equation}
\label{perfect-equation-generator-Pn}
P =
\mathcal{O}_X \oplus \mathcal{O}_X(-1) \oplus \ldots \oplus \mathcal{O}_X(-n)
\end{equation}
は $D_\QCoh(\mathcal{O}_X)$ の完全生成対象である。
（非可換な）$A$-代数
\begin{equation}
\label{perfect-equation-algebra-for-Pn}
R = \Hom_{\mathcal{O}_X}(P, P) =
\left(
\begin{matrix}
S_0 & S_1 & S_2 & \ldots & \ldots \\
0 & S_0 & S_1 & \ldots & \ldots\\
0 & 0 & S_0 & \ldots & \ldots \\
\ldots & \ldots & \ldots & \ldots & \ldots \\
0 & \ldots & \ldots & \ldots & S_0
\end{matrix}
\right)
\end{equation}
を、その明らかな乗法と加法とともに考える。$P$ を通常の仕方で
$\mathcal{O}_X$-加群の複体（すなわち次数 $0$ に $P$ を置き、
他のすべての次数には零を置いたもの）とみなすと
$$
R = \Hom_{\text{Comp}^{dg}(\mathcal{O}_X)}(P, P)
$$
である。ここで右辺では $R$ を、微分が零である $A$ 上の
微分次数付き代数（すなわち次数 $0$ に $R$ を置き、他のすべての
次数には零を置いたもの）とみなしている。
Differential Graded Algebra, Section \ref{dga-section-variant-base-change}
の議論により、導来関手
$$
- \otimes_R^\mathbf{L} P : D(R) \longrightarrow D(\mathcal{O}_X),
$$
を得る。特に Differential Graded Algebra, Lemma
\ref{dga-lemma-tensor-with-complex-derived} を参照されたい。
補題 \ref{perfect-lemma-tensor-with-QCoh-complex} により、この関手の本質的像は
$D_\QCoh(\mathcal{O}_X)$ に含まれる。

\begin{lemma}
\label{perfect-lemma-Pn-module-category}
\begin{reference}
\cite{Beilinson}
\end{reference}
$A$ を環とし、$X = \mathbf{P}^n_A = \text{Proj}(S)$ とおく。
ただし $S = A[X_0, \ldots, X_n]$ である。
(\ref{perfect-equation-generator-Pn}) の $P$ と
(\ref{perfect-equation-algebra-for-Pn}) の $R$ に対して、関手
$$
- \otimes_R^\mathbf{L} P : D(R) \longrightarrow D_\QCoh(\mathcal{O}_X)
$$
は $R$ を $P$ に送る三角圏の $A$-線型同値である。
\end{lemma}

\noindent
言い換えると、射影空間上の擬連接加群の導来圏は
ある（非可換）代数上の加群の導来圏と同値である。
射影空間のこの性質は、$A$ 上の全射影スキームの中では
かなり例外的なものと思われる。

\begin{proof}
この関手が充満忠実であることを示すには、$i \not = 0$ に対して
$\Ext^i_X(P, P)$ が零であり、$i = 0$ に対して $R$ に等しいことを
示せば十分である。Differential Graded Algebra, Lemma
\ref{dga-lemma-fully-faithful-in-compact-case} を参照されたい。
補題 \ref{perfect-lemma-ext-from-perfect-into-bounded-QCoh} の証明と同様に
$$
\Ext^i_X(P, P) = H^i(X, P^\wedge \otimes P) =
\bigoplus\nolimits_{0 \leq a, b \leq n} H^i(X, \mathcal{O}_X(a - b))
$$
を得る。射影空間のコホモロジーの計算
（Cohomology of Schemes, Lemma
\ref{coherent-lemma-cohomology-projective-space-over-ring}）により、
これらの $\Ext$-群は $i = 0$ でない限り零である。
$i = 0$ の場合には $R$ を得る。これは
(\ref{perfect-equation-algebra-for-Pn}) における $R$ の定義そのものである。
Differential Graded Algebra, Lemma
\ref{dga-lemma-tensor-with-complex-hom-adjoint}
により、この関手は右随伴
$R\Hom(P, -) : D_\QCoh(\mathcal{O}_X) \to D(R)$ をもつ。
補題 \ref{perfect-lemma-generator-P1} により $P$ は
$D_\QCoh(\mathcal{O}_X)$ の生成対象なので、
$R\Hom(P, -)$ の核は零である。したがって
Derived Categories, Lemma
\ref{derived-lemma-fully-faithful-adjoint-kernel-zero} により、
この関手は三角圏の同値である。
\end{proof}






\section{コヒーレータ再論}
\label{perfect-section-better-coherator}

\noindent
節 \ref{perfect-section-coherator} では、標準関手
$D(\QCoh(\mathcal{O}_X)) \to D(\mathcal{O}_X)$ の右随伴 $RQ_X$
を構成して調べた。これはコヒーレータ
$Q_X : \textit{Mod}(\mathcal{O}_X) \to \QCoh(\mathcal{O}_X)$ の
右導来延長として構成された。本節では、包含関手
$$
D_\QCoh(\mathcal{O}_X) \longrightarrow D(\mathcal{O}_X)
$$
がいつ右随伴をもつかを調べる。この右随伴が存在するとき、それを
\footnote{これはおそらく標準的な記法ではない。しかし、すでに
コヒーレータには $Q_X$ を、その導来延長には $RQ_X$ を用いている。}
$$
DQ_X :
D(\mathcal{O}_X) \longrightarrow D_\QCoh(\mathcal{O}_X)
$$
と記す。準コンパクトかつ準分離なスキームは、このような右随伴を
もつことが分かる。

\begin{lemma}
\label{perfect-lemma-better-coherator}
$X$ を準コンパクトかつ準分離なスキームとする。
包含関手 $D_\QCoh(\mathcal{O}_X) \to D(\mathcal{O}_X)$ は
右随伴 $DQ_X$ をもつ。
\end{lemma}

\begin{proof}[第一の証明]
Cohomology of Schemes, Lemma \ref{coherent-lemma-induction-principle}
と同じ帰納原理を用いて証明する。
$D(\QCoh(\mathcal{O}_X)) \to D_\QCoh(\mathcal{O}_X)$ が同値なら、
節 \ref{perfect-section-coherator} の関手 $RQ_X$ は
$D(\QCoh(\mathcal{O}_X)) \to D(\mathcal{O}_X)$ の右随伴なので、
補題は成り立つ。特に補題 \ref{perfect-lemma-affine-coherator} により、
アフィン・スキームに対して成り立つ。したがって次を示せば十分である。
$X = U \cup V$ が二つの準コンパクト開集合の和であり、補題が
$U$、$V$、$U \cap V$ に対して成り立つなら、$X$ に対しても成り立つ。

\medskip\noindent
この随伴が存在するための必要十分条件は、$D(\mathcal{O}_X)$ の
すべての対象 $K$ に対して完全三角形
$$
E' \to E \to K \to E'[1]
$$
を $D(\mathcal{O}_X)$ 内で見いだせて、$E'$ が
$D_\QCoh(\mathcal{O}_X)$ に属し、かつ
$D_\QCoh(\mathcal{O}_X)$ のすべての $M$ に対して
$\Hom(M, K) = 0$ となることである。Derived Categories, Lemma
\ref{derived-lemma-right-adjoint} を参照されたい。完全三角形
$$
E \to Rj_{U, *}E|_U \oplus Rj_{V, *}E|_V \to
Rj_{U \cap V, *}E|_{U \cap V} \to E[1]
$$
を考える。これは Cohomology, Lemma
\ref{cohomology-lemma-exact-sequence-j-star} の
$D(\mathcal{O}_X)$ における完全三角形である。
Derived Categories, Lemma \ref{derived-lemma-prepare-adjoint} により、
$Rj_{U, *}E|_U$、$Rj_{V, *}E|_V$、および
$Rj_{U \cap V, *}E|_{U \cap V}$ に対して所望の完全三角形を
構成すれば十分である。これは次の段落で論じる主張に帰着する。

\medskip\noindent
$j : U \to X$ を、補題が成り立つ準コンパクト開集合 $U$ に対応する
開埋入とする。$L$ を $D(\mathcal{O}_U)$ の対象とする。このとき
完全三角形
$$
E' \to Rj_*L \to K \to E'[1]
$$
が $D(\mathcal{O}_X)$ 内に存在し、$E'$ は
$D_\QCoh(\mathcal{O}_X)$ に属し、$D_\QCoh(\mathcal{O}_X)$ の
すべての $M$ に対して $\Hom(M, K) = 0$ である。
これを見るため、完全三角形
$$
L' \to L \to Q \to L'[1]
$$
を $D(\mathcal{O}_U)$ 内で選ぶ。ただし $L'$ は
$D_\QCoh(\mathcal{O}_U)$ に属し、$D_\QCoh(\mathcal{O}_U)$ の
すべての $N$ に対して $\Hom(N, Q) = 0$ である。
これは Derived Categories, Lemma \ref{derived-lemma-right-adjoint}
の主張が必要十分条件だから可能である。完全三角形
$$
Rj_*L' \to Rj_*L \to Rj_*Q \to Rj_*L'[1]
$$
を $D(\mathcal{O}_X)$ 内に得る。補題
\ref{perfect-lemma-quasi-coherence-direct-image} により
$Rj_*L'$ は $D_\QCoh(\mathcal{O}_X)$ に属する。一方、
$M$ が $D_\QCoh(\mathcal{O}_X)$ に属するなら
$$
\Hom(M, Rj_*Q) = \Hom(Lj^*M, Q) = 0
$$
である。実際、補題 \ref{perfect-lemma-quasi-coherence-pullback} により
$Lj^*M$ は $D_\QCoh(\mathcal{O}_U)$ に属する。これで証明が完了する。
\end{proof}

\begin{proof}[第二の証明]
Derived Categories, Proposition \ref{derived-proposition-brown}
により随伴は存在する。実際、その仮定は満たされる。
まず $D_\QCoh(\mathcal{O}_X)$ は直和をもち、直和は包含関手と可換する
（補題 \ref{perfect-lemma-quasi-coherence-direct-sums}）。
一方、$D_\QCoh(\mathcal{O}_X)$ はコンパクト生成される。
なぜなら定理 \ref{perfect-theorem-bondal-van-den-Bergh} の完全生成対象をもち、
命題 \ref{perfect-proposition-compact-is-perfect} により完全対象は
コンパクトだからである。
\end{proof}

\begin{lemma}
\label{perfect-lemma-pushforward-better-coherator}
$f : X \to Y$ をスキームの準コンパクトかつ準分離な射とする。
$X$ と $Y$ に対して包含関手 $D_\QCoh \to D$ の右随伴
$DQ_X$ と $DQ_Y$ が存在するなら
$$
Rf_* \circ DQ_X = DQ_Y \circ Rf_*
$$
\end{lemma}

\begin{proof}
補題 \ref{perfect-lemma-quasi-coherence-direct-image} により $Rf_*$ は
$D_\QCoh(\mathcal{O}_X)$ を $D_\QCoh(\mathcal{O}_Y)$ に送るので、
この主張には意味がある。同様に、補題
\ref{perfect-lemma-quasi-coherence-pullback} により $Lf^*$ は
$D_\QCoh(\mathcal{O}_Y)$ を $D_\QCoh(\mathcal{O}_X)$ に送る。
したがって $Rf_* \circ DQ_X$ と $DQ_Y \circ Rf_*$ はともに
$Lf^* : D_\QCoh(\mathcal{O}_Y) \to D(\mathcal{O}_X)$ の右随伴であり、
主張が従う。
\end{proof}

\begin{remark}
\label{perfect-remark-explain-consequence}
$X$ を準コンパクトかつ準分離なスキームとし、
$X = U \cup V$ とする。ただし $U$ と $V$ は準コンパクト開集合である。
補題 \ref{perfect-lemma-better-coherator} により関手
$DQ_X$、$DQ_U$、$DQ_V$、$DQ_{U \cap V}$ は存在する。さらに、
任意の $K \in D(\mathcal{O}_X)$ に対して標準的な完全三角形
$$
DQ_X(K) \to Rj_{U, *}DQ_U(K|_U) \oplus Rj_{V, *}DQ_V(K|_V)
\to Rj_{U \cap V, *}DQ_{U \cap V}(K|_{U \cap V}) \to
$$
がある。これは Cohomology, Lemma
\ref{cohomology-lemma-exact-sequence-j-star} の完全三角形に
完全関手 $DQ_X$ を適用し、補題
\ref{perfect-lemma-pushforward-better-coherator} を三回用いれば従う。
\end{remark}

\begin{lemma}
\label{perfect-lemma-boundedness-better-coherator}
$X$ を準コンパクトかつ準分離なスキームとする。補題
\ref{perfect-lemma-better-coherator} の関手 $DQ_X$ は次の有界性をもつ。
整数 $N = N(X)$ が存在し、$K$ が $D(\mathcal{O}_X)$ に属して、
$X$ のアフィン開集合 $U$ と $i \not \in [a, b]$ に対して
$H^i(U, K) = 0$ なら、$i \not \in [a, b + N]$ に対して
コホモロジー層 $H^i(DQ_X(K))$ は零である。
\end{lemma}

\begin{proof}
Cohomology of Schemes, Lemma \ref{coherent-lemma-induction-principle}
の帰納原理を用いて証明する。

\medskip\noindent
$X$ がアフィンなら、$RQ_X = DQ_X$ は $R\Gamma(X, K)$ に付随する
擬連接層の複体を取ることで与えられるので、補題は $N = 0$
として成り立つ。補題 \ref{perfect-lemma-affine-compare-bounded} と
\ref{perfect-lemma-affine-coherator} を参照されたい。

\medskip\noindent
$U, V$ を $X$ の準コンパクト開集合とし、補題が $U$、$V$、
$U \cap V$ に対して、それぞれ整数 $N(U)$、$N(V)$、
$N(U \cap V)$ を用いて成り立つと仮定する。$K$ が
$D(\mathcal{O}_X)$ に属し、すべてのアフィン開集合 $W \subset X$ と
すべての $i \not \in [a, b]$ に対して $H^i(W, K) = 0$ とする。
すると $K|_U$、$K|_V$、$K|_{U \cap V}$ も同じ性質をもつ。
したがって $RQ_U(K|_U)$、$RQ_V(K|_V)$、および
$RQ_{U \cap V}(K|_{U \cap V})$ のコホモロジー層は区間
$[a, b + \max(N(U), N(V), N(U \cap V))$ の外で消滅する。
補題 \ref{perfect-lemma-quasi-coherence-direct-image} により、関手
$Rj_{U, *}$、$Rj_{V, *}$、$Rj_{U \cap V, *}$ は $D_\QCoh$ 上で
有限コホモロジー次元をもつ。ゆえに整数 $N$ が存在し、
$Rj_{U, *}DQ_U(K|_U)$、$Rj_{V, *}DQ_V(K|_V)$、および
$Rj_{U \cap V, *}DQ_{U \cap V}(K|_{U \cap V})$ の
コホモロジー層は区間 $[a, b + N]$ の外で消滅する。
最後に注意 \ref{perfect-remark-explain-consequence} の完全三角形から結論する。
\end{proof}

\begin{example}
\label{perfect-example-inverse-limit-quasi-coherent}
$X$ を準コンパクトかつ準分離なスキームとし、
$(\mathcal{F}_n)$ を擬連接層の逆系とする。$DQ_X$ は右随伴なので
直積と可換し、したがって導来極限とも可換する。ゆえに
$$
DQ_X(R\lim \mathcal{F}_n) =
(R\lim\text{ in }D_\QCoh(\mathcal{O}_X))(\mathcal{F}_n)
$$
である。ここで最初の $R\lim$ は $D(\mathcal{O}_X)$ 内で取る。
これを $K = R\lim \mathcal{F}_n$ と書こう。
任意のアフィン開集合 $U \subset X$ に対して
$$
H^i(U, K) =
H^i(R\Gamma(U, R\lim \mathcal{F}_n)) =
H^i(R\lim R\Gamma(U, \mathcal{F}_n)) =
H^i(R\lim \Gamma(U, \mathcal{F}_n))
$$
である。実際、コホモロジーは導来極限と可換し、擬連接層
$\mathcal{F}_n$ はアフィン上で高次コホモロジーをもたない。
アーベル群の圏における $R\lim$ の計算により、
$i \in [0, 1]$ でない限り $H^i(U, K) = 0$ である。
したがって補題 \ref{perfect-lemma-boundedness-better-coherator} により、
$D_\QCoh(\mathcal{O}_X)$ 内の $R\lim$、すなわち上で見た
$DQ_X(K)$ は $D^b_\QCoh(\mathcal{O}_X)$ に属する。
\end{example}













\section{コホモロジーと基底変換 IV}
\label{perfect-section-cohomology-and-base-change-perfect}

\noindent
本節では Cohomology of Schemes, Section
\ref{coherent-section-cohomology-and-base-change-perfect}
の議論を続ける。まず、スキームの準コンパクトかつ準分離な射と、
擬連接なコホモロジー層をもつ複体に対する射影公式の
非常に一般的な形を与える。

\begin{lemma}
\label{perfect-lemma-cohomology-base-change}
$f : X \to Y$ をスキームの準コンパクトかつ準分離な射とする。
$D_\QCoh(\mathcal{O}_X)$ の $E$ と
$D_\QCoh(\mathcal{O}_Y)$ の $K$ に対して、Cohomology, Equation
(\ref{cohomology-equation-projection-formula-map}) で定義された写像
$$
Rf_*(E) \otimes_{\mathcal{O}_Y}^\mathbf{L} K
\longrightarrow
Rf_*(E \otimes_{\mathcal{O}_X}^\mathbf{L} Lf^*K)
$$
は同型である。
\end{lemma}

\begin{proof}
この写像が同型であることは $Y$ 上局所的に確認できる。
したがって $Y$ がアフィンの場合に帰着する。

\medskip\noindent
$K = \bigoplus K_i$ がいくつかの複体
$K_i \in D_\QCoh(\mathcal{O}_Y)$ の直和であるとする。
各 $K_i$ に対して主張が成り立つなら、$K$ に対しても成り立つ。
実際、構成により関手 $Lf^*$ と $\otimes^\mathbf{L}$ は直和を保ち、
補題 \ref{perfect-lemma-quasi-coherence-pushforward-direct-sums} により
$Rf_*$ は（擬連接なコホモロジー層をもつ複体について）直和と可換する。
さらに $K \to L \to M \to K[1]$ が $D_\QCoh(Y)$ の完全三角形なら、
$K, L, M$ のうち二つに対して補題の主張が成り立てば、三つ目に対しても
成り立つ（現れる関手はいずれも三角圏の完全関手だからである）。

\medskip\noindent
$Y$ がアフィン、$Y = \Spec(A)$ であるとする。関手
$\widetilde{\ } : D(A) \to D_\QCoh(\mathcal{O}_Y)$ は同値である
（補題 \ref{perfect-lemma-affine-compare-bounded}）。
$K \in D(A)$ に対し、$\widetilde{K}$ について補題の主張が
成り立つという性質を $T$ とする。上の議論と
More on Algebra, Remark \ref{more-algebra-remark-P-resolution} により、
$A[k]$ が $T$ を満たすことを示せば十分である。構造層のシフトに
対して補題の主張は明らかなので、これで証明が完了する。
\end{proof}

\begin{definition}
\label{perfect-definition-tor-independent}
$S$ をスキームとし、$X$, $Y$ を $S$ 上のスキームとする。
$s \in S$ の同じ点に写るすべての $x \in X$ と $y \in Y$ に対して、
環 $\mathcal{O}_{X, x}$ と $\mathcal{O}_{Y, y}$ が
$\mathcal{O}_{S, s}$ 上 Tor 独立であるとき、$X$ と $Y$ は
{\it $S$ 上 Tor 独立}であるという（More on Algebra, Definition
\ref{more-algebra-definition-tor-independent} を参照）。
\end{definition}

\begin{lemma}
\label{perfect-lemma-tor-independent}
$f : X \to S$ と $g : Y \to S$ をスキームの射とする。
次は同値である。
\begin{enumerate}
\item $X$ と $Y$ は $S$ 上 tor 独立である。
\item $f(U) \subset W$ かつ $g(V) \subset W$ を満たすすべての
アフィン開集合 $U \subset X$, $V \subset Y$, $W \subset S$ に対し、
環 $\mathcal{O}_X(U)$ と $\mathcal{O}_Y(V)$ は
$\mathcal{O}_S(W)$ 上 tor 独立である。
\item アフィン開被覆 $S = \bigcup W_i$ と、各 $i$ に対して
アフィン開被覆 $f^{-1}(W_i) = \bigcup U_{ij}$ および
$g^{-1}(W_i) = \bigcup V_{ik}$ が存在し、すべての $i, j, k$ に対して
環 $\mathcal{O}_X(U_{ij})$ と $\mathcal{O}_Y(V_{ik})$ は
$\mathcal{O}_S(W_i)$ 上 tor 独立である。
\end{enumerate}
\end{lemma}

\begin{proof}
省略する。ヒント：More on Algebra, Lemma
\ref{more-algebra-lemma-tor-independent} を用いる。
\end{proof}

\begin{lemma}
\label{perfect-lemma-flat-base-change-tor-independent}
$X \to S$ と $Y \to S$ をスキームの射とする。
$S' \to S$ をスキームの射とし、$X' = X \times_S S'$ および
$Y' = Y \times_S S'$ と記す。$X$ と $Y$ が $S$ 上 tor 独立で、
$S' \to S$ が平坦なら、$X'$ と $Y'$ は $S'$ 上 tor 独立である。
\end{lemma}

\begin{proof}
省略する。ヒント：補題 \ref{perfect-lemma-tor-independent} を用い、
アフィン開集合上では More on Algebra, Lemma
\ref{more-algebra-lemma-flat-base-change-tor-independent} を用いる。
\end{proof}

\begin{lemma}
\label{perfect-lemma-compare-base-change}
$g : S' \to S$ をスキームの射とし、
$f : X \to S$ を準コンパクトかつ準分離な射とする。
基底変換図式
$$
\xymatrix{
X' \ar[r]_{g'} \ar[d]_{f'} &
X \ar[d]^f \\
S' \ar[r]^g &
S
}
$$
を考える。$X$ と $S'$ が $S$ 上 Tor 独立なら、
すべての $E \in D_\QCoh(\mathcal{O}_X)$ に対して標準射
$Lg^*Rf_*E \to Rf'_*L(g')^*E$ は同型である。
\end{lemma}

\begin{proof}
$D(\mathcal{O}_X)$ の任意の対象 $E$ に対し、Cohomology, Remark
\ref{cohomology-remark-base-change} から標準的な基底変換写像
$Lg^*Rf_*E \to Rf'_*L(g')^*E$ を得る。これが同型であることは
$S'$ 上局所的に確認できる。したがって
$g : S' \to S$ はアフィン・スキームの射と仮定してよい。
特に $g$ はアフィンであり、
$$
Rg_*Lg^*Rf_*E \to Rg_*Rf'_*L(g')^*E = Rf_*(Rg'_* L(g')^* E)
$$
が同型であることを示せば十分である。補題
\ref{perfect-lemma-affine-morphism} を参照されたい
（対象 $Rf'_*L(g')^*E$ と $Lg^*Rf_*E$ が擬連接なコホモロジー層を
もつことを見るには、補題 \ref{perfect-lemma-quasi-coherence-pullback}、
\ref{perfect-lemma-quasi-coherence-tensor-product}、および
\ref{perfect-lemma-quasi-coherence-direct-image} を用いる）。
なお $g'$ もアフィンである（Morphisms, Lemma
\ref{morphisms-lemma-base-change-affine}）。
補題 \ref{perfect-lemma-affine-morphism-pull-push} により、この写像は
$$
Rf_*E \otimes_{\mathcal{O}_S}^\mathbf{L} g_*\mathcal{O}_{S'}
\longrightarrow
Rf_*(E \otimes_{\mathcal{O}_X}^\mathbf{L} g'_*\mathcal{O}_{X'})
$$
となる。アフィン基底変換により
$g'_*\mathcal{O}_{X'} = f^*g_*\mathcal{O}_{S'}$ である
（Cohomology of Schemes, Lemma
\ref{coherent-lemma-affine-base-change} を参照）。
したがって補題 \ref{perfect-lemma-cohomology-base-change} により、
$Lf^*g_*\mathcal{O}_{S'} = f^*g_*\mathcal{O}_{S'}$ を示せば十分である。
これは $X$ と $S'$ が $S$ 上 Tor 独立であるという仮定から従う。
実際、これは $X$ 上局所的に確認できるので、$X$ もアフィンと
仮定してよい。$X = \Spec(A)$、$S = \Spec(R)$、
$S' = \Spec(R')$ と書く。この仮定から $A$ と $R'$ は $R$ 上
Tor 独立である（More on Algebra, Lemma
\ref{more-algebra-lemma-tor-independent}）。すなわち $i > 0$ に対して
$\text{Tor}_i^R(A, R') = 0$ である。言い換えると
$A \otimes_R^\mathbf{L} R' = A \otimes_R R'$ であり、これはまさに
$Lf^*g_*\mathcal{O}_{S'} = f^*g_*\mathcal{O}_{S'}$ を意味する
（補題 \ref{perfect-lemma-quasi-coherence-pullback} を用いる）。
\end{proof}

\noindent
次の補題は双対化複体の章で用いる。

\begin{lemma}
\label{perfect-lemma-affine-morphism-and-hom-out-of-perfect}
準コンパクトかつ準分離なスキームのカルテシアン正方形
$$
\xymatrix{
X' \ar[r]_{g'} \ar[d]_{f'} & X \ar[d]^f \\
S' \ar[r]^g & S
}
$$
を考える。$g$ と $f$ は Tor 独立であり、
$S = \Spec(R)$、$S' = \Spec(R')$ はアフィンであると仮定する。
$M, K \in D(\mathcal{O}_X)$ に対して標準写像
$$
R\Hom_X(M, K) \otimes^\mathbf{L}_R R'
\longrightarrow
R\Hom_{X'}(L(g')^*M, L(g')^*K)
$$
は次の二つの場合に $D(R')$ 内の同型である。
\begin{enumerate}
\item $M \in D(\mathcal{O}_X)$ は完全で、$K \in D_\QCoh(X)$ である。
\item $M \in D(\mathcal{O}_X)$ は擬連接で、
$K \in D_\QCoh^+(X)$ であり、$R'$ は $R$ 上有限 tor 次元をもつ。
\end{enumerate}
\end{lemma}

\begin{proof}
大域 Hom 複体の圏 $D(\Gamma(X, \mathcal{O}_X))$ には標準写像\newline
$R\Hom_X(M, K) \to R\Hom_{X'}(L(g')^*M, L(g')^*K)$
がある。Cohomology, Section \ref{cohomology-section-global-RHom}
を参照されたい。スカラーを制限すれば、これを $D(R)$ の写像と
みなせる。そこでスカラー制限と $- \otimes_R^\mathbf{L} R'$ の
随伴性を用いて、補題に表示された写像を得る。写像を定義したので、
アーベル群の導来圏で同型であることを示せば十分である。

\medskip\noindent
右辺は、補題 \ref{perfect-lemma-affine-morphism-pull-push} により
$$
R\Hom_X(M, R(g')_*L(g')^*K) =
R\Hom_X(M, K \otimes_{\mathcal{O}_X}^\mathbf{L} g'_*\mathcal{O}_{X'})
$$
に等しい。どちらの場合にも補題
\ref{perfect-lemma-quasi-coherence-internal-hom} により複体
$R\SheafHom(M, K)$ は $D_\QCoh(\mathcal{O}_X)$ の対象である。
自然な写像
$$
R\SheafHom(M, K) \otimes_{\mathcal{O}_X}^\mathbf{L} g'_*\mathcal{O}_{X'}
\longrightarrow
R\SheafHom(M, K \otimes_{\mathcal{O}_X}^\mathbf{L} g'_*\mathcal{O}_{X'})
$$
があり、補題 \ref{perfect-lemma-internal-hom-evaluate-tensor-isomorphism}
により、どちらの場合にも同型である。この補題を (2) の場合に
適用できることを見るため、
$g'_*\mathcal{O}_{X'} = Rg'_*\mathcal{O}_{X'} =
Lf^*g_*\mathcal{O}_{S'}$
であることに注意する。第二の等号は補題
\ref{perfect-lemma-compare-base-change} による。補題
\ref{perfect-lemma-tor-dimension-affine} と Cohomology, Lemma
\ref{cohomology-lemma-tor-amplitude-pullback} を用いると、
$g'_*\mathcal{O}_{X'}$ は有限 Tor 次元をもつと分かる。
したがってどちらの場合にも、$K$ を $R\SheafHom(M, K)$ で
置き換えることにより、
$$
R\Gamma(X, K) \otimes^\mathbf{L}_A A' \longrightarrow
R\Gamma(X, K \otimes^\mathbf{L}_{\mathcal{O}_X} g'_*\mathcal{O}_{X'})
$$
が同型であることに帰着する。補題
\ref{perfect-lemma-affine-morphism-pull-push} により、左辺は
$R\Gamma(X', L(g')^*K)$ に等しい。ゆえに結果は補題
\ref{perfect-lemma-compare-base-change} から従う。
\end{proof}

\begin{remark}
\label{perfect-remark-multiplication-map}
補題 \ref{perfect-lemma-affine-morphism-and-hom-out-of-perfect} の記法を用いる。
図式
$$
\xymatrix{
R\Hom_X(M, Rg'_*L) \otimes_R^\mathbf{L} R' \ar[r] \ar[d]_\mu &
R\Hom_{X'}(L(g')^*M, L(g')^*Rg'_*L) \ar[d]^a \\
R\Hom_X(M, R(g')_*L) \ar@{=}[r] &
R\Hom_{X'}(L(g')^*M, L)
}
$$
は可換である。ここで上の水平射は補題の写像、$\mu$ は乗法写像であり、
$a$ は随伴写像 $L(g')^*Rg'_*L \to L$ から来る。乗法写像は、
任意の $K' \in D(R')$ に対する随伴写像
$K' \otimes_R^\mathbf{L} R' \to K'$ である。
\end{remark}

\begin{lemma}
\label{perfect-lemma-tor-independence-and-tor-amplitude}
スキームのカルテシアン正方形
$$
\xymatrix{
X' \ar[r]_{g'} \ar[d]_{f'} & X \ar[d]^f \\
S' \ar[r]^g & S
}
$$
を考える。$g$ と $f$ は Tor 独立であると仮定する。
\begin{enumerate}
\item $E \in D(\mathcal{O}_X)$ が $f^{-1}\mathcal{O}_S$-加群の
複体として $[a, b]$ 内の tor 振幅をもつなら、
$L(g')^*E$ は $f^{-1}\mathcal{O}_{S'}$-加群の複体として
$[a, b]$ 内の tor 振幅をもつ。
\item $\mathcal{G}$ が $S$ 上平坦な $\mathcal{O}_X$-加群なら、
$L(g')^*\mathcal{G} = (g')^*\mathcal{G}$ である。
\end{enumerate}
\end{lemma}

\begin{proof}
tor 次元は茎で計算できる。Cohomology, Lemma
\ref{cohomology-lemma-tor-amplitude-stalk} を参照されたい。
$x' \in X'$ の像を $x \in X$ とすると
$$
(L(g')^*E)_{x'} =
E_x \otimes_{\mathcal{O}_{X, x}}^\mathbf{L} \mathcal{O}_{X', x'}
$$
である。$x'$ と $x$ の像をそれぞれ $s' \in S'$ と $s \in S$
とする。$X$ と $S'$ は $S$ 上 tor 独立なので、
More on Algebra, Lemma \ref{more-algebra-lemma-base-change-comparison}
を適用すると、表示式の右辺は $D(\mathcal{O}_{S', s'})$ 内で
$E_x \otimes_{\mathcal{O}_{S, s}}^\mathbf{L} \mathcal{O}_{S', s'}$
に等しいと分かる。したがって (1) は More on Algebra, Lemma
\ref{more-algebra-lemma-pull-tor-amplitude} から従う。
(2) を見るには、$\mathcal{G}$ の平坦性が
$\mathcal{G}[0]$ が $[0, 0]$ 内の tor 振幅をもつことと同値である
ことに注意し、(1) を適用すればよい。
\end{proof}

\begin{lemma}
\label{perfect-lemma-compare-base-change-closed-immersion}
スキームのカルテシアン図式
$$
\xymatrix{
Z' \ar[r]_{i'} \ar[d]_g & X' \ar[d]^f \\
Z \ar[r]^i & X
}
$$
を考える。ただし $i$ は閉埋入とする。$Z$ と $X'$ が
$X$ 上 tor 独立なら、$D(\mathcal{O}_Z) \to D(\mathcal{O}_{X'})$
の関手として $Ri'_* \circ Lg^* = Lf^* \circ Ri_*$ である。
\end{lemma}

\begin{proof}
この補題はすべての $K \in D(\mathcal{O}_Z)$ に対して
成り立つべき主張であることに注意する。$i_*$ と $i'_*$ は完全関手
なので、$Ri_*$ と $Ri'_*$ は任意の代表に $i_*$ と $i'_*$ を
適用して計算される。したがって基底変換写像
$$
Lf^*(Ri_*(K)) \longrightarrow Ri'_*(Lg^*(K))
$$
を、像が $z \in Z$ である点 $z' \in Z'$ の茎で取ると
$$
K_z \otimes_{\mathcal{O}_{X, z}}^\mathbf{L} \mathcal{O}_{X', z'}
\longrightarrow
K_z \otimes_{\mathcal{O}_{Z, z}}^\mathbf{L} \mathcal{O}_{Z', z'}
$$
となる。この写像は More on Algebra, Lemma
\ref{more-algebra-lemma-base-change-comparison} と仮定した
tor 独立性により同型である。
\end{proof}








\section{K\"unneth 公式 II}
\label{perfect-section-kunneth}

\noindent
基底が体である場合については
Varieties, Section \ref{varieties-section-kunneth} を参照されたい。
スキームのカルテシアン図式
$$
\xymatrix{
& X \times_S Y \ar[ld]^p \ar[rd]_q \ar[dd]^f \\
X \ar[rd]_a & & Y \ar[ld]^b \\
& S
}
$$
を考える。$K \in D(\mathcal{O}_X)$ および
$M \in D(\mathcal{O}_Y)$ とする。標準写像
\begin{equation}
\label{perfect-equation-kunneth}
Ra_*K \otimes_{\mathcal{O}_S}^\mathbf{L} Rb_*M
\longrightarrow
Rf_*(Lp^*K \otimes_{\mathcal{O}_{X \times_S Y}}^\mathbf{L} Lq^*M)
\end{equation}
がある。実際、写像
$Ra_*K \to Ra_*Rp_* Lp^*K = Rf_*Lp^*K$ と
$Rb_*M \to Rb_*Rq_* Lq^*M = Rf_*Lq^*M$ を用い、続いて相対カップ積
（Cohomology, Remark \ref{cohomology-remark-cup-product}）を用いればよい。

\medskip\noindent
$A = \Gamma(S, \mathcal{O}_S)$ とおく。$D(A)$ には大域 K\"unneth 写像
\begin{equation}
\label{perfect-equation-kunneth-global}
R\Gamma(X, K) \otimes_A^\mathbf{L} R\Gamma(Y, M)
\longrightarrow
R\Gamma(X \times_S Y,
Lp^*K \otimes_{\mathcal{O}_{X \times_S Y}}^\mathbf{L} Lq^*M)
\end{equation}
がある。この写像は引き戻し写像
$R\Gamma(X, K) \to R\Gamma(X \times_S Y, Lp^*K)$、
$R\Gamma(Y, M) \to R\Gamma(X \times_S Y, Lq^*M)$、および
Cohomology, Section \ref{cohomology-section-cup-product}
で構成されたカップ積を用いて構成される。

\begin{lemma}
\label{perfect-lemma-kunneth}
上の状況で $a$ と $b$ が準コンパクトかつ準分離であり、
$X$ と $Y$ が $S$ 上 tor 独立なら、
$K \in D_\QCoh(\mathcal{O}_X)$ と
$M \in D_\QCoh(\mathcal{O}_Y)$ に対して
(\ref{perfect-equation-kunneth}) は同型である。さらに
$S = \Spec(A)$ がアフィンなら、写像
(\ref{perfect-equation-kunneth-global}) は同型である。
\end{lemma}

\begin{proof}[第一の証明]
これは次の同型列から従う。
\begin{align*}
Rf_*(Lp^*K \otimes_{\mathcal{O}_{X \times_S Y}}^\mathbf{L} Lq^*M)
& =
Ra_*Rp_*(Lp^*K \otimes_{\mathcal{O}_{X \times_S Y}}^\mathbf{L} Lq^*M) \\
& =
Ra_*(K \otimes_{\mathcal{O}_X}^\mathbf{L} Rp_*Lq^*M) \\
& =
Ra_*(K \otimes_{\mathcal{O}_X}^\mathbf{L} La^*Rb_*M) \\
& =
Ra_*K \otimes_{\mathcal{O}_S}^\mathbf{L} Rb_*M
\end{align*}
第一の等号は $f = a \circ p$ だからである。第二の等号は補題
\ref{perfect-lemma-cohomology-base-change}、第三の等号は補題
\ref{perfect-lemma-compare-base-change}、第四の等号は補題
\ref{perfect-lemma-cohomology-base-change} による。これらの同型の合成が
写像 (\ref{perfect-equation-kunneth}) と一致することの確認は省略する。
$S$ がアフィンなら、射 (\ref{perfect-equation-kunneth-global}) の始域と終域は
(\ref{perfect-equation-kunneth}) の始域と終域に $R\Gamma(S, -)$ を適用して
得られる。これで最後の主張を得る。詳細は省略する。
\end{proof}

\begin{proof}[第二の証明]
相対カップ積の構成から直ちに分かるように、射
(\ref{perfect-equation-kunneth}) の構成は $S$ の開部分スキームへの制限と
両立する。したがって $S$ がアフィンのときに
(\ref{perfect-equation-kunneth}) が同型であることを示せば十分である。

\medskip\noindent
$S = \Spec(A)$ がアフィンであると仮定する。Leray により
$R\Gamma(S, Rf_*K) = R\Gamma(X, K)$ であり、他の場合も同様である。
Cohomology, Lemma \ref{cohomology-lemma-compose-cup-product} により、
写像 (\ref{perfect-equation-kunneth}) に $R\Gamma(S, -)$ を取ると
写像 (\ref{perfect-equation-kunneth-global}) が誘導される。一方、補題
\ref{perfect-lemma-quasi-coherence-direct-image} と
\ref{perfect-lemma-quasi-coherence-tensor-product} により、
写像 (\ref{perfect-equation-kunneth}) の始域と終域は
$D_\QCoh(\mathcal{O}_S)$ に属する。したがって補題
\ref{perfect-lemma-affine-compare-bounded} により、
(\ref{perfect-equation-kunneth-global}) が同型であることを示せば十分である。

\medskip\noindent
$S = \Spec(A)$、$X = \Spec(B)$、$Y = \Spec(C)$ がすべて
アフィンであると仮定する。以下、補題
\ref{perfect-lemma-affine-compare-bounded} は断りなく用いる。この場合、
項が平坦な $B$-加群の K-平坦複体 $K^\bullet$ で、$K$ が
$\widetilde{K}^\bullet$ により表されるものを選べる。同様に、
項が平坦な $C$-加群の K-平坦複体 $M^\bullet$ で、$M$ が
$\widetilde{M}^\bullet$ により表されるものを選べる。
More on Algebra, Lemma \ref{more-algebra-lemma-K-flat-resolution}
を参照されたい。このとき $\widetilde{K}^\bullet$ は
$\mathcal{O}_X$-加群の K-平坦複体であり、
$\widetilde{M}^\bullet$ についても同様である。補題
\ref{perfect-lemma-affine-K-flat} を参照されたい。したがって
$La^*K$ は
$$
a^*\widetilde{K}^\bullet = \widetilde{K^\bullet \otimes_A C}
$$
により表され、$Lb^*M$ についても同様である。さらにこれは、
上で引用した補題と More on Algebra, Lemma
\ref{more-algebra-lemma-base-change-K-flat} により
$\mathcal{O}_{X \times_S Y}$-加群の K-平坦複体である。
したがって
$$
\text{Tot}((K^\bullet \otimes_A C) \otimes_{B \otimes_A C}
(B \otimes_A M^\bullet)) =
\text{Tot}(K^\bullet \otimes_A M^\bullet)
$$
に付随する $\mathcal{O}_{X \times_S Y}$-加群の複体は、
$X \times_S Y$ の導来圏において
$La^*K \otimes_{\mathcal{O}_{X \times_S Y}}^\mathbf{L} Lb^*M$
を表す。大域切断を取ると
$\text{Tot}(K^\bullet \otimes_A M^\bullet)$ を得るが、これはもちろん
$R\Gamma(X, K) \otimes_A^\mathbf{L} R\Gamma(Y, M)$ を表す複体でもある。
この同型がカップ積によって与えられることは、
Cohomology, Section \ref{cohomology-section-cup-product}
における真正のカップ積と素朴なカップ積との関係から従う
（特に Cohomology, Lemma
\ref{cohomology-lemma-cup-compatible-with-naive} とその後の議論を参照）。

\medskip\noindent
$S = \Spec(A)$ と $Y$ がアフィンであると仮定する。
Cohomology of Schemes, Lemma \ref{coherent-lemma-induction-principle}
の帰納原理を用いて主張を証明する。そのためには次を示せばよい。
$X = U \cup V$ が準コンパクトな $U$ と $V$ による開被覆であり、
$S$ 上の $U$ と $Y$ に対する写像 (\ref{perfect-equation-kunneth-global})
$$
R\Gamma(U, K) \otimes_A^\mathbf{L} R\Gamma(Y, M)
\longrightarrow
R\Gamma(U \times_S Y,
Lp^*K \otimes_{\mathcal{O}_{X \times_S Y}}^\mathbf{L} Lq^*M)
$$
$S$ 上の $V$ と $Y$ に対する写像 (\ref{perfect-equation-kunneth-global})
$$
R\Gamma(V, K) \otimes_A^\mathbf{L} R\Gamma(Y, M)
\longrightarrow
R\Gamma(V \times_S Y,
Lp^*K \otimes_{\mathcal{O}_{X \times_S Y}}^\mathbf{L} Lq^*M)
$$
および $S$ 上の $U \cap V$ と $Y$ に対する写像
(\ref{perfect-equation-kunneth-global})
$$
R\Gamma(U \cap V, K) \otimes_A^\mathbf{L} R\Gamma(Y, M)
\longrightarrow
R\Gamma((U \cap V) \times_S Y,
Lp^*K \otimes_{\mathcal{O}_{X \times_S Y}}^\mathbf{L} Lq^*M)
$$
が同型なら、$S$ 上の $X$ と $Y$ に対する写像
(\ref{perfect-equation-kunneth-global}) も同型である。
実際、Cohomology, Lemma \ref{cohomology-lemma-mayer-vietoris-cup}
により、これらの写像は
(\ref{perfect-equation-kunneth-global}) を最後の辺とする完全三角形間の
写像をなす。したがって Derived Categories, Lemma
\ref{derived-lemma-third-isomorphism-triangle} により結論を得る。

\medskip\noindent
$S = \Spec(A)$ がアフィンであると仮定する。証明を終えるには、
$Y$ に Cohomology of Schemes, Lemma
\ref{coherent-lemma-induction-principle} の帰納原理を用いればよい。
実際、上により $Y$ がアフィンなら写像は同型であることが
すでに分かっている。残りの議論は $X$ と $Y$ の役割を交換すれば
前段落とまったく同じである。
\end{proof}

\begin{lemma}
\label{perfect-lemma-cohomology-de-rham-base-change}
$a : X \to S$ をスキームの準コンパクトかつ準分離な射とする。
$\mathcal{F}^\bullet$ を $a^{-1}\mathcal{O}_S$-加群の局所有界複体とする。
すべての $n \in \mathbf{Z}$ に対して、層 $\mathcal{F}^n$ は
平坦な $a^{-1}\mathcal{O}_S$-加群であり、$\mathcal{F}^n$ は与えられた
$a^{-1}\mathcal{O}_S$-加群構造と両立する擬連接
$\mathcal{O}_X$-加群の構造をもつと仮定する
（ただし複体 $\mathcal{F}^\bullet$ の微分は
$\mathcal{O}_X$-線型でなくてもよい）。このとき次が成り立つ。
\begin{enumerate}
\item $Ra_*\mathcal{F}^\bullet$ は局所有界である。
\item $Ra_*\mathcal{F}^\bullet$ は $D_\QCoh(\mathcal{O}_S)$ に属する。
\item $Ra_*\mathcal{F}^\bullet$ は局所的に有限 tor 次元をもつ。
\item $\mathcal{G} \otimes_{\mathcal{O}_S}^\mathbf{L} Ra_*\mathcal{F}^\bullet =
Ra_*(a^{-1}\mathcal{G} \otimes_{a^{-1}\mathcal{O}_S} \mathcal{F}^\bullet)$
が $\mathcal{G} \in \QCoh(\mathcal{O}_S)$ に対して成り立つ。
\item $K \otimes_{\mathcal{O}_S}^\mathbf{L} Ra_*\mathcal{F}^\bullet =
Ra_*(a^{-1}K \otimes_{a^{-1}\mathcal{O}_S}^\mathbf{L} \mathcal{F}^\bullet)$
が $K \in D_\QCoh(\mathcal{O}_S)$ に対して成り立つ。
\end{enumerate}
\end{lemma}

\begin{proof}
(1)、(2)、(3) は $S$ 上局所的なので、$S$ はアフィンであると
仮定してよい。$a$ は準コンパクトだから $X$ は準コンパクトである。
$\mathcal{F}^\bullet$ は局所有界なので
$\mathcal{F}^\bullet$ は有界である。

\medskip\noindent
(1) と (2) には第一スペクトル系列
$R^pa_*\mathcal{F}^q \Rightarrow R^{p + q}a_*\mathcal{F}^\bullet$
を用いる。これは Derived Categories, Lemma
\ref{derived-lemma-two-ss-complex-functor} のスペクトル系列である。
Cohomology of Schemes, Lemma
\ref{coherent-lemma-quasi-coherence-higher-direct-images}
と Homology, Lemma \ref{homology-lemma-biregular-ss-converges}
を合わせれば結論を得る。

\medskip\noindent
(3) を Cohomology of Schemes, Lemma
\ref{coherent-lemma-induction-principle} の帰納原理により証明する。
すなわち、$X$ の準コンパクト開集合 $U$ に対して
$R(a|_U)_*(\mathcal{F}^\bullet|_U)$ が有限 tor 次元をもつという
条件を考える。$U, V$ が $X$ の準コンパクト開集合なら、相対
Mayer--Vietoris 完全三角形
$$
R(a|_{U \cup V})_*\mathcal{F}^\bullet|_{U \cup V} \to
R(a|_U)_*\mathcal{F}^\bullet|_U \oplus
R(a|_V)_*\mathcal{F}^\bullet|_V \to
R(a|_{U \cap V})_*\mathcal{F}^\bullet|_{U \cap V} \to
$$
がある。これは Cohomology, Lemma
\ref{cohomology-lemma-unbounded-relative-mayer-vietoris} による。
完全三角形における tor 振幅の挙動
（Cohomology, Lemma \ref{cohomology-lemma-cone-tor-amplitude} を参照）
により、$U$、$V$、$U \cap V$ に対して結果が分かれば
$U \cup V$ に対しても成り立つ。したがって $X$ がアフィンの場合に
帰着する。この場合
$$
Ra_*\mathcal{F}^\bullet = a_*\mathcal{F}^\bullet
$$
である。これは Leray の非輪状性補題
（Derived Categories, Lemma \ref{derived-lemma-leray-acyclicity}）と、
アフィン射による擬連接加群の高次直像の消滅
（Cohomology of Schemes, Lemma
\ref{coherent-lemma-relative-affine-vanishing}）による。
仮定により $\mathcal{F}^n$ は $S$-平坦で、$X$ はアフィンなので、
すべての $n$ に対して加群 $a_*\mathcal{F}^n$ は平坦である。
ゆえに $a_*\mathcal{F}^\bullet$ は平坦
$\mathcal{O}_S$-加群の有界複体であり、有限 tor 次元をもつ。

\medskip\noindent
(5) を証明する。
$a' : (X, a^{-1}\mathcal{O}_S) \to (S, \mathcal{O}_S)$
を環付き空間の明らかな平坦射と記す。(5) の主張は
$$
K \otimes_{\mathcal{O}_S}^\mathbf{L} Ra'_*\mathcal{F}^\bullet =
Ra'_*(L(a')^*K \otimes_{a^{-1}\mathcal{O}_S}^\mathbf{L}
\mathcal{F}^\bullet)
$$
である。したがって Cohomology, Equation
(\ref{cohomology-equation-projection-formula-map}) は左辺から右辺への
関手的な写像を与え、これが同型であることを示したい。
この問題は $S$ 上局所的なので、$S$ はアフィンと仮定してよい。
残りの証明は、関手
$K \mapsto Ra'_*(L(a')^*K \otimes_{a^{-1}\mathcal{O}_S}^\mathbf{L}
\mathcal{F}^\bullet)$
が直和と可換することも示す必要がある点を除き、補題
\ref{perfect-lemma-cohomology-base-change} の証明と{\it まったく}同じである。
ここで $\mathcal{F}^n$ が擬連接 $\mathcal{O}_X$-加群の構造を
もつことを用いる。実際、関手
$K \mapsto L(a')^*K
\otimes_{a^{-1}\mathcal{O}_S}^\mathbf{L} \mathcal{F}^\bullet$
は任意の直和と可換する。次に $\mathcal{F}^\bullet$ が一つの擬連接
$\mathcal{O}_X$-加群のみからなり、
$\mathcal{F}^\bullet = \mathcal{F}^n[-n]$ なら
$L(a')^*G
\otimes_{a^{-1}\mathcal{O}_S}^\mathbf{L} \mathcal{F}^\bullet =
La^*K \otimes_{\mathcal{O}_X}^\mathbf{L} \mathcal{F}^n[-n]$
である。Cohomology, Lemma
\ref{cohomology-lemma-variant-derived-pullback} を参照されたい。
したがってこの場合、直和との可換性は補題
\ref{perfect-lemma-quasi-coherence-pushforward-direct-sums} から従う。
一般の場合、$S$ はアフィン（したがって $X$ は準コンパクト）で、
$\mathcal{F}^\bullet$ は局所有界なので
$$
\mathcal{F}^\bullet = (\mathcal{F}^a \to \ldots \to \mathcal{F}^b)
$$
は有界である。$b - a$ に関する帰納法を用い、完全三角形
$$
\mathcal{F}^b[-b] \to (\mathcal{F}^a \to \ldots \to \mathcal{F}^b)
\to (\mathcal{F}^a \to \ldots \to \mathcal{F}^{b - 1}) \to
\mathcal{F}^b[-b + 1]
$$
を考えれば、この部分の証明は完了する。いくつかの詳細は省略する。

\medskip\noindent
(4) を証明する。
$a' : (X, a^{-1}\mathcal{O}_S) \to (S, \mathcal{O}_S)$
を上の通りとする。$\mathcal{F}^\bullet$ は平坦な
$a^{-1}\mathcal{O}_S$-加群の局所有界複体なので、複体
$a^{-1}\mathcal{G} \otimes_{a^{-1}\mathcal{O}_S} \mathcal{F}^\bullet$
は
$L(a')^*\mathcal{G}
\otimes_{a^{-1}\mathcal{O}_S}^\mathbf{L}
\mathcal{F}^\bullet$
を $D(a^{-1}\mathcal{O}_S)$ 内で表す。したがって (4) は (5) から従う。
\end{proof}

\begin{lemma}
\label{perfect-lemma-K-flat}
$f : X \to Y$ をスキームの射とし、$Y = \Spec(A)$ はアフィンとする。
$\mathcal{U} : X = \bigcup_{i \in I} U_i$ を有限アフィン開被覆で、
すべての有限交叉
$U_{i_0 \ldots i_p} = U_{i_0} \cap \ldots \cap U_{i_p}$
がアフィンであるものとする。$\mathcal{F}^\bullet$ を
$f^{-1}\mathcal{O}_Y$-加群の有界複体とする。すべての
$n \in \mathbf{Z}$ に対して、層 $\mathcal{F}^n$ は平坦な
$f^{-1}\mathcal{O}_Y$-加群であり、$\mathcal{F}^n$ は与えられた
$p^{-1}\mathcal{O}_Y$-加群構造と両立する擬連接
$\mathcal{O}_X$-加群の構造をもつと仮定する
（ただし複体 $\mathcal{F}^\bullet$ の微分は
$\mathcal{O}_X$-線型でなくてもよい）。このとき複体
$\text{Tot}(\check{\mathcal{C}}^\bullet(\mathcal{U}, \mathcal{F}^\bullet))$
は $A$-加群の複体として K-平坦である。
\end{lemma}

\begin{proof}
次のように書ける。
$$
\mathcal{F}^\bullet = (\mathcal{F}^a \to \ldots \to \mathcal{F}^b)
$$
$b - a$ に関する帰納法を用い、完全三角形
$$
\mathcal{F}^b[-b] \to (\mathcal{F}^a \to \ldots \to \mathcal{F}^b)
\to (\mathcal{F}^a \to \ldots \to \mathcal{F}^{b - 1}) \to
\mathcal{F}^b[-b + 1]
$$
と More on Algebra, Lemma
\ref{more-algebra-lemma-K-flat-two-out-of-three} を用いると、
$\mathcal{F}^\bullet$ が次数 $0$ に置かれた一つの擬連接
$\mathcal{O}_X$-加群 $\mathcal{F}$ だけからなる場合に帰着する。
この場合、$\mathcal{F}$ と $\mathcal{U}$ の {\v C}ech 複体は
交代 {\v C}ech 複体とホモトピー同値である。Cohomology, Lemma
\ref{cohomology-lemma-alternating-usual} を参照されたい。
$U_{i_0 \ldots i_p}$ は常にアフィンなので、
$\mathcal{F}(U_{i_0 \ldots i_p})$ は $A$-平坦である。したがって
$\check{\mathcal{C}}_{alt}^\bullet(\mathcal{U}, \mathcal{F})$
は平坦 $A$-加群の有界複体であり、More on Algebra, Lemma
\ref{more-algebra-lemma-derived-tor-quasi-isomorphism} により
K-平坦である。
\end{proof}

\noindent
$X, Y, S, a, b, p, q, f$ を本節の導入と同じものとする。
$\mathcal{F}$ を $\mathcal{O}_X$-加群、
$\mathcal{G}$ を $\mathcal{O}_Y$-加群とする。
$A = \Gamma(S, \mathcal{O}_S)$ とおく。$D(A)$ の写像
\begin{equation}
\label{perfect-equation-kunneth-single-sheaves}
R\Gamma(X, \mathcal{F})
\otimes_A^\mathbf{L}
R\Gamma(Y, \mathcal{G})
\longrightarrow
R\Gamma(X \times_S Y,
p^*\mathcal{F}
\otimes_{\mathcal{O}_{X \times_S Y}}
q^*\mathcal{G})
\end{equation}
を考える。この写像は引き戻し写像
$R\Gamma(X, \mathcal{F}) \to R\Gamma(X \times_S Y, p^*\mathcal{F})$ と
$R\Gamma(Y, \mathcal{G}) \to R\Gamma(X \times_S Y, q^*\mathcal{G})$、
Cohomology, Section \ref{cohomology-section-cup-product}
で構成されたカップ積、および標準写像
$p^*\mathcal{F}
\otimes_{\mathcal{O}_{X \times_S Y}}^\mathbf{L}
q^*\mathcal{G}
\to
p^*\mathcal{F}
\otimes_{\mathcal{O}_{X \times_S Y}}
q^*\mathcal{G}$
を用いて構成される。

\begin{lemma}
\label{perfect-lemma-kunneth-single-sheaf}
上の状況で $S$ がアフィン、$\mathcal{F}$ と $\mathcal{G}$ が
$S$-平坦かつ擬連接であり、$X$ と $Y$ が準コンパクトで
アフィン対角射をもつなら、写像
(\ref{perfect-equation-kunneth-single-sheaves}) は同型である。
\end{lemma}

\begin{proof}
まず Varieties, Lemma \ref{varieties-lemma-kunneth} の証明を
読むことを強く勧める。有限アフィン開被覆
$\mathcal{U} : X = \bigcup_{i \in I} U_i$ と
$\mathcal{V} : Y = \bigcup_{j \in J} V_j$
を選ぶ。これはアフィン開被覆
$\mathcal{W} : X \times_S Y = \bigcup_{(i, j) \in I \times J} U_i \times_S V_j$.
を定める。$\mathcal{W}$ は $\text{pr}_1^{-1}\mathcal{U}$ と
$\text{pr}_2^{-1}\mathcal{V}$ の細分であることに注意する。
したがって Cohomology, Section
\ref{cohomology-section-cech-cohomology-of-complexes}
の議論により、写像
$$
\check{\mathcal{C}}^\bullet(\mathcal{U}, \mathcal{F})
\to
\check{\mathcal{C}}^\bullet(\mathcal{W}, p^*\mathcal{F})
\quad\text{かつ}\quad
\check{\mathcal{C}}^\bullet(\mathcal{V}, \mathcal{G})
\to
\check{\mathcal{C}}^\bullet(\mathcal{W}, q^*\mathcal{G})
$$
を得る。これらはホモトピーを除いてよく定義され、コホモロジー上の
引き戻し写像と両立する。Cohomology, Equation
(\ref{cohomology-equation-needs-signs}) では複体の写像
$$
\text{Tot}(
\check{\mathcal{C}}^\bullet(\mathcal{W}, p^*\mathcal{F})
\otimes_A
\check{\mathcal{C}}^\bullet(\mathcal{W}, q^*\mathcal{G}))
\longrightarrow
\check{\mathcal{C}}^\bullet(\mathcal{W},
p^*\mathcal{F} \otimes_{\mathcal{O}_{X \times_S Y}}
q^*\mathcal{G})
$$
を構成した。これは Cohomology, Lemma
\ref{cohomology-lemma-diagrams-commute} により
コホモロジー上のカップ積と両立する。以上を合わせると複体の写像
\begin{equation}
\label{perfect-equation-kunneth-on-cech}
\text{Tot}(
\check{\mathcal{C}}^\bullet(\mathcal{U}, \mathcal{F})
\otimes_A
\check{\mathcal{C}}^\bullet(\mathcal{V}, \mathcal{G}))
\to
\check{\mathcal{C}}^\bullet(\mathcal{W},
p^*\mathcal{F}
\otimes_{\mathcal{O}_{X \times_S Y}}
q^*\mathcal{G})
\end{equation}
を得る。これは補題の主張に現れる写像、すなわちこの射の始域と終域は
(\ref{perfect-equation-kunneth-single-sheaves}) の始域と終域に一致すると
主張する。実際、Cohomology of Schemes, Lemma
\ref{coherent-lemma-quasi-coherent-affine-cohomology-zero}
と Cohomology, Lemma
\ref{cohomology-lemma-cech-complex-complex-computes} により、標準写像
$$
\check{\mathcal{C}}^\bullet(\mathcal{U}, \mathcal{F})
\to
R\Gamma(X, \mathcal{F}),
\quad
\check{\mathcal{C}}^\bullet(\mathcal{V}, \mathcal{G})
\to
R\Gamma(Y, \mathcal{G})
$$
および
$$
\check{\mathcal{C}}^\bullet(\mathcal{W},
p^*\mathcal{F} \otimes_{\mathcal{O}_{X \times_S Y}} q^*\mathcal{G})
\to
R\Gamma(X \times_S Y,
p^*\mathcal{F} \otimes_{\mathcal{O}_{X \times_S Y}}
q^*\mathcal{G})
$$
は同型である。一方、複体
$\check{\mathcal{C}}^\bullet(\mathcal{U}, \mathcal{F})$
は補題 \ref{perfect-lemma-K-flat} により K-平坦なので、
$\text{Tot}(
\check{\mathcal{C}}^\bullet(\mathcal{U}, \mathcal{F})
\otimes_A
\check{\mathcal{C}}^\bullet(\mathcal{V}, \mathcal{G}))$
は導来テンソル積
$R\Gamma(X, \mathcal{F}) \otimes_A^\mathbf{L} R\Gamma(Y, \mathcal{G})$
を表す。主張が従う。

\medskip\noindent
残るのは (\ref{perfect-equation-kunneth-on-cech}) が擬同型であることを
示すことである。次元移動を用いる。
$d(\mathcal{F}) = \max \{d \mid H^d(X, \mathcal{F}) \not = 0\}$ とおき、
$d(\mathcal{F}) > 0$ と仮定する。
$U = \coprod\nolimits_{i \in I} U_i$ とおく。$I$ は有限なので
これはアフィン・スキームである。$j : U \to X$ を、各 $U_i$ 上で
包含 $U_i \to X$ となる射と記す。$X$ の対角射はアフィンなので
$j$ はアフィンである。Morphisms, Lemma
\ref{morphisms-lemma-affine-permanence} を参照されたい。
したがって $\mathcal{F}' = j_*j^*\mathcal{F}$ は $S$-平坦である。
Morphisms, Lemma \ref{morphisms-lemma-pushforward-flat-affine}
を参照されたい。また Cohomology of Schemes, Lemmas
\ref{coherent-lemma-relative-affine-cohomology} と
\ref{coherent-lemma-quasi-coherent-affine-cohomology-zero}.
を合わせると $d(\mathcal{F}') = 0$ も従う。
すべての $x \in X$ に対して
$\mathcal{F}_x  \to \mathcal{F}'_x$ は直和因子の包含である。
実際、$x \in U_i$ なら
$\mathcal{F}' \to (U_i \to X)_*\mathcal{F}|_{U_i}$ が分裂を与える。
したがって $\mathcal{F} \to \mathcal{F}'$ は単射であり、
$\mathcal{F}'' = \mathcal{F}'/\mathcal{F}$
も $S$-平坦である。平坦な擬連接 $\mathcal{O}_X$-加群の短完全列
$0 \to \mathcal{F} \to \mathcal{F}' \to \mathcal{F}'' \to 0$
は複体の短完全列
$$
0 \to
\text{Tot}(
\check{\mathcal{C}}^\bullet(\mathcal{U}, \mathcal{F})
\otimes_A
\check{\mathcal{C}}^\bullet(\mathcal{V}, \mathcal{G})) \to
\text{Tot}(
\check{\mathcal{C}}^\bullet(\mathcal{U}, \mathcal{F}')
\otimes_A
\check{\mathcal{C}}^\bullet(\mathcal{V}, \mathcal{G})) \to
\text{Tot}(
\check{\mathcal{C}}^\bullet(\mathcal{U}, \mathcal{F}'')
\otimes_A
\check{\mathcal{C}}^\bullet(\mathcal{V}, \mathcal{G})) \to 0
$$
および複体の短完全列
$$
0 \to
\check{\mathcal{C}}^\bullet(\mathcal{W},
p^*\mathcal{F}
\otimes_{\mathcal{O}_{X \times_S Y}}
q^*\mathcal{G}) \to
\check{\mathcal{C}}^\bullet(\mathcal{W},
p^*\mathcal{F}'
\otimes_{\mathcal{O}_{X \times_S Y}}
q^*\mathcal{G}) \to
\check{\mathcal{C}}^\bullet(\mathcal{W},
p^*\mathcal{F}''
\otimes_{\mathcal{O}_{X \times_S Y}}
q^*\mathcal{G}) \to 0
$$
を生じる。さらに、これらの間の写像
(\ref{perfect-equation-kunneth-on-cech}) はこの短完全列と両立する。
したがって $\mathcal{F}'$ と $\mathcal{F}''$ に対して
(\ref{perfect-equation-kunneth-on-cech}) が同型であることを示せば十分である。
最後に $d(\mathcal{F}'') < d(\mathcal{F})$ である。
このようにして $d(\mathcal{F}) = 0$ の場合に帰着する。

\medskip\noindent
同じ仕方で $\mathcal{G}$ について論じれば、
$\mathcal{F}$ と $\mathcal{G}$ はともに次数 $0$ にしか
非零コホモロジーをもたないと仮定してよい。これは
$\Gamma(X, \mathcal{F})$ が、次数 $\geq 0$ に位置する
$A$-加群の $K$-平坦複体
$\check{\mathcal{C}}^\bullet(\mathcal{U}, \mathcal{F})$
と擬同型であることを意味する。したがって
$\Gamma(X, \mathcal{F})$ は平坦 $A$-加群である
（この加群に対する高次 Tor は {\v C}ech 複体とのテンソル積で
計算できる）。$V \subset Y$ をアフィン開集合とし、アフィン開被覆
$\mathcal{U}_V : X \times_S V = \bigcup_{i \in I} U_i \times_S V$.
を考える。直ちに
$$
\check{\mathcal{C}}^\bullet(\mathcal{U}, \mathcal{F})
\otimes_A \mathcal{G}(V) =
\check{\mathcal{C}}^\bullet(\mathcal{U}_V,
p^*\mathcal{F} \otimes_{\mathcal{O}_{X \times Y}}
q^*\mathcal{G})
$$
である（複体の等号）。$\mathcal{G}(V)$ の $A$ 上の平坦性により
$\Gamma(X, \mathcal{F}) \otimes_A \mathcal{G}(V) \to
\check{\mathcal{C}}^\bullet(\mathcal{U}, \mathcal{F})
\otimes_A \mathcal{G}(V)$
は擬同型である。前層
$V \mapsto \check{\mathcal{C}}^\bullet(\mathcal{U}_V,
p^*\mathcal{F} \otimes_{\mathcal{O}_{X \times Y}}
q^*\mathcal{G})$
の層化は
$Rq_*(p^*\mathcal{F} \otimes_{\mathcal{O}_{X \times Y}} q^*\mathcal{G})$
を表す。これは Cohomology of Schemes, Lemma
\ref{coherent-lemma-separated-case-relative-cech}
による。したがって $Y$ 上
$$
Rq_*(p^*\mathcal{F} \otimes_{\mathcal{O}_{X \times Y}} q^*\mathcal{G})
\cong
\Gamma(X, \mathcal{F}) \otimes_A \mathcal{G}
$$
である。ここで右辺の記法は加群
$$
b^*\widetilde{\Gamma(X, \mathcal{F})} \otimes_{\mathcal{O}_Y} \mathcal{G}
$$
を表す。$q$ に対する Leray スペクトル系列を用いると
$$
H^n(X \times_S Y, p^*\mathcal{F} \otimes_{\mathcal{O}_{X \times Y}}
q^*\mathcal{G}) =
H^n(Y,
b^*\widetilde{\Gamma(X, \mathcal{F})} \otimes_{\mathcal{O}_Y} \mathcal{G})
$$
を得る。射 $b : Y \to S = \Spec(A)$ に補題
\ref{perfect-lemma-cohomology-base-change} を用い、さらに
$\Gamma(X, \mathcal{F})$ が $A$-平坦であることを用いると、
$H^n(X \times_S Y, p^*\mathcal{F} \otimes_{\mathcal{O}_{X \times Y}}
q^*\mathcal{G})$
は $n > 0$ に対して零であり、$n = 0$ に対して
$H^0(X, \mathcal{F}) \otimes_A H^0(Y, \mathcal{G})$ と同型である。
もちろん、ここでは $\mathcal{G}$ が次数 $0$ にしか
コホモロジーをもたないことも用いている。これで証明は完了する
（ただし、この同型が実際に次数 $0$ のカップ積で与えられることを
確認すべきである。その確認は省略する）。
\end{proof}

\begin{remark}
\label{perfect-remark-annoying-compatibility}
$S = \Spec(A)$ をアフィン・スキームとし、
$a : X \to S$ と $b : Y \to S$ をスキームの射とする。
$\mathcal{F}$、$\mathcal{G}$ を擬連接 $\mathcal{O}_X$-加群、
$\mathcal{E}$ を擬連接 $\mathcal{O}_Y$-加群とする。
$\xi \in H^i(X, \mathcal{G})$ とし、その引き戻しを
$p^*\xi \in H^i(X \times_S Y, p^*\mathcal{G})$ とする。
このとき次の図式は可換である。
$$
\xymatrix{
R\Gamma(X, \mathcal{F})[-i] \otimes_A^\mathbf{L} R\Gamma(Y, \mathcal{E})
\ar[d] \ar[rr]_-{\xi \otimes \text{id}} & &
R\Gamma(X, \mathcal{G} \otimes_{\mathcal{O}_X} \mathcal{F})
\otimes_A^\mathbf{L} R\Gamma(Y, \mathcal{E}) \ar[d] \\
R\Gamma(X \times_S Y, p^*\mathcal{F} \otimes q^*\mathcal{E})[-i]
\ar[rr]^-{p^*\xi} & &
R\Gamma(X \times_S Y,
p^*(\mathcal{G} \otimes_{\mathcal{O}_X} \mathcal{F}) \otimes q^*\mathcal{E})
}
$$
ここで下付きのないテンソル積は
$\mathcal{O}_{X \times_S Y}$ 上のものである。水平射は
Cohomology, Remark
\ref{cohomology-remark-cup-with-element-map-total-cohomology}
から来ており、垂直射は (\ref{perfect-equation-kunneth-global}) である。
したがって垂直射は引き戻しの後に $X \times_S Y$ 上のカップ積を
取ることで与えられる。図式が可換なのは、大域カップ積
（$X \times_S Y$ 上で層 $p^*\mathcal{G}$、
$p^*\mathcal{F}$、$q^*\mathcal{E}$ に対して取るもの）が結合的だから
である。Cohomology, Lemma
\ref{cohomology-lemma-cup-product-associative} を参照されたい。
\end{remark}







\section{K\"unneth 公式 III}
\label{perfect-section-kunneth-complexes}

\noindent
$X, Y, S, a, b, p, q, f$ を節 \ref{perfect-section-kunneth} の導入と
同じものとする。本節では、$\mathcal{O}_X$-加群 $\mathcal{F}$ と
$\mathcal{O}_Y$-加群 $\mathcal{G}$ に対して
$$
\mathcal{F} \boxtimes \mathcal{G} =
p^*\mathcal{F} \otimes_{\mathcal{O}_{X \times_S Y}} q^*\mathcal{G}
$$
とおく。後の節で行うこととは異なり、ここでは非導来テンソル積を
取ることに注意する。

\medskip\noindent
$X$ 上で $\mathcal{F}^\bullet$ をアーベル群の層の複体とし、
その各項は擬連接 $\mathcal{O}_X$-加群で、微分
$d^i_\mathcal{F} : \mathcal{F}^i \to \mathcal{F}^{i + 1}$
は $X/S$ 上の有限階微分作用素であるとする。Morphisms, Section
\ref{morphisms-section-differential-operators} を参照されたい。
同様に、$Y$ 上で $\mathcal{G}^\bullet$ をアーベル群の層の複体とし、
その各項は擬連接 $\mathcal{O}_Y$-加群で、微分
$d^j_\mathcal{G} : \mathcal{G}^j \to \mathcal{G}^{j + 1}$
は $Y/S$ 上の有限階微分作用素であるとする。
Morphisms, Lemma
\ref{morphisms-lemma-base-change-differential-operators}
の構成を適用すると、二重複体
$$
\xymatrix{
\ldots &
\ldots &
\ldots &
\ldots \\
\ldots \ar[r] &
\mathcal{F}^i \boxtimes \mathcal{G}^{j + 1}
\ar[r]^{d_1^{i, j + 1}} \ar[u] &
\mathcal{F}^{i + 1} \boxtimes \mathcal{G}^{j + 1} \ar[r] \ar[u] &
\ldots \\
\ldots \ar[r] &
\mathcal{F}^i \boxtimes \mathcal{G}^j
\ar[r]^{d_1^{i, j}} \ar[u]^{d_2^{i, j}} &
\mathcal{F}^{i + 1} \boxtimes \mathcal{G}^j \ar[r] \ar[u]_{d_2^{i + 1, j}} &
\ldots \\
\ldots &
\ldots \ar[u] &
\ldots \ar[u] &
\ldots
}
$$
を得る。その各項は擬連接加群であり、写像は
$X \times_S Y / S$ 上の有限階微分作用素である。
Morphisms, Remark
\ref{morphisms-remark-base-change-differential-operators}
および Homology, Example
\ref{homology-example-double-complex-as-tensor-product-of}
の議論を参照されたい。明示的には
$$
d_1^{i, j} = d^i_\mathcal{F} \boxtimes 1
\quad\text{かつ}\quad
d_2^{i, j} = 1 \boxtimes d^j_\mathcal{G}
$$
とおく。以下の議論で記法
$$
\text{Tot}(\mathcal{F}^\bullet \boxtimes \mathcal{G}^\bullet)
$$
はこの二重複体に付随する全複体を表す。この複体の各項は擬連接
$\mathcal{O}_{X \times_S Y}$-加群であり、その微分は
$X \times_S Y / S$ 上の有限階微分作用素である。

\medskip\noindent
上の状況では「相対カップ積」写像
\begin{equation}
\label{perfect-equation-relative-de-rham-kunneth}
Ra_*(\mathcal{F}^\bullet)
\otimes_{\mathcal{O}_S}^\mathbf{L}
Rb_*(\mathcal{G}^\bullet)
\longrightarrow
Rf_*\left(\text{Tot}(\mathcal{F}^\bullet \boxtimes \mathcal{G}^\bullet)\right)
\end{equation}
が存在する。実際、次の写像を合成して構成できる。
\begin{enumerate}
\item $Ra_*(\mathcal{F}^\bullet) \to Rf_*(p^{-1}\mathcal{F}^\bullet)$,
\item $Rb_*(\mathcal{G}^\bullet) \to Rf_*(q^{-1}\mathcal{G}^\bullet)$,
\item
$Rf_*(p^{-1}\mathcal{F}^\bullet)
\otimes_{\mathcal{O}_S}^\mathbf{L}
Rf_*(q^{-1}\mathcal{G}^\bullet)
\to
Rf_*(p^{-1}\mathcal{F}^\bullet \otimes_{f^{-1}\mathcal{O}_S}^\mathbf{L}
q^{-1}\mathcal{G}^\bullet)$,
\item
$p^{-1}\mathcal{F}^\bullet \otimes_{f^{-1}\mathcal{O}_S}^\mathbf{L}
q^{-1}\mathcal{G}^\bullet \to
\text{Tot}(p^{-1}\mathcal{F}^\bullet \otimes_{f^{-1}\mathcal{O}_S}
q^{-1}\mathcal{G}^\bullet)$
\item
$\text{Tot}(p^{-1}\mathcal{F}^\bullet \otimes_{f^{-1}\mathcal{O}_S}
q^{-1}\mathcal{G}^\bullet) \to \text{Tot}(\mathcal{F}^\bullet \boxtimes
\mathcal{G}^\bullet)$.
\end{enumerate}
(1) と (2) は引き戻し写像、(3) は相対カップ積である。
Cohomology, Remark \ref{cohomology-remark-cup-product}
を参照されたい。(4) は導来テンソル積と非導来テンソル積を比較し、
(5) は基礎にある二重複体上の明らかな写像
$p^{-1}\mathcal{F}^i \otimes_{f^{-1}\mathcal{O}_S} q^{-1}\mathcal{G}^j
\to \mathcal{F}^i \boxtimes \mathcal{G}^j$
によって与えられる。

\medskip\noindent
$A = \Gamma(S, \mathcal{O}_S)$ とおく。「大域カップ積」写像
\begin{equation}
\label{perfect-equation-de-rham-kunneth}
R\Gamma(X, \mathcal{F}^\bullet)
\otimes_A^\mathbf{L}
R\Gamma(Y, \mathcal{G}^\bullet)
\longrightarrow
R\Gamma(X \times_S Y,
\text{Tot}(\mathcal{F}^\bullet \boxtimes \mathcal{G}^\bullet))
\end{equation}
が $D(A)$ 内に存在する。これは上の相対カップ積と同様に、
次の写像を用いて構成される。
\begin{enumerate}
\item $R\Gamma(X, \mathcal{F}^\bullet) \to
R\Gamma(X \times_S Y, p^{-1}\mathcal{F}^\bullet)$
\item $R\Gamma(Y, \mathcal{G}^\bullet) \to
R\Gamma(X \times_S Y, q^{-1}\mathcal{G}^\bullet)$,
\item $R\Gamma(X \times_S Y, p^{-1}\mathcal{F}^\bullet)
\otimes_A^\mathbf{L} R\Gamma(X \times_S Y, q^{-1}\mathcal{G}^\bullet) \to
R\Gamma(X \times_S Y,
p^{-1}\mathcal{F}^\bullet \otimes_{f^{-1}\mathcal{O}_S}^\mathbf{L}
q^{-1}\mathcal{G}^\bullet)$,
\item
$p^{-1}\mathcal{F}^\bullet \otimes_{f^{-1}\mathcal{O}_S}^\mathbf{L}
q^{-1}\mathcal{G}^\bullet \to
\text{Tot}(p^{-1}\mathcal{F}^\bullet \otimes_{f^{-1}\mathcal{O}_S}
q^{-1}\mathcal{G}^\bullet)$
\item
$\text{Tot}(p^{-1}\mathcal{F}^\bullet \otimes_{f^{-1}\mathcal{O}_S}
q^{-1}\mathcal{G}^\bullet) \to \text{Tot}(\mathcal{F}^\bullet \boxtimes
\mathcal{G}^\bullet)$.
\end{enumerate}
ここで (1) と (2) は引き戻し写像、(3) は Cohomology, Section
\ref{cohomology-section-cup-product} で構成されたカップ積である。
(4) と (5) は前段落で述べたものと同じである。

\begin{lemma}
\label{perfect-lemma-kunneth-special}
上の状況で次の仮定が成り立つなら、カップ積
(\ref{perfect-equation-de-rham-kunneth}) は $D(A)$ 内の同型である。
\begin{enumerate}
\item $S = \Spec(A)$ はアフィンである。
\item $X$ と $Y$ は準コンパクトで、アフィン対角射をもつ。
\item $\mathcal{F}^\bullet$ は有界である。
\item $\mathcal{G}^\bullet$ は下に有界である。
\item $\mathcal{F}^n$ は $S$-平坦である。
\item $\mathcal{G}^m$ は $S$-平坦である。
\end{enumerate}
\end{lemma}

\begin{proof}
Morphisms, Remark
\ref{morphisms-remark-base-change-differential-operators}
で導入された記法 $\mathcal{A}_{X/S}$ と $\mathcal{A}_{Y/S}$ を用いる。
複体の写像
$$
\mathcal{F}_1^\bullet \to
\mathcal{F}_2^\bullet \to
\mathcal{F}_3^\bullet \to
\mathcal{F}_1^\bullet[1]
$$
が圏 $\mathcal{A}_{X/S}$ 内にあるとする。
このときカップ積の関手性により可換図式
$$
\xymatrix{
R\Gamma(X, \mathcal{F}_1^\bullet)
\otimes_A^\mathbf{L}
R\Gamma(Y, \mathcal{G}^\bullet)
\ar[r] \ar[d] &
R\Gamma(X \times_S Y,
\text{Tot}(\mathcal{F}_1^\bullet \boxtimes \mathcal{G}^\bullet)) \ar[d] \\
R\Gamma(X, \mathcal{F}_2^\bullet)
\otimes_A^\mathbf{L}
R\Gamma(Y, \mathcal{G}^\bullet)
\ar[r] \ar[d] &
R\Gamma(X \times_S Y,
\text{Tot}(\mathcal{F}_2^\bullet \boxtimes \mathcal{G}^\bullet)) \ar[d] \\
R\Gamma(X, \mathcal{F}_3^\bullet)
\otimes_A^\mathbf{L}
R\Gamma(Y, \mathcal{G}^\bullet)
\ar[r] \ar[d] &
R\Gamma(X \times_S Y,
\text{Tot}(\mathcal{F}_3^\bullet \boxtimes \mathcal{G}^\bullet)) \ar[d] \\
R\Gamma(X, \mathcal{F}_1^\bullet[1])
\otimes_A^\mathbf{L}
R\Gamma(Y, \mathcal{G}^\bullet)
\ar[r] &
R\Gamma(X \times_S Y,
\text{Tot}(\mathcal{F}_1^\bullet[1] \boxtimes \mathcal{G}^\bullet))
}
$$
を得る。もとの写像が $\mathcal{A}_{X/S}$ のホモトピー圏で
完全三角形をなすなら、この図式の各列は $D(A)$ 内で
完全三角形をなす。

\medskip\noindent
補題の状況で、$n < i$ に対して $\mathcal{F}^n = 0$ と仮定する。
すると項ごとに分裂する複体の短完全列
$$
0 \to \sigma_{\geq i + 1}\mathcal{F}^\bullet \to
\mathcal{F}^\bullet \to \mathcal{F}^i[-i] \to 0
$$
を考えられる。ここで切断は Homology, Section
\ref{homology-section-truncations} のものである。これは
完全三角形
$$
\sigma_{\geq i + 1}\mathcal{F}^\bullet \to
\mathcal{F}^\bullet \to
\mathcal{F}^i[-i] \to
(\sigma_{\geq i + 1}\mathcal{F}^\bullet)[1]
$$
を $\mathcal{A}_{X/S}$ のホモトピー圏に生じる。ここで最後の射は
境界写像 $\mathcal{F}^i \to \mathcal{F}^{i + 1}$ によって与えられる。
上の議論により、$\mathcal{F}^i[-i]$ と
$\sigma_{\geq i + 1}\mathcal{F}^\bullet$ に対して補題を
示せば十分である。$\sigma_{\geq i + 1}\mathcal{F}^\bullet$ は
非零な項がより少ないので、帰納法により、
$\mathcal{F}^\bullet$ が一つの次数でのみ非零の場合に補題を
示せればよい。したがって $\mathcal{F}^\bullet$ は一つの次数でのみ
非零であると仮定してよい。

\medskip\noindent
$\mathcal{F}^\bullet$ は、次数 $0$ に $S$-平坦な擬連接
$\mathcal{O}_X$-加群 $\mathcal{F}$ をもち、他の次数では零である
複体と仮定する。例えば補題
\ref{perfect-lemma-cohomology-de-rham-base-change} により
$R\Gamma(X, \mathcal{F})$ は有限 tor 次元をもつ。
その tor 振幅を $[i, j]$ 内とする。$N \gg 0$ を選び、
完全三角形
$$
\sigma_{\geq N + 1}\mathcal{G}^\bullet \to
\mathcal{G}^\bullet \to
\sigma_{\leq N}\mathcal{G}^\bullet \to
(\sigma_{\geq N + 1}\mathcal{G}^\bullet)[1]
$$
を $\mathcal{A}_{Y/S}$ のホモトピー圏で考える。このとき
$$
R\Gamma(X, \mathcal{F})
\otimes_A^\mathbf{L}
R\Gamma(Y, \sigma_{\geq N + 1}\mathcal{G}^\bullet)
\quad\text{かつ}\quad
R\Gamma(X \times_S Y,
\text{Tot}(\mathcal{F} \boxtimes \sigma_{\geq N + 1}\mathcal{G}^\bullet))
$$
はともに次数 $\leq N + i$ でコホモロジーが消滅する。
したがって上の議論を用いると、所与の次数で主張を証明する際には
$\mathcal{G}^\bullet$ が有界であると仮定してよい。
上の議論をもう一度繰り返せば、$\mathcal{G}^\bullet$ も
一つの次数でのみ非零であると仮定してよい。この場合は補題
\ref{perfect-lemma-kunneth-single-sheaf} で扱われている。
\end{proof}







\section{Ext に対する K\"unneth 公式}
\label{perfect-section-kunneth-Ext}

\noindent
スキームのカルテシアン図式
$$
\xymatrix{
& X \times_S Y \ar[ld]^p \ar[rd]_q \ar[dd]^f \\
X \ar[rd]_a & & Y \ar[ld]^b \\
& S
}
$$
を考える。本節では $K \in D(\mathcal{O}_X)$ と
$M \in D(\mathcal{O}_Y)$ に対して
$$
K \boxtimes M =
Lp^*K \otimes_{\mathcal{O}_{X \times_S Y}}^\mathbf{L} Lq^*M
$$
と定義する。$K, K' \in D(\mathcal{O}_X)$ と
$M, M' \in D(\mathcal{O}_Y)$ に対して標準写像
\begin{equation}
\label{perfect-equation-kunneth-ext}
Ra_*R\SheafHom(K, K')
\otimes_{\mathcal{O}_S}^\mathbf{L}
Rb_*R\SheafHom(M, M')
\longrightarrow
Rf_*(R\SheafHom(K \boxtimes M, K' \boxtimes M'))
\end{equation}
があると主張する。実際、次の写像に随伴する写像を取ればよい。
$$
\begin{matrix}
Lf^*\left(Ra_*R\SheafHom(K, K')
\otimes_{\mathcal{O}_S}^\mathbf{L}
Rb_* R\SheafHom(M, M')\right) = \\
Lf^* Ra_* R\SheafHom(K, K')
\otimes_{\mathcal{O}_{X \times_S Y}}^\mathbf{L}
Lf^* Rb_* R\SheafHom(M, M') = \\
Lp^* La^* Ra_* R\SheafHom(K, K')
\otimes_{\mathcal{O}_{X \times_S Y}}^\mathbf{L}
Lq^* Lb^* Rb_* R\SheafHom(M, M') \to \\
Lp^* R\SheafHom(K, K')
\otimes_{\mathcal{O}_{X \times_S Y}}^\mathbf{L}
Lq^* R\SheafHom(M, M') \to \\
R\SheafHom(Lp^*K, Lp^*K')
\otimes_{\mathcal{O}_{X \times_S Y}}^\mathbf{L}
R\SheafHom(Lq^*M, Lq^*M') \to \\
R\SheafHom(K \boxtimes M, K' \boxtimes M')
\end{matrix}
$$
ここで第一の等号は引き戻しとテンソル積の両立性による。
Cohomology, Lemma \ref{cohomology-lemma-pullback-tensor-product}
を参照されたい。第二の等号は $f = a \circ p = b \circ q$ と
引き戻しの合成による。Cohomology, Lemma
\ref{cohomology-lemma-derived-pullback-composition} を参照されたい。
第一の射は、直像と引き戻しが随伴であることによる随伴写像
$La^* Ra_* \to \text{id}$ と $Lb^* Rb_* \to \text{id}$ から与えられる。
Cohomology, Lemma \ref{cohomology-lemma-adjoint} を参照されたい。
第二の射は Cohomology, Remark
\ref{cohomology-remark-prepare-fancy-base-change} により、
第三かつ最後の射は Cohomology, Remark
\ref{cohomology-remark-tensor-internal-hom} により与えられる。
この簡単な特別場合として次を得る。

\begin{lemma}
\label{perfect-lemma-kunneth-Ext}
上の状況で、$a$ と $b$ は準コンパクトかつ準分離であり、
$X$ と $Y$ は $S$ 上 tor 独立であると仮定する。
$K$ が完全、$K' \in D_\QCoh(\mathcal{O}_X)$、
$M$ が完全、$M' \in D_\QCoh(\mathcal{O}_Y)$ なら、
(\ref{perfect-equation-kunneth-ext}) は同型である。
\end{lemma}

\begin{proof}
この場合 $R\SheafHom(K, K') = K' \otimes^\mathbf{L} K^\vee$、
$R\SheafHom(M, M') = M' \otimes^\mathbf{L} M^\vee$、および
$$
R\SheafHom(K \boxtimes M, K' \boxtimes M') =
(K' \otimes^\mathbf{L} K^\vee) \boxtimes
(M' \otimes^\mathbf{L} M^\vee)
$$
である。Cohomology, Lemma
\ref{cohomology-lemma-dual-perfect-complex} を参照されたい。
また完全であることが引き戻しとテンソル積で保たれることも用いる。
したがってこの場合は補題 \ref{perfect-lemma-kunneth} から従う
（これらの同一視のもとで同じ写像が得られることの確認は省略する）。
\end{proof}








\section{コホモロジーと基底変換 V}
\label{perfect-section-cohomology-base-change}

\noindent
節 \ref{perfect-section-cohomology-and-base-change-perfect} では、射が tor 独立なら
基底変換定理が成り立つことを見た。アフィンの場合でさえ、
このような条件なしには基底変換定理は成り立ち得ない。
More on Algebra, Section
\ref{more-algebra-section-tor-independence} を参照されたい。
本節では「一度に一つの複体」について基底変換の結果が
いつ得られるかを解析する。

\medskip\noindent
可換図式
$$
\xymatrix{
X' \ar[r]_{g'} \ar[d]_{f'} &
X \ar[d]^f \\
S' \ar[r]^g &
S
}
$$
がスキームについて与えられているとする（通常はカルテシアンと仮定する）。
$K \in D_\QCoh(\mathcal{O}_X)$ とし、
$L(g')^*K \to K'$ を $D_\QCoh(\mathcal{O}_{X'})$ の写像とする。
点 $x' \in X'$ に対して $x = g'(x') \in X$、
$s' = f'(x') \in S'$、$s = f(x) = g(s')$ とおく。このとき写像
$$
K_x \otimes_{\mathcal{O}_{S, s}}^\mathbf{L} \mathcal{O}_{S', s'} \to
K_x \otimes_{\mathcal{O}_{X, x}}^\mathbf{L} \mathcal{O}_{X', x'} \to
K'_{x'}
$$
を考えられる。第一の射は More on Algebra, Equation
(\ref{more-algebra-equation-comparison-map}) であり、第二の射は
$(L(g')^*K)_{x'} =
K_x \otimes_{\mathcal{O}_{X, x}}^\mathbf{L} \mathcal{O}_{X', x'}$
と与えられた写像 $L(g')^*K \to K'$ から来る。
各 $i \in \mathbf{Z}$ に対して
$H^i(K_x \otimes_{\mathcal{O}_{S, s}}^\mathbf{L} \mathcal{O}_{S', s'})$
には
$\mathcal{O}_{X, x} \otimes_{\mathcal{O}_{S, s}} \mathcal{O}_{S', s'}$-加群
の構造が入る。以上をまとめると
$\mathcal{O}_{X', x'}$-加群の標準写像
\begin{equation}
\label{perfect-equation-bc}
H^i(K_x \otimes_{\mathcal{O}_{S, s}}^\mathbf{L} \mathcal{O}_{S', s'})
\otimes_{(\mathcal{O}_{X, x} \otimes_{\mathcal{O}_{S, s}} \mathcal{O}_{S', s'})}
\mathcal{O}_{X', x'}
\longrightarrow
H^i(K'_{x'})
\end{equation}
を得る。

\begin{lemma}
\label{perfect-lemma-single-complex-base-change-condition}
カルテシアン図式
$$
\xymatrix{
X' \ar[r]_{g'} \ar[d]_{f'} &
X \ar[d]^f \\
S' \ar[r]^g &
S
}
$$
をスキームについて考える。$K \in D_\QCoh(\mathcal{O}_X)$ とし、
$L(g')^*K \to K'$ を $D_\QCoh(\mathcal{O}_{X'})$ の写像とする。
次は同値である。
\begin{enumerate}
\item 任意の $x' \in X'$ と $i \in \mathbf{Z}$ に対して
写像 (\ref{perfect-equation-bc}) は同型である。
\item $U \subset X$ と $V' \subset S'$ がともにアフィン開集合で、
アフィン開集合 $V \subset S$ に写り、
$U' = V' \times_V U$ であるとき、合成
$$
R\Gamma(U, K) \otimes_{\mathcal{O}_S(U)}^\mathbf{L} \mathcal{O}_{S'}(V')
\to
R\Gamma(U, K) \otimes_{\mathcal{O}_X(U)}^\mathbf{L} \mathcal{O}_{X'}(U')
\to
R\Gamma(U', K')
$$
は $D(\mathcal{O}_{S'}(V'))$ 内の同型である。
\item (2) のような四つ組 $U_i, V_i', V_i, U_i'$（$i \in I$）の
集合 $I$ で、$X' = \bigcup U'_i$ を満たすものがある。
\end{enumerate}
\end{lemma}

\begin{proof}
(2) の第二の射は、補題 \ref{perfect-lemma-quasi-coherence-pullback} の等式
$$
R\Gamma(U, K) \otimes_{\mathcal{O}_X(U)}^\mathbf{L} \mathcal{O}_{X'}(U') =
R\Gamma(U', L(g')^*K)
$$
と与えられた射 $L(g')^*K \to K'$ から来る。(2) の第一の射は
More on Algebra, Equation
(\ref{more-algebra-equation-comparison-map}) である。
(2) が (3) を含意することは明らかである。(1) は $X'$ 上局所的である。
したがって $X$、$S$、$S'$、$X'$ がアフィンなら、(1) が
次の写像が同型であるという条件と同値であることを示せば十分である。
$$
R\Gamma(X, K) \otimes_{\mathcal{O}_S(S)}^\mathbf{L} \mathcal{O}_{S'}(S')
\to
R\Gamma(X, K) \otimes_{\mathcal{O}_X(X)}^\mathbf{L} \mathcal{O}_{X'}(X')
\to
R\Gamma(X', K')
$$
ここで同型は $D(\mathcal{O}_{S'}(S'))$ 内で考える。
$S = \Spec(R)$、$X = \Spec(A)$、$S' = \Spec(R')$、
$X' = \Spec(A')$ と書き、$K$ は $A$-加群の複体 $M^\bullet$ に、
$K'$ は $A'$-加群の複体 $N^\bullet$ に対応するとする。
$A' = A \otimes_R R'$ であることに注意する。上の条件は合成
$$
M^\bullet \otimes_R^\mathbf{L} R' \to
M^\bullet \otimes_A^\mathbf{L} A' \to
N^\bullet
$$
が $D(R')$ 内の同型であることである。同値な言い方をすれば、
すべての $i \in \mathbf{Z}$ に対して写像
$$
H^i(M^\bullet \otimes_R^\mathbf{L} R') \to
H^i(M^\bullet \otimes_A^\mathbf{L} A') \to
H^i(N^\bullet)
$$
が同型であることである。これは $A \otimes_R R'$-加群、すなわち
$A'$-加群の写像であることに注意する。一方、(1) は両立する素イデアル
$\mathfrak q' \subset A'$, $\mathfrak q \subset A$,
$\mathfrak p' \subset R'$, $\mathfrak p \subset R$
に対して、合成
$$
H^i(M^\bullet_\mathfrak q \otimes_{R_\mathfrak p}^\mathbf{L} R'_{\mathfrak p'})
\otimes_{(A_\mathfrak q \otimes_{R_\mathfrak p} R'_{\mathfrak p'})}
A'_{\mathfrak q'} \to
H^i(M^\bullet_{\mathfrak q}
\otimes_{A_\mathfrak q}^\mathbf{L} A'_{\mathfrak q'})
\to H^i(N^\bullet_{\mathfrak q'})
$$
が同型であることを要求する。ところで
$$
H^i(M^\bullet_\mathfrak q \otimes_{R_\mathfrak p}^\mathbf{L} R'_{\mathfrak p'})
\otimes_{(A_\mathfrak q \otimes_{R_\mathfrak p} R'_{\mathfrak p'})}
A'_{\mathfrak q'} =
H^i(M^\bullet \otimes_R^\mathbf{L} R') \otimes_{A'} A'_{\mathfrak q'}
$$
は $\mathfrak q'$ における局所化であるから、Algebra, Lemma
\ref{algebra-lemma-characterize-zero-local} により
二つの条件は同値である。
\end{proof}

\begin{lemma}
\label{perfect-lemma-single-complex-base-change}
カルテシアン図式
$$
\xymatrix{
X' \ar[r]_{g'} \ar[d]_{f'} &
X \ar[d]^f \\
S' \ar[r]^g &
S
}
$$
をスキームについて考える。$K \in D_\QCoh(\mathcal{O}_X)$ とし、
$L(g')^*K \to K'$ を $D_\QCoh(\mathcal{O}_{X'})$ の写像とする。
次を仮定する。
\begin{enumerate}
\item 補題 \ref{perfect-lemma-single-complex-base-change-condition}
の同値な条件が成り立つ。
\item $f$ は準コンパクトかつ準分離である。
\end{enumerate}
このとき合成 $Lg^*Rf_*K \to Rf'_*L(g')^*K \to Rf'_*K'$
は同型である。
\end{lemma}

\begin{proof}
補題 \ref{perfect-lemma-compare-base-change} の証明と同じ方法でも証明できるが、
ここでは帰納原理と相対 Mayer--Vietoris を用いて証明する。

\medskip\noindent
写像が同型であることは $S'$ 上局所的に確認できる。
したがって $g : S' \to S$ はアフィン・スキームの射と仮定してよい。
特に $X$ は準コンパクトかつ準分離なスキームである。
Cohomology of Schemes, Lemma
\ref{coherent-lemma-induction-principle} の帰納原理を用いて、
任意の準コンパクト開集合 $U \subset X$ に対して同様に構成される写像
$Lg^*R(U \to S)_*K|_U \to R(U' \to S')_*K'|_{U'}$
が同型であることを示す。ここで $U' = (g')^{-1}(U)$ である。

\medskip\noindent
$U \subset X$ がアフィン開集合なら、仮定により結果は成り立つ。
補題 \ref{perfect-lemma-single-complex-base-change-condition} の (2) と、
補題 \ref{perfect-lemma-affine-compare-bounded} および
\ref{perfect-lemma-quasi-coherence-pullback} による代数への翻訳を参照されたい。

\medskip\noindent
帰納段階を示す。$X = U \cup V$ を開被覆とし、
$U$、$V$、$U \cap V$ は準コンパクトで、
$U$、$V$、$U \cap V$ に対して結果が
成り立つと仮定する。$a = f|_U$、$b = f|_V$、
$c = f|_{U \cap V}$ と記す。$a$、$b$、$c$ の基底変換をそれぞれ
$a' : U' \to S'$、$b' : V' \to S'$、
$c' : U' \cap V' \to S'$ とする。相対 Mayer--Vietoris の完全三角形
（Cohomology, Lemma
\ref{cohomology-lemma-unbounded-relative-mayer-vietoris}）を用いると、
可換図式
$$
\xymatrix{
Lg^*Rf_*K \ar[r] \ar[d] &
Rf'_* K' \ar[d] \\
Lg^*Ra_* K|_U \oplus
Lg^*Rb_* K|_V \ar[r] \ar[d] &
Ra'_* K'|_{U'} \oplus
Rb'_* K'|_{V'} \ar[d] \\
Lg^*Rc_* K|_{U \cap V} \ar[r] \ar[d] &
Rc'_* K'|_{U' \cap V'} \ar[d] \\
Lg^*Rf_* K[1] \ar[r] &
Rf'_* K'[1]
}
$$
を得る。第二と第三の水平射は同型なので、第一の水平射も同型である
（Derived Categories, Lemma
\ref{derived-lemma-third-isomorphism-triangle}）。これで補題の証明が
完了する。
\end{proof}

\begin{lemma}
\label{perfect-lemma-single-complex-base-change-condition-inherited}
カルテシアン図式
$$
\xymatrix{
X' \ar[r]_{g'} \ar[d]_{f'} &
X \ar[d]^f \\
S' \ar[r]^g &
S
}
$$
をスキームについて考える。$K \in D_\QCoh(\mathcal{O}_X)$ とし、
$L(g')^*K \to K'$ を $D_\QCoh(\mathcal{O}_{X'})$ の写像とする。
補題 \ref{perfect-lemma-single-complex-base-change-condition}
の同値な条件が成り立つなら、次が成り立つ。
\begin{enumerate}
\item $E \in D_\QCoh(\mathcal{O}_X)$ に対して、写像
$L(g')^*(E \otimes^\mathbf{L} K) \to L(g')^*E \otimes^\mathbf{L} K'$
は補題 \ref{perfect-lemma-single-complex-base-change-condition}
の同値な条件を満たす。
\item $D(\mathcal{O}_X)$ の $E$ が完全なら、写像
$L(g')^*R\SheafHom(E, K) \to R\SheafHom(L(g')^*E, K')$
は補題 \ref{perfect-lemma-single-complex-base-change-condition}
の同値な条件を満たす。
\item $K$ が下に有界で、$D(\mathcal{O}_X)$ の $E$ が擬連接なら、
写像 $L(g')^*R\SheafHom(E, K) \to R\SheafHom(L(g')^*E, K')$
は補題 \ref{perfect-lemma-single-complex-base-change-condition}
の同値な条件を満たす。
\end{enumerate}
\end{lemma}

\begin{proof}
現れる複体は、補題
\ref{perfect-lemma-quasi-coherence-pullback}、
\ref{perfect-lemma-quasi-coherence-tensor-product}、
\ref{perfect-lemma-quasi-coherence-internal-hom}、および
Cohomology, Lemmas \ref{cohomology-lemma-pseudo-coherent-pullback} と
\ref{cohomology-lemma-perfect-pullback}
により擬連接なコホモロジー層をもつので、主張には意味がある。
その上で (1) の写像 (\ref{perfect-equation-bc}) が同型であることは、
テンソル積の推移性を用いて (\ref{perfect-equation-bc}) の始域と終域を
計算すれば確認できる。More on Algebra, Lemma
\ref{more-algebra-lemma-triple-tensor-product} を参照されたい。
(2) と (3) の写像は合成
$$
L(g')^*R\SheafHom(E, K) \to R\SheafHom(L(g')^*E, L(g')^*K)
\to R\SheafHom(L(g')^*E, K')
$$
である。第一の射は Cohomology, Remark
\ref{cohomology-remark-prepare-fancy-base-change} であり、
第二の射は与えられた写像 $L(g')^*K \to K'$ から来る。
写像 (\ref{perfect-equation-bc}) が同型であることを示すには、(2) では
$E_x$ を有限生成射影 $\mathcal{O}_{X. x}$-加群の有界複体で、
(3) では有限階自由加群の上に有界な複体で表し、射の始域と終域を
計算すればよい。いくつかの詳細は省略する。
\end{proof}

\begin{lemma}
\label{perfect-lemma-base-change-tensor}
$f : X \to S$ をスキームの準コンパクトかつ準分離な射とする。
$E \in D_\QCoh(\mathcal{O}_X)$ とする。$\mathcal{G}^\bullet$ を
$S$ 上平坦な擬連接 $\mathcal{O}_X$-加群の上に有界な複体とする。
このとき
$$
Rf_*(E \otimes^\mathbf{L}_{\mathcal{O}_X} \mathcal{G}^\bullet)
$$
の形成は任意の基底変換と可換する
（精密な主張については証明を参照）。
\end{lemma}

\begin{proof}
主張の意味は次の通りである。$g : S' \to S$ をスキームの射とし、
基底変換図式
$$
\xymatrix{
X' \ar[r]_{g'} \ar[d]_{f'} &
X \ar[d]^f \\
S' \ar[r]^g &
S
}
$$
を考える。すなわち $X' = S' \times_S X$ である。補題は
$$
Lg^*Rf_*(E \otimes^\mathbf{L}_{\mathcal{O}_X} \mathcal{G}^\bullet)
\longrightarrow
Rf'_*\left(
L(g')^*E \otimes^\mathbf{L}_{\mathcal{O}_{X'}} (g')^*\mathcal{G}^\bullet
\right)
$$
が同型であると主張する。右辺では $\mathcal{G}^\bullet$ に
導来引き戻しを{\bf 用いない}ことに注意する。これを示すため、
補題 \ref{perfect-lemma-single-complex-base-change} と
\ref{perfect-lemma-single-complex-base-change-condition-inherited} を適用すると、
標準写像
$$
L(g')^*\mathcal{G}^\bullet \to (g')^*\mathcal{G}^\bullet
$$
が補題 \ref{perfect-lemma-single-complex-base-change-condition}
の同値な条件を満たすことを示せば十分である。これは条件を茎で
確認すれば従う。実際、複体に関する仮定により
$\mathcal{G}^\bullet_x \otimes_{\mathcal{O}_{S, s}} \mathcal{O}_{S', s'}$
が導来テンソル積を計算することが、複体
$\mathcal{G}^\bullet$ に関する仮定から直ちに従う。
\end{proof}

\begin{lemma}
\label{perfect-lemma-base-change-RHom}
$f : X \to S$ をスキームの準コンパクトかつ準分離な射とする。
$E$ を $D(\mathcal{O}_X)$ の対象とし、$\mathcal{G}^\bullet$ を
擬連接 $\mathcal{O}_X$-加群の複体とする。次のいずれかを仮定する。
\begin{enumerate}
\item $E$ は完全、$\mathcal{G}^\bullet$ は上に有界で、
$\mathcal{G}^n$ は $S$ 上平坦である。
\item $E$ は擬連接、$\mathcal{G}^\bullet$ は有界で、
$\mathcal{G}^n$ は $S$ 上平坦である。
\end{enumerate}
このとき
$$
Rf_*R\SheafHom(E, \mathcal{G}^\bullet)
$$
の形成は任意の基底変換と可換する
（精密な主張については証明を参照）。
\end{lemma}

\begin{proof}
主張の意味は次の通りである。$g : S' \to S$ をスキームの射とし、
基底変換図式
$$
\xymatrix{
X' \ar[r]_{g'} \ar[d]_{f'} &
X \ar[d]^f \\
S' \ar[r]^g &
S
}
$$
を考える。すなわち $X' = S' \times_S X$ である。補題は
$$
Lg^*Rf_*R\SheafHom(E, \mathcal{G}^\bullet)
\longrightarrow
R(f')_*R\SheafHom(L(g')^*E, (g')^*\mathcal{G}^\bullet)
$$
が同型であると主張する。右辺では $\mathcal{G}^\bullet$ に
導来引き戻しを{\bf 用いない}ことに注意する。これを示すため、
補題 \ref{perfect-lemma-single-complex-base-change} と
\ref{perfect-lemma-single-complex-base-change-condition-inherited} を適用すると、
標準写像
$$
L(g')^*\mathcal{G}^\bullet \to (g')^*\mathcal{G}^\bullet
$$
が補題 \ref{perfect-lemma-single-complex-base-change-condition}
の同値な条件を満たすことを示せば十分である。これは補題
\ref{perfect-lemma-base-change-tensor} の証明で示した。
\end{proof}





\section{完全複体の生成}
\label{perfect-section-producing-perfect}

\noindent
次の補題は完全複体を生成するための主要な技術的道具である。
後で与えるこの結果の変種は、ネーター近似によってこれに帰着する。
節 \ref{perfect-section-cohomology-and-base-change-final} を参照されたい。

\begin{lemma}
\label{perfect-lemma-perfect-direct-image}
$S$ をネーター・スキームとする。$f : X \to S$ を局所有限型な
スキームの射とする。$E \in D(\mathcal{O}_X)$ が次を満たすとする。
\begin{enumerate}
\item $E \in D^b_{\textit{Coh}}(\mathcal{O}_X)$,
\item すべての $i$ に対して $H^i(E)$ の台は $S$ 上固有である。
\item $E$ は $D(f^{-1}\mathcal{O}_S)$ の対象として有限 tor 次元をもつ。
\end{enumerate}
このとき $Rf_*E$ は $D(\mathcal{O}_S)$ の完全対象である。
\end{lemma}

\begin{proof}
補題 \ref{perfect-lemma-direct-image-coherent} により $Rf_*E$ は
$D^b_{\textit{Coh}}(\mathcal{O}_S)$ の対象である。したがって
$Rf_*E$ は擬連接である（補題
\ref{perfect-lemma-identify-pseudo-coherent-noetherian}）。
ゆえに $Rf_*E$ が有限 tor 次元をもつことを示せば十分である。
Cohomology, Lemma \ref{cohomology-lemma-perfect} を参照されたい。
補題 \ref{perfect-lemma-tor-qc-qs} により、$S$ 上のすべての擬連接層
$\mathcal{F}$ に対して
$Rf_*(E) \otimes_{\mathcal{O}_S}^\mathbf{L} \mathcal{F}$
が一様有界なコホモロジーをもつことを確認すれば十分である。
$Rf_*(E)$ は上に有界なので、上からの有界性は明らかである。
$T \subset X$ をすべての $i$ に対する $H^i(E)$ の台の和とする。
仮定 (1) と (2) により $T$ は $S$ 上固有である。Cohomology of
Schemes, Lemma
\ref{coherent-lemma-union-closed-proper-over-base}.
を参照されたい。特に $T$ を含む準コンパクト開集合
$X' \subset X$ が存在する。$f' = f|_{X'}$ とおくと、
$E$ は $X \setminus T$ 上で零に制限されるので
$Rf_*(E) = Rf'_*(E|_{X'})$ である。したがって $X$ を $X'$ で
置き換え、$f$ は準コンパクトと仮定してよい。さらに Morphisms,
Lemma \ref{morphisms-lemma-finite-type-Noetherian-quasi-separated}
により $f$ は準分離である。ここで
$$
Rf_*(E) \otimes_{\mathcal{O}_S}^\mathbf{L} \mathcal{F} =
Rf_*\left(E \otimes_{\mathcal{O}_X}^\mathbf{L} Lf^*\mathcal{F}\right) =
Rf_*\left(E \otimes_{f^{-1}\mathcal{O}_S}^\mathbf{L} f^{-1}\mathcal{F}\right)
$$
である。これは補題 \ref{perfect-lemma-cohomology-base-change} と
Cohomology, Lemma \ref{cohomology-lemma-variant-derived-pullback}
による。仮定 (3) により、複体
$E \otimes_{f^{-1}\mathcal{O}_S}^\mathbf{L} f^{-1}\mathcal{F}$
のコホモロジー層は、ある一定の有限区間、例えば $[a, b]$ にある。
するとその $Rf_*$ は区間 $[a, \infty)$ にコホモロジーをもち、
結論を得る。
\end{proof}

\begin{lemma}
\label{perfect-lemma-tensor-perfect}
$S$ をネーター・スキームとし、$f : X \to S$ を局所有限型な
スキームの射とする。$E \in D(\mathcal{O}_X)$ を完全とする。
$\mathcal{G}^\bullet$ を、$S$ 上平坦で台が $S$ 上固有である
連接 $\mathcal{O}_X$-加群の有界複体とする。このとき
$K = Rf_*(E \otimes_{\mathcal{O}_X}^\mathbf{L} \mathcal{G}^\bullet)$
は $D(\mathcal{O}_S)$ の完全対象である。
\end{lemma}

\begin{proof}
補題 \ref{perfect-lemma-perfect-direct-image} により対象 $K$ は完全である。
補題を適用できることを確認する。局所的に $E$ は有限階自由
$\mathcal{O}_X$-加群の有限複体と同型である。したがって局所的に
$E \otimes^\mathbf{L}_{\mathcal{O}_X} \mathcal{G}^\bullet$ は、
各項が
$$
\bigoplus\nolimits_{i = a, \ldots, b} (\mathcal{G}^i)^{\oplus r_i}
$$
という形である有限複体と同型である。ただし
$a, b, r_a, \ldots, r_b$ は整数である。これにより
コホモロジー層
$H^i(E \otimes^\mathbf{L}_{\mathcal{O}_X} \mathcal{G})$
が連接であることが直ちに従う。$\mathcal{G}^i$ は
$f^{-1}\mathcal{O}_S$ 上平坦なので、tor 次元に関する仮定も従う。
\end{proof}

\begin{lemma}
\label{perfect-lemma-ext-perfect}
$S$ をネーター・スキームとし、$f : X \to S$ を局所有限型な
スキームの射とする。$E \in D(\mathcal{O}_X)$ を完全とする。
$\mathcal{G}^\bullet$ を、$S$ 上平坦で台が $S$ 上固有である
連接 $\mathcal{O}_X$-加群の有界複体とする。このとき
$K = Rf_*R\SheafHom(E, \mathcal{G}^\bullet)$ は
$D(\mathcal{O}_S)$ の完全対象である。
\end{lemma}

\begin{proof}
$E$ は完全複体なので、双対完全複体 $E^\vee$ が存在する。
Cohomology, Lemma \ref{cohomology-lemma-dual-perfect-complex}
を参照されたい。
$R\SheafHom(E, \mathcal{G}^\bullet) =
E^\vee \otimes^\mathbf{L}_{\mathcal{O}_X} \mathcal{G}^\bullet$
であることに注意する。したがって $K$ の完全性は補題
\ref{perfect-lemma-tensor-perfect} から従う。
\end{proof}

\noindent
次の補題は、補題
\ref{perfect-lemma-flat-proper-perfect-direct-image-general} で一般の基底上の
平坦固有射へ、また More on Morphisms, Lemma
\ref{more-morphisms-lemma-perfect-proper-perfect-direct-image}
で完全固有射へ一般化する。

\begin{lemma}
\label{perfect-lemma-flat-proper-perfect-direct-image}
$S$ をネーター・スキームとする。$f : X \to S$ をスキームの
平坦固有射とし、$E \in D(\mathcal{O}_X)$ を完全とする。
このとき $Rf_*E$ は $D(\mathcal{O}_S)$ の完全対象である。
\end{lemma}

\begin{proof}
補題 \ref{perfect-lemma-perfect-direct-image} が適用できると主張する。
条件 (1) と (2) は直ちに成り立つ。条件 (3) は $X$ 上局所的である。
したがって $X$ と $S$ はアフィンで、$E$ は
$\mathcal{O}_X$-加群の厳密完全複体で表されると仮定してよい。
$\mathcal{O}_X$ は $f^{-1}\mathcal{O}_S$-加群の層として
平坦なので、条件 (3) も満たされる。
\end{proof}






\section{Ext に対する射影公式}
\label{perfect-section-ext}

\noindent
補題 \ref{perfect-lemma-compute-ext}（または文献中の類似の結果）は、
Quot 関手、Hilbert スキーム、および他のモジュライ問題に対する
Artin の判定条件の一つを確認するために用いられることがある。
$f : X \to S$ をスキームの固有、平坦、有限表示な射とし、
$E \in D(\mathcal{O}_X)$ を完全とする。この場合、補題は
$$
\Ext^i_X(E, f^*\mathcal{F}) =
\Ext^i_S((Rf_*E^\vee)^\vee, \mathcal{F})
$$
と述べている。ただし $\mathcal{F}$ は $S$ 上擬連接である。
このように書くと Ext に対する射影公式のように見え、実際この結果は
補題 \ref{perfect-lemma-cohomology-base-change} からかなり容易に従う。

\begin{lemma}
\label{perfect-lemma-compute-tensor-perfect}
仮定と記法を補題 \ref{perfect-lemma-tensor-perfect} と同じものとする。
このとき関手的な同型
$$
H^i(S, K \otimes^\mathbf{L}_{\mathcal{O}_S} \mathcal{F})
\longrightarrow
H^i(X, E \otimes_{\mathcal{O}_X}^\mathbf{L}
(\mathcal{G}^\bullet \otimes_{\mathcal{O}_X} f^*\mathcal{F}))
$$
が $S$ 上擬連接な $\mathcal{F}$ に対して存在し、境界写像と両立する
（証明を参照）。
\end{lemma}

\begin{proof}
等式
$$
\mathcal{G}^\bullet \otimes_{\mathcal{O}_X}^\mathbf{L} Lf^*\mathcal{F} =
\mathcal{G}^\bullet \otimes_{f^{-1}\mathcal{O}_S}^\mathbf{L} f^{-1}\mathcal{F} =
\mathcal{G}^\bullet \otimes_{f^{-1}\mathcal{O}_S} f^{-1}\mathcal{F} =
\mathcal{G}^\bullet \otimes_{\mathcal{O}_X} f^*\mathcal{F}
$$
が成り立つ。第一の等号は Cohomology, Lemma
\ref{cohomology-lemma-variant-derived-pullback}、第二の等号は
$\mathcal{G}^n$ が平坦な $f^{-1}\mathcal{O}_S$-加群であること、
第三の等号は引き戻しの定義による。したがって
\begin{align*}
H^i(X, E \otimes^\mathbf{L}_{\mathcal{O}_X}
(\mathcal{G}^\bullet \otimes_{\mathcal{O}_X} f^*\mathcal{F}))
& =
H^i(X, E \otimes^\mathbf{L}_{\mathcal{O}_X} \mathcal{G}^\bullet
\otimes_{\mathcal{O}_X}^\mathbf{L} Lf^*\mathcal{F}) \\
& =
H^i(S,
Rf_*(E \otimes^\mathbf{L}_{\mathcal{O}_X} \mathcal{G}^\bullet
\otimes^\mathbf{L}_{\mathcal{O}_X} Lf^*\mathcal{F})) \\
& =
H^i(S, Rf_*(E \otimes^\mathbf{L}_{\mathcal{O}_X} \mathcal{G}^\bullet)
\otimes^\mathbf{L}_{\mathcal{O}_S} \mathcal{F}) \\
& =
H^i(S, K \otimes^\mathbf{L}_{\mathcal{O}_S} \mathcal{F}) 
\end{align*}
を得る。第一の等号は上の議論、第二の等号は Leray
（Cohomology, Lemma \ref{cohomology-lemma-before-Leray}）、
第三の等号は補題 \ref{perfect-lemma-cohomology-base-change} による。
境界写像に関する主張の意味は次の通りである。
擬連接 $\mathcal{O}_S$-加群の短完全列
$0 \to \mathcal{F}_1 \to \mathcal{F}_2 \to \mathcal{F}_3 \to 0$
が与えられると、同型は可換図式
$$
\xymatrix{
H^i(S, K \otimes^\mathbf{L}_{\mathcal{O}_S} \mathcal{F}_3)
\ar[r] \ar[d]_\delta &
H^i(X, E \otimes^\mathbf{L}_{\mathcal{O}_X}
(\mathcal{G}^\bullet \otimes_{\mathcal{O}_X} f^*\mathcal{F}_3))
\ar[d]^\delta \\
H^{i + 1}(S, K \otimes^\mathbf{L}_{\mathcal{O}_S} \mathcal{F}_1)
\ar[r] &
H^{i + 1}(X, E \otimes^\mathbf{L}_{\mathcal{O}_X}
(\mathcal{G}^\bullet \otimes_{\mathcal{O}_X} f^*\mathcal{F}_1))
}
$$
に入る。ここで境界写像は完全三角形
$$
K \otimes^\mathbf{L}_{\mathcal{O}_S} \mathcal{F}_1 \to
K \otimes^\mathbf{L}_{\mathcal{O}_S} \mathcal{F}_2 \to
K \otimes^\mathbf{L}_{\mathcal{O}_S} \mathcal{F}_3 \to
K \otimes^\mathbf{L}_{\mathcal{O}_S} \mathcal{F}_1[1]
$$
と、短完全列
$$
0 \to
\mathcal{G}^\bullet \otimes_{\mathcal{O}_X} f^*\mathcal{F}_1 \to
\mathcal{G}^\bullet \otimes_{\mathcal{O}_X} f^*\mathcal{F}_2 \to
\mathcal{G}^\bullet \otimes_{\mathcal{O}_X} f^*\mathcal{F}_3 \to 0
$$
に付随する $D(\mathcal{O}_X)$ 内の完全三角形から来る。
これは $\mathcal{O}_X$-加群の複体の短完全列である。
$\mathcal{G}^n$ は $S$ 上平坦なので、この列は完全である。
表示した図式の可換性の確認は省略する。
\end{proof}

\begin{lemma}
\label{perfect-lemma-compute-ext-perfect}
仮定と記法を補題 \ref{perfect-lemma-ext-perfect} と同じものとする。
このとき関手的な同型
$$
H^i(S, K \otimes^\mathbf{L}_{\mathcal{O}_S} \mathcal{F})
\longrightarrow
\Ext^i_{\mathcal{O}_X}(E,
\mathcal{G}^\bullet \otimes_{\mathcal{O}_X} f^*\mathcal{F})
$$
が $S$ 上擬連接な $\mathcal{F}$ に対して存在し、境界写像と両立する
（証明を参照）。
\end{lemma}

\begin{proof}
補題 \ref{perfect-lemma-ext-perfect} の証明と同様に、$E^\vee$ を
双対完全複体とする。また
$K = Rf_*(E^\vee \otimes_{\mathcal{O}_X}^\mathbf{L} \mathcal{G}^\bullet)$.
であることを思い出す。さらに
$$
\Ext^i_{\mathcal{O}_X}(E,
\mathcal{G}^\bullet \otimes_{\mathcal{O}_X} f^*\mathcal{F})
=
H^i(X, E^\vee \otimes^\mathbf{L}_{\mathcal{O}_X}
(\mathcal{G}^\bullet \otimes_{\mathcal{O}_X} f^*\mathcal{F}))
$$
である。これは $E^\vee$ の構成による。したがって同型の存在は、
$E^\vee$ と $\mathcal{G}^\bullet$ に補題
\ref{perfect-lemma-compute-tensor-perfect} を適用すれば従う。
境界写像に関する主張の意味は次の通りである。
短完全列
$0 \to \mathcal{F}_1 \to \mathcal{F}_2 \to \mathcal{F}_3 \to 0$
が与えられると、同型は可換図式
$$
\xymatrix{
H^i(S, K \otimes^\mathbf{L}_{\mathcal{O}_S} \mathcal{F}_3)
\ar[r] \ar[d]_\delta &
\Ext^i_{\mathcal{O}_X}(E,
\mathcal{G}^\bullet \otimes_{\mathcal{O}_X} f^*\mathcal{F}_3) \ar[d]^\delta \\
H^{i + 1}(S, K \otimes^\mathbf{L}_{\mathcal{O}_S} \mathcal{F}_1)
\ar[r] &
\Ext^{i + 1}_{\mathcal{O}_X}(E,
\mathcal{G}^\bullet \otimes_{\mathcal{O}_X} f^*\mathcal{F}_1)
}
$$
に入る。ここで境界写像は完全三角形
$$
K \otimes^\mathbf{L}_{\mathcal{O}_S} \mathcal{F}_1 \to
K \otimes^\mathbf{L}_{\mathcal{O}_S} \mathcal{F}_2 \to
K \otimes^\mathbf{L}_{\mathcal{O}_S} \mathcal{F}_3 \to
K \otimes^\mathbf{L}_{\mathcal{O}_S} \mathcal{F}_1[1]
$$
と、短完全列
$$
0 \to
\mathcal{G}^\bullet \otimes_{\mathcal{O}_X} f^*\mathcal{F}_1 \to
\mathcal{G}^\bullet \otimes_{\mathcal{O}_X} f^*\mathcal{F}_2 \to
\mathcal{G}^\bullet \otimes_{\mathcal{O}_X} f^*\mathcal{F}_3 \to 0
$$
に付随する $D(\mathcal{O}_X)$ 内の完全三角形から来る。
これは複体の短完全列である。$\mathcal{G}$ は $S$ 上平坦なので
この列は完全である。表示した図式の可換性の確認は省略する。
\end{proof}

\begin{lemma}
\label{perfect-lemma-compute-ext}
$f : X \to S$ をスキームの射、$E \in D(\mathcal{O}_X)$ とし、
$\mathcal{G}^\bullet$ を $\mathcal{O}_X$-加群の複体とする。
次を仮定する。
\begin{enumerate}
\item $S$ はネーターである。
\item $f$ は局所有限型である。
\item $E \in D^-_{\textit{Coh}}(\mathcal{O}_X)$,
\item $\mathcal{G}^\bullet$ は、$S$ 上平坦で台が $S$ 上固有である
連接 $\mathcal{O}_X$-加群の有界複体である。
\end{enumerate}
このとき次の二つの主張が成り立つ。
\begin{enumerate}
\item[(A)] すべての $m \in \mathbf{Z}$ に対して、
$D(\mathcal{O}_S)$ の完全対象 $K$ と関手的な写像
$$
\alpha^i_\mathcal{F} :
\Ext^i_{\mathcal{O}_X}(E,
\mathcal{G}^\bullet \otimes_{\mathcal{O}_X} f^*\mathcal{F})
\longrightarrow
H^i(S, K \otimes^\mathbf{L}_{\mathcal{O}_S} \mathcal{F})
$$
が $S$ 上擬連接な $\mathcal{F}$ に対して存在する。これらは境界写像と
両立し（証明を参照）、$i \leq m$ に対して
$\alpha^i_\mathcal{F}$ は同型である。
\item[(B)] 擬連接な $L \in D(\mathcal{O}_S)$ と関手的な同型
$$
\Ext^i_{\mathcal{O}_S}(L, \mathcal{F}) \longrightarrow
\Ext^i_{\mathcal{O}_X}(E,
\mathcal{G}^\bullet \otimes_{\mathcal{O}_X} f^*\mathcal{F})
$$
が $S$ 上擬連接な $\mathcal{F}$ に対して存在し、境界写像と両立する。
\end{enumerate}
\end{lemma}

\begin{proof}
(A) を証明する。$\mathcal{G}^i$ は $i \in [a, b]$ に対してのみ
非零であるとする。$X$ を $\mathcal{G}^i$ の台の和の
準コンパクト開近傍で置き換えてよい。したがって $X$ は
ネーターと仮定してよい。この場合 $X$ と $f$ は準コンパクトかつ
準分離である。完全複体 $P$ による $(X, E, -m - 1 + a)$ の近似
$P \to E$ を選ぶ（定理 \ref{perfect-theorem-approximation} により可能）。
このとき誘導写像
$$
\Ext^i_{\mathcal{O}_X}(E,
\mathcal{G}^\bullet \otimes_{\mathcal{O}_X} f^*\mathcal{F})
\longrightarrow
\Ext^i_{\mathcal{O}_X}(P,
\mathcal{G}^\bullet \otimes_{\mathcal{O}_X} f^*\mathcal{F})
$$
は $i \leq m$ に対して同型である。実際、この写像の核、resp.\ 余核は
$$
\Ext^i_{\mathcal{O}_X}(C,
\mathcal{G}^\bullet \otimes_{\mathcal{O}_X} f^*\mathcal{F})
\quad\text{それぞれ}\quad
\Ext^{i + 1}_{\mathcal{O}_X}(C,
\mathcal{G}^\bullet \otimes_{\mathcal{O}_X} f^*\mathcal{F})
$$
の商、resp.\ 部分加群である。ここで $C$ は $P \to E$ の錐である。
$C$ のコホモロジー層は次数 $\geq -m - 1 + a$ で消滅するので、
Derived Categories, Lemma \ref{derived-lemma-negative-exts}
により、これらの $\Ext$-群は $i \leq m + 1$ に対して零である。
したがって $E$ が完全複体の場合、すなわち補題
\ref{perfect-lemma-compute-ext-perfect} に帰着する。境界写像に関する主張は
補題 \ref{perfect-lemma-compute-ext-perfect} の証明で説明されている。

\medskip\noindent
(B) を証明する。(A) の証明と同様に $X$ はネーターと仮定してよい。
補題 \ref{perfect-lemma-identify-pseudo-coherent-noetherian} により
$E$ は擬連接であることに注意する。補題
\ref{perfect-lemma-pseudo-coherent-hocolim} により、
$E = \text{hocolim} E_n$ と書ける。ここで $E_n$ は完全で、
$E_n \to E$ は切断 $\tau_{\geq -n}$ 上の同型を誘導する。
$E_n^\vee$ を双対完全複体とする
（Cohomology, Lemma \ref{cohomology-lemma-dual-perfect-complex}）。
完全対象の逆系
$\ldots \to E_3^\vee \to E_2^\vee \to E_1^\vee$ を得る。
これはさらに逆系
$$
\ldots \to K_3 \to K_2 \to K_1\quad\text{ここで}\quad
K_n = Rf_*(E_n^\vee \otimes_{\mathcal{O}_X}^\mathbf{L} \mathcal{G}^\bullet)
$$
を生じる。その各項は $S$ 上完全である。補題
\ref{perfect-lemma-tensor-perfect} を参照されたい。補題
\ref{perfect-lemma-compute-ext-perfect} とその証明、および前段落の議論
（$P = E_n$ とする）により、$S$ 上の任意の擬連接
$\mathcal{F}$ に対して関手的な標準写像
$$
\xymatrix{
& \Ext^i_{\mathcal{O}_X}(E,
\mathcal{G}^\bullet \otimes_{\mathcal{O}_X} f^*\mathcal{F})
\ar[ld] \ar[rd] \\
H^i(S, K_{n + 1} \otimes_{\mathcal{O}_S}^\mathbf{L} \mathcal{F})
\ar[rr] & &
H^i(S, K_n \otimes_{\mathcal{O}_S}^\mathbf{L} \mathcal{F})
}
$$
を得る。これらは $i \leq n + a$ に対して同型である。
$L_n = K_n^\vee$ を双対完全複体とする。このとき
$L_1 \to L_2 \to L_3 \to \ldots$ は $D(\mathcal{O}_S)$ の
完全対象の系であり、$S$ 上の任意の擬連接 $\mathcal{F}$ に対して
写像
$$
\Ext^i_{\mathcal{O}_S}(L_{n + 1}, \mathcal{F})
\longrightarrow
\Ext^i_{\mathcal{O}_S}(L_n, \mathcal{F})
$$
は $i \leq n + a - 1$ に対して同型である。これは
$L_n \to L_{n + 1}$ が切断 $\tau_{\geq -n - a + 2}$ 上の同型を
誘導することを意味する（ヒント：$L_n \to L_{n + 1}$ の錐を取り、
その最後の非消滅コホモロジー層を見る）。
したがって $L = \text{hocolim} L_n$ は擬連接である。
補題 \ref{perfect-lemma-pseudo-coherent-hocolim} を参照されたい。
ホモトピー余極限の写像性質により
$\Ext^i_{\mathcal{O}_S}(L, \mathcal{F}) =
\Ext^i_{\mathcal{O}_S}(L_n, \mathcal{F})$
が $i \leq n + a - 3$ に対して成り立ち、証明が完了する。
\end{proof}

\begin{remark}
\label{perfect-remark-base-change-of-L}
補題 \ref{perfect-lemma-compute-ext} の (B) の擬連接複体 $L$ は、
この状況に標準的に付随する。例えば (B) の $L$ の形成は
基底変換と両立する。言い換えると、スキームのカルテシアン図式
$$
\xymatrix{
X' \ar[r]_{g'} \ar[d]_{f'} &
X \ar[d]^f \\
S' \ar[r]^g &
S
}
$$
が与えられると、標準的かつ関手的な同型
$$
\Ext^i_{\mathcal{O}_{S'}}(Lg^*L, \mathcal{F}') \longrightarrow
\Ext^i_{\mathcal{O}_X}(L(g')^*E,
(g')^*\mathcal{G}^\bullet \otimes_{\mathcal{O}_{X'}} (f')^*\mathcal{F}')
$$
が $S'$ 上擬連接な $\mathcal{F}'$ に対して存在する。右辺では
$\mathcal{G}^\bullet$ に導来引き戻しを{\bf 用いない}ことに注意する。
これが必要になれば、ここに精密な結果を定式化して詳しい証明を与える。
\end{remark}






\section{極限と導来圏}
\label{perfect-section-limits}

\noindent
この節では、スキームの逆系の極限となるスキームの導来圏に関する
いくつかの結果をまとめる。より正確には、次の設定で議論する。

\begin{situation}
\label{perfect-situation-descent}
$S = \lim_{i \in I} S_i$ を、アフィン遷移射
$f_{i'i} : S_{i'} \to S_i$ をもつスキームの有向系の極限とする。
すべての $i \in I$ に対して $S_i$ は準コンパクトかつ準分離であると仮定する。
射影を $f_i : S \to S_i$ と記す。また、元 $0 \in I$ を一つ固定する。
\end{situation}

\begin{lemma}
\label{perfect-lemma-descend-homomorphisms}
状況 \ref{perfect-situation-descent} において、$E_0$ と $K_0$ を
$D(\mathcal{O}_{S_0})$ の対象とする。
$i \geq 0$ に対して $E_i = Lf_{i0}^*E_0$ および
$K_i = Lf_{i0}^*K_0$ とおき、さらに $E = Lf_0^*E_0$ および
$K = Lf_0^*K_0$ とおく。このとき、次のいずれかが成り立てば、写像
$$
\colim_{i \geq 0} \Hom_{D(\mathcal{O}_{S_i})}(E_i, K_i)
\longrightarrow
\Hom_{D(\mathcal{O}_S)}(E, K)
$$
は同型である。
\begin{enumerate}
\item $E_0$ は完全であり、$K_0 \in D_\QCoh(\mathcal{O}_{S_0})$ である。
\item $E_0$ は擬連接であり、
$K_0 \in D_\QCoh(\mathcal{O}_{S_0})$ は有限 tor 次元をもつ。
\end{enumerate}
\end{lemma}

\begin{proof}
各開部分 $U_0 \subset S_0$ に対し、標準写像
$$
\colim_{i \geq 0} \Hom_{D(\mathcal{O}_{U_i})}(E_i|_{U_i}, K_i|_{U_i})
\longrightarrow
\Hom_{D(\mathcal{O}_U)}(E|_U, K|_U)
$$
が同型であるという条件を $P$ とする。ここで
$U = f_0^{-1}(U_0)$ および $U_i = f_{i0}^{-1}(U_0)$ である。
Cohomology of Schemes, Lemma \ref{coherent-lemma-induction-principle}
の帰納原理により、すべての準コンパクト開部分 $U_0$ に対して
$P$ が成り立つことを示す。この補題の条件 (2) は、導来圏における
hom の Mayer--Vietoris、すなわち Cohomology, Lemma
\ref{cohomology-lemma-mayer-vietoris-hom} から直ちに従う。
したがって、$S_0$ がアフィンの場合に補題を証明すれば十分である。

\medskip\noindent
$S_0$ がアフィンであると仮定し、$S_0 = \Spec(A_0)$、
$S_i = \Spec(A_i)$、および $S = \Spec(A)$ と書く。
以下では補題 \ref{perfect-lemma-affine-compare-bounded} を断りなく用いる。

\medskip\noindent
(1) の場合、対象 $E_0^\bullet$ は有限生成射影 $A_0$-加群からなる
有限複体に対応する。補題 \ref{perfect-lemma-perfect-affine} を参照されたい。
対象 $K_0$ は $A_0$-加群の K-flat 複体 $K_0^\bullet$ で表せる。
この状況で示すべきことは
$$
\colim_{i \geq 0} \Hom_{D(A_i)}(E_0^\bullet \otimes_{A_0} A_i,
K_0^\bullet \otimes_{A_0} A_i)
\longrightarrow
\Hom_{D(A)}(E_0^\bullet \otimes_{A_0} A, K_0^\bullet \otimes_{A_0} A)
$$
である。$E_0^\bullet$ は射影加群からなる上に有界な複体なので、
これは
$$
\colim_{i \geq 0} \Hom_{K(A_0)}(E_0^\bullet,
K_0^\bullet \otimes_{A_0} A_i)
\longrightarrow
\Hom_{K(A_0)}(E_0^\bullet, K_0^\bullet \otimes_{A_0} A)
$$
と書き直せる。零でない加群 $E_0^n$ は有限個しかなく、しかもそれらは
すべて有限表示加群なので、この写像は同型である。

\medskip\noindent
(2) の場合、対象 $E_0$ は有限階数自由 $A_0$-加群からなる上に有界な
複体 $E_0^\bullet$ に対応する。補題
\ref{perfect-lemma-pseudo-coherent-affine} を参照されたい。
$K_0$ は平坦 $A_0$-加群からなる有限複体 $K_0^\bullet$ で表せる。
補題 \ref{perfect-lemma-tor-dimension-affine} および More on Algebra, Lemma
\ref{more-algebra-lemma-tor-amplitude} を参照されたい。
特に $K_0^\bullet$ は K-flat であり、前と同様に論じると写像
$$
\colim_{i \geq 0} \Hom_{K(A_0)}(E_0^\bullet,
K_0^\bullet \otimes_{A_0} A_i)
\longrightarrow
\Hom_{K(A_0)}(E_0^\bullet, K_0^\bullet \otimes_{A_0} A)
$$
を得る。この写像が同型であることは明らかである
（$K_0^\bullet$ は有界なので、関与する項は有限個にすぎない）。
\end{proof}

\begin{lemma}
\label{perfect-lemma-descend-perfect}
状況 \ref{perfect-situation-descent} において、$D(\mathcal{O}_S)$ の完全対象の圏は、
$D(\mathcal{O}_{S_i})$ の完全対象の圏の余極限である。
\end{lemma}

\begin{proof}
各開部分 $U_0 \subset S_0$ に対し、関手
$$
\colim_{i \geq 0} D_{perf}(\mathcal{O}_{U_i})
\longrightarrow
D_{perf}(\mathcal{O}_U)
$$
が同値であるという条件を $P$ とする。ここで ${}_{perf}$ は完全対象からなる
充満部分圏を表し、$U = f_0^{-1}(U_0)$ および
$U_i = f_{i0}^{-1}(U_0)$ である。Cohomology of Schemes, Lemma
\ref{coherent-lemma-induction-principle} の帰納原理により、すべての
準コンパクト開部分 $U_0$ に対して $P$ が成り立つことを示す。
まず、補題 \ref{perfect-lemma-descend-homomorphisms} により、この関手がすでに
充満忠実であることが分かっている。したがって本質的全射性を示せば十分である。

\medskip\noindent
まず帰納原理の条件 (2) を確認する。そこで $S_0 = U_0 \cup V_0$ であり、
$U_0$、$V_0$、および $U_0 \cap V_0$ に対して $P$ が成り立つと仮定する。
$E$ を $D(\mathcal{O}_S)$ の完全対象とする。ある $i \geq 0$ と、
$U_i$ 上の完全対象 $E_{U, i}$ および $V_i$ 上の完全対象 $E_{V, i}$ で、
$U$ および $V$ への引き戻しがそれぞれ $E|_U$ および $E|_V$ と
同型になるものを取れる。写像
$$
a : E_{U, i} \to (Rf_{i, *}E)|_{U_i}
\quad\text{かつ}\quad
b : E_{V, i} \to (Rf_{i, *}E)|_{V_i}
$$
を、同型 $Lf_i^*E_{U, i} \to E|_U$ および
$Lf_i^*E_{V, i} \to E|_V$ に随伴する写像とする。
充満忠実性により、$i$ を大きくすれば、引き戻したときに同一視
に一致する同型
$c : E_{U, i}|_{U_i \cap V_i} \to E_{V, i}|_{U_i \cap V_i}$
を取れる：
$$
Lf_i^*E_{U, i}|_{U \cap V} \to E|_{U \cap V} \to Lf_i^*E_{V, i}|_{U \cap V}.
$$
Cohomology, Lemma \ref{cohomology-lemma-glue} を適用すると、
$S_i$ 上の対象 $E_i$ と写像 $d : E_i \to Rf_{i, *}E$ で、
$U_i$ および $V_i$ 上への制限がそれぞれ $a$ および $b$ となるものを得る。
このとき $E_i$ が完全であり、$d$ が同型 $Lf_i^*E_i \to E$ に
随伴していることは明らかである。

\medskip\noindent
最後に帰納原理の条件 (1)、すなわち $S_0$ がアフィンの場合に
補題が成り立つことを確認する。$S_0 = \Spec(A_0)$、
$S_i = \Spec(A_i)$、および $S = \Spec(A)$ と書く。
補題 \ref{perfect-lemma-affine-compare-bounded} および
\ref{perfect-lemma-perfect-affine} を用いると、示すべきことは
$$
D_{perf}(A) = \colim D_{perf}(A_i)
$$
である。これは、環上の完全複体が有限生成射影（したがって有限表示）加群からなる
有限複体によって与えられることから明らかである。詳しくは More on Algebra,
Lemma \ref{more-algebra-lemma-colimit-perfect-complexes} を参照されたい。
\end{proof}





\section{コホモロジーと基底変換 VI}
\label{perfect-section-cohomology-and-base-change-final}

\noindent
コホモロジーと基底変換に関する最後の節として、Sections
\ref{perfect-section-cohomology-and-base-change-perfect},
\ref{perfect-section-cohomology-base-change}, および
\ref{perfect-section-producing-perfect} の議論を続ける。
理解しやすい特別な場合を Remark
\ref{perfect-remark-explain-perfect-direct-image} に示す。

\begin{lemma}
\label{perfect-lemma-base-change-tensor-perfect}
$f : X \to S$ を有限表示射とする。
$E \in D(\mathcal{O}_X)$ を完全対象とする。$\mathcal{G}^\bullet$ を、
$S$ 上平坦で、その台が $S$ 上固有である有限表示
$\mathcal{O}_X$-加群からなる有界複体とする。このとき
$$
K = Rf_*(E \otimes_{\mathcal{O}_X}^\mathbf{L} \mathcal{G}^\bullet)
$$
は $D(\mathcal{O}_S)$ の完全対象であり、その形成は任意の基底変換と可換する。
\end{lemma}

\begin{proof}
基底変換に関する主張は補題 \ref{perfect-lemma-base-change-tensor} である。
したがって、$K$ が完全対象であることを示せば十分である。
$S$ がネーターなら、これは補題 \ref{perfect-lemma-tensor-perfect} から従う。
ネーター近似によりこの場合に帰着する。読者には以下の証明を
読み飛ばすことを勧める。

\medskip\noindent
問題は $S$ 上局所的なので、$S$ はアフィンと仮定してよい。
$S = \Spec(R)$ と書く。この環をネーター環 $R_i$ のフィルター余極限
$R = \colim R_i$ と書く。Limits, Lemma
\ref{limits-lemma-descend-finite-presentation}
により、ある $i$ と $R_i$ 上有限表示なスキーム $X_i$ で、
$R$ への基底変換が $X$ となるものが存在する。Limits, Lemma
\ref{limits-lemma-descend-modules-finite-presentation} により、
$i$ を大きくすれば、有限表示 $\mathcal{O}_{X_i}$-加群からなる有界複体
$\mathcal{G}_i^\bullet$ で、$X$ への引き戻しが
$\mathcal{G}^\bullet$ となるものが存在すると仮定してよい。
$i$ を大きくすれば、$\mathcal{G}_i^n$ は $R_i$ 上平坦であると
仮定してよい。Limits, Lemma
\ref{limits-lemma-descend-module-flat-finite-presentation} を参照されたい。
さらに $i$ を大きくすれば、$\mathcal{G}_i^n$ の台は $R_i$ 上固有であると
仮定してよい。Limits, Lemma
\ref{limits-lemma-eventually-proper-support} および Cohomology of Schemes,
Lemma \ref{coherent-lemma-module-support-proper-over-base} を参照されたい。
最後に補題 \ref{perfect-lemma-descend-perfect} により、$i$ を大きくすれば、
$D(\mathcal{O}_{X_i})$ の完全対象 $E_i$ で、$X$ への引き戻しが
$E$ となるものが存在すると仮定してよい。$X_i \to \Spec(R_i)$、
$E_i$、$\mathcal{G}_i^\bullet$ に補題 \ref{perfect-lemma-tensor-perfect} を適用し、
すでに示した基底変換性を用いれば結果を得る。
\end{proof}

\begin{remark}
\label{perfect-remark-explain-perfect-direct-image}
$R$ を環とする。$X$ を $R$ 上有限表示なスキームとする。
$\mathcal{G}$ を、$R$ 上平坦で、その台が $R$ 上固有である有限表示
$\mathcal{O}_X$-加群とする。補題
\ref{perfect-lemma-base-change-tensor-perfect} により、有限生成射影
$R$-加群からなる有限複体 $M^\bullet$ で
$$
R\Gamma(X_{R'}, \mathcal{G}_{R'}) = M^\bullet \otimes_R R'
$$
が $R$-代数 $R'$ に関して関手的に成り立つものが存在する。
\end{remark}

\begin{lemma}
\label{perfect-lemma-base-change-tensor-pseudo-coherent}
$f : X \to S$ を有限表示射とする。
$E \in D(\mathcal{O}_X)$ を擬連接対象とする。
$\mathcal{G}^\bullet$ を、$S$ 上平坦で、その台が $S$ 上固有である
有限表示 $\mathcal{O}_X$-加群からなる上に有界な複体とする。このとき
$$
K = Rf_*(E \otimes_{\mathcal{O}_X}^\mathbf{L} \mathcal{G}^\bullet)
$$
は $D(\mathcal{O}_S)$ の擬連接対象であり、その形成は任意の基底変換と可換する。
\end{lemma}

\begin{proof}
基底変換に関する主張は補題 \ref{perfect-lemma-base-change-tensor} である。
したがって、$K$ が擬連接対象であることを示せば十分である。
これは完全複体による近似を用いて補題
\ref{perfect-lemma-base-change-tensor-perfect} から従う。読者には以下の証明を
読み飛ばすことを勧める。

\medskip\noindent
問題は $S$ 上局所的なので、$S$ はアフィンと仮定してよい。
すると $X$ は準コンパクトかつ準分離である。さらに、補題
\ref{perfect-lemma-quasi-coherence-direct-image} で説明したように、全直像
$Rf_* : D_\QCoh(\mathcal{O}_X) \to D_\QCoh(\mathcal{O}_S)$
のコホモロジー次元が $N$ となる整数 $N$ が存在する。
$i > b$ に対して $\mathcal{G}^i = 0$ となる整数 $b$ を選ぶ。
すべての $m$ に対して $K$ が $m$-擬連接であることを示せば十分である。
完全複体 $P$ による $(X, E, m - N - 1 - b)$ の近似
$P \to E$ を選ぶ。これは Theorem \ref{perfect-theorem-approximation} により可能である。
$D_\QCoh(\mathcal{O}_X)$ における識別三角
$$
P \to E \to C \to P[1]
$$
を選ぶ。$C$ のコホモロジー層は次数 $\geq m - N - 1 - b$ で零である。
したがって $C \otimes^\mathbf{L} \mathcal{G}^\bullet$ の
コホモロジー層は次数 $\geq m - N - 1$ で零である。ゆえに
$Rf_*(C \otimes^\mathbf{L} \mathcal{G}^\bullet)$
のコホモロジー層は次数 $\geq m - 1$ で零である。したがって
$$
Rf_*(P \otimes^\mathbf{L} \mathcal{G}^\bullet) \to
Rf_*(E \otimes^\mathbf{L} \mathcal{G}^\bullet)
$$
は次数 $\geq m$ のコホモロジー層上で同型である。
次に、$i > a$ に対して $H^i(P) = 0$ と仮定する。このとき
$
P \otimes^\mathbf{L} \sigma_{\geq m - N - 1 - a}\mathcal{G}^\bullet
\longrightarrow
P \otimes^\mathbf{L} \mathcal{G}^\bullet
$
は次数 $\geq m - N - 1$ のコホモロジー層上で同型である。
したがって再び
$$
Rf_*(P \otimes^\mathbf{L} \sigma_{\geq m - N - 1 - a}\mathcal{G}^\bullet) \to
Rf_*(P \otimes^\mathbf{L} \mathcal{G}^\bullet)
$$
は次数 $\geq m$ のコホモロジー層上で同型である。補題
\ref{perfect-lemma-base-change-tensor-perfect} により始域は完全複体である。
したがって、所望どおり $K$ は $m$-擬連接である。
\end{proof}

\begin{lemma}
\label{perfect-lemma-flat-proper-perfect-direct-image-general}
$S$ をスキームとし、$f : X \to S$ を有限表示固有射とする。
\begin{enumerate}
\item $E \in D(\mathcal{O}_X)$ は完全であり、$f$ は平坦であるとする。
このとき $Rf_*E$ は $D(\mathcal{O}_S)$ の完全対象であり、
その形成は任意の基底変換と可換する。
\item $\mathcal{G}$ を $S$ 上平坦な有限表示 $\mathcal{O}_X$-加群とする。
このとき $Rf_*\mathcal{G}$ は $D(\mathcal{O}_S)$ の完全対象であり、
その形成は任意の基底変換と可換する。
\end{enumerate}
\end{lemma}

\begin{proof}
補題 \ref{perfect-lemma-base-change-tensor-perfect} において、
(1) $\mathcal{G}^\bullet$ を次数 $0$ に置いた $\mathcal{O}_X$ とし、
(2) $E = \mathcal{O}_X$ かつ $\mathcal{G}^\bullet$ を次数 $0$ に置いた
$\mathcal{G}$ からなる複体とした特別な場合である。
\end{proof}

\begin{lemma}
\label{perfect-lemma-flat-proper-pseudo-coherent-direct-image-general}
$S$ をスキームとする。$f : X \to S$ を有限表示な平坦固有射とする。
$E \in D(\mathcal{O}_X)$ を擬連接とする。このとき $Rf_*E$ は
$D(\mathcal{O}_S)$ の擬連接対象であり、その形成は任意の基底変換と可換する。
\end{lemma}

\noindent
より一般に、$f : X \to S$ が固有であり、$X$ 上の $E$ が $S$ に関して
擬連接であれば（More on Morphisms, Definition
\ref{more-morphisms-definition-relative-pseudo-coherence}),
$Rf_*E$ は擬連接である（ただし、この一般性では形成は基底変換と可換しない）。
\cite{Kiehl} を参照されたい。

\begin{proof}
補題 \ref{perfect-lemma-base-change-tensor-pseudo-coherent} において、
$\mathcal{G}^\bullet$ を次数 $0$ に置いた $\mathcal{O}_X$ とした特別な場合である。
\end{proof}

\begin{lemma}
\label{perfect-lemma-pullback-and-limits}
$R$ を環とする。$X$ をスキームとし、$f : X \to \Spec(R)$ を
有限表示な平坦固有射とする。$(M_n)$ を、遷移写像が全射である
$R$-加群の逆系とする。このとき標準写像
$$
\mathcal{O}_X \otimes_R (\lim M_n)
\longrightarrow
\lim \mathcal{O}_X \otimes_R M_n
$$
は、始域から終域に $DQ_X$ を適用したものへの同型を誘導する。
\end{lemma}

\begin{proof}
この主張は、$D_\QCoh(\mathcal{O}_X)$ の任意の対象 $E$ に対して、誘導写像
$$
\Hom(E, \mathcal{O}_X \otimes_R (\lim M_n))
\longrightarrow
\Hom(E, \lim \mathcal{O}_X \otimes_R M_n)
$$
が同型であることを意味する。$D_\QCoh(\mathcal{O}_X)$ は完全生成元をもつので
（Theorem \ref{perfect-theorem-bondal-van-den-Bergh}）、完全な $E$ に対して
これを確認すれば十分である。補題 \ref{perfect-lemma-Rlim-quasi-coherent} により
$\lim \mathcal{O}_X \otimes_R M_n = R\lim \mathcal{O}_X \otimes_R M_n$.
Cohomology, Section \ref{cohomology-section-global-RHom} の完全関手
$R\Hom_X(E, -) : D_\QCoh(\mathcal{O}_X) \to D(R)$
は積と、したがって導来極限と可換する。ゆえに
$$
R\Hom_X(E, \lim \mathcal{O}_X \otimes_R M_n) =
R\lim R\Hom_X(E, \mathcal{O}_X \otimes_R M_n)
$$
$E^\vee$ を双対完全複体とする。Cohomology, Lemma
\ref{cohomology-lemma-dual-perfect-complex} を参照されたい。補題
\ref{perfect-lemma-cohomology-base-change} により
$$
R\Hom_X(E, \mathcal{O}_X \otimes_R M_n) =
R\Gamma(X, E^\vee \otimes_{\mathcal{O}_X}^\mathbf{L} Lf^*M_n) =
R\Gamma(X, E^\vee) \otimes_R^\mathbf{L} M_n
$$
である。補題 \ref{perfect-lemma-flat-proper-perfect-direct-image-general} により
$R\Gamma(X, E^\vee)$ は $R$-加群の完全複体である。特にこれは擬連接複体であり、
More on Algebra, Lemma \ref{more-algebra-lemma-pseudo-coherent-tensor-limit}
により
$$
R\lim R\Gamma(X, E^\vee) \otimes_R^\mathbf{L} M_n =
R\Gamma(X, E^\vee) \otimes_R^\mathbf{L} \lim M_n
$$
を得る。これは所望の等式である。
\end{proof}

\begin{lemma}
\label{perfect-lemma-base-change-RHom-perfect}
$f : X \to S$ を有限表示射とする。
$E \in D(\mathcal{O}_X)$ を完全対象とする。$\mathcal{G}^\bullet$ を、
$S$ 上平坦で、その台が $S$ 上固有である有限表示
$\mathcal{O}_X$-加群からなる有界複体とする。このとき
$$
K = Rf_*R\SheafHom(E, \mathcal{G}^\bullet)
$$
は $D(\mathcal{O}_S)$ の完全対象であり、その形成は任意の基底変換と可換する。
\end{lemma}

\begin{proof}
基底変換に関する主張は補題 \ref{perfect-lemma-base-change-RHom} である。
したがって、$K$ が完全対象であることを示せば十分である。
$S$ がネーターなら、これは補題 \ref{perfect-lemma-ext-perfect} から従う。
ネーター近似によりこの場合に帰着する。読者には以下の証明を
読み飛ばすことを勧める。

\medskip\noindent
問題は $S$ 上局所的なので、$S$ はアフィンと仮定してよい。
$S = \Spec(R)$ と書く。この環をネーター環 $R_i$ のフィルター余極限
$R = \colim R_i$ と書く。Limits, Lemma
\ref{limits-lemma-descend-finite-presentation}
により、ある $i$ と $R_i$ 上有限表示なスキーム $X_i$ で、
$R$ への基底変換が $X$ となるものが存在する。Limits, Lemma
\ref{limits-lemma-descend-modules-finite-presentation} により、
$i$ を大きくすれば、有限表示 $\mathcal{O}_{X_i}$-加群からなる有界複体
$\mathcal{G}_i^\bullet$ で、$X$ への引き戻しが $\mathcal{G}^\bullet$
となるものが存在すると仮定してよい。$i$ を大きくすれば、
$\mathcal{G}_i^n$ は $R_i$ 上平坦であると仮定してよい。Limits, Lemma
\ref{limits-lemma-descend-module-flat-finite-presentation} を参照されたい。
さらに $i$ を大きくすれば、$\mathcal{G}_i^n$ の台は $R_i$ 上固有であると
仮定してよい。Limits, Lemma \ref{limits-lemma-eventually-proper-support}
および Cohomology of Schemes, Lemma
\ref{coherent-lemma-module-support-proper-over-base} を参照されたい。
最後に補題 \ref{perfect-lemma-descend-perfect} により、$i$ を大きくすれば、
$D(\mathcal{O}_{X_i})$ の完全対象 $E_i$ で、$X$ への引き戻しが
$E$ となるものが存在すると仮定してよい。$X_i \to \Spec(R_i)$、
$E_i$、$\mathcal{G}_i^\bullet$ に補題 \ref{perfect-lemma-ext-perfect} を適用し、
すでに示した基底変換性を用いれば結果を得る。
\end{proof}









\section{完全複体}
\label{perfect-section-perfect-complexes}

\noindent
まず、完全複体の Betti 数の飛躍軌跡について述べる。
スキーム $X$ 上の複体 $E$ と $X$ の点 $x$ が与えられたとき、
より正確な $Li_x^*E$ の代わりに、しばしば
$E \otimes_{\mathcal{O}_X}^\mathbf{L} \kappa(x)$ と書く。
ここで $i_x : x \to X$ は標準射である。

\begin{lemma}
\label{perfect-lemma-jump-loci}
$X$ をスキームとし、$E \in D(\mathcal{O}_X)$ を擬連接
（例えば完全）とする。任意の $i \in \mathbf{Z}$ に対して関数
$$
\beta_i : X \longrightarrow \{0, 1, 2, \ldots\},\quad
x \longmapsto
\dim_{\kappa(x)}
H^i(E \otimes_{\mathcal{O}_X}^\mathbf{L} \kappa(x))
$$
を考える。このとき次が成り立つ。
\begin{enumerate}
\item $\beta_i$ の形成は任意の基底変換と可換する。
\item 関数 $\beta_i$ は上半連続である。
\item $\beta_i$ の等位集合は $X$ において局所構成可能である。
\end{enumerate}
\end{lemma}

\begin{proof}
スキームの射 $f : Y \to X$ と点 $y \in Y$ を考える。
$x$ を $y$ の像とし、可換図式
$$
\xymatrix{
y \ar[r]_j \ar[d]_g & Y \ar[d]^f \\
x \ar[r]^i & X
}
$$
を考える。このとき $Lg^* \circ Li^* = Lj^* \circ Lf^*$ である。
したがって、擬連接複体 $Lf^*E$ に付随する関数 $\beta'_i$ は
関数 $\beta_i$ の引き戻しである。式で書けば
$\beta'_i = \beta_i \circ f$ である。これが (1) の意味である。

\medskip\noindent
$i$ を固定し、$x \in X$ とする。(2) と (3) は $x$ のある開近傍で
証明すれば十分なので、$X$ はアフィンと仮定してよい。
このとき $E$ は有限階数自由加群からなる上に有界な複体
$\mathcal{F}^\bullet$ で表せる（補題
\ref{perfect-lemma-lift-pseudo-coherent}）。すると
$P = \sigma_{\geq i - 1}\mathcal{F}^\bullet$ は完全対象であり、
$P \to E$ はすべての $x' \in X$ に対して同型
$$
H^i(P \otimes_{\mathcal{O}_X}^\mathbf{L} \kappa(x')) \to
H^i(E \otimes_{\mathcal{O}_X}^\mathbf{L} \kappa(x'))
$$
を誘導する。したがって $E$ は完全であると仮定してよい。この場合、
More on Algebra, Lemma
\ref{more-algebra-lemma-lift-perfect-from-residue-field}
により、$x$ のアフィン開近傍 $U$ と $a \leq b$ が存在し、
$E|_U$ は複体
$$
\ldots \to 0 \to \mathcal{O}_U^{\oplus \beta_a(x)}
\to \mathcal{O}_U^{\oplus \beta_{a + 1}(x)} \to
\ldots \to
\mathcal{O}_U^{\oplus \beta_{b - 1}(x)} \to
\mathcal{O}_U^{\oplus \beta_b(x)} \to 0
\to \ldots
$$
で表される（ここでは問題を代数に帰着するための先行結果、例えば補題
\ref{perfect-lemma-affine-compare-bounded} および
\ref{perfect-lemma-perfect-affine} も用いている）。
すべての $x' \in U$ に対して $\beta_i(x') \leq \beta_i(x)$ であることが
直ちに従う。これにより $\beta_i$ が上半連続であることが分かる。

\medskip\noindent
(3) を証明するため、$X$ はアフィンであり、$E$ は有限階数自由
$\mathcal{O}_X$-加群からなる複体で与えられていると仮定してよい
（例えば前段落と同様に論じるか、Cohomology, Lemma
\ref{cohomology-lemma-perfect-on-locally-ringed} を用いる）。
したがって、複体
$$
\mathcal{O}_X^{\oplus a} \to
\mathcal{O}_X^{\oplus b} \to
\mathcal{O}_X^{\oplus c}
$$
が与えられたとき、点 $x \in X$ に
$\kappa_x^{\oplus a} \to \kappa_x^{\oplus b} \to \kappa_x^{\oplus c}$
の中央のコホモロジーの次元を対応させる関数の等位集合が構成可能であることを
示せばよい。$A \in \text{Mat}(a \times b, \Gamma(X, \mathcal{O}_X))$
を最初の矢印の行列とする。$\text{Mat}(a \times b, \kappa(x))$ における
$A$ の像の階数が $r$ であることは、$A$ のすべての
$(r + 1) \times (r + 1)$-小行列式が $x$ で消え、かつ $A$ の
$r \times r$-小行列式で $x$ で消えないものが存在することと同値である。したがって階数が
$r$ となる点の集合は構成可能な局所閉集合である。二番目の矢印についても
同様に論じ、それらを合わせれば所望の結果を得る。
\end{proof}

\begin{lemma}
\label{perfect-lemma-chi-locally-constant}
$X$ をスキームとし、$E \in D(\mathcal{O}_X)$ を完全とする。関数
$$
\chi_E : X \longrightarrow \mathbf{Z},\quad
x \longmapsto \sum (-1)^i
\dim_{\kappa(x)} H^i(E \otimes_{\mathcal{O}_X}^\mathbf{L} \kappa(x))
$$
は $X$ 上局所定数である。
\end{lemma}

\begin{proof}
Cohomology, Lemma
\ref{cohomology-lemma-perfect-on-locally-ringed}
により、$X$ 上局所的に、$E$ は有限階数自由 $\mathcal{O}_X$-加群からなる
有限複体 $\mathcal{E}^\bullet$ で表せる。そのような開集合上で関数
$\chi_E$ は値 $\sum (-1)^i \text{rank}(\mathcal{E}^i)$ をとる定数関数である。
\end{proof}

\begin{lemma}
\label{perfect-lemma-open-where-cohomology-in-degree-i-rank-r}
$X$ をスキームとし、$E \in D(\mathcal{O}_X)$ を完全とする。
$i, r \in \mathbf{Z}$ が与えられたとき、次の性質で特徴づけられる
開部分スキーム $U \subset X$ が存在する。
\begin{enumerate}
\item $E|_U \cong H^i(E|_U)[-i]$ であり、$H^i(E|_U)$ は階数 $r$ の
局所自由 $\mathcal{O}_U$-加群である。
\item 射 $f : Y \to X$ が $U$ を経由するための必要十分条件は、
$Lf^*E$ が次数 $i$ に置かれた階数 $r$ の局所自由加群と同型になることである。
\end{enumerate}
\end{lemma}

\begin{proof}
$j \in \mathbf{Z}$ に対し、$\beta_j : X \to \{0, 1, 2, \ldots\}$ を
補題 \ref{perfect-lemma-jump-loci} の関数とする。このとき集合
$$
W = \{x \in X \mid \beta_j(x) \leq 0\text{ for all }j \not = i\}
$$
は $X$ の開集合であり、その形成は $X$ 上の任意の $Y$ への引き戻しと
可換する。これは、どの点の近傍でも非零な $\beta_j$ は先験的に有限個しか
ないことを用いれば、この補題から従う。したがって $X$ を $W$ で置き換え、
すべての $x \in X$ と $j \not = i$ に対して $\beta_j(x) = 0$ と
仮定してよい。この場合 $H^i(E)$ は有限局所自由加群であり、
$E \cong H^i(E)[-i]$ である。例えば More on Algebra, Lemma
\ref{more-algebra-lemma-lift-perfect-from-residue-field} を参照されたい。
したがって $X$ は、$H^i(E)$ の階数が一定である開部分スキームの
非交和であり、結果が従う。
\end{proof}

\begin{lemma}
\label{perfect-lemma-locally-closed-where-H0-locally-free}
$X$ をスキームとする。$E \in D(\mathcal{O}_X)$ を、ある
$a, b \in \mathbf{Z}$ に対して $[a, b]$ 内の tor 振幅をもつ完全対象とする。
$r \geq 0$ とする。このとき、次の性質で特徴づけられる局所閉部分スキーム
$j : Z \to X$ が存在する。
\begin{enumerate}
\item $H^a(Lj^*E)$ は階数 $r$ の局所自由 $\mathcal{O}_Z$-加群である。
\item 射 $f : Y \to X$ が $Z$ を経由するための必要十分条件は、
すべての射 $g : Y' \to Y$ に対して $\mathcal{O}_{Y'}$-加群
$H^a(L(f \circ g)^*E)$ が階数 $r$ の局所自由加群になることである。
\end{enumerate}
さらに $j : Z \to X$ は有限表示であり、次が成り立つ。
\begin{enumerate}
\item[(3)] $f : Y \to X$ が $Y \xrightarrow{g} Z \to X$ と分解すれば
$H^a(Lf^*E) = g^*H^a(Lj^*E)$,
\item[(4)] すべての $x \in X$ に対して $\beta_a(x) \leq r$ なら、
$j$ は閉埋入であり、$f : Y \to X$ が与えられたとき次は同値である。
\begin{enumerate}
\item $f : Y \to X$ は $Z$ を経由する。
\item $H^a(Lf^*E)$ は階数 $r$ の局所自由 $\mathcal{O}_Y$-加群である。
\end{enumerate}
さらに $r = 1$ なら、これらは次とも同値である。
\begin{enumerate}
\item[(c)] $\mathcal{O}_Y \to \SheafHom_{\mathcal{O}_Y}(H^a(Lf^*E), H^a(Lf^*E))$
は単射である。
\end{enumerate}
\end{enumerate}
\end{lemma}

\begin{proof}
まず $U \subset X$ を、補題 \ref{perfect-lemma-jump-loci} の関数 $\beta_a$ が
値 $\leq r$ をとる局所構成可能な開部分スキームとする。
$f : Y \to X$ を (2) の射とする。このとき任意の $y \in Y$ に対して
$\beta_a(Lf^*E) = r$ なので、補題 \ref{perfect-lemma-jump-loci} により $y$ は
$U$ に写る。したがって (2) の $f$ は $U$ を経由する。
よって $X$ を $U$ で置き換え、すべての $x \in X$ に対して
$\beta_a(x) \in \{0, 1, \ldots, r\}$ と仮定してよい。
この場合、有限型準連接イデアルで切り出される閉部分スキーム
$Z \subset X$ が存在し、(4) の (a) と (b)、さらに $r = 1$ なら
(4)(c) の同値性によって特徴づけられ、かつ (3) が成り立つことを示す。
これにより (2) の射が $Z$ を経由することもなおさら従うので、
証明が完了する。

\medskip\noindent
$x \in X$ かつ $\beta_a(x) < r$ なら、$\beta_a < r$ となる
$x$ の開近傍が存在する（補題 \ref{perfect-lemma-jump-loci}）。
このようにして、少なくとも集合論的には $Z$ が閉部分集合であることが分かる。

\medskip\noindent
スキーム論的構造を得るため、$\beta_a(x) = r$ となる点 $x \in X$ を考え、
$\beta = \beta_{a + 1}(x)$ とおく。More on Algebra, Lemma
\ref{more-algebra-lemma-lift-perfect-from-residue-field}
により、$x$ のアフィン開近傍 $U$ が存在し、$K|_U$ は複体
$$
\ldots \to 0 \to \mathcal{O}_U^{\oplus r}
\xrightarrow{(f_{ij})} \mathcal{O}_U^{\oplus \beta} \to
\ldots \to
\mathcal{O}_U^{\oplus \beta_{b - 1}(x)} \to
\mathcal{O}_U^{\oplus \beta_b(x)} \to 0
\to \ldots
$$
で表される（ここでも問題を代数に帰着するための先行結果、例えば補題
\ref{perfect-lemma-affine-compare-bounded} および
\ref{perfect-lemma-perfect-affine} を用いている）。ここで $g : Y \to U$ を
スキームの任意の射とし、ある組 $i, j$ に対して $g^\sharp(f_{ij})$ が
非零であるとする。このとき $H^a(Lg^*E)$ は階数 $r$ の局所自由
$\mathcal{O}_Y$-加群ではない。More on Algebra, Lemma
\ref{more-algebra-lemma-coker-injective-free} を参照されたい。
すべての $i, j$ に対して $g^\sharp(f_{ij}) = 0$ なら、
$H^a(Lg^*E)$ が局所自由 $\mathcal{O}_Y$-加群であることは明らかである。
したがって $U$ 上では、すべての $f_{ij}$ で切り出される閉部分スキームが
(3) を満たし、(4)(a) と (b) は同値である。$Z$ の特徴づけにより、
局所的に構成した断片は貼り合わさる（詳細は省略する）。最後に $r = 1$ なら、
この場合局所的に $H^a(Lg^*E) \subset \mathcal{O}_Y$ は元
$g^\sharp(f_{ij})$ が生成するイデアルの零化イデアルなので、
(4)(c) と (4)(b) は同値である。
\end{proof}





\section{応用}
\label{perfect-section-applications}

\noindent
主としてコホモロジーと基底変換の応用を扱う。将来、これらの結果を
補題 \ref{perfect-lemma-base-change-tensor-perfect} で論じた状況へ
一般化するかもしれない。

\begin{lemma}
\label{perfect-lemma-jump-loci-geometric}
$f : X \to S$ を有限表示固有射とする。$\mathcal{F}$ を $S$ 上平坦な
有限表示 $\mathcal{O}_X$-加群とする。固定した $i \in \mathbf{Z}$ に対して関数
$$
\beta_i : S \to \{0, 1, 2, \ldots\},\quad
s \longmapsto \dim_{\kappa(s)} H^i(X_s, \mathcal{F}_s)
$$
を考える。このとき次が成り立つ。
\begin{enumerate}
\item $\beta_i$ の形成は任意の基底変換と可換する。
\item 関数 $\beta_i$ は上半連続である。
\item $\beta_i$ の等位集合は $S$ において局所構成可能である。
\end{enumerate}
\end{lemma}

\begin{proof}
コホモロジーと基底変換（より正確には補題
\ref{perfect-lemma-flat-proper-perfect-direct-image-general}）により、対象
$K = Rf_*\mathcal{F}$ は $S$ の導来圏の完全対象であり、その形成は
任意の基底変換と可換する。特に
$$
H^i(X_s, \mathcal{F}_s) = H^i(K \otimes_{\mathcal{O}_S}^\mathbf{L} \kappa(s))
$$
である。したがって補題 \ref{perfect-lemma-jump-loci} から主張が従う。
\end{proof}

\begin{lemma}
\label{perfect-lemma-chi-locally-constant-geometric}
$f : X \to S$ を有限表示固有射とする。$\mathcal{F}$ を $S$ 上平坦な
有限表示 $\mathcal{O}_X$-加群とする。関数
$$
s \longmapsto \chi(X_s, \mathcal{F}_s)
$$
は $S$ 上局所定数である。この関数の形成は基底変換と可換する。
\end{lemma}

\begin{proof}
コホモロジーと基底変換（より正確には補題
\ref{perfect-lemma-flat-proper-perfect-direct-image-general}）により、対象
$K = Rf_*\mathcal{F}$ は $S$ の導来圏の完全対象であり、その形成は
任意の基底変換と可換する。したがって写像
$$
s \longmapsto \sum (-1)^i \dim_{\kappa(s)}
H^i(K \otimes^\mathbf{L}_{\mathcal{O}_S} \kappa(s))
$$
が $S$ 上局所定数であることを示せばよい。これは補題
\ref{perfect-lemma-chi-locally-constant} である。
\end{proof}

\begin{lemma}
\label{perfect-lemma-open-where-cohomology-in-degree-i-rank-r-geometric}
$f : X \to S$ を有限表示固有射とする。$\mathcal{F}$ を $S$ 上平坦な
有限表示 $\mathcal{O}_X$-加群とする。$i, r \in \mathbf{Z}$ を固定する。
このとき、次の性質をもつ開部分スキーム $U \subset S$ が存在する。
射 $T \to S$ が $U$ を経由するための必要十分条件は、
$Rf_{T, *}\mathcal{F}_T$ が次数 $i$ に置かれた階数 $r$ の
有限局所自由加群と同型になることである。
\end{lemma}

\begin{proof}
コホモロジーと基底変換（より正確には補題
\ref{perfect-lemma-flat-proper-perfect-direct-image-general}）により、対象
$K = Rf_*\mathcal{F}$ は $S$ の導来圏の完全対象であり、その形成は
任意の基底変換と可換する。したがってこの補題は補題
\ref{perfect-lemma-open-where-cohomology-in-degree-i-rank-r} から直ちに従う。
\end{proof}

\begin{lemma}
\label{perfect-lemma-vanishing-implies-locally-free}
$f : X \to S$ を有限表示射とする。$\mathcal{F}$ を、その台が $S$ 上固有である
$S$ 上平坦な有限表示 $\mathcal{O}_X$-加群とする。$i > 0$ に対して
$R^if_*\mathcal{F} = 0$ なら、$f_*\mathcal{F}$ は局所自由であり、その形成は
任意の基底変換と可換する（説明は証明を参照されたい）。
\end{lemma}

\begin{proof}
補題 \ref{perfect-lemma-base-change-tensor-perfect} により、$D(\mathcal{O}_S)$ の対象
$E = Rf_*\mathcal{F}$ は完全であり、その形成は任意の基底変換と可換する。
すなわち、スキームの任意のカルテシアン図式
$$
\xymatrix{
X' \ar[r]_{g'} \ar[d]_{f'} &
X \ar[d]^f \\
S' \ar[r]^g &
S
}
$$
に対して $Rf'_*(g')^*\mathcal{F} = Lg^*E$ である。
次数 $< 0$ にはコホモロジーが決して存在しないので、ある $b$ に対して
$E$ は（局所的に）$[0, b]$ 内の tor 振幅をもつ。
$i > 0$ に対して $H^i(E) = R^if_*\mathcal{F} = 0$ なら、
$E$ は $[0, 0]$ 内の tor 振幅をもつ。したがって $E = H^0(E)[0]$ である。
More on Algebra, Lemma \ref{more-algebra-lemma-perfect}
（および Algebra, Lemma \ref{algebra-lemma-finite-projective} における
有限射影加群の特徴づけ）により、$H^0(E) = f_*\mathcal{F}$ は有限局所自由である。
基底変換との可換性とは、上のような図式に対して
$g^*f_*\mathcal{F} = f'_*(g')^*\mathcal{F}$ となることを意味し、これは
$E$ についてすでに確立した基底変換との可換性から従う。
\end{proof}

\begin{lemma}
\label{perfect-lemma-proper-flat-h0}
$f : X \to S$ をスキームの射とする。次を仮定する。
\begin{enumerate}
\item $f$ は固有、平坦、かつ有限表示である。
\item すべての $s \in S$ に対して
$\kappa(s) = H^0(X_s, \mathcal{O}_{X_s})$ である。
\end{enumerate}
このとき次が成り立つ。
\begin{enumerate}
\item[(a)] $f_*\mathcal{O}_X = \mathcal{O}_S$ であり、
これは任意の基底変換後にも成り立つ。
\item[(b)] $S$ 上局所的に、$D(\mathcal{O}_S)$ において
$$
Rf_*\mathcal{O}_X = \mathcal{O}_S \oplus P
$$
と書ける。ここで $P$ は $[1, \infty)$ 内の tor 振幅をもつ完全対象である。
\end{enumerate}
\end{lemma}

\begin{proof}
コホモロジーと基底変換（補題
\ref{perfect-lemma-flat-proper-perfect-direct-image-general}）により、複体
$E = Rf_*\mathcal{O}_X$ は完全であり、その形成は任意の基底変換と可換する。
第一に、これは $E$ が $[0, \infty)$ 内の tor 振幅をもつことを含意する。
第二に、$s \in S$ に対して
$H^0(E \otimes^\mathbf{L} \kappa(s)) =
H^0(X_s, \mathcal{O}_{X_s}) = \kappa(s)$.
であることを含意する。したがって写像
$\mathcal{O}_S \to Rf_*\mathcal{O}_X = E$ は同型
$\mathcal{O}_S \otimes \kappa(s) \to H^0(E \otimes^\mathbf{L} \kappa(s))$.
を誘導する。ゆえに
$H^0(E) \otimes \kappa(s) \to H^0(E \otimes^\mathbf{L} \kappa(s))$
は全射である。More on Algebra, Lemma
\ref{more-algebra-lemma-better-cut-complex-in-two} を適用すると、$S$ を
$s$ のアフィン開近傍で置き換えた後、分解
$E = H^0(E) \oplus \tau_{\geq 1}E$ を得る。ここで
$\tau_{\geq 1}E$ は $[1, \infty)$ 内の tor 振幅をもつ完全対象である。
$E$ は $[0, \infty)$ 内の tor 振幅をもつので、$H^0(E)$ は平坦な
$\mathcal{O}_S$-加群である。したがって $H^0(E)$ は平坦かつ完全な
$\mathcal{O}_S$-加群であり、ゆえに有限局所自由である。More on Algebra,
Lemma \ref{more-algebra-lemma-perfect}（および有限射影加群が有限局所自由であるという
Algebra, Lemma \ref{algebra-lemma-finite-projective}）を参照されたい。
写像 $\mathcal{O}_S \to H^0(E)$ が同型であることは、剰余体と
テンソルした後に確認できるので従う（Algebra, Lemma
\ref{algebra-lemma-cokernel-flat}）。
\end{proof}

\begin{lemma}
\label{perfect-lemma-proper-flat-geom-red-connected}
$f : X \to S$ をスキームの射とする。次を仮定する。
\begin{enumerate}
\item $f$ は固有、平坦、かつ有限表示である。
\item $f$ の幾何学的ファイバーは被約かつ連結である。
\end{enumerate}
このとき $f_*\mathcal{O}_X = \mathcal{O}_S$ であり、これは
任意の基底変換後にも成り立つ。
\end{lemma}

\begin{proof}
補題 \ref{perfect-lemma-proper-flat-h0} により、すべての $s \in S$ に対して
$\kappa(s) = H^0(X_s, \mathcal{O}_{X_s})$ であることを示せば十分である。
これは Varieties, Lemma
\ref{varieties-lemma-proper-geometrically-reduced-global-sections} と、
$X_s$ が幾何学的に連結かつ幾何学的に被約であることから従う。
\end{proof}

\begin{lemma}
\label{perfect-lemma-proper-idempotent-on-fibre}
$f : X \to S$ をスキームの固有射とする。$s \in S$ とし、
$e \in H^0(X_s, \mathcal{O}_{X_s})$ を冪等元とする。このとき $e$ は写像
$(f_*\mathcal{O}_X)_s \to H^0(X_s, \mathcal{O}_{X_s})$ の像に属する。
\end{lemma}

\begin{proof}
$X_s = T_1 \amalg T_2$ を $e$ に対応する非交和分解とする。ここで
$T_1$ と $T_2$ は空でなく、$X_s$ において開かつ閉であり、
$e$ は $T_1$ 上恒等的に $1$、$T_2$ 上恒等的に $0$ である。

\medskip\noindent
$S$ がネーターであると仮定する。形式関数定理を Cohomology of Schemes,
Lemma \ref{coherent-lemma-formal-functions-stalk} の形で用いる。これによれば
$$
(f_*\mathcal{O}_X)_s^\wedge = \lim_n H^0(X_n, \mathcal{O}_{X_n})
$$
である。ここで $X_n$ は $X_s$ の第 $n$ 無限小近傍である。
$X_n$ の台位相空間は $X_s$ の台位相空間と等しいので、すべての $n $ に対して
スキームの非交和分解 $X_n = T_{1, n} \amalg T_{2, n}$ を得る。ここで
$i = 1, 2$ に対し $T_{i, n}$ の台位相空間は $T_i$ である。したがって
$H^0(X_n, \mathcal{O}_{X_n})$ は非自明な冪等元 $e_n$、すなわち
$T_{1, n}$ 上恒等的に $1$ であり、$T_{2, n}$ 上恒等的に $0$ である関数を含む。
$e_{n + 1}$ の $X_n$ への制限が $e_n$ であることは明らかである。
ゆえに $e_\infty = \lim e_n$ は極限の非自明な冪等元である。
したがって $e_\infty$ は $(f_*\mathcal{O}_X)_s$ の完備化の元であり、
$H^0(X_s, \mathcal{O}_{X_s})$ の $e$ に写る。写像
$(f_*\mathcal{O}_X)_s^\wedge \to H^0(X_s, \mathcal{O}_{X_s})$ は
$(f_*\mathcal{O}_X)^\wedge_s / \mathfrak m_s (f_*\mathcal{O}_X)_s^\wedge =
(f_*\mathcal{O}_X)_s / \mathfrak m_s (f_*\mathcal{O}_X)_s$
を経由するので（Algebra, Lemma
\ref{algebra-lemma-hathat-finitely-generated}）、$e$ は所望どおり写像
$(f_*\mathcal{O}_X)_s \to H^0(X_s, \mathcal{O}_{X_s})$ の像に属する。

\medskip\noindent
一般の場合：極限の議論により一般の場合をネーターの場合へ帰着する。
読者にはこの証明を読み飛ばすことを勧める。$S$ を $s$ のアフィン開近傍で
置き換えてよい。したがって $S$ はアフィンと仮定する。Limits, Lemma
\ref{limits-lemma-proper-limit-of-proper-finite-presentation-noetherian}
により、$(f : X \to S) = \lim (f_i : X_i \to S_i)$ と書ける。ここで
$f_i$ は固有であり、$S_i$ はネーターである。$s$ の像を $s_i \in S_i$ と記す。
このとき $s = \lim s_i$ である。Limits, Lemma
\ref{limits-lemma-inverse-limit-irreducibles} を参照されたい。
極限は極限と可換するので、
$X_s = X \times_S s = \lim X_i \times_{S_i} s_i = \lim X_{i, s_i}$ である
（Categories, Lemma \ref{categories-lemma-colimits-commute}）。
したがって、Limits, Lemma \ref{limits-lemma-descend-section} により、
$e$ はある冪等元
$e_i \in H^0(X_{i, s_i}, \mathcal{O}_{X_{i, s_i}})$
の像である。ネーターの場合により、茎
$(f_{i, *}\mathcal{O}_{X_i})_{s_i}$ の元 $\tilde e_i$ で $e_i$ に写るものが存在する。
$\tilde e_i$ を引き戻すと、$(f_*\mathcal{O}_X)_s$ の元 $\tilde e$ で
$e$ に写るものを得る。これで証明が完了する。
\end{proof}

\begin{lemma}
\label{perfect-lemma-proper-flat-geom-red}
$f : X \to S$ をスキームの射とし、$s \in S$ とする。次を仮定する。
\begin{enumerate}
\item $f$ は固有、平坦、かつ有限表示である。
\item ファイバー $X_s$ は幾何学的に被約である。
\end{enumerate}
このとき、$S$ を $s$ の開近傍で置き換えた後、$D(\mathcal{O}_S)$ において
直和分解
$Rf_*\mathcal{O}_X = f_*\mathcal{O}_X \oplus P$
が存在する。ここで $f_*\mathcal{O}_X$ は有限エタール
$\mathcal{O}_S$-代数であり、$P$ は $[1, \infty)$ 内の tor 振幅をもつ完全対象である。
\end{lemma}

\begin{proof}
この補題の証明は補題 \ref{perfect-lemma-proper-flat-h0} の証明と同様であり、
読者にはまずそちらを読むことを勧める。コホモロジーと基底変換
（補題 \ref{perfect-lemma-flat-proper-perfect-direct-image-general}）により、複体
$E = Rf_*\mathcal{O}_X$ は完全であり、その形成は任意の基底変換と可換する。
これはまず、$E$ が $[0, \infty)$ 内の tor 振幅をもつことを含意する。

\medskip\noindent
次を主張する。$S$ を $s$ の開近傍で置き換えた後、$D(\mathcal{O}_S)$ において
直和分解 $E = H^0(E) \oplus \tau_{\geq 1}E$ を取ることができ、
$\tau_{\geq 1}E$ は $[1, \infty)$ 内の tor 振幅をもつ。
ひとまずこの主張が成り立つとし、置き換えを行って直和分解を得たと仮定する。
$E$ は $[0, \infty)$ 内の tor 振幅をもつので、$H^0(E)$ は平坦な
$\mathcal{O}_S$-加群である。ゆえに $H^0(E)$ は平坦かつ完全な
$\mathcal{O}_S$-加群であり、したがって有限局所自由である。More on Algebra,
Lemma \ref{more-algebra-lemma-perfect}（および有限射影加群が有限局所自由であるという
Algebra, Lemma \ref{algebra-lemma-finite-projective}）を参照されたい。
もちろん $H^0(E) = f_*\mathcal{O}_X$ は $\mathcal{O}_S$-代数である。
コホモロジーと基底変換により
$H^0(E) \otimes \kappa(s) = H^0(X_s, \mathcal{O}_{X_s})$.
Varieties, Lemma
\ref{varieties-lemma-proper-geometrically-reduced-global-sections}
と $X_s$ が幾何学的に被約であるという仮定により、
$\kappa(s) \to H^0(E) \otimes \kappa(s)$ は有限エタールである。
有限局所自由射 $\underline{\Spec}_S(H^0(E)) \to S$ に Morphisms, Lemma
\ref{morphisms-lemma-set-points-where-fibres-etale}
を適用すると、$S$ を縮小した後、$\mathcal{O}_S$-代数 $H^0(E)$ は
有限エタールであることが分かる。

\medskip\noindent
残るのは主張の証明である。そのためには写像
$$
(f_*\mathcal{O}_X)_s
\longrightarrow
H^0(X_s, \mathcal{O}_{X_s}) = H^0(E \otimes^\mathbf{L} \kappa(s))
$$
が全射であることを示せば十分である。More on Algebra, Lemma
\ref{more-algebra-lemma-better-cut-complex-in-two} を参照されたい。
$A$ の剰余体 $k$ が代数閉となる平坦局所環準同型
$\mathcal{O}_{S, s} \to A$ を選ぶ。Algebra, Lemma
\ref{algebra-lemma-flat-local-given-residue-field} を参照されたい。
平坦基底変換（Cohomology of Schemes, Lemma
\ref{coherent-lemma-flat-base-change-cohomology}）により
$H^0(X_A, \mathcal{O}_{X_A}) =
(f_*\mathcal{O}_X)_s \otimes_{\mathcal{O}_{S, s}} A$
および $H^0(X_k, \mathcal{O}_{X_k}) =
H^0(X_s, \mathcal{O}_{X_s}) \otimes_{\kappa(s)} k$.
したがって
$H^0(X_A, \mathcal{O}_{X_A}) \to H^0(X_k, \mathcal{O}_{X_k})$
が全射であることを示せば十分である。$X_k$ は $k$ 上の被約固有スキームであり、
$k$ は代数閉なので、すでに用いた Varieties, Lemma
\ref{varieties-lemma-proper-geometrically-reduced-global-sections} により、
$H^0(X_k, \mathcal{O}_{X_k})$ は $k$ のコピーの有限積である。
補題 \ref{perfect-lemma-proper-idempotent-on-fibre} により、この $k$-代数の冪等元は
$H^0(X_A, \mathcal{O}_{X_A}) \to H^0(X_k, \mathcal{O}_{X_k})$ の像に
属するので、結論が従う。
\end{proof}





\section{その他の応用}
\label{perfect-section-other-applications}

\noindent
この節では、上で展開した理論から導けるいくつかの結果を述べ、証明する。

\begin{lemma}
\label{perfect-lemma-cohomology-over-coherent-ring}
$R$ をコヒーレント環とする。$X$ を $R$ 上有限表示なスキームとする。
$\mathcal{G}$ を、$R$ 上平坦で、その台が $R$ 上固有である有限表示
$\mathcal{O}_X$-加群とする。このとき $H^i(X, \mathcal{G})$ は
コヒーレント $R$-加群である。
\end{lemma}

\begin{proof}
補題 \ref{perfect-lemma-base-change-tensor-perfect} と More on Algebra, Lemmas
\ref{more-algebra-lemma-coherent-pseudo-coherent} および
\ref{more-algebra-lemma-perfect} を組み合わせればよい。
\end{proof}

\begin{lemma}
\label{perfect-lemma-countable-cohomology}
$X$ を準コンパクトかつ準分離なスキームとする。
$K$ を $D_\QCoh(\mathcal{O}_X)$ の対象で、そのコホモロジー層
$H^i(K)$ がアフィン開集合上で可算な切断集合をもつものとする。
このとき、任意の準コンパクト開集合 $U \subset X$ と
$D(\mathcal{O}_X)$ の任意の完全対象 $E$ に対して、集合
$$
H^i(U, K \otimes^\mathbf{L} E),\quad \Ext^i(E|_U, K|_U)
$$
は可算である。
\end{lemma}

\begin{proof}
Cohomology, Lemma \ref{cohomology-lemma-dual-perfect-complex} を用いると、
群 $H^i(U, K \otimes^\mathbf{L} E)$ について結果を証明すれば十分である。
帰納原理を用いて補題を証明する。Cohomology of Schemes, Lemma
\ref{coherent-lemma-induction-principle} を参照されたい。

\medskip\noindent
まず $U = \Spec(A)$ がアフィンの場合に成り立つことを示す。
$K$ は $A$-加群の複体 $K^\bullet$ で表せ、$E$ は有限生成射影
$A$-加群からなる有限複体 $P^\bullet$ で表せる。補題
\ref{perfect-lemma-affine-compare-bounded}、\ref{perfect-lemma-perfect-affine}、および
$A$-加群の完全複体の定義（More on Algebra, Definition
\ref{more-algebra-definition-perfect}）を参照されたい。
このとき $(E \otimes^\mathbf{L} K)|_U$ は二重複体
$P^\bullet \otimes_A K^\bullet$
に付随する全複体で表される（補題
\ref{perfect-lemma-quasi-coherence-tensor-product}）。複体 $P^\bullet$ の長さに関する
帰納法（または適切なスペクトル系列）を用いると、各 $a$ に対して
$H^i(P^a \otimes_A K^\bullet)$ が可算であることを示せば十分である。
$P^a$ はある $n$ に対して $A^{\oplus n}$ の直和因子なので、これは
コホモロジー群 $H^i(K^\bullet)$ が可算であるという仮定から従う。

\medskip\noindent
証明を終えるには、$U = V \cup W$ であり、$V$、$W$、および
$V \cap W$ に対して結果が成り立つなら $U$ に対しても成り立つことを
示せば十分である。これは Mayer--Vietoris 系列から直ちに従う。
Cohomology, Lemma \ref{cohomology-lemma-unbounded-mayer-vietoris} を
参照されたい。
\end{proof}

\begin{lemma}
\label{perfect-lemma-countable}
$X$ を準コンパクトかつ準分離なスキームで、アフィン開集合上の
$\mathcal{O}_X$ の切断集合が可算であるものとする。
$K$ を $D_\QCoh(\mathcal{O}_X)$ の対象とする。次は同値である。
\begin{enumerate}
\item $K = \text{hocolim} E_n$ であり、$E_n$ は $D(\mathcal{O}_X)$ の
完全対象である。
\item コホモロジー層 $H^i(K)$ はアフィン開集合上で可算な切断集合をもつ。
\end{enumerate}
\end{lemma}

\begin{proof}
(1) が成り立つなら (2) も成り立つ。実際、ホモトピー余極限を取る操作は
コホモロジー層を取る操作と可換し（Derived Categories, Lemma
\ref{derived-lemma-cohomology-of-hocolim}）、完全複体は局所的に有限階数自由
$\mathcal{O}_X$-加群からなる有限複体と同型なので、$X$ に関する仮定により
(2) を満たす。

\medskip\noindent
(2) を仮定する。$K$ を表す K-injective 複体 $\mathcal{K}^\bullet$ を選ぶ。
$D_\QCoh(\mathcal{O}_X)$ の完全生成元 $E$ を選び、K-injective 複体
$\mathcal{I}^\bullet$ で表す。Theorem \ref{perfect-theorem-DQCoh-is-Ddga} とその証明により、
三角圏の同値 $F : D_\QCoh(\mathcal{O}_X) \to D(A, \text{d})$ が存在する。
ここで $(A, \text{d})$ は微分次数付き代数
$$
(A, \text{d}) =
\Hom_{\text{Comp}^{dg}(\mathcal{O}_X)}
(\mathcal{I}^\bullet, \mathcal{I}^\bullet)
$$
であり、これは $K$ を微分次数付き加群
$$
M = \Hom_{\text{Comp}^{dg}(\mathcal{O}_X)}
(\mathcal{I}^\bullet, \mathcal{K}^\bullet)
$$
に写す。$H^i(A) = \Ext^i(E, E)$ および
$H^i(M) = \Ext^i(E, K)$ であることに注意する。
さらに $F$ は同値なので、この関手とその準逆はホモトピー余極限と可換する。
したがって $M$ を $D(A, \text{d})$ のコンパクト対象のホモトピー余極限として
書けば十分である。Differential Graded Algebra, Lemma
\ref{dga-lemma-countable} により、各 $i$ に対して
$\Ext^i(E, E)$ と $\Ext^i(E, K)$ が可算であることを示せば十分である。
これは補題 \ref{perfect-lemma-countable-cohomology} から従う。
\end{proof}

\begin{lemma}
\label{perfect-lemma-computing-sections-as-colim}
$A$ を環とする。$X$ を $A$ 上有限表示なスキームとする。
$f : U \to X$ を有限表示平坦射とする。このとき次が成り立つ。
\begin{enumerate}
\item $D(\mathcal{O}_X)$ の完全対象 $L_n$ の逆系で、
$$
R\Gamma(U, Lf^*K) = \text{hocolim}\ R\Hom_X(L_n, K)
$$
が $D(A)$ において、$D_\QCoh(\mathcal{O}_X)$ の $K$ に関して関手的に
成り立つものが存在する。
\item $D(\mathcal{O}_X)$ の完全対象 $E_n$ の系で、
$$
R\Gamma(U, Lf^*K) = \text{hocolim}\ R\Gamma(X, E_n \otimes^\mathbf{L} K)
$$
が $D(A)$ において、$D_\QCoh(\mathcal{O}_X)$ の $K$ に関して関手的に
成り立つものが存在する。
\end{enumerate}
\end{lemma}

\begin{proof}
補題 \ref{perfect-lemma-cohomology-base-change} により、$K$ に関して関手的に
$$
R\Gamma(U, Lf^*K) = R\Gamma(X, Rf_*\mathcal{O}_U \otimes^\mathbf{L} K)
$$
である。補題 \ref{perfect-lemma-quasi-coherence-pushforward-direct-sums} により
$R\Gamma(X, -)$ は直和と可換するので、ホモトピー余極限とも可換する。
同様に、$- \otimes^\mathbf{L} K$ は導来余極限と可換する。実際、
$- \otimes^\mathbf{L} K$ は直和と可換する
（$D(\mathcal{O}_X)$ における直和は、表現複体の直和によって与えられる）。
したがって (2) を証明するには、$Rf_*\mathcal{O}_U = \text{hocolim} E_n$ を
$D(\mathcal{O}_X)$ の完全対象 $E_n$ の系について書けば十分である。
これができれば、$L_n = E_n^\vee$ とおくことで (1) を得る。
Cohomology, Lemma \ref{cohomology-lemma-dual-perfect-complex} を参照されたい。

\medskip\noindent
$A = \colim A_i$ と書く。ここで $A_i$ は $\mathbf{Z}$ 上有限型である。
Limits, Lemma \ref{limits-lemma-descend-finite-presentation} により、
ある $i$ と有限表示射 $U_i \to X_i \to \Spec(A_i)$ で、
$\Spec(A)$ への基底変換が $U \to X \to \Spec(A)$ を復元するものを取れる。
$i$ を大きくすれば、$f_i : U_i \to X_i$ は平坦であると仮定してよい。
Limits, Lemma \ref{limits-lemma-descend-flat-finite-presentation} を参照されたい。
補題 \ref{perfect-lemma-compare-base-change} により、$g : X \to X_i$ による
$Rf_{i, *}\mathcal{O}_{U_i}$ の導来引き戻しは $Rf_*\mathcal{O}_U$ に等しい。
$Lg^*$ は導来余極限と可換するので、$f_i$ について所望の主張を
証明すれば十分である。したがって $U$ と $X$ は $\mathbf{Z}$ 上有限型と
仮定してよい。

\medskip\noindent
$f : U \to X$ を $\mathbf{Z}$ 上有限型なスキームの射と仮定する。
証明を終えるため、$Rf_*\mathcal{O}_U$ が完全複体のホモトピー余極限であることを
示す。これには補題 \ref{perfect-lemma-countable} を適用する。したがって
$R^if_*\mathcal{O}_U$ がアフィン開集合上で可算な切断集合をもつことを
示せば十分である。これは構造層に補題
\ref{perfect-lemma-countable-cohomology} を適用すれば従う。
\end{proof}




\section{擬連接複体の特徴づけ II}
\label{perfect-section-pseudo-coherent}

\noindent
この節は Section \ref{perfect-section-pseudo-coherent-hocolim} の続きである。
ここでは、コホモロジーによる擬連接複体の特徴づけを論じる。
この種の結果は More on Morphisms, Section
\ref{more-morphisms-section-characterize-pseudo-coherent} にさらにある。

\begin{lemma}
\label{perfect-lemma-pseudo-coherent-over-algebra}
$A$ を環とする。$R$ を（非可換でもよい）$A$-代数で、$A$-加群として
有限自由であるものとする。このとき $D(A)$ において擬連接である
$D(R)$ の任意の対象 $M$ は、有限自由（右）$R$-加群からなる
上に有界な複体で表せる。
\end{lemma}

\begin{proof}
$M$ を表す右 $R$-加群の複体 $M^\bullet$ を選ぶ。$M$ は擬連接なので、
十分大きい $i$ に対して $H^i(M) = 0$ である。
$H^m(M)$ が非零となる最小の添字を $m$ とする。More on Algebra, Lemma
\ref{more-algebra-lemma-finite-cohomology} により $H^m(M)$ は有限
$A$-加群である。したがって有限自由 $R$-加群 $F^m$ と写像
$F^m \to M^m$ で、$F^m \to M^m \to M^{m + 1}$ が零であり、かつ
$F^m \to H^m(M)$ が全射となるものを選べる。図式で書けば
$$
\xymatrix{
& F^m \ar[d]^\alpha \ar[r] & 0 \ar[d] \ar[r] & \ldots \\
M^{m - 1} \ar[r] & M^m \ar[r] & M^{m + 1} \ar[r] & \ldots
}
$$
$n \leq m$ に関する降下帰納法により、$i \geq n$ に対する有限自由
$R$-加群 $F^i$、$i \geq n$ に対する微分
$d^i : F^i \to F^{i + 1}$、および微分と両立する写像
$\alpha : F^i \to K^i$ を構成する。ただし、
(1) $H^i(\alpha)$ は $i > n$ に対して同型、$i = n$ に対して全射であり、
(2) $i > m$ に対して $F^i = 0$ である。図式は
$$
\xymatrix{
& F^n \ar[r] \ar[d]^\alpha & F^{n + 1} \ar[d]^\alpha \ar[r] & \ldots \ar[r]
& F^i \ar[d]^\alpha \ar[r] & 0 \ar[d] \ar[r] & \ldots \\
M^{n - 1} \ar[r] & M^n \ar[r] & M^{n + 1} \ar[r] & \ldots \ar[r] &
M^i \ar[r] & M^{i + 1} \ar[r] & \ldots
}
$$
基底の場合 $n = m$ は上で済んでいる。帰納段階に進む。
$C^\bullet$ を $\alpha$ の錐とする（Derived Categories, Definition
\ref{derived-definition-cone}）。コホモロジーの長完全系列により、
$i \geq n$ に対して $H^i(C^\bullet) = 0$ である。
$R$ は $A$-加群として有限自由なので、$F^\bullet$ は $A$-加群の
複体として擬連接であることに注意する。したがって More on Algebra,
Lemma \ref{more-algebra-lemma-cone-pseudo-coherent} により、$C^\bullet$ は
$A$-加群の複体として $(n - 1)$-擬連接である。More on Algebra,
Lemma \ref{more-algebra-lemma-finite-cohomology} により、
$H^{n - 1}(C^\bullet)$ は有限 $A$-加群である。有限自由 $R$-加群
$F^{n - 1}$ と写像 $\beta : F^{n - 1} \to C^{n - 1}$ で、合成
$F^{n - 1} \to C^{n - 1} \to C^n$ が零であり、かつ $F^{n - 1}$ が
$H^{n - 1}(C^\bullet)$ へ全射となるものを選ぶ。
$C^{n - 1} = M^{n - 1} \oplus F^n$ なので、
$\beta = (\alpha^{n - 1}, -d^{n - 1})$ と書ける。合成
$F^{n - 1} \to C^{n - 1} \to C^n$ が消えることから、これらの写像は
複体の射
$$
\xymatrix{
& F^{n - 1} \ar[d]^{\alpha^{n - 1}} \ar[r]_{d^{n - 1}} &
F^n \ar[r] \ar[d]^\alpha &
F^{n + 1} \ar[d]^\alpha \ar[r] & \ldots \\
\ldots \ar[r] &
M^{n - 1} \ar[r] & M^n \ar[r] & M^{n + 1} \ar[r] & \ldots
}
$$
さらに、これらの写像は識別三角の射
$$
\xymatrix{
(F^n \to \ldots) \ar[r] \ar[d] &
(F^{n - 1} \to \ldots) \ar[r] \ar[d] &
F^{n - 1} \ar[r] \ar[d]_\beta &
(F^n \to \ldots)[1] \ar[d] \\
(F^n \to \ldots) \ar[r] &
M^\bullet \ar[r] &
C^\bullet \ar[r] &
(F^n \to \ldots)[1]
}
$$
を定める。したがって $\beta$ の選び方により、複体の写像
$(F^{n - 1} \to \ldots) \to M^\bullet$ は次数 $\geq n$ の
コホモロジー上で同型を、次数 $n - 1$ で全射を誘導する。
これで補題の証明が完了する。
\end{proof}

\begin{lemma}
\label{perfect-lemma-pseudo-coherent-on-projective-space}
$A$ を環とし、$n \geq 0$ とする。
$K \in D_\QCoh(\mathcal{O}_{\mathbf{P}^n_A})$ とする。次は同値である。
\begin{enumerate}
\item $K$ は擬連接である。
\item $D(\mathcal{O}_{\mathbf{P}^n_A})$ の各擬連接対象 $E$ に対して、
$R\Gamma(\mathbf{P}^n_A, E \otimes^\mathbf{L} K)$ は $D(A)$ の
擬連接対象である。
\item $D(\mathcal{O}_{\mathbf{P}^n_A})$ の各完全対象 $E$ に対して、
$R\Gamma(\mathbf{P}^n_A, E \otimes^\mathbf{L} K)$ は $D(A)$ の
擬連接対象である。
\item $D(\mathcal{O}_{\mathbf{P}^n_A})$ の各完全対象 $E$ に対して、
$R\Hom_{\mathbf{P}^n_A}(E, K)$ は $D(A)$ の擬連接対象である。
\item $R\Gamma(\mathbf{P}^n_A,
K \otimes^\mathbf{L} \mathcal{O}_{\mathbf{P}^n_A}(d))$ は
$d = 0, 1, \ldots, n$ に対して $D(A)$ の擬連接対象である。
\end{enumerate}
\end{lemma}

\begin{proof}
定義により
$$
R\Hom_{\mathbf{P}^n_A}(E, K) =
R\Gamma(\mathbf{P}^n_A, R\SheafHom_{\mathcal{O}_{\mathbf{P}^n_A}}(E, K))
$$
であることを思い出そう。Cohomology, Section
\ref{cohomology-section-global-RHom} を参照されたい。したがって
Cohomology, Lemma \ref{cohomology-lemma-dual-perfect-complex} により
(4) と (3) は同値である。

\medskip\noindent
すべての完全複体は擬連接なので、(2) が (3) を含意することは明らかである。

\medskip\noindent
(1) が成り立つと仮定する。このとき各擬連接 $E$ に対して
$E \otimes^\mathbf{L} K$ は擬連接である。Cohomology, Lemma
\ref{cohomology-lemma-tensor-pseudo-coherent} を参照されたい。
補題 \ref{perfect-lemma-flat-proper-pseudo-coherent-direct-image-general} により、
そのような擬連接複体の直像は擬連接であり、(2) が成り立つ。

\medskip\noindent
$d = 0, 1, \ldots, n$ に対して $E = \mathcal{O}_{\mathbf{P}^n_A}(d)$ と
取れるので、(3) は (5) を含意する。

\medskip\noindent
証明を終えるには、(5) が (1) を含意することを示さなければならない。
$P$ を (\ref{perfect-equation-generator-Pn}) のもの、$R$ を
(\ref{perfect-equation-algebra-for-Pn}) のものとする。補題
\ref{perfect-lemma-Pn-module-category} により同値
$$
- \otimes^\mathbf{L}_R P :
D(R) \longrightarrow D_\QCoh(\mathcal{O}_{\mathbf{P}^n_A})
$$
$M \otimes^\mathbf{L} P = K$ となる対象 $M \in D(R)$ を取る。
Differential Graded Algebra, Lemma
\ref{dga-lemma-upgrade-tensor-with-complex-derived}
により、$D(A)$ における同型
$$
R\Hom(R, M) = R\Hom_{\mathbf{P}^n_A}(P, K)
$$
が存在する。上と同様に論じると
$$
R\Hom_{\mathbf{P}^n_A}(P, K) = R\Gamma(\mathbf{P}^n_A,
R\SheafHom_{\mathcal{O}_{\mathbf{P}^n_A}}(E, K)) =
R\Gamma(\mathbf{P}^n_A,
P^\vee \otimes^\mathbf{L}_{\mathcal{O}_{\mathbf{P}^n_A}} K).
$$
$P^\vee$ が $d = 0, 1, \ldots, n$ に対する
$\mathcal{O}_{\mathbf{P}^n_A}(d)$ の直和であることと (5) を用いると、
$R\Hom(R, M)$ は $A$-加群の複体として擬連接である。
もちろん $D(A)$ において $M = R\Hom(R, M)$ である。したがって
$M$ は $A$-加群の複体として擬連接である。補題
\ref{perfect-lemma-pseudo-coherent-over-algebra} により、$M$ は有限自由
$R$-加群からなる上に有界な複体 $F^\bullet$ で表せる。このとき
$F^\bullet = \bigcup_{p \geq 0} \sigma_{\geq p}F^\bullet$
は、$F^\bullet$ が性質 (P) をもつ微分次数付き $R$-加群であることを示す
フィルトレーションである。Differential Graded Algebra, Section
\ref{dga-section-P-resolutions} を参照されたい。したがって
$K = M \otimes^\mathbf{L}_R P$ は $F^\bullet \otimes_R P$ で表される
（これは導来テンソル関手の構成から従う。例えば Differential Graded Algebra,
Lemma \ref{dga-lemma-tensor-with-complex-derived} の証明を参照されたい）。
$F^\bullet \otimes_R P$ は各項が $P$ のコピーの直和である上に有界な
複体なので、補題が成り立つ。
\end{proof}

\begin{lemma}
\label{perfect-lemma-perfect-enough}
$A$ を環とする。$X$ を $A$ 上の準コンパクトかつ準分離なスキームとする。
$K \in D^-_\QCoh(\mathcal{O}_X)$ とする。$D(\mathcal{O}_X)$ のすべての
完全対象 $E$ に対して $R\Gamma(X, E \otimes^\mathbf{L} K)$ が
$D(A)$ において擬連接なら、$D(\mathcal{O}_X)$ のすべての擬連接対象
$E$ に対しても $R\Gamma(X, E \otimes^\mathbf{L} K)$ は
$D(A)$ において擬連接である。
\end{lemma}

\begin{proof}
補題 \ref{perfect-lemma-quasi-coherence-direct-image} で説明したように、
$R\Gamma(X, -) : D_\QCoh(\mathcal{O}_X) \to D(A)$
のコホモロジー次元が $N$ となる整数 $N$ が存在する。
$i > b$ に対して $H^i(K) = 0$ となる $b \in \mathbf{Z}$ を取る。
$E$ を $X$ 上擬連接とする。すべての $m$ に対して
$R\Gamma(X, E \otimes^\mathbf{L} K)$ が $m$-擬連接であることを
示せば十分である。完全複体 $P$ による
$(X, E, m - N - 1 - b)$ の近似 $P \to E$ を選ぶ。これは Theorem
\ref{perfect-theorem-approximation} により可能である。
$D_\QCoh(\mathcal{O}_X)$ における識別三角
$$
P \to E \to C \to P[1]
$$
を選ぶ。$C$ のコホモロジー層は次数 $\geq m - N - 1 - b$ で零である。
したがって $C \otimes^\mathbf{L} K$ のコホモロジー層は次数
$\geq m - N - 1$ で零である。ゆえに
$R\Gamma(X, C \otimes^\mathbf{L} K)$ のコホモロジーは次数
$\geq m - 1$ で零である。したがって
$$
R\Gamma(X, P \otimes^\mathbf{L} K) \to R\Gamma(X, E \otimes^\mathbf{L} K)
$$
は次数 $\geq m$ のコホモロジー上で同型である。仮定により始域は擬連接である。
したがって所望どおり $R\Gamma(X, E \otimes^\mathbf{L} K)$ は
$m$-擬連接である。
\end{proof}









\section{相対的完全対象}
\label{perfect-section-relatively-perfect}

\noindent
この節では \cite{lieblich-complexes} の概念を導入する。

\begin{definition}
\label{perfect-definition-relatively-perfect}
$f : X \to S$ を平坦かつ局所有限表示なスキームの射とする。
$D(\mathcal{O}_X)$ の対象 $E$ が擬連接であり
（Cohomology, Definition \ref{cohomology-definition-pseudo-coherent}）、
かつ $E$ が局所的に $D(f^{-1}\mathcal{O}_S)$ の対象として有限 tor 次元を
もつとき（Cohomology, Definition
\ref{cohomology-definition-tor-amplitude}）、$E$ は
{\it $S$ に関して完全}、または {\it $S$-完全} であるという。
\end{definition}

\noindent
この概念の議論については Remark \ref{perfect-remark-discuss-rel-perfect} を参照されたい。

\begin{example}
\label{perfect-example-relatively-perfect-field}
$k$ を体とする。$X$ を $k$ 上有限表示なスキームとする
（特に $X$ は準コンパクトである）。このとき $D(\mathcal{O}_X)$ の対象 $E$ が
$k$-完全であるための必要十分条件は、それが有界かつ擬連接であることであり
（定義による）、これは対象が $D^b_{\textit{Coh}}(X)$ に属することと同値である
（補題 \ref{perfect-lemma-identify-pseudo-coherent-noetherian} による）。したがって
相対的完全であることは「ファイバー上完全」であることを意味 {\bf しない}。
\end{example}

\noindent
対応する代数的概念は More on Algebra, Section
\ref{more-algebra-section-relatively-perfect} で研究されている。
スキームに対する概念と代数的概念は次のように結びつく。

\begin{lemma}
\label{perfect-lemma-affine-locally-rel-perfect}
$f : X \to S$ を平坦かつ局所有限表示なスキームの射とし、$E$ を
$D_\QCoh(\mathcal{O}_X)$ の対象とする。次は同値である。
\begin{enumerate}
\item $E$ は $S$-完全である。
\item アフィン開集合 $V \subset S$ に写る任意のアフィン開集合
$U \subset X$ に対して、複体 $R\Gamma(U, E)$ は
$\mathcal{O}_S(V)$-完全である。
\item アフィン開被覆 $S = \bigcup V_i$ と、各 $i$ に対するアフィン開被覆
$f^{-1}(V_i) = \bigcup U_{ij}$ で、複体 $R\Gamma(U_{ij}, E)$ が
$\mathcal{O}_S(V_i)$-完全となるものが存在する。
\end{enumerate}
\end{lemma}

\begin{proof}
擬連接であることは局所的性質であり、「局所的に有限 tor 次元をもつ」ことも
局所的性質である。したがってこの補題は直ちに次の主張へ帰着する：
$X$ と $S$ がアフィンなら、$E$ が $S$-完全であるための必要十分条件は、
$K = R\Gamma(X, E)$ が $\mathcal{O}_S(S)$-完全であることである。
$X = \Spec(A)$、$S = \Spec(R)$ と書き、$E$ が $K \in D(A)$ に対応するとする。
すなわち $K = R\Gamma(X, E)$ である。補題
\ref{perfect-lemma-affine-compare-bounded} を参照されたい。

\medskip\noindent
$K$ が $R$-完全であるための必要十分条件は、$K$ が擬連接で、かつ
$R$-加群の複体として有限 tor 次元をもつことであることに注意する
（More on Algebra, Definition
\ref{more-algebra-definition-relatively-perfect}）。補題
\ref{perfect-lemma-pseudo-coherent-affine} により、$E$ が擬連接であるための
必要十分条件は $K$ が擬連接であることである。補題
\ref{perfect-lemma-tor-dimension-rel-affine} により、$E$ が
$f^{-1}\mathcal{O}_S$ 上有限 tor 次元をもつための必要十分条件は、
$K$ が $R$-加群の複体として有限 tor 次元をもつことである。
\end{proof}

\begin{lemma}
\label{perfect-lemma-triangulated}
$f : X \to S$ を平坦かつ局所有限表示なスキームの射とする。
$D(\mathcal{O}_X)$ の $S$-完全対象からなる充満部分圏は、飽和
\footnote{Derived Categories, Definition
\ref{derived-definition-saturated}.} 三角部分圏である。
\end{lemma}

\begin{proof}
これは Cohomology, Lemmas
\ref{cohomology-lemma-cone-pseudo-coherent},
\ref{cohomology-lemma-summands-pseudo-coherent},
\ref{cohomology-lemma-cone-tor-amplitude}, および
\ref{cohomology-lemma-summands-tor-amplitude} から従う。
\end{proof}

\begin{lemma}
\label{perfect-lemma-perfect-relatively-perfect}
$f : X \to S$ を平坦かつ局所有限表示なスキームの射とする。
$D(\mathcal{O}_X)$ の完全対象は $S$-完全である。
$K, M \in D(\mathcal{O}_X)$ とする。$K$ が完全で $M$ が $S$-完全なら、
$K \otimes_{\mathcal{O}_X}^\mathbf{L} M$ は $S$-完全である。
\end{lemma}

\begin{proof}
第一の証明：補題 \ref{perfect-lemma-affine-locally-rel-perfect} により
アフィンの場合へ帰着し、More on Algebra, Lemma
\ref{more-algebra-lemma-perfect-relatively-perfect} を適用する。
\end{proof}

\begin{lemma}
\label{perfect-lemma-base-change-relatively-perfect}
$f : X \to S$ を平坦かつ局所有限表示なスキームの射とする。
$g : S' \to S$ をスキームの射とする。$X' = S' \times_S X$ とおき、
射影を $g' : X' \to X$ と記す。
$K \in D(\mathcal{O}_X)$ が $S$-完全なら、$L(g')^*K$ は $S'$-完全である。
\end{lemma}

\begin{proof}
第一の証明：補題 \ref{perfect-lemma-affine-locally-rel-perfect} により
アフィンの場合へ帰着し、More on Algebra, Lemma
\ref{more-algebra-lemma-base-change-relatively-perfect} を適用する。

\medskip\noindent
第二の証明：Cohomology, Lemma
\ref{cohomology-lemma-pseudo-coherent-pullback} により $L(g')^*K$ は
擬連接であり、tor 次元の有界性は補題
\ref{perfect-lemma-tor-independence-and-tor-amplitude} から従う。
\end{proof}

\begin{situation}
\label{perfect-situation-relative-descent}
$S = \lim_{i \in I} S_i$ を、アフィン遷移射
$g_{i'i} : S_{i'} \to S_i$ をもつスキームの有向系の極限とする。
すべての $i \in I$ に対して $S_i$ は準コンパクトかつ準分離であると仮定する。
射影を $g_i : S \to S_i$ と記す。元 $0 \in I$ と有限表示平坦射
$X_0 \to S_0$ を固定する。$X_i = S_i \times_{S_0} X_0$ および
$X = S \times_{S_0} X_0$ とおき、遷移射を
$f_{i'i} : X_{i'} \to X_i$、射影を $f_i : X \to X_i$ と記す。
\end{situation}

\begin{lemma}
\label{perfect-lemma-relative-descend-homomorphisms}
状況 \ref{perfect-situation-relative-descent} において、$K_0$ と $L_0$ を
$D(\mathcal{O}_{X_0})$ の対象とする。$i \geq 0$ に対して
$K_i = Lf_{i0}^*K_0$ および $L_i = Lf_{i0}^*L_0$ とおき、さらに
$K = Lf_0^*K_0$ および $L = Lf_0^*L_0$ とおく。このとき、$K_0$ が
擬連接であり、$L_0 \in D_\QCoh(\mathcal{O}_{X_0})$ が
$D((X_0 \to S_0)^{-1}\mathcal{O}_{S_0})$ の対象として（局所的に）
有限 tor 次元をもつなら、写像
$$
\colim_{i \geq 0} \Hom_{D(\mathcal{O}_{X_i})}(K_i, L_i)
\longrightarrow
\Hom_{D(\mathcal{O}_X)}(K, L)
$$
は同型である。
\end{lemma}

\begin{proof}
各準コンパクト開集合 $U_0 \subset X_0$ に対し、写像
$$
\colim_{i \geq 0} \Hom_{D(\mathcal{O}_{U_i})}(K_i|_{U_i}, L_i|_{U_i})
\longrightarrow
\Hom_{D(\mathcal{O}_U)}(K|_U, L|_U)
$$
が同型であるという条件を $P$ とする。ここで $U = f_0^{-1}(U_0)$ および
$U_i = f_{i0}^{-1}(U_0)$ である。$U_0$、$V_0$、および
$U_0 \cap V_0$ に対して $P$ が成り立つなら、導来圏における hom の
Mayer--Vietoris により $U_0 \cup V_0$ に対しても成り立つ。
Cohomology, Lemma \ref{cohomology-lemma-mayer-vietoris-hom} を参照されたい。

\medskip\noindent
与えられた射を $\pi_0 : X_0 \to S_0$ と記す。まず、準コンパクト開集合
$W_0 \subset S_0$ に対する $U_0 = \pi_0^{-1}(W_0)$ を考える。
$S_0$ の準コンパクト開集合に Cohomology of Schemes, Lemma
\ref{coherent-lemma-induction-principle} の帰納原理を適用し、上の注意を用いると、
$W_0$ がアフィンである $U_0 = \pi_0^{-1}(W_0)$ に対して $P$ を
証明すれば十分である。すなわち $S_0$ がアフィンの場合へ帰着した。
次に、今度は $X_0$ のすべての準コンパクトかつ準分離な開集合に
再び帰納原理を適用し、$X_0$ もアフィンである場合へ帰着する。

\medskip\noindent
$X_0$ と $S_0$ がアフィンなら、結果は More on Algebra, Lemma
\ref{more-algebra-lemma-colimit-relatively-perfect} から従う。
実際、補題 \ref{perfect-lemma-pseudo-coherent} および
\ref{perfect-lemma-affine-compare-bounded} により、主張は加群の導来圏における
hom の計算へ移される。補題 \ref{perfect-lemma-pseudo-coherent-affine} により、
$K_0$ に対応する加群の複体は擬連接である。補題
\ref{perfect-lemma-tor-dimension-rel-affine} により、$L_0$ に対応する加群の複体は
$\mathcal{O}_{S_0}(S_0)$ 上有限 tor 次元をもつ。したがって More on Algebra,
Lemma \ref{more-algebra-lemma-colimit-relatively-perfect} の仮定が満たされ、
結果が従う。
\end{proof}

\begin{lemma}
\label{perfect-lemma-descend-relatively-perfect}
状況 \ref{perfect-situation-relative-descent} において、$D(\mathcal{O}_X)$ の
$S$-完全対象の圏は、$D(\mathcal{O}_{X_i})$ の $S_i$-完全対象の圏の
余極限である。
\end{lemma}

\begin{proof}
各準コンパクト開集合 $U_0 \subset X_0$ に対し、関手
$$
\colim_{i \geq 0} D_{S_i\text{-perfect}}(\mathcal{O}_{U_i})
\longrightarrow
D_{S\text{-perfect}}(\mathcal{O}_U)
$$
が同値であるという条件を $P$ とする。ここで $U = f_0^{-1}(U_0)$ および
$U_i = f_{i0}^{-1}(U_0)$ である。補題
\ref{perfect-lemma-relative-descend-homomorphisms} により、この関手がすでに
充満忠実であることが分かっている。したがって本質的全射性を示せば十分である。

\medskip\noindent
$X_0$ の準コンパクト開集合 $U_0$ と $V_0$ に対して $P$ が成り立つとする。
$U_0 \cup V_0$ に対して $P$ が成り立つことを主張する。記法
$U_i = f_{i0}^{-1}U_0$、$U = f_0^{-1}U_0$、$V_i = f_{i0}^{-1}V_0$、
および $V = f_0^{-1}V_0$ を用いる。また、射 $U \to U_i$、$V \to V_i$、
$U \cap V \to U_i \cap V_i$、および $U \cup V \to U_i \cup V_i$ の
すべてを濫用して $f_i$ と記す。$E$ を $D(\mathcal{O}_{U \cup V})$ の
$S$-完全対象とする。目標は、$E$ が関手の本質像に属することを示すことである。
仮定により、ある $i \geq 0$、$U_i$ 上の $S_i$-完全対象 $E_{U, i}$、
$V_i$ 上の $S_i$-完全対象 $E_{V, i}$、および同型
$Lf_i^*E_{U, i} \to E|_U$ と $Lf_i^*E_{V, i} \to E|_V$ を取れる。
写像
$$
a : E_{U, i} \to (Rf_{i, *}E)|_{U_i}
\quad\text{かつ}\quad
b : E_{V, i} \to (Rf_{i, *}E)|_{V_i}
$$
を、同型 $Lf_i^*E_{U, i} \to E|_U$ および
$Lf_i^*E_{V, i} \to E|_V$ に随伴する写像とする。
充満忠実性により、$i$ を大きくすれば、引き戻したときに同一視
に一致する同型
$c : E_{U, i}|_{U_i \cap V_i} \to E_{V, i}|_{U_i \cap V_i}$
を取れる：
$$
Lf_i^*E_{U, i}|_{U \cap V} \to E|_{U \cap V} \to Lf_i^*E_{V, i}|_{U \cap V}.
$$
Cohomology, Lemma \ref{cohomology-lemma-glue} を適用すると、
$U_i \cup V_i$ 上の対象 $E_i$ と写像 $d : E_i \to Rf_{i, *}E$ で、
$U_i$ および $V_i$ 上への制限がそれぞれ $a$ および $b$ となるものを得る。
このとき $E_i$ が $S_i$-完全であり（相対的完全であることは局所的性質なので）、
$d$ が同型 $Lf_i^*E_i \to E$ に随伴することは明らかである。

\medskip\noindent
補題 \ref{perfect-lemma-relative-descend-homomorphisms} の証明で用いたものと
まったく同じ帰納原理の議論（Cohomology of Schemes, Lemma
\ref{coherent-lemma-induction-principle}）により、$X_0$ と $S_0$ がともに
アフィンである場合へ帰着する（まず $S_0$ の開集合を扱って $S_0$ が
アフィンの場合へ帰着し、次に $X_0$ の開集合を扱って $X_0$ が
アフィンの場合へ帰着する）。アフィンの場合、結果は More on Algebra,
Lemma \ref{more-algebra-lemma-colimit-relatively-perfect} から従う。
代数への翻訳は補題 \ref{perfect-lemma-affine-locally-rel-perfect} による。
\end{proof}

\begin{lemma}
\label{perfect-lemma-derived-pushforward-rel-perfect}
$f : X \to S$ を有限表示な平坦固有射とする。
$E \in D(\mathcal{O}_X)$ を $S$-完全とする。このとき $Rf_*E$ は
$D(\mathcal{O}_S)$ の完全対象であり、その形成は任意の基底変換と可換する。
\end{lemma}

\begin{proof}
基底変換に関する主張は補題 \ref{perfect-lemma-compare-base-change} である。
したがって $Rf_*E$ が完全対象であることを示せば十分である。
極限の議論により $S$ がネーターかつアフィンである場合へ帰着する。

\medskip\noindent
問題は $S$ 上局所的なので、$S$ はアフィンと仮定してよい。
$S = \Spec(R)$ と書く。$R = \colim R_i$ をネーター環 $R_i$ の
フィルター余極限として書く。Limits, Lemma
\ref{limits-lemma-descend-finite-presentation}
により、ある $i$ と $R_i$ 上有限表示なスキーム $X_i$ で、$R$ への
基底変換が $X$ となるものが存在する。Limits, Lemmas
\ref{limits-lemma-eventually-proper} および
\ref{limits-lemma-descend-flat-finite-presentation}
により、$X_i$ は $R_i$ 上固有かつ平坦であると仮定してよい。
補題 \ref{perfect-lemma-descend-relatively-perfect} により、
$D(\mathcal{O}_{X_i})$ の $R_i$-完全対象 $E_i$ で、$X$ への引き戻しが
$E$ となるものが存在すると仮定してよい。$X_i \to \Spec(R_i)$ と
$E_i$ に補題 \ref{perfect-lemma-perfect-direct-image} を適用し、すでに示した
基底変換性を用いれば結果を得る。
\end{proof}

\begin{lemma}
\label{perfect-lemma-compute-ext-rel-perfect}
$f : X \to S$ をスキームの射とし、$E, K \in D(\mathcal{O}_X)$ とする。
次を仮定する。
\begin{enumerate}
\item $S$ は準コンパクトかつ準分離である。
\item $f$ は固有、平坦、かつ有限表示である。
\item $E$ は $S$-完全である。
\item $K$ は擬連接である。
\end{enumerate}
このとき擬連接な $L \in D(\mathcal{O}_S)$ で
$$
Rf_*R\SheafHom(K, E) = R\SheafHom(L, \mathcal{O}_S)
$$
を満たすものが存在し、同じことが任意の基底変換後にも成り立つ。すなわち
$$
\vcenter{
\xymatrix{
X' \ar[r]_{g'} \ar[d]_{f'} &
X \ar[d]^f \\
S' \ar[r]^g &
S
}
}
\quad\quad
\begin{matrix}
\text{がカルテシアンなら } \\
Rf'_*R\SheafHom(L(g')^*K, L(g')^*E) \\
= R\SheafHom(Lg^*L, \mathcal{O}_{S'})
\end{matrix}
$$
\end{lemma}

\begin{proof}
$S$ は準コンパクトかつ準分離なので、$X$ もそうである。
補題 \ref{perfect-lemma-pseudo-coherent-hocolim} により
$K = \text{hocolim} K_n$ と書ける。ここで $K_n$ は完全であり、
$K_n \to K$ は切断 $\tau_{\geq -n}$ 上の同型を誘導する。
$K_n^\vee$ を双対完全複体とする（Cohomology, Lemma
\ref{cohomology-lemma-dual-perfect-complex}）。完全対象の逆系
$\ldots \to K_3^\vee \to K_2^\vee \to K_1^\vee$ を得る。
補題 \ref{perfect-lemma-perfect-relatively-perfect} により
$K_n^\vee \otimes_{\mathcal{O}_X} E$ は $S$-完全である。したがって
$K_n^\vee \otimes_{\mathcal{O}_X} E$ に補題
\ref{perfect-lemma-derived-pushforward-rel-perfect} を適用すると、$S$ 上の
完全複体の逆系
$$
\ldots \to M_3 \to M_2 \to M_1
$$
で、
$$
M_n = Rf_*(K_n^\vee \otimes_{\mathcal{O}_X}^\mathbf{L} E) =
Rf_*R\SheafHom(K_n, E)
$$
を満たすものを得る。さらに、これらの複体の形成は任意の基底変換と可換する。
すなわち $Lg^*M_n =
Rf'_*((L(g')^*K_n)^\vee \otimes_{\mathcal{O}_{X'}}^\mathbf{L} L(g')^*E) =
Rf'_*R\SheafHom(L(g')^*K_n, L(g')^*E)$.

\medskip\noindent
$K_n \to K$ は $\tau_{\geq -n}$ 上の同型を誘導するので、
$K_n \to K_{n + 1}$ も $\tau_{\geq -n}$ 上の同型を誘導する。
$K_n^\vee = R\SheafHom(K_n, \mathcal{O}_X)$ なので、
$K_{n + 1}^\vee \to K_n^\vee$ は $\tau_{\leq n}$ 上の同型を誘導する。
$E$ が $f^{-1}\mathcal{O}_Y$-加群の複体として $[a, b]$ 内の tor 振幅を
もつと仮定する。このとき任意の基底変換後にも同じことが成り立つ。
補題 \ref{perfect-lemma-tor-independence-and-tor-amplitude} を参照されたい。
したがって
$K_{n + 1}^\vee \otimes_{\mathcal{O}_X} E \to
K_n^\vee \otimes_{\mathcal{O}_X} E$
は $\tau_{\leq n + a}$ 上の同型を誘導し、任意の基底変換後にも
同じことが成り立つ。右導来関手 $Rf_*$ を適用すると、写像
$M_{n + 1} \to M_n$ は $\tau_{\leq n + a}$ 上の同型を誘導し、
任意の基底変換後にも同じことが成り立つ。識別三角
$$
M_{n + 1} \to M_n \to C_n \to M_{n + 1}[1]
$$
を選ぶ。$S'$ を点 $s \in S$ の剰余体のスペクトルとし、引き戻すと、
$C_n \otimes_{\mathcal{O}_S}^\mathbf{L} \kappa(s)$ は次数
$\geq n + a$ にのみ非零なコホモロジーをもつことが分かる。
More on Algebra, Lemma
\ref{more-algebra-lemma-lift-perfect-from-residue-field}
により、完全複体 $C_n$ はある整数 $m_n$ に対して
$[n + a, m_n]$ 内の tor 振幅をもつ。特に、双対完全複体 $C_n^\vee$ は
$[-m_n, -n - a]$ 内の tor 振幅をもつ。

\medskip\noindent
$L_n = M_n^\vee$ を双対完全複体とする。前段落の議論から、
$L_n \to L_{n + 1}$ は $\tau_{\geq -n - a}$ 上の同型を誘導する。
したがって $L = \text{hocolim} L_n$ は擬連接である。補題
\ref{perfect-lemma-pseudo-coherent-hocolim} を参照されたい。また
$$
R\SheafHom(K, E) = R\SheafHom(\text{hocolim} K_n, E) =
R\lim R\SheafHom(K_n, E) = R\lim K_n^\vee \otimes_{\mathcal{O}_X} E
$$
であり（Cohomology, Lemma
\ref{cohomology-lemma-colim-and-lim-of-duals}）、$R\lim$ は $Rf_*$ と
可換するので
$$
Rf_*R\SheafHom(K, E) = R\lim M_n = R\lim R\SheafHom(L_n, \mathcal{O}_S) =
R\SheafHom(L, \mathcal{O}_S)
$$
を得る。これで $S$ 上の公式が証明された。$M_n$ の構成は基底変換と
両立するので、この公式は任意の基底変換後にも成り立つ。
\end{proof}

\begin{remark}
\label{perfect-remark-compare-L}
読者は補題 \ref{perfect-lemma-compute-ext-rel-perfect} と補題
\ref{perfect-lemma-compute-ext} の類似に気づいたかもしれない。実際、補題
\ref{perfect-lemma-compute-ext-rel-perfect} の擬連接複体 $L$ は、関手的同型
$$
\Ext^i_{\mathcal{O}_S}(L, \mathcal{F}) \longrightarrow
\Ext^i_{\mathcal{O}_X}(K,
E \otimes_{\mathcal{O}_X}^\mathbf{L} Lf^*\mathcal{F})
$$
が $\QCoh(\mathcal{O}_S)$ を動く $\mathcal{F}$ に対して存在し、境界写像と
両立するような $S$ 上の一意な擬連接複体として特徴づけられる。
これが必要になれば、ここに精密な結果を定式化して詳しい証明を与える。
\end{remark}

\begin{lemma}
\label{perfect-lemma-bounded-on-fibres}
$f : X \to S$ を平坦かつ局所有限表示なスキームの射とする。
$E$ を $D(\mathcal{O}_X)$ の擬連接対象とする。次は同値である。
\begin{enumerate}
\item $E$ は $S$-完全である。
\item $E$ は局所的に下に有界であり、すべての点 $s \in S$ に対して
$D(\mathcal{O}_{X_s})$ の対象 $L(X_s \to X)^*E$ は局所的に下に有界である。
\end{enumerate}
\end{lemma}

\begin{proof}
すべてが局所的なので、補題 \ref{perfect-lemma-affine-locally-rel-perfect} により
$X$ と $S$ がアフィンである場合へ直ちに帰着する。
$X \to S$ が $\Spec(A) \to \Spec(R)$ に対応し、$E$ が $D(A)$ の
$K$ に対応するとする。$s$ が素イデアル $\mathfrak p \subset R$ に対応するなら、
$R \to A$ は平坦なので、$L(X_s \to X)^*E$ は
$K \otimes_R^\mathbf{L} \kappa(\mathfrak p)$ に対応する。例えば補題
\ref{perfect-lemma-compare-base-change} を参照されたい。したがって補題は対応する
代数的結果 More on Algebra, Lemma
\ref{more-algebra-lemma-bounded-on-fibres-relatively-perfect} から従う。
\end{proof}

\begin{remark}
\label{perfect-remark-discuss-rel-perfect}
相対的完全複体に関する Definition \ref{perfect-definition-relatively-perfect} は、
本書の定義が適用できる場合には \cite{lieblich-complexes} の定義と同値である
\footnote{これを見るには、補題 \ref{perfect-lemma-affine-locally-rel-perfect} と
More on Algebra, Lemma \ref{more-algebra-lemma-structure-relatively-perfect} を用いる。}。
次に、$f : X \to S$ について局所有限型であることだけを仮定する
（平坦または局所有限表示とは限らない）。上で引用した論文の定義では、
$E \in D(\mathcal{O}_X)$ が相対的完全であるとは、次を満たすことである。
\begin{enumerate}
\item[(A)] $X$ 上局所的に、対象 $E$ は $S$ 上平坦な有限表示
$\mathcal{O}_X$-加群からなる有限複体と準同型である。
\end{enumerate}
一方、Definition \ref{perfect-definition-relatively-perfect} の自然な一般化は次である。
\begin{enumerate}
\item[(B)] $E$ は $S$ に関して擬連接であり（More on Morphisms, Definition
\ref{more-morphisms-definition-relative-pseudo-coherence}）、かつ $E$ は局所的に
$D(f^{-1}\mathcal{O}_S)$ の対象として有限 tor 次元をもつ
（Cohomology, Definition \ref{cohomology-definition-tor-amplitude}）。
\end{enumerate}
条件 (B) の利点は、それが明らかに $D(\mathcal{O}_X)$ の三角部分圏を
定めることであるが、条件 (A) についてはそうでないと予想される。
条件 (A) の利点は、特に極限に関して扱いやすいことである。
\end{remark}









\section{分解性質}
\label{perfect-section-resolution-property}

\noindent
この概念は論文 \cite{totaro_resolution} で論じられ、その議論は
\cite{Gross-thesis}、\cite{Gross-surface}、および \cite{Gross-stack}
で続けられている。体上の固有スキームが常に分解性質をもつのか、
それともそうでないのかは現在知られていない。この問いの答えを知っているなら、
\href{mailto:stacks.project@gmail.com}{stacks.project@gmail.com} へ連絡されたい。

\medskip\noindent
一般のスキームについて考える意味はほとんどないが、次の定義を置ける。

\begin{definition}
\label{perfect-definition-resolution-property}
$X$ をスキームとする。有限型準連接 $\mathcal{O}_X$-加群がすべて有限局所自由
$\mathcal{O}_X$-加群の商であるとき、$X$ は {\it 分解性質} をもつという。
\end{definition}

\noindent
$X$ が準コンパクトかつ準分離なら、有限表示 $\mathcal{O}_X$-加群
（自動的に準連接）がすべて有限局所自由 $\mathcal{O}_X$-加群の商であることを
確認すれば十分である。Properties, Lemma
\ref{properties-lemma-finite-directed-colimit-surjective-maps} を参照されたい。
$X$ がネーターなら、有限型準連接加群はちょうどコヒーレント
$\mathcal{O}_X$-加群である。Cohomology of Schemes, Lemma
\ref{coherent-lemma-coherent-Noetherian} を参照されたい。

\begin{lemma}
\label{perfect-lemma-resolution-property-ample}
$X$ をスキームとする。$X$ が豊富な可逆 $\mathcal{O}_X$-加群をもつなら、
$X$ は分解性質をもつ。
\end{lemma}

\begin{proof}
Properties, Proposition \ref{properties-proposition-characterize-ample}
の直ちに分かる帰結である。
\end{proof}

\begin{lemma}
\label{perfect-lemma-resolution-property-ample-relative}
$f : X \to Y$ をスキームの射とする。次を仮定する。
\begin{enumerate}
\item $Y$ は準コンパクトかつ準分離であり、分解性質をもつ。
\item $X$ 上に $f$-豊富な可逆加群が存在する。
\end{enumerate}
このとき $X$ は分解性質をもつ。
\end{lemma}

\begin{proof}
$\mathcal{F}$ を有限型準連接 $\mathcal{O}_X$-加群とし、$\mathcal{L}$ を
$f$-豊富な可逆加群とする。アフィン開被覆
$Y = V_1 \cup \ldots \cup V_m$ を選び、$U_j = f^{-1}(V_j)$ とおく。
Properties, Proposition
\ref{properties-proposition-characterize-ample}
により、各 $j$ に対して有限個の写像
$s_{j, i} : \mathcal{L}^{\otimes n_{j, i}}|_{U_j} \to \mathcal{F}|_{U_j}$
で、合わせて全射となるものが存在する。準連接 $\mathcal{O}_Y$-加群
$$
\mathcal{H}_n =
f_*(\mathcal{F} \otimes_{\mathcal{O}_X} \mathcal{L}^{\otimes n})
$$
を考える。$s_{j, i}$ は層 $\mathcal{H}_{-n_{j, i}}$ の $V_j$ 上の切断と
みなせる。有限局所自由 $\mathcal{O}_Y$-加群 $\mathcal{E}_{i, j}$ と写像
$\mathcal{E}_{i, j} \to \mathcal{H}_{-n_{j, i}}$ で、$s_{j, i}$ が像に
属するものを取れたと仮定する。このとき対応する写像
$$
f^*\mathcal{E}_{i, j}
\otimes_{\mathcal{O}_X}
\mathcal{L}^{\otimes n_{i, j}} \longrightarrow \mathcal{F}
$$
は合わせて全射となり、補題が証明される。Properties, Lemma
\ref{properties-lemma-quasi-coherent-colimit-finite-type} により、各 $i, j$ に
対して有限型準連接部分加群
$\mathcal{H}'_{i, j} \subset  \mathcal{H}_{-n_{j, i}}$
で、$V_j$ 上の切断 $s_{i, j}$ を含むものを取れる。したがって $Y$ の
分解性質により全射 $\mathcal{E}_{i, j} \to \mathcal{H}'_{j, i}$ を得て、
結論が従う。
\end{proof}

\begin{lemma}
\label{perfect-lemma-resolution-property-goes-up-affine}
$f : X \to Y$ をアフィンまたは準アフィンなスキームの射とし、$Y$ は
準コンパクトかつ準分離であるとする。$Y$ が分解性質をもつなら、
$X$ も分解性質をもつ。
\end{lemma}

\begin{proof}
Morphisms, Lemma \ref{morphisms-lemma-quasi-affine-O-ample} により、
これは補題 \ref{perfect-lemma-resolution-property-ample-relative} の特別な場合である。
\end{proof}

\noindent
分解性質が下へ移ることを証明できる場合を一つ挙げる。

\begin{lemma}
\label{perfect-lemma-resolution-property-goes-down-finite-flat}
$f : X \to Y$ を全射な有限局所自由スキーム射とする。
$X$ が分解性質をもつなら、$Y$ も分解性質をもつ。
\end{lemma}

\begin{proof}
この条件は、$f$ がアフィンであり、$f_*\mathcal{O}_X$ が正の階数をもつ
有限局所自由 $\mathcal{O}_Y$-加群であることを意味する。
$\mathcal{G}$ を有限型準連接 $\mathcal{O}_Y$-加群とする。仮定により、
ある有限局所自由 $\mathcal{O}_X$-加群 $\mathcal{E}$ からの全射
$\mathcal{E} \to f^*\mathcal{G}$ が存在する。$f_*$ は準連接加群上完全なので
（Cohomology of Schemes, Lemma
\ref{coherent-lemma-relative-affine-vanishing}）、全射
$$
f_*\mathcal{E}
\longrightarrow
f_*f^*\mathcal{G} = \mathcal{G} \otimes_{\mathcal{O}_Y} f_*\mathcal{O}_X
$$
双対を取ると全射
$$
f_*\mathcal{E}
\otimes_{\mathcal{O}_Y}
\SheafHom_{\mathcal{O}_Y}(f_*\mathcal{O}_X, \mathcal{O}_Y)
\longrightarrow
\mathcal{G}
$$
を得る。$f_*\mathcal{E}$ は有限局所自由なので\footnote{実際、
$A \to B$ が有限局所自由な環準同型で、$N$ が有限局所自由
$B$-加群なら、$N$ は有限局所自由 $A$-加群である。これを見るには、まず
$N$ が $B$ 上有限局所自由であることから、$N$ は $B$-加群として
平坦かつ有限表示であることに注意する。Algebra, Lemma
\ref{algebra-lemma-finite-projective} を参照されたい。次に $N$ は
Algebra, Lemma \ref{algebra-lemma-finite-finitely-presented-extension} により
有限表示 $A$-加群であり、Algebra, Lemma
\ref{algebra-lemma-composition-flat} により平坦 $A$-加群である。
最後に $A$ 上で Algebra, Lemma \ref{algebra-lemma-finite-projective}
を用いればよい。}、結論が従う。
\end{proof}

\begin{lemma}
\label{perfect-lemma-resolution-property-finite-number}
$X$ をスキームとする。次が与えられているとする。
\begin{enumerate}
\item 有限アフィン開被覆 $X = U_1 \cup \ldots \cup U_m$。
\item $V(\mathcal{I}_j) = X \setminus U_j$ を満たす有限型準連接イデアル
$\mathcal{I}_j$。
\end{enumerate}
このとき $X$ が分解性質をもつための必要十分条件は、
$j = 1, \ldots, m$ に対して $\mathcal{I}_j$ が有限局所自由
$\mathcal{O}_X$-加群の商となることである。
\end{lemma}

\begin{proof}
補題の一方向は自明である。逆方向について、$\mathcal{E}_j$ が有限局所自由で、
$\mathcal{E}_j \to \mathcal{I}_j$ が全射であるとする。次の段落で
ネーターの場合へ帰着するが、読者には読み飛ばすことを勧める。

\medskip\noindent
まず、$U_j \cap U_{j'}$ は準コンパクトである。実際、これはアフィンスキーム
$U_j$ 上の有限型準連接イデアル $\mathcal{I}_{j'}|{U_j}$ の零点スキームの
補集合である。Properties, Lemma
\ref{properties-lemma-quasi-coherent-finite-type-ideals} を参照されたい。
したがって $X$ は準コンパクトかつ準分離である。Schemes, Lemma
\ref{schemes-lemma-characterize-quasi-separated} を参照されたい。
Limits, Proposition \ref{limits-proposition-approximate} により、$X = \lim X_i$ を
アフィン遷移射をもつネータースキームの有向系の極限として書ける。各 $j$ に
対して、ある $i$ と有限局所自由 $\mathcal{O}_{X_i}$-加群
$\mathcal{E}_{i, j}$ で、引き戻しが $\mathcal{E}_j$ となるものを取れる。
Limits, Lemma \ref{limits-lemma-descend-invertible-modules} を参照されたい。
$i$ を大きくすれば、合成 $\mathcal{E}_j \to \mathcal{I}_j \to \mathcal{O}_X$
は写像 $\mathcal{E}_{i, j} \to \mathcal{O}_{X_i}$ の引き戻しであると
仮定してよい。Limits, Lemma
\ref{limits-lemma-descend-modules-finite-presentation} を参照されたい。
この写像の像を $\mathcal{I}_{i, j} \subset \mathcal{O}_{X_i}$ と記す。
これはネータースキーム $X_i$ 上の準連接イデアル層であり、$X$ への引き戻しは
$\mathcal{I}_j$ である。補開集合を $U_{i, j} \subset X_i$ と記すと、
すべての $i, j$ に対してこれらはアフィンであると仮定してよい。
Limits, Lemma \ref{limits-lemma-limit-affine} を参照されたい。
$X_i$ 上の開集合 $U_{i, j}$ とイデアル層 $\mathcal{I}_{i, j}$ に対して
補題を証明できれば、$X$ は $X_i$ 上アフィンなので、補題
\ref{perfect-lemma-resolution-property-goes-up-affine} により分解性質をもつ。
このようにしてネータースキームの場合へ帰着する。

\medskip\noindent
$X$ がネーターであると仮定する。各コヒーレント加群 $\mathcal{F}$ に対して、
$\mathcal{F}$ の $U_j$ への制限を生成する有限個の切断
$s_{jk} \in \mathcal{F}(U_j)$、$k = 1, \ldots, e_j$ を選べる。
Cohomology of Schemes, Lemma \ref{coherent-lemma-homs-over-open} により、
ある $n_{jk} \geq 1$ に対する写像
$s'_{jk} : \mathcal{I}_i^{n_{jk}} \to \mathcal{F}$ へ $s_{jk}$ を延長できる。
そこで合成
$$
\mathcal{E}_j^{\otimes n_{jk}} \to \mathcal{I}_j^{n_{jk}} \to \mathcal{F}
$$
を考えれば結論が従う。
\end{proof}

\begin{lemma}
\label{perfect-lemma-resolution-property-ample-family}
$X$ をスキームとする。$X$ が可逆加群の豊富族をもつなら
（Morphisms, Definition
\ref{morphisms-definition-family-ample-invertible-modules}）、
$X$ は分解性質をもつ。
\end{lemma}

\begin{proof}
$X$ は準コンパクトなので、$n$ と組 $(\mathcal{L}_i, s_i)$、
$i = 1, \ldots, n$ で、$\mathcal{L}_i$ が可逆 $\mathcal{O}_X$-加群、
$s_i \in \Gamma(X, \mathcal{L}_i)$ が切断であり、$s_i$ が消えない点の集合
$U_i \subset X$ がアフィンで、$X = U_1 \cup \ldots \cup U_n$ となるものが
存在する。$\mathcal{I}_i \subset \mathcal{O}_X$ を
$s_i : \mathcal{L}_i^{\otimes -1} \to \mathcal{O}_X$ の像とする。
補題 \ref{perfect-lemma-resolution-property-finite-number} を適用すると、
$X$ が分解性質をもつことが分かる。
\end{proof}

\begin{lemma}
\label{perfect-lemma-regular-resolution-property}
$X$ をアフィン対角射をもつ準コンパクト正則スキームとする。
このとき $X$ は分解性質をもつ。
\end{lemma}

\begin{proof}
Divisors, Lemma \ref{divisors-lemma-regular-ample-family} と上の補題
\ref{perfect-lemma-resolution-property-ample-family} を組み合わせればよい。
\end{proof}

\begin{lemma}
\label{perfect-lemma-resolution-property-descends}
$X = \lim X_i$ を、アフィン遷移射をもつ準コンパクトかつ準分離な
スキームの有向系の極限とする。このとき $X$ が分解性質をもつための
必要十分条件は、ある $i$ に対して $X_i$ が分解性質をもつことである。
\end{lemma}

\begin{proof}
$X_i$ が分解性質をもつなら、補題
\ref{perfect-lemma-resolution-property-goes-up-affine} により $X$ もそうである。
$X$ が分解性質をもつと仮定する。$i \in I$ を選び、有限アフィン開被覆
$X_i = U_{i, 1} \cup \ldots \cup U_{i, m}$ を選ぶ。各 $j$ に対して、
$X_i \setminus V(\mathcal{I}_{i, j}) = U_{i, j}$ を満たす有限型準連接
イデアル層 $\mathcal{I}_{i, j} \subset \mathcal{O}_{X_i}$ を選ぶ。
Properties, Lemma \ref{properties-lemma-quasi-coherent-finite-type-ideals}
を参照されたい。$U_{i, j}$ の逆像を $U_j \subset X$ と記し、
$\mathcal{I}_{i, j}$ の引き戻しを $\mathcal{I}_j \subset \mathcal{O}_X$ と記す。
$X$ は分解性質をもつので、有限局所自由 $\mathcal{O}_X$-加群
$\mathcal{E}_j$ と全射 $\mathcal{E}_j \to \mathcal{I}_j$ を選べる。
Limits, Lemmas \ref{limits-lemma-descend-invertible-modules} および
\ref{limits-lemma-descend-modules-finite-presentation} により、$i$ を
大きくすれば、有限局所自由 $\mathcal{O}_{X_i}$-加群
$\mathcal{E}_{i, j}$ と写像 $\mathcal{E}_{i, j} \to \mathcal{O}_{X_i}$ で、
$X$ への基底変換が合成
$\mathcal{E}_j \to \mathcal{I}_j \to \mathcal{O}_X$、
$j = 1, \ldots, m$ を復元するものを取れる。有限表示
$\mathcal{O}_{X_i}$-加群
$\Coker(\mathcal{E}_{i, j} \to \mathcal{O}_{X_i})$ と
$\mathcal{O}_{X_i}/\mathcal{I}_{i, j}$ の $X$ への引き戻しは、
$\mathcal{O}_X$ の商として一致する。したがって Limits, Lemma
\ref{limits-lemma-descend-modules-finite-presentation} により、これらが
一致すると仮定してよい。言い換えると、
$\mathcal{E}_{i, j} \to \mathcal{O}_{X_i}$ の像は
$\mathcal{I}_{i, j}$ に等しい。補題
\ref{perfect-lemma-resolution-property-finite-number} により、$X_i$ は
分解性質をもつ。
\end{proof}

\begin{lemma}
\label{perfect-lemma-resolution-property-affine-diagonal}
\begin{reference}
\cite[Proposition 1.3]{totaro_resolution} の特別な場合である。
\end{reference}
$X$ を分解性質をもつ準コンパクトかつ準分離なスキームとする。
このとき $X$ の対角射はアフィンである。
\end{lemma}

\begin{proof}
Limits, Proposition \ref{limits-proposition-approximate} と補題
\ref{perfect-lemma-resolution-property-descends} を組み合わせると、$X$ が
ネーターである場合へ帰着する（細部を一つ省略する）。$X$ がネーターであると
仮定する。$X \times X$ は、$X$ のアフィン開集合 $U$, $V$ に対する
アフィン開集合 $U \times V$ で被覆されることを思い出そう。Schemes, Section
\ref{schemes-section-fibre-products} を参照されたい。したがって対角射
$\Delta : X \to X \times X$ がアフィンであることを示すには、$X$ の
すべてのアフィン開集合 $U$, $V$ に対して
$U \cap V = \Delta^{-1}(U \times V)$ がアフィンであることを示せば十分である。
Morphisms, Lemma \ref{morphisms-lemma-characterize-affine} を参照されたい。
特に、$U$ が $X$ のアフィン開集合なら、包含射 $j : U \to X$ がアフィンである
ことを示せば十分である。Cohomology of Schemes, Lemma
\ref{coherent-lemma-criterion-affine-morphism} により、任意の準連接
$\mathcal{O}_U$-加群 $\mathcal{G}$ に対して $R^1j_*\mathcal{G} = 0$ を
示せば十分である。Proposition \ref{perfect-proposition-Noetherian}
（ここでネーターの場合へ帰着したことを用いる）により、
$Rj_*\mathcal{G}$ は準連接 $\mathcal{O}_X$-加群からなる複体
$\mathcal{H}^\bullet$ で表せる。背理法のため
$$
H^1(\mathcal{H}^\bullet) =
\Ker(\mathcal{H}^1 \to \mathcal{H}^2)/\Im(\mathcal{H}^0 \to \mathcal{H}^1)
$$
が非零であると仮定する。このときコヒーレント $\mathcal{O}_X$-加群
$\mathcal{F}$ と写像
$$
\mathcal{F} \longrightarrow
\Ker(\mathcal{H}^1 \to \mathcal{H}^2)
$$
で、$H^1(\mathcal{H}^\bullet)$ への射影との合成が非零となるものを取れる。
実際、Properties, Lemma
\ref{properties-lemma-quasi-coherent-colimit-finite-type} により
$\Ker(\mathcal{H}^1 \to \mathcal{H}^2)$ をそのコヒーレント部分加群の
フィルター合併として書け、そのうち一つが条件を満たす。次に $X$ の
分解性質を用いて、有限局所自由 $\mathcal{O}_X$-加群 $\mathcal{E}$ と
全射 $\mathcal{E} \to \mathcal{F}$ を選ぶ。これにより導来圏の写像
$$
\mathcal{E}[-1] \longrightarrow Rj_*\mathcal{G}
$$
を得る。この写像はコホモロジー層上で非零なので $D(\mathcal{O}_X)$ において
非零である。随伴により、これは写像
$$
j^*\mathcal{E}[-1] \to \mathcal{G}
$$
が $D(\mathcal{O}_U)$ において非零であることと同じである。
$\mathcal{E}$ は有限局所自由なので、これは
$$
H^1(U, j^*\mathcal{E}^\vee \otimes_{\mathcal{O}_U} \mathcal{G})
$$
の非零元と同じである。ここで
$\mathcal{E}^\vee = \SheafHom_{\mathcal{O}_X}(\mathcal{E}, \mathcal{O}_X)$
は双対有限局所自由加群である。しかし、この群は Cohomology of Schemes,
Lemma
\ref{coherent-lemma-quasi-coherent-affine-cohomology-zero}
により零であり、所望の矛盾を得る（双対を用いる段階に疑問があれば、
より一般の Cohomology, Lemma
\ref{cohomology-lemma-dual-perfect-complex} を参照されたい）。
\end{proof}






\section{分解性質と完全複体}
\label{perfect-section-resolution-property-perfect}

\noindent
本節では、分解性質をもつスキーム上の完全複体と厳密完全複体との関係を論じる。

\begin{lemma}
\label{perfect-lemma-construct-strictly-perfect}
$X$ を分解性質をもつ準コンパクトかつ準分離なスキームとする。
$\mathcal{F}^\bullet$ を $D(\mathcal{O}_X)$ の完全対象を表す、準連接
$\mathcal{O}_X$-加群の下に有界な複体とする。このとき、有限局所自由
$\mathcal{O}_X$-加群の有界複体 $\mathcal{E}^\bullet$ および擬同型
$\mathcal{E}^\bullet \to \mathcal{F}^\bullet$ が存在する。
\end{lemma}

\begin{proof}
$\mathcal{F}^\bullet$ が $[a, b]$ に Tor 振幅をもち、かつ $n < a$ に対して
$\mathcal{F}^n = 0$ となるような整数 $a, b \in \mathbf{Z}$ を取る。
そのような整数の組が存在することは Cohomology, Lemma
\ref{cohomology-lemma-perfect} と $X$ の準コンパクト性から従う。
$b < a$ なら $\mathcal{F}^\bullet$ は導来圏で零であり、補題は成立する。
$b - a \geq 0$ に関する帰納法により、複体
$\mathcal{E}^a \to \ldots \to \mathcal{E}^b$
で各 $\mathcal{E}^i$ が有限局所自由であり、擬同型
$\mathcal{E}^\bullet \to \mathcal{F}^\bullet$ をもつものが存在することを示す。

\medskip\noindent
基底の場合は $b - a = 0$ の場合である。このとき
$H^b(\mathcal{F}^\bullet) = H^a(\mathcal{F}^\bullet) =
\Ker(\mathcal{F}^a \to \mathcal{F}^{a + 1})$
は有限局所自由である。実際、これは Tor 次元 $0$ の有限表示
$\mathcal{O}_X$-加群であり、したがって有限局所自由である。Cohomology,
Lemmas \ref{cohomology-lemma-perfect} および
\ref{cohomology-lemma-finite-cohomology} および
Properties, Lemma \ref{properties-lemma-finite-locally-free}
を参照されたい。したがって $\mathcal{E}^\bullet$ として、次数 $b$ に置いた
$H^b(\mathcal{F}^\bullet)$ を取れる。証明の残りでは帰納段階を扱う。

\medskip\noindent
$b > a$ と仮定する。ここで
$$
H^b(\mathcal{F}^\bullet) =
\Ker(\mathcal{F}^b \to \mathcal{F}^{b + 1})/
\Im(\mathcal{F}^{b - 1} \to \mathcal{F}^b)
$$
は有限型準連接 $\mathcal{O}_X$-加群である。Cohomology, Lemmas
\ref{cohomology-lemma-perfect} および
\ref{cohomology-lemma-finite-cohomology} を参照されたい。したがって、有限型
準連接 $\mathcal{O}_X$-加群 $\mathcal{F}$ と写像
$$
\mathcal{F} \longrightarrow
\Ker(\mathcal{F}^b \to \mathcal{F}^{b + 1})
$$
で、$H^b(\mathcal{F}^\bullet)$ への射影との合成が全射となるものを取れる。
実際、Properties, Lemma
\ref{properties-lemma-quasi-coherent-colimit-finite-type} により
$\Ker(\mathcal{F}^b \to \mathcal{F}^{b + 1})$ をその有限型準連接部分加群の
フィルター合併として書け、そのうち一つが条件を満たす。次に $X$ の分解性質を
用いて、有限局所自由 $\mathcal{O}_X$-加群 $\mathcal{E}^b$ と全射
$\mathcal{E}^b \to \mathcal{F}$ を選ぶ。複体の写像
$$
\alpha : \mathcal{E}^b[-b] \to \mathcal{F}^\bullet
$$
とその錐 $C(\alpha)^\bullet$ を考える。Derived Categories, Definition
\ref{derived-definition-cone} を参照されたい。$C(\alpha)^\bullet$ は次数
$\geq a$ でのみ非零であり、Cohomology, Lemma
\ref{cohomology-lemma-cone-tor-amplitude} により $[a, b]$ に Tor 振幅をもち、
構成により $H^b(C(\alpha)^\bullet) = 0$ である。したがって実際には
$C(\alpha)^\bullet$ は $[a, b - 1]$ に Tor 振幅をもつ。それゆえ帰納法の仮定を
$C(\alpha)^\bullet$ に適用でき、複体の写像
$$
(\mathcal{E}^a \to \ldots \to \mathcal{E}^{b - 1})
\longrightarrow
C(\alpha)^\bullet
$$
で帰納法の仮定に述べた性質をもつものを得る。錐の定義を展開すると、これは
可換図式
$$
\xymatrix{
\ldots \ar[r] &
\mathcal{E}^{b - 2} \ar[r] \ar[d] &
\mathcal{E}^{b - 1} \ar[r] \ar[d] &
0 \ar[r] \ar[d] &
\ldots \\
\ldots \ar[r] &
\mathcal{F}^{b - 2} \ar[r] &
\mathcal{F}^{b - 1} \oplus \mathcal{E}^b \ar[r] &
\mathcal{F}^b \ar[r] &
\ldots
}
$$
を与える。複体の写像
$(\mathcal{E}^a \to \ldots \to \mathcal{E}^b) \to \mathcal{F}^\bullet$
が得られることは明らかである。この写像が擬同型であることの確認は省略する。
\end{proof}

\begin{lemma}
\label{perfect-lemma-resolution-property-perfect-complex}
$X$ を分解性質をもつ準コンパクトかつ準分離なスキームとする。このとき
$D(\mathcal{O}_X)$ の任意の完全対象は、有限局所自由
$\mathcal{O}_X$-加群の有界複体で表せる。
\end{lemma}

\begin{proof}
$E$ を $D(\mathcal{O}_X)$ の完全対象とする。補題
\ref{perfect-lemma-resolution-property-affine-diagonal} により、$X$ の対角射は
アフィンである。したがって Proposition
\ref{perfect-proposition-quasi-compact-affine-diagonal} により、$E$ は準連接
$\mathcal{O}_X$-加群の複体 $\mathcal{F}^\bullet$ で表せる。
$X$ は準コンパクトなので $E$ は $D^b(\mathcal{O}_X)$ に属することに注意する。
それゆえ $n \ll 0$ なら $\tau_{\geq n}\mathcal{F}^\bullet$ は $E$ を表す
準連接 $\mathcal{O}_X$-加群の下に有界な複体である。よって補題
\ref{perfect-lemma-construct-strictly-perfect} を適用して結論を得る。
\end{proof}

\begin{lemma}
\label{perfect-lemma-resolution-property-map-perfect-complex}
$X$ を分解性質をもつ準コンパクトかつ準分離なスキームとする。
$\mathcal{E}^\bullet$ と $\mathcal{F}^\bullet$ を有限局所自由
$\mathcal{O}_X$-加群の有限複体とする。このとき任意の
$\alpha \in \Hom_{D(\mathcal{O}_X)}(\mathcal{E}^\bullet, \mathcal{F}^\bullet)$
は図式
$$
\mathcal{E}^\bullet \leftarrow \mathcal{G}^\bullet \to \mathcal{F}^\bullet
$$
で表せる。ここで $\mathcal{G}^\bullet$ は有限局所自由
$\mathcal{O}_X$-加群の有界複体であり、
$\mathcal{G}^\bullet \to \mathcal{E}^\bullet$ は擬同型である。
\end{lemma}

\begin{proof}
補題 \ref{perfect-lemma-resolution-property-affine-diagonal} により、$X$ の対角射は
アフィンである。したがって Proposition
\ref{perfect-proposition-quasi-compact-affine-diagonal} により、$\alpha$ は図式
$$
\mathcal{E}^\bullet \leftarrow \mathcal{H}^\bullet \to \mathcal{F}^\bullet
$$
で表せる。ここで $\mathcal{H}^\bullet$ は準連接
$\mathcal{O}_X$-加群の複体であり、
$\mathcal{H}^\bullet \to \mathcal{E}^\bullet$
は擬同型である。$n \ll 0$ に対して写像
$\mathcal{H}^\bullet \to \mathcal{E}^\bullet$ および
$\mathcal{H}^\bullet \to \mathcal{F}^\bullet$
は擬同型
$\mathcal{H}^\bullet \to \tau_{\geq n}\mathcal{H}^\bullet$
を経由する。これは単に $\mathcal{E}^\bullet$ と $\mathcal{F}^\bullet$ が
有界複体だからである。したがって $\mathcal{H}^\bullet$ を
$\tau_{\geq n}\mathcal{H}^\bullet$ で置き換え、$\mathcal{H}^\bullet$ が
下に有界であると仮定してよい。そこで補題
\ref{perfect-lemma-construct-strictly-perfect} を適用して結論を得る。
\end{proof}

\begin{lemma}
\label{perfect-lemma-resolution-property-homotopy-map-perfect-complex}
$X$ を分解性質をもつ準コンパクトかつ準分離なスキームとする。
$\mathcal{E}^\bullet$ と $\mathcal{F}^\bullet$ を有限局所自由
$\mathcal{O}_X$-加群の有限複体とする。
$\alpha^\bullet, \beta^\bullet :\mathcal{E}^\bullet \to \mathcal{F}^\bullet$
を $D(\mathcal{O}_X)$ で同じ写像を定める二つの複体の写像とする。このとき擬同型
$\gamma^\bullet : \mathcal{G}^\bullet \to \mathcal{E}^\bullet$
が存在する。ここで $\mathcal{G}^\bullet$ は有限局所自由
$\mathcal{O}_X$-加群の有界複体であり、
$\alpha^\bullet \circ \gamma^\bullet$ と
$\beta^\bullet \circ \gamma^\bullet$ は複体の写像としてホモトピックである。
\end{lemma}

\begin{proof}
補題 \ref{perfect-lemma-resolution-property-affine-diagonal} により、$X$ の対角射は
アフィンである。したがって Proposition
\ref{perfect-proposition-quasi-compact-affine-diagonal}（および導来圏の定義）により擬同型
$\gamma^\bullet : \mathcal{G}^\bullet \to \mathcal{E}^\bullet$
で、$\mathcal{G}^\bullet$ が準連接 $\mathcal{O}_X$-加群の複体であり、
$\alpha^\bullet \circ \gamma^\bullet$ と
$\beta^\bullet \circ \gamma^\bullet$ が複体の写像としてホモトピックとなるものが
存在する。この事実を与えるホモトピー
$h^i : \mathcal{G}^i \to \mathcal{F}^{i - 1}$ を選ぶ。$n \ll 0$ を選ぶ。
このとき写像 $\gamma^\bullet$ は、$\mathcal{E}^\bullet$ が下に有界なので商写像
$\mathcal{G}^\bullet \to \tau_{\geq n}\mathcal{G}^\bullet$ を標準的に経由する。
まったく同じ理由により、写像 $h^i$ は全射
$\mathcal{G}^i \to (\tau_{\geq n}\mathcal{G})^i$.
を経由する。したがって $\mathcal{G}^\bullet$ を
$\tau_{\geq n}\mathcal{G}^\bullet$ で置き換えてよい。そこで補題
\ref{perfect-lemma-construct-strictly-perfect} を適用して結論を得る。
\end{proof}

\begin{proposition}
\label{perfect-proposition-perfect-resolution-property}
$X$ を分解性質をもつ準コンパクトかつ準分離なスキームとする。次のように記す。
\begin{enumerate}
\item $\mathcal{A}$ は有限局所自由 $\mathcal{O}_X$-加群の加法圏、
\item $K^b(\mathcal{A})$ は $\mathcal{A}$ の有界複体のホモトピー圏
（Derived Categories, Section \ref{derived-section-homotopy} を参照）、
\item $D_{perf}(\mathcal{O}_X)$ は完全対象からなる $D(\mathcal{O}_X)$ の
厳密充満・飽和・三角部分圏。
\end{enumerate}
この記法のもとで、明らかな関手
$$
K^b(\mathcal{A}) \longrightarrow D_{perf}(\mathcal{O}_X)
$$
は三角圏の完全関手であり、三角圏の同値
$S^{-1}K^b(\mathcal{A}) \to D_{perf}(\mathcal{O}_X)$ を経由する。ここで
$S$ は $K^b(\mathcal{A})$ における擬同型からなる飽和乗法系である。
\end{proposition}

\begin{proof}
命題の主張を読解できるなら、この最初の段落は読み飛ばしてよい。用いた定義の
いくつかについては、Derived Categories, Definition
\ref{derived-definition-triangulated-subcategory}
（三角部分圏）、
Derived Categories, Definition \ref{derived-definition-saturated}
（飽和三角部分圏）、
Derived Categories, Definition \ref{derived-definition-localization}
（三角構造と両立する乗法系）、および Categories,
Definition \ref{categories-definition-saturated-multiplicative-system}
（飽和乗法系）を参照されたい。Cohomology, Lemmas
\ref{cohomology-lemma-two-out-of-three-perfect} および
\ref{cohomology-lemma-summands-perfect} により
$D_{perf}(\mathcal{O}_X)$ は $D(\mathcal{O}_X)$ の飽和三角部分圏である。
また $K^b(\mathcal{A})$ は三角圏であることにも注意する。Derived Categories,
Lemma \ref{derived-lemma-bounded-triangulated-subcategories} を参照されたい。

\medskip\noindent
この関手が区別三角を区別三角へ送る、すなわち完全であることは明らかである。
すると Derived Categories, Lemma
\ref{derived-lemma-triangle-functor-localize} により、$S$ は
$K^b(\mathcal{A})$ の三角構造と両立する飽和乗法系である。したがって Derived
Categories, Proposition \ref{derived-proposition-construct-localization} により、
局所化 $S^{-1}K^b(\mathcal{A})$ が存在し、三角圏となる。Derived Categories,
Lemma \ref{derived-lemma-universal-property-localization} により完全な分解
$S^{-1}K^b(\mathcal{A}) \to D_{perf}(\mathcal{O}_X)$ を得る。補題
\ref{perfect-lemma-resolution-property-perfect-complex},
\ref{perfect-lemma-resolution-property-map-perfect-complex}, および
\ref{perfect-lemma-resolution-property-homotopy-map-perfect-complex}
により、この関手は同値である。最後に関手
$S^{-1}K^b(\mathcal{A}) \to D_{perf}(\mathcal{O}_X)$
は Derived Categories, Lemma \ref{derived-lemma-exact-equivalence} により、
三角圏の同値（区別三角が対応するという意味で）である。
\end{proof}






\section{K 群}
\label{perfect-section-K-groups}

\noindent
スキームに付随するさまざまな圏の $K_0$ について少しだけ述べる。先行する内容は
Algebra, Section \ref{algebra-section-K-groups}、
Homology, Section \ref{homology-section-K-groups}、
Derived Categories, Section \ref{derived-section-K-groups}、および
More on Algebra, Lemma \ref{more-algebra-lemma-perfect-to-K-group-universal}
にある。

\medskip\noindent
Algebra, Section \ref{algebra-section-K-groups} と同様に、任意のネーター・
スキーム $X$ に対して二つの $K$ 群 $K'_0(X)$ と $K_0(X)$ を定義する。
前者ではコヒーレント $\mathcal{O}_X$-加群を、後者では有限局所自由
$\mathcal{O}_X$-加群を用いる。

\begin{lemma}
\label{perfect-lemma-Noetherian-Kprime}
$X$ をネーター・スキームとする。このとき
$$
K_0(\textit{Coh}(\mathcal{O}_X)) =
K_0(D^b(\textit{Coh}(\mathcal{O}_X)) =
K_0(D^b_{\textit{Coh}}(\mathcal{O}_X))
$$
\end{lemma}

\begin{proof}
最初の等式は Derived Categories, Lemma
\ref{derived-lemma-K-bounded-derived} である。Derived Categories, Lemma
\ref{derived-lemma-DBA-map-K} により
$K_0(\textit{Coh}(\mathcal{O}_X)) =
K_0(D^b_{\textit{Coh}}(\mathcal{O}_X))$ である。これで補題が示された。
（Proposition \ref{perfect-proposition-DCoh} による等式
$D^b(\textit{Coh}(\mathcal{O}_X)) = D^b_{\textit{Coh}}(\mathcal{O}_X)$
を用いて第二の等式を示すこともできる。）
\end{proof}

\noindent
定義は次のとおりである。

\begin{definition}
\label{perfect-definition-K-group}
$X$ をスキームとする。
\begin{enumerate}
\item $K_0(X)$ を {\it $X$ のグロタンディーク群} と呼ぶ。これは完全対象からなる
$D(\mathcal{O}_X)$ の厳密充満・飽和・三角部分圏
$D_{perf}(\mathcal{O}_X)$ の第0 K 群である。式で書けば
$$
K_0(X) = K_0(D_{perf}(\mathcal{O}_X))
$$
\item $X$ が局所ネーターなら、$K'_0(X)$ を
{\it $X$ 上のコヒーレント層のグロタンディーク群} と呼ぶ。これはコヒーレント
$\mathcal{O}_X$-加群のアーベル圏の第0 $K$ 群である。式で書けば
$$
K'_0(X) = K_0(\textit{Coh}(\mathcal{O}_X))
$$
\end{enumerate}
\end{definition}

\noindent
$X$ が分解性質をもつ場合（例えば $X$ が豊富な可逆加群をもつ場合）、この
$K_0(X)$ の定義が有限局所自由加群による通常の定義と一致することを示す。
補題 \ref{perfect-lemma-K-is-old-K} を参照されたい。

\begin{lemma}
\label{perfect-lemma-K-agrees-affine}
$X = \Spec(R)$ をアフィン・スキームとする。このとき $K_0(X) = K_0(R)$ であり、
$R$ がネーターなら $K'_0(X) = K'_0(R)$ である。
\end{lemma}

\begin{proof}
$K'_0(R)$ と $K_0(R)$ は Algebra, Section
\ref{algebra-section-K-groups} で定義されたことを思い出そう。

\medskip\noindent
More on Algebra, Lemma \ref{more-algebra-lemma-perfect-to-K-group-universal}
により $K_0(R) = K_0(D_{perf}(R))$ である。補題
\ref{perfect-lemma-perfect-affine} および \ref{perfect-lemma-affine-compare-bounded} により
$D_{perf}(R) = D_{perf}(\mathcal{O}_X)$ である。これで等式
$K_0(R) = K_0(X)$ が示された。

\medskip\noindent
Cohomology of Schemes, Lemma
\ref{coherent-lemma-coherent-Noetherian} により
$\textit{Coh}(\mathcal{O}_X)$ は有限 $R$-加群の圏と同値なので、等式
$K'_0(R) = K'_0(X)$ が成り立つ。また定義から、$K'_0(R)$ が有限
$R$-加群の圏の第0 K 群であることは明らかである。
\end{proof}

\noindent
$X$ をネーター・スキームとする。このとき $K'_0(X)$ と $K_0(X)$ はともに
定義され、この場合には標準写像
$$
K_0(X) = K_0(D_{perf}(\mathcal{O}_X))
\longrightarrow
K_0(D^b_{\textit{Coh}}(\mathcal{O}_X)) = K'_0(X)
$$
がある。実際、補題 \ref{perfect-lemma-identify-pseudo-coherent-noetherian} により
完全複体は $D^b_{\textit{Coh}}(\mathcal{O}_X)$ に属し、包含関手
$D_{perf}(\mathcal{O}_X) \to D^b_{\textit{Coh}}(\mathcal{O}_X)$
は第0 $K$ 群上の写像を誘導し（Derived Categories, Lemma
\ref{derived-lemma-map-K}）、右辺の等式は補題
\ref{perfect-lemma-Noetherian-Kprime} による。

\begin{lemma}
\label{perfect-lemma-Kprime-K}
$X$ をネーター正則スキームとする。このとき写像
$K_0(X) \to K'_0(X)$ は同型である。
\end{lemma}

\begin{proof}
補題 \ref{perfect-lemma-perfect-on-regular} と上での写像
$K_0(X) \to K'_0(X)$ の構成から直ちに従う。
\end{proof}

\noindent
$X$ をスキームとする。有限局所自由 $\mathcal{O}_X$-加群の圏を
$\textit{Vect}(X)$ と記す。一般に $\textit{Vect}(X)$ はアーベル圏ではないが、
$\textit{Vect}(X)$ の短完全列が何を意味するかは明らかである。次の性質をもつ
一意なアーベル群を $K_0(\textit{Vect}(X))$ と記す\footnote{ここで正しい
一般性は任意の完全圏に対して $K_0$ を定義することである。Injectives,
Remark \ref{injectives-remark-embed-exact-category} を参照されたい。}。
\begin{enumerate}
\item 各有限局所自由 $\mathcal{O}_X$-加群 $\mathcal{E}$ に対して、
$K_0(\textit{Vect}(X))$ の元 $[\mathcal{E}]$ が与えられる。
\item 有限局所自由 $\mathcal{O}_X$-加群の各短完全列
$0 \to \mathcal{E}' \to \mathcal{E} \to \mathcal{E}'' \to 0$
に対して、$K_0(\textit{Vect}(X))$ で関係式
$[\mathcal{E}] = [\mathcal{E}'] + [\mathcal{E}'']$ が成り立つ。
\item 群 $K_0(\textit{Vect}(X))$ は元 $[\mathcal{E}]$ で生成される。
\item 生成元 $[\mathcal{E}]$ の間の $K_0(\textit{Vect}(X))$ におけるすべての
関係式は、上のような完全列から生じる関係式の $\mathbf{Z}$-線型結合である。
\end{enumerate}
$K_0(\textit{Vect}(X))$ の詳しい構成は省略する。自然な写像
$$
K_0(\textit{Vect}(X)) \longrightarrow K_0(X)
$$
がある。実際、有限局所自由 $\mathcal{O}_X$-加群 $\mathcal{E}$ に対し、次数
$0$ に $\mathcal{E}$ をもち、他の次数では零である $X$ 上の完全複体を
$\mathcal{E}[0]$ と記す。有限局所自由 $\mathcal{O}_X$-加群の短完全列
$0 \to \mathcal{E} \to \mathcal{E}' \to \mathcal{E}'' \to 0$ から区別三角
$\mathcal{E}[0] \to \mathcal{E}'[0] \to \mathcal{E}''[0] \to \mathcal{E}[1]$,
を得る。Derived Categories, Section
\ref{derived-section-canonical-delta-functor} を参照されたい。したがって写像
$K_0(\textit{Vect}(X)) \to K_0(D_{perf}(\mathcal{O}_X)) = K_0(X)$
を $[\mathcal{E}]$ を $[\mathcal{E}[0]]$ に送ることで得る。このひどい記法は
ご容赦願いたい。

\begin{lemma}
\label{perfect-lemma-K-is-old-K}
$X$ を分解性質をもつ準コンパクトかつ準分離なスキームとする。このとき写像
$K_0(\textit{Vect}(X)) \to K_0(X)$ は同型である。
\end{lemma}

\begin{proof}
この補題は、補題 \ref{perfect-lemma-resolution-property-perfect-complex},
\ref{perfect-lemma-resolution-property-map-perfect-complex}, および
\ref{perfect-lemma-resolution-property-homotopy-map-perfect-complex}
から直接的に従う。以下ではこれらの結果を断りなく用いる。逆写像
$$
c : K_0(X) = K_0(D_{perf}(\mathcal{O}_X)) \longrightarrow
K_0(\textit{Vect}(X))
$$
を構成しよう。$D_{perf}(\mathcal{O}_X)$ の任意の対象は有限局所自由
$\mathcal{O}_X$-加群の有界複体 $\mathcal{E}^\bullet$ で表せる。そこで
$$
c([\mathcal{E}^\bullet]) = \sum (-1)^i[\mathcal{E}^i]
$$
と置く。もちろん、これが良定義であることを示さなければならない。ひとまず $c$ を
有限局所自由 $\mathcal{O}_X$-加群の有界複体上で定義された写像とみなす。

\medskip\noindent
$\mathcal{E}^\bullet \to \mathcal{F}^\bullet$ を有限局所自由
$\mathcal{O}_X$-加群の有界複体の全射とする。その核を
$\mathcal{K}^\bullet$ とする。このとき $\mathcal{O}_X$-加群の短完全列
$$
0 \to \mathcal{K}^n \to \mathcal{E}^n \to \mathcal{F}^n \to 0
$$
を得る。これは $\mathcal{F}^n$ が有限局所自由なので局所的に分裂する。
したがって $\mathcal{K}^\bullet$ も有限局所自由 $\mathcal{O}_X$-加群の
有界複体であり、
$c(\mathcal{E}^\bullet) = c(\mathcal{K}^\bullet) + c(\mathcal{F}^\bullet)$
が $K_0(\textit{Vect}(X))$ で成り立つ。

\medskip\noindent
有限局所自由 $\mathcal{O}_X$-加群の有界かつ非輪状な複体
$\mathcal{E}^\bullet$ が与えられたとする。$n \not \in [a, b]$ に対して
$\mathcal{E}^n = 0$ とする。このとき $\mathcal{E}^\bullet$ を短完全列
$$
\begin{matrix}
0 \to \mathcal{E}^a \to \mathcal{E}^{a + 1} \to \mathcal{F}^{a + 1} \to 0,\\
0 \to \mathcal{F}^{a + 1} \to \mathcal{E}^{a + 2} \to
\mathcal{F}^{a + 3} \to 0, \\
\ldots \\
0 \to \mathcal{F}^{b - 3} \to \mathcal{E}^{b - 2} \to
\mathcal{F}^{b - 2} \to 0, \\
0 \to \mathcal{F}^{b - 2} \to \mathcal{E}^{b - 1} \to \mathcal{E}^b \to 0
\end{matrix}
$$
に分解できる。降下帰納法により
$\mathcal{F}^{b - 2}, \ldots, \mathcal{F}^{a + 1}$
は有限局所自由 $\mathcal{O}_X$-加群であり、
$$
c(\mathcal{E}^\bullet) = \sum (-1)[\mathcal{E}^n] =
\sum (-1)^n([\mathcal{F}^{n - 1}] + [\mathcal{F}^n]) = 0
$$
となることが分かる。したがってこの構成は非輪状複体上で零を与える。

\medskip\noindent
前二段落の結果から $c$ が良定義であることが従う。実際、有限局所自由
$\mathcal{O}_X$-加群の有界複体 $\mathcal{E}^\bullet$ と
$\mathcal{F}^\bullet$ が $D(\mathcal{O}_X)$ の同じ対象を表すとする。このとき擬同型
$a : \mathcal{G}^\bullet \to \mathcal{E}^\bullet$ および
$b : \mathcal{G}^\bullet \to \mathcal{F}^\bullet$
で、$\mathcal{G}^\bullet$ が有限局所自由 $\mathcal{O}_X$-加群の有界複体である
ものを取れる。複体の短完全列
$$
0 \to \mathcal{E}^\bullet \to C(a)^\bullet \to \mathcal{G}^\bullet[1] \to 0
$$
を得る。Derived Categories, Definition \ref{derived-definition-cone}
を参照されたい。$a$ は擬同型なので、錐 $C(a)^\bullet$ は非輪状である
（例えば Derived Categories, Section
\ref{derived-section-canonical-delta-functor} の議論から従う）。したがって
$$
0 = c(C(f)^\bullet) = c(\mathcal{E}^\bullet) + c(\mathcal{G}^\bullet[1]) =
c(\mathcal{E}^\bullet) - c(\mathcal{G}^\bullet)
$$
となり、望む結果を得る。$b$ を用いた同じ議論により
$0 = c(\mathcal{F}^\bullet) - c(\mathcal{G}^\bullet)$ である。したがって
$c(\mathcal{E}^\bullet) = c(\mathcal{F}^\bullet)$ であり、$c$ は良定義である。

\medskip\noindent
有限局所自由 $\mathcal{O}_X$-加群の有界複体の写像
$\mathcal{E}^\bullet \to \mathcal{F}^\bullet$ の錐を用いる同様の議論により、
$X \to Y \to Z$ が $D_{perf}(\mathcal{O}_X)$ の区別三角なら
$c(Y) = c(X) + c(Z)$ である。細部は省略する。こうして望みのアーベル群の
準同型 $c : K_0(X) \to K_0(\textit{Vect}(X))$ を得る。

\medskip\noindent
合成
$K_0(\textit{Vect}(X)) \to K_0(X) \to K_0(\textit{Vect}(X))$
が恒等写像であることは明らかである。一方、$\mathcal{E}^\bullet$ を有限局所自由
$\mathcal{O}_X$-加群の有界複体とする。このとき「愚直な切断」の区別三角
（Homology, Section \ref{homology-section-truncations} を参照）
$$
\sigma_{\geq n}\mathcal{E}^\bullet \to
\sigma_{\geq n - 1}\mathcal{E}^\bullet \to
\mathcal{E}^{n - 1}[-n + 1] \to
(\sigma_{\geq n}\mathcal{E}^\bullet)[1]
$$
の存在と帰納法により
$$
[\mathcal{E}^\bullet] = \sum (-1)^i[\mathcal{E}^i[0]]
$$
が $K_0(X) = K_0(D_{perf}(\mathcal{O}_X))$ で成り立つ。この記法はご容赦願いたい。
したがって写像
$K_0(\textit{Vect}(X)) \to K_0(D_{perf}(\mathcal{O}_X)) = K_0(X)$
は全射であり、証明が完了する。
\end{proof}

\begin{remark}
\label{perfect-remark-K-ring}
$X$ をスキームとする。K 群 $K_0(X)$ は標準的に可換環である。実際、導来
テンソル積
$$
\otimes = \otimes^\mathbf{L}_{\mathcal{O}_X} :
D_{perf}(\mathcal{O}_X) \times D_{perf}(\mathcal{O}_X)
\longrightarrow
D_{perf}(\mathcal{O}_X)
$$
と Derived Categories, Lemma \ref{derived-lemma-bilinear-map-K} により
双線型な積を得る。$K \otimes L \cong L \otimes K$ なのでこの積は可換である。
また
$(K \otimes L) \otimes M = K \otimes (L \otimes M)$
なのでこの積は結合的である。最後に、$K_0(X)$ の単位元は
$1 = [\mathcal{O}_X]$ である。

\medskip\noindent
$\textit{Vect}(X)$ と $K_0(\textit{Vect}(X))$ を上のとおりとすると、
$K_0(\textit{Vect}(X))$ も明らかに環構造をもつ。すなわち、
$\mathcal{E}$ と $\mathcal{F}$ が有限局所自由 $\mathcal{O}_X$-加群なら
$$
[\mathcal{E}] \cdot [\mathcal{F}] =
[\mathcal{E} \otimes_{\mathcal{O}_X} \mathcal{F}]
$$
と置く。これが実際に双線型・可換・結合的な積を定めることは容易に確認できる。
細部は省略する。上で構成した写像
$$
K_0(\textit{Vect}(X)) \longrightarrow K_0(X)
$$
は、これらの定義のもとで環準同型である。

\medskip\noindent
次に $X$ がネーターであると仮定する。導来テンソル積は写像
$$
\otimes = \otimes^\mathbf{L}_{\mathcal{O}_X} :
D_{perf}(\mathcal{O}_X) \times D^b_{\textit{Coh}}(\mathcal{O}_X)
\longrightarrow
D^b_{\textit{Coh}}(\mathcal{O}_X)
$$
も与える。再び Derived Categories, Lemma
\ref{derived-lemma-bilinear-map-K} を用いると、補題
\ref{perfect-lemma-Noetherian-Kprime} により
$K'_0(X) = K_0(D^b_{\textit{Coh}}(\mathcal{O}_X))$ なので、双線型な積
$K_0(X) \times K'_0(X) \to K'_0(X)$ を得る。これが $K'_0(X)$ に環
$K_0(X)$ 上の加群構造を与えることは容易に示せる。
\end{remark}

\begin{remark}
\label{perfect-remark-pushforward-K}
$f : X \to Y$ を局所ネーター・スキームの固有射とする。写像
$$
f_* : K'_0(X) \longrightarrow K'_0(Y)
$$
があり、$[\mathcal{F}]$ を
$$
[\bigoplus\nolimits_{i \geq 0} R^{2i}f_*\mathcal{F}] -
[\bigoplus\nolimits_{i \geq 0} R^{2i + 1}f_*\mathcal{F}]
$$
に送る。これは次の理由で良定義である。層 $R^if_*\mathcal{F}$ はコヒーレントで
あり（Cohomology of Schemes, Proposition
\ref{coherent-proposition-proper-pushforward-coherent}）、局所的には有限個しか
非零でなく、さらに $X$ 上のコヒーレント層の短完全列は $Y$ 上の
$R^if_*$ の長完全列を与える。$Y$ が準コンパクトなら（実際に最もよく使うのは
この場合だけである）、上式は
$$
f_*[\mathcal{F}] = \sum (-1)^i[R^if_*\mathcal{F}] = [Rf_*\mathcal{F}]
$$
と書き直せる。ここでは補題 \ref{perfect-lemma-Noetherian-Kprime} の等式
$K'_0(Y) = K_0(D^b_{\textit{Coh}}(Y))$ を用いた。
\end{remark}

\begin{lemma}
\label{perfect-lemma-projection-formula}
$f : X \to Y$ を局所ネーター・スキームの固有射とする。このとき
$\alpha \in K'_0(X)$ および $\beta \in K_0(Y)$ に対して
$f_*(\alpha \cdot f^*\beta) = f_*\alpha \cdot \beta$ が成り立つ。
\end{lemma}

\begin{proof}
補題 \ref{perfect-lemma-cohomology-base-change}、Remark
\ref{perfect-remark-pushforward-K} の議論、および Remark \ref{perfect-remark-K-ring}
における積 $K'_0(X) \times K_0(X) \to K'_0(X)$ の定義から従う。
\end{proof}

\begin{remark}
\label{perfect-remark-perf-Z}
$X$ をスキームとし、$Z \subset X$ を閉部分スキームとする。厳密充満・飽和・
三角部分圏
$$
D_{Z, perf}(\mathcal{O}_X) \subset D(\mathcal{O}_X)
$$
を考える。これは、コホモロジー層が集合論的に $Z$ 上に台をもつ
$\mathcal{O}_X$-加群の完全複体からなる。この三角圏の第0 $K$ 群
$K_0(D_{Z, perf}(\mathcal{O}_X))$ は $K_Z(X)$ または $K_{0, Z}(X)$ と
記されることもある。Remark \ref{perfect-remark-K-ring} とまったく同様に導来
テンソル積を用いると、$K_0(D_{Z, perf}(\mathcal{O}_X))$ は結合的かつ可換な
積をもつことが分かるが、一般には $K_0(D_{Z, perf}(\mathcal{O}_X))$ は
単位元をもたない。
\end{remark}













\section{複体の行列式}
\label{perfect-section-det-two-terms}

\noindent
本節は More on Algebra, Section
\ref{more-algebra-section-determinants-complexes} の続きである。任意の環付き空間
$(X, \mathcal{O}_X)$ に対して関手
$$
\det :
\left\{
\begin{matrix}
\text{完全複体の圏} \\
\text{射は同型}
\end{matrix}
\right\}
\longrightarrow
\left\{
\begin{matrix}
\text{可逆加群の圏} \\
\text{射は同型}
\end{matrix}
\right\}
$$
さらに、フィルター付き導来圏 $DF(\mathcal{O}_X)$ の対象 $(L, F)$ で、
フィルトレーションが有限かつ次数付き部分が完全複体であるものに対して、標準同型
$\det(\text{gr}L) \to \det(L)$ がある。原論文の解説については
\cite{determinant} を参照されたい。この内容は後で追加する（将来の参照を挿入）。

\medskip\noindent
ひとまず、$X$ がスキームであり、$D(\mathcal{O}_X)$ の Tor 振幅が
$[-1, 0]$ にある完全対象 $L$ を考える場合の個別的な構成を与える。

\begin{lemma}
\label{perfect-lemma-determinant-two-term-complexes}
$X$ をスキームとする。関手
$$
\det :
\left\{
\begin{matrix}
\text{完全複体の圏} \\
\text{Tor 振幅が }[-1, 0]\text{ にある} \\
\text{射は同型}
\end{matrix}
\right\}
\longrightarrow
\left\{
\begin{matrix}
\text{可逆加群の圏} \\
\text{射は同型}
\end{matrix}
\right\}
$$
がある。さらに、$D(\mathcal{O}_X)$ の階数 $0$、Tor 振幅 $[-1, 0]$ の完全対象
$L$ に対して標準元 $\delta(L) \in \Gamma(X, \det(L))$ があり、
$D(\mathcal{O}_X)$ の任意の同型 $a : L \to K$ に対して
$\det(a)(\delta(L)) = \delta(K)$ が成り立つ。またこの構成はアフィン局所的には
More on Algebra, Section
\ref{more-algebra-section-determinants-complexes} の構成で与えられる。
\end{lemma}

\begin{proof}
$L$ を左辺の対象とする。$\Spec(A) = U \subset X$ がアフィン開なら、$L|_U$ は
Tor 振幅が $[-1, 0]$ にある $A$-加群の完全複体 $L^\bullet$ に対応する。補題
\ref{perfect-lemma-affine-compare-bounded},
\ref{perfect-lemma-tor-dimension-affine}, および
\ref{perfect-lemma-perfect-affine} を参照されたい。そこで More on Algebra, Lemma
\ref{more-algebra-lemma-determinant-two-term-complexes} で構成した可逆
$A$-加群 $\det(L^\bullet)$ を考えられる。$\Spec(B) = V \subset U$ が $U$ に
含まれる別のアフィン開なら、
$\det(L^\bullet) \otimes_A B = \det(L^\bullet \otimes_A B)$ であり、したがって
この構成は制限写像と両立する（補題 \ref{perfect-lemma-quasi-coherence-pullback} を
参照し、$A \to B$ が平坦であることに注意する）。よってこれらの可逆加群を
貼り合わせ、$X$ 上の可逆加群 $\det(L)$ を得られる。関手性と標準切断も
まったく同様に構成する。細部は省略する。
\end{proof}

\begin{remark}
\label{perfect-remark-functorial-det}
補題 \ref{perfect-lemma-determinant-two-term-complexes} の構成は引き戻しと両立する。
より正確には、スキームの射 $f : X \to Y$ と、$D(\mathcal{O}_Y)$ の
Tor 振幅が $[-1, 0]$ にある完全対象 $K$ が与えられると、$Lf^*K$ は
$D(\mathcal{O}_X)$ の Tor 振幅が $[-1, 0]$ にある完全対象 $K$ であり、標準同一視
$$
f^*\det(K) \longrightarrow \det(Lf^*K)
$$
がある。さらに $K$ の階数が $0$ なら、この写像により $\delta(K)$ は
$\delta(Lf^*K)$ へ引き戻される。これは行列式のアフィン局所的構成から明らかである。
\end{remark}






\section{有界性の検出}
\label{perfect-section-detecting-boundedness}

\noindent
本節では、Section \ref{perfect-section-generating} で構成した準コンパクトかつ準分離な
スキームの $D_\QCoh$ のコンパクト生成対象が特別な性質をもつことを示す。
その節は本節と非常によく似ているので、先に読むことを勧める。

\begin{lemma}
\label{perfect-lemma-orthogonal-koszul-first-variant}
Situation \ref{perfect-situation-complex} において開埋め込みを $j : U \to X$ と記し、
$K$ を $A$ 上の $f_1, \ldots, f_r$ に関する Koszul 複体に対応する
$D(\mathcal{O}_X)$ の完全対象とする。$E \in D_\QCoh(\mathcal{O}_X)$ および
$a \in \mathbf{Z}$ とする。次の条件を考える。
\begin{enumerate}
\item 標準写像 $\tau_{\geq a}E \to \tau_{\geq a} Rj_*(E|_U)$ は同型である。
\item すべての $n \geq a$ に対して
$\Hom_{D(\mathcal{O}_X)}(K[-n], E) = 0$ である。
\end{enumerate}
このとき (2) は (1) を含意し、(1) は $a$ を $a + 1$ で置き換えた (2) を
含意する。
\end{lemma}

\begin{proof}
区別三角 $N \to E \to Rj_*(E|_U) \to N[1]$ を選ぶ。このとき (1) は
$\tau_{\geq a + 1} N = 0$ を含意し、$\tau_{\geq a}N = 0$ は (1) を含意する。
ここで
$$
\Hom_{D(\mathcal{O}_X)}(K[-n], Rj_*(E|_U)) =
\Hom_{D(\mathcal{O}_U)}(K|_U[-n], E) = 0
$$
は $K|_U = 0$ なので、すべての $n$ に対して成り立つ。したがって (2) は、
すべての $n \geq a$ に対して
$\Hom_{D(\mathcal{O}_X)}(K[-n], N) = 0$ であることと同値である。また Koszul
複体の区別三角
$$
K^\bullet(f_1^{e_1}, \ldots, f_i^{e'_i}, \ldots, f_r^{e_r}) \to
K^\bullet(f_1^{e_1}, \ldots, f_i^{e'_i + e''_i}, \ldots, f_r^{e_r}) \to
K^\bullet(f_1^{e_1}, \ldots, f_i^{e''_i}, \ldots, f_r^{e_r}) \to \ldots
$$
がある。More on Algebra, Lemma \ref{more-algebra-lemma-koszul-mult} を
参照されたい。ゆえに、すべての $n \geq a$ に対する
$\Hom_{D(\mathcal{O}_X)}(K[-n], N) = 0$ は、補題
\ref{perfect-lemma-represent-cohomology-class-on-closed} の $K_e$ を用いて、すべての
$n \geq a$ およびすべての $e \geq 1$ に対する
$\Hom_{D(\mathcal{O}_X)}(K_e[-n], N) = 0$ と同値である。$N|_U = 0$ なので、
同補題によりこれはさらに $n \geq a$ に対する $H^n(X, N) = 0$ と同値である。
$N$ は $A$-加群の複体 $R\Gamma(X, N)$ で決まるので、補題
\ref{perfect-lemma-affine-compare-bounded} により (2) は
$\tau_{\geq a}N = 0$ と同値である。これで補題が示された。
\end{proof}

\begin{lemma}
\label{perfect-lemma-orthogonal-koszul-second-variant}
Situation \ref{perfect-situation-complex} において開埋め込みを $j : U \to X$ と記し、
$K$ を $A$ 上の $f_1, \ldots, f_r$ に関する Koszul 複体に対応する
$D(\mathcal{O}_X)$ の完全対象とする。$E \in D_\QCoh(\mathcal{O}_X)$ および
$a \in \mathbf{Z}$ とする。次の条件を考える。
\begin{enumerate}
\item 標準写像 $\tau_{\leq a}E \to \tau_{\leq a} Rj_*(E|_U)$ は同型である。
\item すべての $n \leq a$ に対して
$\Hom_{D(\mathcal{O}_X)}(K[-n], E) = 0$ である。
\end{enumerate}
このとき (2) は (1) を含意し、(1) は $a$ を $a - 1$ で置き換えた (2) を
含意する。
\end{lemma}

\begin{proof}
区別三角 $E \to Rj_*(E|_U) \to N \to E[1]$ を選ぶ。このとき (1) は
$\tau_{\leq a - 1}N = 0$ を含意し、$\tau_{\leq a}N = 0$ は (1) を含意する。
ここで
$$
\Hom_{D(\mathcal{O}_X)}(K[-n], Rj_*(E|_U)) =
\Hom_{D(\mathcal{O}_U)}(K|_U[-n], E) = 0
$$
は $K|_U = 0$ なので、すべての $n$ に対して成り立つ。したがって (2) は、
すべての $n \leq a$ に対して
$\Hom_{D(\mathcal{O}_X)}(K[-n], N) = 0$ であることと同値である。また Koszul
複体の区別三角
$$
K^\bullet(f_1^{e_1}, \ldots, f_i^{e'_i}, \ldots, f_r^{e_r}) \to
K^\bullet(f_1^{e_1}, \ldots, f_i^{e'_i + e''_i}, \ldots, f_r^{e_r}) \to
K^\bullet(f_1^{e_1}, \ldots, f_i^{e''_i}, \ldots, f_r^{e_r}) \to \ldots
$$
がある。More on Algebra, Lemma \ref{more-algebra-lemma-koszul-mult} を
参照されたい。ゆえに、すべての $n \leq a$ に対する
$\Hom_{D(\mathcal{O}_X)}(K[-n], N) = 0$ は、補題
\ref{perfect-lemma-represent-cohomology-class-on-closed} の $K_e$ を用いて、すべての
$n \leq a$ およびすべての $e \geq 1$ に対する
$\Hom_{D(\mathcal{O}_X)}(K_e[-n], N) = 0$ と同値である。$N|_U = 0$ なので、
同補題によりこれはさらに $n \leq a$ に対する $H^n(X, N) = 0$ と同値である。
$N$ は $A$-加群の複体 $R\Gamma(X, N)$ で決まるので、補題
\ref{perfect-lemma-affine-compare-bounded} により (2) は
$\tau_{\leq a}N = 0$ と同値である。これで補題が示された。
\end{proof}

\begin{lemma}
\label{perfect-lemma-bounded-truncation}
$X$ を準コンパクトかつ準分離なスキームとする。
$P \in D_{perf}(\mathcal{O}_X)$、$E \in D_{\QCoh}(\mathcal{O}_X)$、
$a \in \mathbf{Z}$ とする。次は同値である。
\begin{enumerate}
\item $i \gg 0$ に対して $\Hom_{D(\mathcal{O}_X)}(P[-i], E) = 0$ である。
\item $i \gg 0$ に対して
$\Hom_{D(\mathcal{O}_X)}(P[-i], \tau_{\geq a} E) = 0$ である。
\end{enumerate}
\end{lemma}

\begin{proof}
三角 $ \tau_{< a} E \to E \to \tau_{\geq a} E \to$ を用いると、同値性を
示すには
$$
\Hom_{D(\mathcal{O}_X)}(P[-i], \tau_{< a} E) =
\Hom_{D(\mathcal{O}_X)}(P, (\tau_{< a} E)[i]) = 0 
$$
が $i \gg 0$ に対して成り立つことを示せば十分である。$P$ は完全なので、これは
補題 \ref{perfect-lemma-ext-from-perfect-into-bounded-QCoh} により成り立つ。
\end{proof}

\begin{lemma}
\label{perfect-lemma-bounded-below-truncation}
$X$ を準コンパクトかつ準分離なスキームとする。
$P \in D_{perf}(\mathcal{O}_X)$、$E \in D_{\QCoh}(\mathcal{O}_X)$、
$a \in \mathbf{Z}$ とする。次は同値である。
\begin{enumerate}
\item $i \ll 0$ に対して $\Hom_{D(\mathcal{O}_X)}(P[-i], E) = 0$ である。
\item $i \ll 0$ に対して
$\Hom_{D(\mathcal{O}_X)}(P[-i], \tau_{\leq a} E) = 0$ である。
\end{enumerate}
\end{lemma}

\begin{proof}
三角 $ \tau_{\leq a} E \to E \to \tau_{> a} E \to$ を用いると、同値性を
示すには
$$
\Hom_{D(\mathcal{O}_X)}(P[-i], \tau_{> a} E) =
\Hom_{D(\mathcal{O}_X)}(P, (\tau_{> a} E)[i]) = 0
$$
が $i \ll 0$ に対して成り立つことを示せば十分である。$P$ は完全なので、これは
補題 \ref{perfect-lemma-ext-from-perfect-into-bounded-QCoh} により成り立つ。
\end{proof}

\begin{proposition}
\label{perfect-proposition-detecting-bounded-above}
$X$ を準コンパクトかつ準分離なスキームとする。
$G \in D_{perf}(\mathcal{O}_X)$ を $D_\QCoh (\mathcal{O}_X)$ を生成する完全複体とし、
$E \in D_\QCoh (\mathcal{O}_X)$ とする。次は同値である。
\begin{enumerate}
\item $E \in D^-_\QCoh (\mathcal{O}_X)$,
\item $i \gg 0$ に対して $\Hom_{D(\mathcal{O}_X)}(G[-i], E) = 0$,
\item $i \gg 0$ に対して $\Ext^i_X(G, E) = 0$,
\item $R\Hom_X(G, E)$ は $D^-(\mathbf{Z})$ に属する。
\item $H^i(X, G^\vee \otimes_{\mathcal{O}_X}^\mathbf{L} E) = 0$
が $i \gg 0$ に対して成り立つ。
\item $R\Gamma(X, G^\vee \otimes_{\mathcal{O}_X}^\mathbf{L} E)$
は $D^-(\mathbf{Z})$ に属する。
\item $D(\mathcal{O}_X)$ のすべての完全対象 $P$ に対して、
\begin{enumerate}
\item $G$ を $P$ で置き換えた主張 (2), (3), (4) が成り立つ。
\item $i \gg 0$ に対して
$H^i(X, P \otimes_{\mathcal{O}_X}^\mathbf{L} E) = 0$ である。
\item $R\Gamma(X, P \otimes_{\mathcal{O}_X}^\mathbf{L} E)$
は $D^-(\mathbf{Z})$ に属する。
\end{enumerate}
\end{enumerate}
\end{proposition}

\begin{proof}
(1) を仮定する。
$\Hom_{D(\mathcal{O}_X)}(G[-i], E) = \Hom_{D(\mathcal{O}_X)}(G, E[i])$
であるから、補題 \ref{perfect-lemma-ext-from-perfect-into-bounded-QCoh} によりこれは
$i \gg 0$ に対して零である。よって (1) は (2) を含意する。

\medskip\noindent
Cohomology, Section \ref{cohomology-section-global-RHom} の議論により
(2), (3), (4) は同値である。定義により
$H^i(X, -) = H^i(R\Gamma(X, -))$ なので (5) と (6) は同値である。
Cohomology, Lemma \ref{cohomology-lemma-dual-perfect-complex} と等式
$(G[-i])^\vee = G^\vee[i]$ により、同値な条件 (2), (3), (4) は同値な条件
(5), (6) と同値である。

\medskip\noindent
(7) が (2) を含意することは明らかである。逆に、同値な条件 (2) -- (6) が
(7) を含意することを示そう。Remark \ref{perfect-remark-classical-generator} により
$G$ は $D_{perf}(\mathcal{O}_X)$ の古典的生成対象であることを思い出そう。
$P \in D_{perf}(\mathcal{O}_X)$ に対し、$T(P)$ を
$R\Hom_X(P, E)$ が $D^-(\mathbf{Z})$ に属するという主張とする。明らかに $T$ は
直和と直和因子へ遺伝し、区別三角について二つから三つの性質を満たし、シフトで
保たれる。したがって Derived Categories, Remark
\ref{derived-remark-check-on-generator} により、(4) から $T$ が
$D_{perf}(\mathcal{O}_X)$ 全体で成り立つことが従う。他のすべての性質にも同じ議論が
使えるが、(7)(b) と (7)(c) については $P \mapsto P^\vee$ が
$D_{perf}(\mathcal{O}_X)$ の自己同値であることも用いる。細部を一つ省略する。

\medskip\noindent
Cohomology of Schemes, Lemma \ref{coherent-lemma-induction-principle} の
帰納原理を用いて、同値な条件 (2) -- (7) が (1) を含意することを示す。

\medskip\noindent
まず $X$ がアフィンの場合に (2) -- (7) $\Rightarrow$ (1) を示す。
$P = \mathcal{O}_X[0]$ と置く。(7) から $i \gg 0$ に対して
$H^i (X, E) = 0$ を得る。$E$ は $R\Gamma (X, E)$ で決まるので (1) が従う。
補題 \ref{perfect-lemma-affine-compare-bounded} を参照されたい。

\medskip\noindent
次に $X = U \cup V$ とし、$U$ は $X$ の準コンパクト開、$V$ はアフィン開で、
(2) -- (7) $\Rightarrow$ (1) がスキーム $U$, $V$, $U \cap V$ に対して既知であると
仮定する。$E \in D_\QCoh(\mathcal{O}_X)$ が (2) -- (7) を満たすとする。補題
\ref{perfect-lemma-direct-summand-of-a-restriction} と Theorem
\ref{perfect-theorem-bondal-van-den-Bergh} により、$Q|_U$ が
$D_\QCoh (\mathcal{O}_U)$ を生成するような $X$ 上の完全複体 $Q$ が存在する。
$V$ の部分集合として $V \setminus U = V(f_1, \dots , f_r)$ となる
$f_1, \dots , f_r \in \Gamma (V, \mathcal{O}_V)$ を取る。
$K \in D_{perf}(\mathcal{O}_V)$ を $f_1, \dots , f_r$ に関する Koszul 複体に
対応する対象とする。$K' \in D_{perf}(\mathcal{O}_X)$ を
\begin{equation}
\label{perfect-equation-detecting-bounded-above}
K' = R (V \to X)_* K = R (V \to X)_! K,
\end{equation}
とする。Cohomology, Lemmas \ref{cohomology-lemma-pushforward-restriction}
および \ref{cohomology-lemma-pushforward-perfect} を参照されたい。これは閉集合
$X \setminus U \subset V$ 上に台をもち、$V$ 上で $K$ と同型な $X$ 上の完全複体で
ある。仮定により $R\Hom_{\mathcal{O}_X}(Q, E)$ と
$R\Hom_{\mathcal{O}_X}(K', E)$ は上に有界である。

\medskip\noindent
(\ref{perfect-equation-detecting-bounded-above}) における $K'$ の第二の記述により
$$
\Hom_{D(\mathcal{O}_V)}(K[-i], E|_V) = \Hom_{D(\mathcal{O}_X)}(K'[-i], E) = 0
$$
が $i \gg 0$ に対して成り立つ。したがって補題
\ref{perfect-lemma-orthogonal-koszul-first-variant} を $E|_V$ に適用して、
$\tau_{\geq a}(E|_V) = \tau_{\geq a} R (U \cap V \to V)_* (E|_{U \cap V})$.
となる整数 $a$ を得る。このとき
$\tau_{\geq a} E = \tau_{\geq a} R (U \to X)_* (E |_U)$ である
（標準写像が $U$ と $V$ への制限後に同型となることを確認すればよい）。
したがって補題 \ref{perfect-lemma-bounded-truncation} を二度用いると
$$
\Hom_{D(\mathcal{O}_U)}(Q|_U [-i], E|_U) =
\Hom_{D(\mathcal{O}_X)}(Q[-i], R (U \to X)_* (E|_U)) = 0
$$
が $i \gg 0$ に対して成り立つ。命題は $U$ と生成対象 $Q|_U$ に対して成り立つので、
$E|_U \in D^-_\QCoh(\mathcal{O}_U)$ である。一方、補題
\ref{perfect-lemma-quasi-coherence-direct-image} により関手 $R (U \to X)_*$ は
$D^-_\QCoh$ を保つので、$\tau_{\geq a}E \in D^-_\QCoh(\mathcal{O}_X)$ を得る。
したがって $E \in D^-_\QCoh (\mathcal{O}_X)$ である。
\end{proof}

\begin{proposition}
\label{perfect-proposition-detecting-bounded-below}
$X$ を準コンパクトかつ準分離なスキームとする。
$G \in D_{perf}(\mathcal{O}_X)$ を $D_\QCoh (\mathcal{O}_X)$ を生成する完全複体とし、
$E \in D_\QCoh (\mathcal{O}_X)$ とする。次は同値である。
\begin{enumerate}
\item $E \in D^+_\QCoh (\mathcal{O}_X)$,
\item $i \ll 0$ に対して $\Hom_{D(\mathcal{O}_X)}(G[-i], E) = 0$,
\item $i \ll 0$ に対して $\Ext^i_X(G, E) = 0$,
\item $R\Hom_X(G, E)$ は $D^+(\mathbf{Z})$ に属する。
\item $H^i(X, G^\vee \otimes_{\mathcal{O}_X}^\mathbf{L} E) = 0$
が $i \ll 0$ に対して成り立つ。
\item $R\Gamma(X, G^\vee \otimes_{\mathcal{O}_X}^\mathbf{L} E)$
は $D^+(\mathbf{Z})$ に属する。
\item $D(\mathcal{O}_X)$ のすべての完全対象 $P$ に対して、
\begin{enumerate}
\item $G$ を $P$ で置き換えた主張 (2), (3), (4) が成り立つ。
\item $i \ll 0$ に対して
$H^i(X, P \otimes_{\mathcal{O}_X}^\mathbf{L} E) = 0$ である。
\item $R\Gamma(X, P \otimes_{\mathcal{O}_X}^\mathbf{L} E)$
は $D^+(\mathbf{Z})$ に属する。
\end{enumerate}
\end{enumerate}
\end{proposition}

\begin{proof}
(1) を仮定する。
$\Hom_{D(\mathcal{O}_X)}(G[-i], E) = \Hom_{D(\mathcal{O}_X)}(G, E[i])$
であるから、補題 \ref{perfect-lemma-ext-from-perfect-into-bounded-QCoh} によりこれは
$i \ll 0$ に対して零である。よって (1) は (2) を含意する。

\medskip\noindent
Cohomology, Section \ref{cohomology-section-global-RHom} の議論により
(2), (3), (4) は同値である。定義により
$H^i(X, -) = H^i(R\Gamma(X, -))$ なので (5) と (6) は同値である。
Cohomology, Lemma \ref{cohomology-lemma-dual-perfect-complex} と等式
$(G[-i])^\vee = G^\vee[i]$ により、同値な条件 (2), (3), (4) は同値な条件
(5), (6) と同値である。

\medskip\noindent
(7) が (2) を含意することは明らかである。逆に、同値な条件 (2) -- (6) が
(7) を含意することを示そう。Remark \ref{perfect-remark-classical-generator} により
$G$ は $D_{perf}(\mathcal{O}_X)$ の古典的生成対象であることを思い出そう。
$P \in D_{perf}(\mathcal{O}_X)$ に対し、$T(P)$ を
$R\Hom_X(P, E)$ が $D^+(\mathbf{Z})$ に属するという主張とする。明らかに $T$ は
直和と直和因子へ遺伝し、区別三角について二つから三つの性質を満たし、シフトで
保たれる。したがって Derived Categories, Remark
\ref{derived-remark-check-on-generator} により、(4) から $T$ が
$D_{perf}(\mathcal{O}_X)$ 全体で成り立つことが従う。他のすべての性質にも同じ議論が
使えるが、(7)(b) と (7)(c) については $P \mapsto P^\vee$ が
$D_{perf}(\mathcal{O}_X)$ の自己同値であることも用いる。細部を一つ省略する。

\medskip\noindent
Cohomology of Schemes, Lemma \ref{coherent-lemma-induction-principle} の
帰納原理を用いて、同値な条件 (2) -- (7) が (1) を含意することを示す。

\medskip\noindent
まず $X$ がアフィンの場合に (2) -- (7) $\Rightarrow$ (1) を示す。
$P = \mathcal{O}_X[0]$ と置く。(7) から $i \ll 0$ に対して
$H^i (X, E) = 0$ を得る。$E$ は $R\Gamma (X, E)$ で決まるので (1) が従う。
補題 \ref{perfect-lemma-affine-compare-bounded} を参照されたい。

\medskip\noindent
次に $X = U \cup V$ とし、$U$ は $X$ の準コンパクト開、$V$ はアフィン開で、
(2) -- (7) $\Rightarrow$ (1) がスキーム $U$, $V$, $U \cap V$ に対して既知であると
仮定する。$E \in D_\QCoh(\mathcal{O}_X)$ が (2) -- (7) を満たすとする。補題
\ref{perfect-lemma-direct-summand-of-a-restriction} と Theorem
\ref{perfect-theorem-bondal-van-den-Bergh} により、$Q|_U$ が
$D_\QCoh (\mathcal{O}_U)$ を生成するような $X$ 上の完全複体 $Q$ が存在する。
$V$ の部分集合として $V \setminus U = V(f_1, \dots , f_r)$ となる
$f_1, \dots , f_r \in \Gamma (V, \mathcal{O}_V)$ を取る。
$K \in D_{perf}(\mathcal{O}_V)$ を $f_1, \dots , f_r$ に関する Koszul 複体に
対応する対象とする。$K' \in D_{perf}(\mathcal{O}_X)$ を
\begin{equation}
\label{perfect-equation-detecting-bounded-below}
K' = R (V \to X)_* K = R (V \to X)_! K,
\end{equation}
とする。Cohomology, Lemmas \ref{cohomology-lemma-pushforward-restriction}
および \ref{cohomology-lemma-pushforward-perfect} を参照されたい。これは閉集合
$X \setminus U \subset V$ 上に台をもち、$V$ 上で $K$ と同型な $X$ 上の完全複体で
ある。仮定により $R\Hom_{\mathcal{O}_X}(Q, E)$ と
$R\Hom_{\mathcal{O}_X}(K', E)$ は下に有界である。

\medskip\noindent
(\ref{perfect-equation-detecting-bounded-below}) における $K'$ の第二の記述により
$$
\Hom_{D(\mathcal{O}_V)}(K[-i], E|_V) = \Hom_{D(\mathcal{O}_X)}(K'[-i], E) = 0
$$
が $i \ll 0$ に対して成り立つ。したがって補題
\ref{perfect-lemma-orthogonal-koszul-second-variant} を $E|_V$ に適用して、
$\tau_{\leq a}(E|_V) = \tau_{\leq a} R(U \cap V \to V)_*(E|_{U \cap V})$.
となる整数 $a$ を得る。このとき
$\tau_{\leq a} E = \tau_{\leq a} R(U \to X)_*(E|_U)$ である
（標準写像が $U$ と $V$ への制限後に同型となることを確認すればよい）。
したがって補題 \ref{perfect-lemma-bounded-below-truncation} を二度用いると
$$
\Hom_{D(\mathcal{O}_U)}(Q|_U [-i], E|_U) =
\Hom_{D(\mathcal{O}_X)}(Q[-i], R (U \to X)_* (E|_U)) = 0
$$
が $i \ll 0$ に対して成り立つ。命題は $U$ と生成対象 $Q|_U$ に対して成り立つので、
$E|_U \in D^+_\QCoh (\mathcal{O}_U)$ である。一方、Cohomology, Section
\ref{cohomology-section-derived-functors} により関手 $R(U \to X)_*$ は下に有界な
対象を保つので、$\tau_{\leq a} E \in D^+_\QCoh(\mathcal{O}_X)$ を得る。
したがって $E \in D^+_\QCoh (\mathcal{O}_X)$ である。
\end{proof}










\section{導来圏の準連接対象}
\label{perfect-section-QC}

\noindent
$X$ をスキームとする。$X_{affine, Zar}$ は $X$ のアフィン開集合の圏で、
標準 Zariski 被覆が定める位相を備えたものを表すことを思い出そう。Topologies,
Definition \ref{topologies-definition-big-small-Zariski} を参照されたい。
$X_{affine, Zar}$ のトポスは $X$ の小 Zariski トポスである。Topologies, Lemma
\ref{topologies-lemma-alternative-zariski} を参照されたい。サイト
$X_{affine, Zar}$ は構造層 $\mathcal{O}$ をもち、環付きトポスの同値
$$
(\Sh(X_{affine, Zar}), \mathcal{O})
\longrightarrow
(\Sh(X_{Zar}), \mathcal{O})
$$
がある。Descent, Equation (\ref{descent-equation-alternative-small-ringed})
および、少し異なる記法を用いた Descent, Section
\ref{descent-section-alternative-quasi-coherent} のその周辺の議論を
参照されたい。

\medskip\noindent
本節では $X_{affine, Zar}$ の基礎圏に乱雑位相、すなわち層と前層が一致する
位相を入れたものを $X_{affine}$ と記す。特に構造層 $\mathcal{O}$ は
$X_{affine}$ 上でも層となる。環付きサイトの射
$$
\epsilon :
(X_{affine, Zar}, \mathcal{O})
\longrightarrow
(X_{affine}, \mathcal{O})
$$
を得る。Cohomology on Sites, Section
\ref{sites-cohomology-section-compare} を参照されたい。本節では
$D_\QCoh(\mathcal{O}_X)$ を Cohomology on Sites, Section
\ref{sites-cohomology-section-modules-cohomology} で導入された圏
$\mathit{QC}(X_{affine}, \mathcal{O})$ と同一視する。

\begin{lemma}
\label{perfect-lemma-DQCoh-alternative-small}
上の状況において、次の三角圏の間には標準的な完全同値がある。
\begin{enumerate}
\item $D_\QCoh(\mathcal{O}_X)$,
\item $D_\QCoh(X_{Zar}, \mathcal{O})$,
\item $D_\QCoh(X_{affine, Zar}, \mathcal{O})$,
\item $D_\QCoh(X_{affine}, \mathcal{O}_X)$、
\item $\mathit{QC}(X_{affine}, \mathcal{O})$.
\end{enumerate}
\end{lemma}

\begin{proof}
$U \subset V \subset X$ がアフィン開なら、環準同型
$\mathcal{O}(V) \to \mathcal{O}(U)$ は平坦である。したがって (4) と (5) の
同値は Cohomology on Sites, Lemma
\ref{sites-cohomology-lemma-cartesian-flat-quasi-coherent} の特別な場合である
（その証明は主張の意味も明らかにする）。

\medskip\noindent
Descent, Remark \ref{descent-remark-Zariski-site-space} により、環付きサイト
$(X_{Zar}, \mathcal{O})$ と環付き空間 $(X, \mathcal{O}_X)$ は同じ加群圏をもつ。
Descent, Proposition
\ref{descent-proposition-equivalence-quasi-coherent} により、この同値のもとで
準連接加群が対応する。したがって (1) と (2) の三角圏の間に標準的な完全同値を
得る。

\medskip\noindent
補題の前の議論により、環付きトポスの同値
$(\Sh(X_{affine, Zar}), \mathcal{O}) \to (\Sh(X_{Zar}), \mathcal{O})$
があり、したがって加群のアーベル圏の間に同値がある。準連接加群の概念は内在的
なので（Modules on Sites, Lemma
\ref{sites-modules-lemma-special-locally-free}）、この同値は準連接加群の
部分圏を保つ。よって (2) と (3) の三角圏の間に標準的な完全同値を得る。

\medskip\noindent
(3) と (4) の三角圏の間の完全同値を得るため、Cohomology on Sites, Lemma
\ref{sites-cohomology-lemma-compare-topologies-derived-adequate-modules} を
上の射 $\epsilon : (X_{affine, Zar}, \mathcal{O}) \to
(X_{affine}, \mathcal{O})$ に適用する。$\mathcal{B} = \Ob(X_{affine})$ とし、
$\mathcal{A} \subset \textit{PMod}(X_{affine}, \mathcal{O})$ を、
$\mathcal{F}(V) \otimes_{\mathcal{O}(V)} \mathcal{O}(U)
\to \mathcal{F}(U)$ が同型となる前層 $\mathcal{F}$ の充満部分圏とする。
Descent, Lemma \ref{descent-lemma-quasi-coherent-alternative-small} により、
$\mathcal{A}$ の対象はちょうど Zariski 位相で層をなし、
$(X_{affine, Zar}, \mathcal{O})$ 上の準連接加群であるものに等しい。一方、
Modules on Sites, Lemma \ref{sites-modules-lemma-chaotic-quasi-coherent} により、
$\mathcal{A}$ の対象はちょうど $(X_{affine}, \mathcal{O})$ 上、すなわち乱雑位相での
準連接加群に等しい。したがって Cohomology on Sites, Lemma
\ref{sites-cohomology-lemma-compare-topologies-derived-adequate-modules} が
適用できることを示せば、$\epsilon^*$ と $R\epsilon_*$ を用いて実際に (3) と (4) の
圏の間の標準同値を得る。

\medskip\noindent
次の 4 条件を確認しなければならない。
\begin{enumerate}
\item $\mathcal{A}$ のすべての対象は Zariski 位相の層である。これは上で見た。
\item $\mathcal{A}$ は $\textit{Mod}(X_{affine, Zar}, \mathcal{O})$ の弱 Serre
部分圏である。上で $\mathcal{A} = \QCoh(X_{affine, Zar}, \mathcal{O})$ を見た。
また同値 $\textit{Mod}(X_{affine, Zar}, \mathcal{O}) =
\textit{Mod}(X, \mathcal{O}_X)$ のもとで、これらが $X$ 上の準連接加群に対応する
ことも見た。したがって結果は Schemes, Section
\ref{schemes-section-quasi-coherent} の議論から従う。
\item $X_{affine}$ のすべての対象は、その要素が $\mathcal{B}$ に属する乱雑位相での
被覆をもつ。$\mathcal{B}$ はすべての対象を含むので、これは成り立つ。
\item $X_{affine}$ のすべての対象 $U$ と $\mathcal{A}$ のすべての
$\mathcal{F}$ に対し、$p > 0$ なら $H^p_{Zar}(U, \mathcal{F}) = 0$ である。
これはアフィン上の準連接加群のコホモロジーの消滅から従う。Cohomology of Schemes,
Lemma \ref{coherent-lemma-quasi-coherent-affine-cohomology-zero} を参照されたい。
\end{enumerate}
これで証明が完了する。
\end{proof}

\begin{remark}
\label{perfect-remark-QC-compare}
$S$ をスキームとする。後で $\mathit{QC}((\textit{Aff}/S), \mathcal{O})$ も
$D_\QCoh(\mathcal{O}_S)$ と標準的に同値であることを示す。Sheaves on Stacks,
Proposition \ref{stacks-sheaves-proposition-QC-compare} を参照されたい。
\end{remark}
